\documentclass[10pt,twoside,openright]{book}
\usepackage{../styles/neet_book}

\begin{document}

% =========================================================================
% FRONT MATTER
% =========================================================================
\frontmatter

% --- Half-Title Page ---
\thispagestyle{empty}
\vspace*{4cm}
\begin{center}
  {\Huge\bfseries\color{primarydark} NEET UG PHYSICS}\\[0.6cm]
  {\Large\bfseries\color{accentblue} MASTER TEXTBOOK}\\[1.5cm]
  {\normalsize\color{gray} NEET UG 2027 Master Reference Series $\cdot$ Complete Authoritative Edition}
\end{center}
\cleardoublepage

% --- Full Title Page ---
\thispagestyle{empty}
\vspace*{2cm}
\begin{center}
  {\Huge\bfseries\color{primarydark} NEET UG PHYSICS}\\[0.4cm]
  {\LARGE\bfseries\color{primarydark} MASTER TEXTBOOK}\\[0.8cm]
  {\Large\bfseries\color{accentblue} Mechanics, Thermodynamics, Electrodynamics, Optics,}\\[0.2cm]
  {\Large\bfseries\color{accentblue} Modern Physics, Experimental Skills \& High-Yield Formulas}\\[2cm]
  
  {\large\textbf{Authored by}}\\[0.3cm]
  {\LARGE\bfseries\color{primarydark} Satya Prakash}\\[1.5cm]

  {\normalsize\textbf{Strictly Calibrated to the Revised NMC Syllabus \& Rationalized NCERT Textbooks}}\\[0.3cm]
  {\small Target Examination: National Eligibility cum Entrance Test (NEET UG 2027)}\\[2.5cm]

  {\footnotesize Published for Independent Medical Aspirants, Educators \& Academic Study Circles}\\[0.2cm]
  {\footnotesize Comprehensive Academic Edition $\cdot$ Global Print-on-Demand}
\end{center}
\cleardoublepage

% --- Professional Copyright Page ---
\thispagestyle{empty}
\vspace*{\fill}
\begin{flushleft}
{\small
\textbf{NEET UG PHYSICS MASTER TEXTBOOK}\\[2pt]
\textit{Comprehensive Theory, Rigorous Derivations, Experimental Skills \& Formula Compendium}\\[8pt]
\textbf{Author:} Satya Prakash\\
\textbf{Editorial Correspondence \& Academic Feedback:} \href{mailto:spdedsec@hotmail.com}{spdedsec@hotmail.com}\\
\textbf{Source Repository \& Official Errata:} \url{https://github.com/spdedsec/neet-study-materials}\\[10pt]

\copyright\ 2026--2027 Satya Prakash. All rights reserved.\\[8pt]

\textbf{First Edition:} October 2026\\
\textbf{Publishing Imprint:} Independent Academic Publishing / Satya Prakash Editorial Imprint\\
\textbf{Target Cycle:} NEET UG 2027 (Undergraduate Medical Admissions Entrance Examination)\\
\textbf{ISBN-13 (Paperback):} 979-8-XXXXXXXXX (Placeholder for Official Assigned ISBN)\\
\textbf{ISBN-13 (Hardcover):} 979-8-XXXXXXXXX (Placeholder for Official Assigned ISBN)\\[10pt]

\textbf{All Rights Reserved.} No part of this publication may be reproduced, stored in a retrieval system, or transmitted in any form or by any means---electronic, mechanical, photocopying, recording, scanning, or otherwise---without the prior written permission of the author, except for brief quotations embodied in critical reviews or scholarly analysis as permitted under copyright law.\\[8pt]

\textbf{Academic \& Legal Disclaimer:} This textbook is an independent preparation manual authored by Satya Prakash to assist candidates preparing for the National Eligibility cum Entrance Test (NEET UG). While every effort has been made to verify scientific accuracy and strict concordance with the National Medical Commission (NMC) curriculum and rationalized NCERT textbooks, the author and publisher make no official claims of affiliation with or endorsement by the National Testing Agency (NTA), the National Medical Commission (NMC), or the Ministry of Education, Government of India. The author disclaims any liability for errors or omissions.\\[8pt]

\textbf{Typesetting \& Typography:} Typeset in \LaTeXe\ with the Tectonic XeTeX engine using Latin Modern typefaces. Mathematical typesetting follows ISO 80000-2 standards. Produced and bound under rigorous academic publishing standards.\\[4pt]
Printed in the United States of America / United Kingdom / India.
}
\end{flushleft}
\cleardoublepage

% --- Dedication Page ---
\thispagestyle{empty}
\vspace*{5cm}
\begin{center}
  \itshape\large
  To my parents, whose quiet sacrifices and endless belief built the foundation of my life;\\
  \vspace{1.2cm}
  To the teachers and mentors who showed me that physics is not a collection of arbitrary formulas,\\
  but the poetry of how nature actually works;\\
  \vspace{1.2cm}
  And to you, the medical aspirant reading this in the quiet hours of the night---\\
  may every solved problem and mastered derivation bring you one step closer\\
  to the stethoscope, the white coat, and a lifetime of healing.
\end{center}
\cleardoublepage

% --- Acknowledgements ---
\chapter*{Acknowledgements}
\addcontentsline{toc}{chapter}{Acknowledgements}

Creating a complete, mathematically uncompromising textbook calibrated to a competitive examination as fiercely contested as NEET UG is never a solitary effort. I owe an immense debt of gratitude to several people and institutions whose influence echoes through these chapters.

First and foremost, I thank my family for their patience, understanding, and unwavering encouragement during the countless late nights and weekends spent drafting, proofreading, and verifying derivations.

I express deep gratitude to the mentors and physics educators whose classroom discussions shaped my conviction that physics must be taught from fundamental principles rather than mechanical memorization. Their insistence on clear vector diagrams, rigorous free-body analysis, and boundary condition checks inspired the structure of every chapter in this book.

I am especially grateful to the students and peer study groups whose probing questions, skepticism, and frequent exam mistakes illuminated the exact cognitive hurdles aspirants face. The "NEET Trap" boxes in this book were born directly from real student doubts and mock exam error logs.

Finally, I acknowledge the vibrant open-source scientific typesetting community, specifically the developers of the \LaTeXe\ ecosystem and the Tectonic compiler, which made it possible to produce a publication of this typographical caliber.

\vspace{1cm}
\begin{flushright}
  \textbf{Satya Prakash}\\
  \textit{October 2026}
\end{flushright}
\cleardoublepage

% --- Author's Preface ---
\chapter*{Author's Preface: Straight Talk Before We Begin}
\addcontentsline{toc}{chapter}{Author's Preface: Straight Talk Before We Begin}

Let's be completely frank with each other. If you are preparing for NEET UG, you already know that physics is where ranks are decided. Biology gives you the foundation, Chemistry keeps you in the game, but Physics is what separates a student struggling for a cutoff from one who walks into a top government medical college.

Yet for most aspirants, physics feels like a daily war. You sit down with coaching modules running over 800 pages, and half the material is either filled with fluff or contains obsolete topics that the National Medical Commission (NMC) eliminated years ago. You waste precious study hours memorizing Carnot refrigerator formulas, potentiometer wire balancing equations, and Earth magnetism tables that will not appear on your exam paper. Meanwhile, when the actual exam twists a rolling friction problem on an inclined plane or asks an Assertion-Reasoning question on Gauss's law with non-symmetric charges, formulas from memory suddenly fail.

I wrote this book because I wanted to hand you the exact companion I wished I had when studying: honest, rigorous, beautifully organized, and talking to you straight without any fake "miracle tricks" or empty motivation. 

Here is how this book is built, and how you should use it:

\begin{enumerate}[leftmargin=*]
  \item \textbf{Bedrock Principles Over Rote Memorization:} Every major chapter starts from first principles. When we look at work-energy, fluid dynamics, or electromagnetic induction, you will see the complete mathematical derivation. Why? Because when you understand where a formula comes from and what assumptions were made to reach it, you become immune to the traps exam setters lay for unsuspecting students.
  
  \item \textbf{The Five Tactical Callout Systems:} Throughout every chapter, you will encounter five distinct color-coded boxes. Treat them like field markers:
  \begin{itemize}
    \item \textbf{Core Principle Boxes (\texttt{conceptbox}):} The foundational laws and definitions. Pay close attention to these---NTA draws multi-statement and factual questions directly from these definitions.
    \item \textbf{Master Formula Boxes (\texttt{formulabox}):} Every critical governing equation with explicit variable definitions, limits of integration, and SI units.
    \item \textbf{Application Regime Boxes (\texttt{usecasebox}):} The real-world conditions where formulas work, where they break down, and shortcuts for common numerical situations.
    \item \textbf{NEET Trap Boxes (\texttt{trapbox}):} The exact spots where thousands of students lose 5 marks (-1 penalty plus the 4 marks they missed). These cover sign convention mistakes, hidden assumptions, and misleading exceptions.
    \item \textbf{Direct NCERT Lines (\texttt{ncertbox}):} Verbatim textual extractions and table data that frequently form the basis of Assertion-Reasoning items.
  \end{itemize}

  \item \textbf{Strict NMC Syllabus Calibration:} Everything in this book conforms strictly to the revised NMC syllabus for NEET UG 2027. Obsolete topics like Carnot engines, free and damped oscillations, potentiometers, and transistor amplifiers have been flagged and excised from the core theory. You will not waste a single hour on deleted material.

  \item \textbf{Experimental Skills (Unit 20) Integration:} The NMC added explicit weightage to practical physics. Chapter 6 is dedicated entirely to the 18 prescribed experimental skills---from Vernier calipers and screw gauges to sonometer wires, focal length determination, and meter bridges.
\end{enumerate}

A personal piece of advice: do not read physics like a novel. If you are sitting on your bed flipping pages without a pen in your hand, you are not studying physics. Pick up a notebook, derive the equations alongside the text, draw every free-body diagram yourself, and work through the problem archetypes with the solutions covered.

Mistakes are going to happen. When you get a question wrong in a practice test, do not get demoralized. An error made in practice is a gift: it reveals a blind spot today so you don't repeat it on exam day.

You have the capability and the time to master this subject. Stay consistent, stay disciplined, and keep your focus on the prize. I'll see you on the other side with that white coat.

\vspace{1cm}
\begin{flushright}
  \textbf{Satya Prakash}\\
  Author \& Dedicated Mentor\\
  \href{mailto:spdedsec@hotmail.com}{spdedsec@hotmail.com}\\
  \textit{October 2026}
\end{flushright}
\cleardoublepage

% --- Syllabus Blueprint & NMC Rationalization Guide ---
\chapter*{Syllabus Blueprint \& NMC Rationalization Guide}
\addcontentsline{toc}{chapter}{Syllabus Blueprint \& NMC Rationalization Guide}

The National Eligibility cum Entrance Test (NEET UG 2027) follows the curriculum prescribed by the Under Graduate Medical Education Board (UGMEB) of the National Medical Commission (NMC). This section outlines the structure of the examination, subject weightage, and critical syllabus deletions.

\section*{Official Physics Examination Architecture}

\begin{table}[htbp]
\centering
\small
\renewcommand{\arraystretch}{1.2}
\begin{tabularx}{\textwidth}{l c c c c Y}
\toprule
\textbf{Component} & \textbf{Section A} & \textbf{Section B} & \textbf{Total Questions} & \textbf{Marks} & \textbf{Question Typology} \\
\midrule
\textbf{Physics} & 35 & 15 (Attempt 10) & 45 & 180 & Single Choice, MCQs, \\
\textbf{Exam Total} & 140 & 60 (Attempt 40) & 180 & 720 & Assertion-Reason, Matrix Match \\
\bottomrule
\end{tabularx}
\caption{NEET UG Physics Examination Scheme (+4 for correct response, $-1$ for incorrect response, 0 for unattempted).}
\end{table}

\section*{Official 20-Unit Physics Syllabus Structure}

\begin{enumerate}[leftmargin=*]
  \item \textbf{Unit 1: Physics and Measurement} (Units, dimensions, error propagation, significant figures, Vernier calipers, screw gauge).
  \item \textbf{Unit 2: Kinematics} (Frame of reference, rectilinear motion, relative velocity, 2D projectile motion, uniform circular motion).
  \item \textbf{Unit 3: Laws of Motion} (Newton's laws, momentum conservation, equilibrium, friction dynamics, circular dynamics and banking).
  \item \textbf{Unit 4: Work, Energy and Power} (Work-energy theorem, conservative forces, potential energy of springs, vertical circular motion, collisions).
  \item \textbf{Unit 5: Rotational Motion} (Centre of mass, moment of inertia, parallel and perpendicular axis theorems, torque, angular momentum, pure rolling).
  \item \textbf{Unit 6: Gravitation} (Kepler's laws, universal gravitation, acceleration due to gravity variation, gravitational potential and energy, orbital velocity, escape speed).
  \item \textbf{Unit 7: Properties of Solids and Liquids} (Elasticity, Hooke's law, Young's/bulk/shear moduli, Pascal's law, Archimedes principle, surface tension, capillarity, viscosity, Stokes' law, Bernoulli's theorem).
  \item \textbf{Unit 8: Thermodynamics} (Thermal equilibrium, zeroth law, heat, work, first law of thermodynamics, isothermal, adiabatic, isobaric, isochoric processes, second law).
  \item \textbf{Unit 9: Kinetic Theory of Gases} (Equation of state of perfect gas, kinetic postulates, pressure formulation, kinetic interpretation of temperature, rms speed, degrees of freedom, equipartition of energy, mean free path).
  \item \textbf{Unit 10: Oscillations and Waves} (Periodic motion, SHM kinematics and dynamics, spring-mass systems, simple pendulum, wave motion, longitudinal and transverse waves, superposition, standing waves in strings and pipes, beats).
  \item \textbf{Unit 11: Electrostatics} (Coulomb's law, electric field, electric dipole, Gauss's law and applications, electric potential, equipotential surfaces, capacitors, dielectrics).
  \item \textbf{Unit 12: Current Electricity} (Electric current, drift velocity, Ohm's law, resistivity, temperature dependence, Kirchhoff's laws, Wheatstone bridge, internal resistance).
  \item \textbf{Unit 13: Magnetic Effects of Current and Magnetism} (Biot-Savart law, Ampère's circuital law, force on moving charges and current-carrying conductors, magnetic dipole, moving coil galvanometer).
  \item \textbf{Unit 14: Electromagnetic Induction and Alternating Currents} (Faraday's laws, Lenz's law, eddy currents, self and mutual inductance, AC circuits, impedance, LCR series circuits, resonance, power in AC, transformers).
  \item \textbf{Unit 15: Electromagnetic Waves} (Displacement current, Maxwell's equations characteristics, electromagnetic spectrum from radio waves to gamma rays).
  \item \textbf{Unit 16: Optics} (Ray optics: reflection, refraction, total internal reflection, lens maker's formula, prisms, microscopes, telescopes; Wave optics: Huygens' principle, interference, Young's double slit, diffraction at single slit).
  \item \textbf{Unit 17: Dual Nature of Matter and Radiation} (Photoelectric effect, Hertz and Lenard's observations, Einstein's photoelectric equation, de Broglie relation).
  \item \textbf{Unit 18: Atoms and Nuclei} (Alpha-particle scattering, Rutherford model, Bohr model, energy levels, hydrogen spectrum, atomic nuclei, mass defect, binding energy per nucleon).
  \item \textbf{Unit 19: Electronic Devices} (Semiconductors, energy bands, intrinsic and extrinsic semiconductors, p-n junction diode, I-V characteristics, diode as rectifier, logic gates: OR, AND, NOT, NAND, NOR).
  \item \textbf{Unit 20: Experimental Skills} (18 prescribed experiments including Vernier calipers, screw gauge, simple pendulum, Metre Scale principle of moments, Young's modulus by Searle's method, surface tension by capillary rise, viscosity of liquid, sonometer, resonance tube, Ohm's law, meter bridge, p-n junction diode characteristics).
\end{enumerate}

\begin{trapbox}[Zero-Tolerance NMC Exclusions: Topics Strictly Deleted from Physics]
Candidates must note that legacy textbooks and coaching modules contain several topics that have been \textbf{completely eliminated} by the NMC:
\begin{itemize}
  \item \textbf{Carnot Engine, Carnot Cycle \& Refrigerators} (Completely eliminated from Thermodynamics).
  \item \textbf{Free, Damped \& Forced Oscillations, Resonance} (Eliminated from Oscillations).
  \item \textbf{Doppler Effect in Sound} (Eliminated from Waves).
  \item \textbf{Potentiometer and its applications} (Replaced in practicals by Meter Bridge and Multimeter).
  \item \textbf{Earth's Magnetism and Magnetic Elements} (Eliminated from Magnetism).
  \item \textbf{Cyclotron} (Eliminated from Moving Charges).
  \item \textbf{Radioactive Decay Law, Alpha/Beta/Gamma Decay, Half-Life} (Eliminated from Nuclear Physics).
  \item \textbf{Transistor and its characteristics, Transistor as Amplifier/Oscillator} (Eliminated from Semiconductors).
\end{itemize}
\textbf{Preparation Directive:} Do not waste valuable study time practicing numericals or derivations on these excluded topics.
\end{trapbox}

\section*{Empirical Question Weightage Matrix}

\begin{table}[htbp]
\centering
\small
\renewcommand{\arraystretch}{1.2}
\begin{tabularx}{\textwidth}{L c c}
\toprule
\textbf{Physics Domain} & \textbf{Typical Questions} & \textbf{Marks Weightage} \\
\midrule
Classical Mechanics (Kinematics, Dynamics, Work-Energy, Rotation, Gravitation, Fluids) & 16--18 & 64--72 Marks \\
Thermal Physics, Thermodynamics, Kinetic Theory \& Oscillations & 6--8 & 24--32 Marks \\
Electrodynamics, Magnetism, EMI \& AC Circuits & 11--13 & 44--52 Marks \\
Optics (Ray \& Wave Optics) & 4--5 & 16--20 Marks \\
Modern Physics \& Semiconductors (Dual Nature, Atoms, Nuclei, Diodes/Gates) & 6--8 & 24--32 Marks \\
Experimental Skills (Unit 20 Laboratory Practicals) & 2--3 & 8--12 Marks \\
\midrule
\textbf{Total Physics Paper} & \textbf{45 Attempted / 50} & \textbf{180 Marks} \\
\bottomrule
\end{tabularx}
\caption{Empirical Question Distribution in Recent NEET UG Papers.}
\end{table}

\cleardoublepage

% --- Table of Contents ---
\setcounter{tocdepth}{2}
\setcounter{secnumdepth}{2}
\tableofcontents
\cleardoublepage

% =========================================================================
% MAIN MATTER
% =========================================================================
\mainmatter

% =========================================================================
% PART I: CLASSICAL MECHANICS & FLUID DYNAMICS
% =========================================================================
\part{Classical Mechanics \& Fluid Dynamics}

\chapter{Foundations of Mechanics: Kinematics, Dynamics, Work, Energy \& Collisions}

This chapter establishes the core mechanics principles for NEET UG. We begin with the mathematical foundation of physical measurements, dimensional homogeneity, and error calculus, progressing into rectilinear and two-dimensional kinematics, Newton's laws of motion with constrained multi-block dynamics, microscopic friction mechanisms, vertical and banked circular motion, and the calculus of work, energy, and collision impulse.

\section{Physical Measurement, Dimensional Analysis \& Calculus of Errors}

\begin{conceptbox}[Foundational Standards of Measurement \& Dimensional Homogeneity]
Physical quantities are classified into base (fundamental) and derived quantities. The International System of Units (SI) defines seven base quantities: Mass $[\mathrm{M}]$, Length $[\mathrm{L}]$, Time $[\mathrm{T}]$, Electric Current $[\mathrm{I}]$ or $[\mathrm{A}]$, Thermodynamic Temperature $[\mathrm{K}]$, Amount of Substance $[\mathrm{mol}]$, and Luminous Intensity $[\mathrm{cd}]$.
\begin{enumerate}[leftmargin=*]
  \item \textbf{Principle of Dimensional Homogeneity:} An equation representing physical law $A + B = C$ is physically valid only if $[A] = [B] = [C]$. Transcendental functions (trigonometric $\sin\theta$, exponential $e^x$, logarithmic $\ln x$) must have strictly dimensionless arguments: $[x] = \mathrm{M}^0\mathrm{L}^0\mathrm{T}^0 = 1$.
  \item \textbf{Dimensional Analysis Capabilities \& Inherent Limitations:}
  \begin{itemize}
    \item Useful for deducing power-law dependencies among independent variables, checking algebraic consistency, and converting numerical values between unit systems.
    \item Cannot determine dimensionless numerical constants (e.g., $2\pi$ or $1/2$), cannot solve equations involving more variables than the number of fundamental dimensions (typically 3 in mechanics: $\mathrm{M, L, T}$), and fails for relations involving vector addition or sum of terms.
  \end{itemize}
\end{enumerate}
\end{conceptbox}

\begin{formulabox}[Calculus of Error Propagation \& Uncertainty Formulations]
Let a derived quantity $Z$ depend on independently measured physical variables $A, B, C$ via the functional relation $Z = f(A, B, C)$. By taking the total differential:
\begin{equation}
  dZ = \left(\frac{\partial f}{\partial A}\right) dA + \left(\frac{\partial f}{\partial B}\right) dB + \left(\frac{\partial f}{\partial C}\right) dC
\end{equation}
For experimental uncertainty analysis, maximum absolute error $\Delta Z$ is determined by summing the absolute magnitudes of the partial deviations:
\begin{align}
  \text{Sum / Difference: } Z = A \pm B &\implies \Delta Z = \Delta A + \Delta B \\
  \text{Product / Quotient: } Z = \frac{A \cdot B}{C} &\implies \frac{\Delta Z}{Z} = \frac{\Delta A}{A} + \frac{\Delta B}{B} + \frac{\Delta C}{C} \\
  \text{Generalized Power Formulation: } Z = \frac{A^p B^q}{C^r} &\implies \frac{\Delta Z}{Z} = p\left(\frac{\Delta A}{A}\right) + q\left(\frac{\Delta B}{B}\right) + r\left(\frac{\Delta C}{C}\right)
\end{align}
The percentage fractional uncertainty is given by:
\begin{equation}
  \% \Delta Z = \left[ p\left(\frac{\Delta A}{A}\right) + q\left(\frac{\Delta B}{B}\right) + r\left(\frac{\Delta C}{C}\right) \right] \times 100\%
\end{equation}
\end{formulabox}

\begin{usecasebox}[Experimental Least Count \& Boundary Application Rules]
\begin{itemize}
  \item \textbf{Vernier Calipers Least Count ($\mathrm{LC}$):} Given $1\,\mathrm{MSD}$ (Main Scale Division) and $n$ Vernier divisions coinciding with $(n-1)\,\mathrm{MSD}$:
  \begin{equation}
    \mathrm{LC} = 1\,\mathrm{MSD} - 1\,\mathrm{VSD} = \frac{1\,\mathrm{MSD}}{n}
  \end{equation}
  \item \textbf{Screw Gauge Least Count:} For pitch $p$ and total circular scale divisions $N$:
  \begin{equation}
    \mathrm{LC} = \frac{\text{Pitch } p}{N}, \qquad \text{True Reading} = \text{MSR} + (\text{CSR} \times \mathrm{LC}) - (\pm \text{Zero Error})
  \end{equation}
  \item \textbf{Rule for Significant Figures in Calculus:} For multiplication and division, the final computed value retains the same number of significant figures as the least precise factor. For addition and subtraction, the result retains the least number of decimal places.
\end{itemize}
\end{usecasebox}

\begin{trapbox}[NEET Trap: Sign of Zero Error \& Subtraction of Uncertainty]
\begin{itemize}
  \item \textbf{Zero Error Subtraction:} True Reading = Observed Reading $-$ (Zero Error with appropriate algebraic sign). If a screw gauge has a positive zero error of $+0.04\,\mathrm{mm}$, that quantity must be \textbf{subtracted}. If it has a negative zero error of $-0.03\,\mathrm{mm}$, it is added: $\text{True} = \text{Observed} - (-0.03) = \text{Observed} + 0.03\,\mathrm{mm}$.
  \item \textbf{Never Subtract Errors:} In fractional error propagation for $Z = A/B$ or $Z = A - B$, errors never cancel. Even though variables subtract, their individual experimental uncertainties \textbf{additively compound}: $\Delta Z = \Delta A + \Delta B$.
\end{itemize}
\end{trapbox}

\begin{ncertbox}[NCERT Critical Extraction: Systematic vs Random Errors]
\textbf{NCERT Highlight:} Systematic errors are unidirectional (consistently positive or consistently negative) and arise from instrumental imperfections, zero errors, or experimental methodology flaws. They can be eliminated by calibration. Random errors arise from uncontrollable fluctuations in experimental conditions and follow Gaussian distributions; their influence is minimized by taking the arithmetic mean of $N$ repetitive measurements, reducing the random error magnitude by a factor of $1/\sqrt{N}$.
\end{ncertbox}

% =========================================================================
% SECTION 2: RECTILINEAR KINEMATICS (MOTION IN 1D)
% =========================================================================
\section{Rectilinear Kinematics (Motion in One Dimension)}

\begin{conceptbox}[Calculus Formulations of Rectilinear Motion]
Consider a point particle moving along a coordinate axis $x(t)$.
\begin{enumerate}[leftmargin=*]
  \item \textbf{Position and Displacement:} Position vector $\mathbf{r}(t) = x(t)\hat{\mathbf{i}}$. Displacement over time interval $[t_1, t_2]$ is $\Delta x = x(t_2) - x(t_1)$, which is path-independent. Total distance $s$ is path-dependent: $s = \int_{t_1}^{t_2} |v(t)|\,dt \ge |\Delta x|$.
  \item \textbf{Differential Velocity and Acceleration:}
  \begin{equation}
    v(t) = \frac{dx}{dt}, \qquad a(t) = \frac{dv}{dt} = \frac{d^2 x}{dt^2} = v\frac{dv}{dx}
  \end{equation}
  \item \textbf{Integral Kinematic Relations (Variable Acceleration):}
  \begin{equation}
    \Delta x = \int_{t_1}^{t_2} v(t)\,dt, \qquad \Delta v = \int_{t_1}^{t_2} a(t)\,dt, \qquad \int_{v_1}^{v_2} v\,dv = \int_{x_1}^{x_2} a(x)\,dx
  \end{equation}
\end{enumerate}
\end{conceptbox}

\begin{formulabox}[Rigorous Calculus Derivation of Uniformly Accelerated Kinematic Equations]
For constant acceleration $a = \text{constant}$:
\begin{enumerate}[leftmargin=*]
  \item \textbf{First Kinematic Relation ($v - t$):}
  \begin{equation}
    a = \frac{dv}{dt} \implies \int_{u}^{v} dv = a\int_{0}^{t} dt \implies [v]_u^v = a[t]_0^t \implies v = u + at
  \end{equation}
  \item \textbf{Second Kinematic Relation ($x - t$):}
  \begin{equation}
    v = \frac{dx}{dt} = u + at \implies \int_{0}^{s} dx = \int_{0}^{t} (u + at)\,dt \implies s = ut + \frac{1}{2}at^2
  \end{equation}
  \item \textbf{Third Kinematic Relation ($v - x$):}
  \begin{equation}
    a = v\frac{dv}{dx} \implies \int_{u}^{v} v\,dv = a\int_{0}^{s} dx \implies \left[\frac{v^2}{2}\right]_u^v = as \implies v^2 = u^2 + 2as
  \end{equation}
  \item \textbf{Displacement in the $n^{\text{th}}$ Second ($s_n$):}
  \begin{align}
    s_n &= s(n) - s(n-1) = \left[un + \frac{1}{2}an^2\right] - \left[u(n-1) + \frac{1}{2}a(n-1)^2\right] \nonumber \\
    &= u[n - (n-1)] + \frac{1}{2}a[n^2 - (n^2 - 2n + 1)] = u + \frac{a}{2}(2n - 1)
  \end{align}
\end{enumerate}
\end{formulabox}

\begin{usecasebox}[Motion Under Gravity: Symmetries \& Galileo's Odd Integer Theorem]
For vertical throw under uniform gravitational acceleration $a = -g$ (taking upward as positive):
\begin{itemize}
  \item \textbf{Maximum Height Attained:} Setting $v = 0$ in $v^2 = u^2 - 2gH \implies H = \frac{u^2}{2g}$.
  \item \textbf{Time of Ascent and Descent:} $t_{\text{ascent}} = \frac{u}{g}$, Total flight time $T = \frac{2u}{g}$. The speed of the particle returning to the release plane equals its launch speed: $v_{\text{return}} = -u$.
  \item \textbf{Galileo's Law of Odd Numbers:} For a body released from rest ($u=0$), the distances traversed in successive equal time intervals $\Delta t$ follow the ratio of odd integers:
  \begin{equation}
    s_1 : s_2 : s_3 : \dots : s_n = 1 : 3 : 5 : \dots : (2n - 1)
  \end{equation}
  \item \textbf{Stopping Distance \& Reaction Time:} If driver reacts after time $t_r$ and braking deceleration is $a$:
  \begin{equation}
    d_{\text{stop}} = d_{\text{reaction}} + d_{\text{braking}} = u \cdot t_r + \frac{u^2}{2a}
  \end{equation}
\end{itemize}
\end{usecasebox}

\begin{trapbox}[NEET Trap: Distance vs Displacement in Turning Point Trajectories]
When computing the total distance traveled by an object undergoing deceleration (e.g., $u > 0, a < 0$), the kinematic formula $s = ut + \frac{1}{2}at^2$ yields the \textbf{net displacement}, which can be zero or negative.
\begin{itemize}
  \item \textbf{Verification Step:} First solve $v(t^*) = u + at^* = 0 \implies t^* = -u/a$.
  \item If the specified evaluation time $t > t^*$, the particle has reversed direction. The total path distance is:
  \begin{equation}
    \text{Distance } = |x(t^*) - x(0)| + |x(t) - x(t^*)| = \frac{u^2}{2|a|} + \frac{1}{2}|a|(t - t^*)^2
  \end{equation}
  Using $s = ut + \frac{1}{2}at^2$ directly will cause a catastrophic error in NEET problems!
\end{itemize}
\end{trapbox}

\begin{ncertbox}[NCERT High-Yield Extraction: Geometric Meanings of Motion Graphs]
\textbf{NCERT Highlight:}
\begin{itemize}
  \item In a Position--Time ($x-t$) graph: Slope of the tangent at time $t$ gives instantaneous velocity $v = \frac{dx}{dt}$. An inflection point represents zero acceleration.
  \item In a Velocity--Time ($v-t$) graph: Slope gives instantaneous acceleration $a = \frac{dv}{dt}$. The signed definite integral (area between curve and time axis) represents net displacement: $\Delta x = \int v\,dt$. Total area with positive magnitudes represents total distance: $s = \int |v|\,dt$.
  \item In an Acceleration--Time ($a-t$) graph: Area under the curve yields change in velocity: $\Delta v = \int a\,dt$.
\end{itemize}
\end{ncertbox}

% =========================================================================
% SECTION 3: TWO-DIMENSIONAL KINEMATICS \& PROJECTILE DYNAMICS
% =========================================================================
\section{Two-Dimensional Kinematics, Projectile Motion \& Relative Velocity}

\begin{conceptbox}[Orthogonal Vector Decomposition in Planar Motion]
Planar motion under constant acceleration vector $\mathbf{a} = a_x \hat{\mathbf{i}} + a_y \hat{\mathbf{j}}$ decomposes into two completely independent, uncoupled rectilinear motions along perpendicular Cartesian axes:
\begin{equation}
  x(t) = x_0 + u_x t + \frac{1}{2}a_x t^2, \qquad y(t) = y_0 + u_y t + \frac{1}{2}a_y t^2
\end{equation}
The instantaneous velocity vector is $\mathbf{v}(t) = (u_x + a_x t)\hat{\mathbf{i}} + (u_y + a_y t)\hat{\mathbf{j}}$, with speed $|\mathbf{v}| = \sqrt{v_x^2 + v_y^2}$ and trajectory inclination angle $\tan\phi(t) = \frac{v_y(t)}{v_x(t)}$.
\end{conceptbox}

\begin{formulabox}[Complete Mathematical Derivation of Oblique Projectile Dynamics]
Consider a projectile launched with speed $u$ at angle $\theta$ to the horizontal over level ground ($a_x = 0, a_y = -g$):
\begin{align}
  u_x &= u\cos\theta, \qquad v_x(t) = u\cos\theta, \qquad x(t) = (u\cos\theta)t \\
  u_y &= u\sin\theta, \qquad v_y(t) = u\sin\theta - gt, \qquad y(t) = (u\sin\theta)t - \frac{1}{2}gt^2
\end{align}
\textbf{1. Time of Flight ($T$):} Setting vertical displacement $y(T) = 0$:
\begin{equation}
  T\left(u\sin\theta - \frac{1}{2}gT\right) = 0 \implies T = \frac{2u\sin\theta}{g}
\end{equation}
\textbf{2. Maximum Height Attained ($H$):} At the trajectory apex, $v_y = 0 \implies t_{\text{apex}} = \frac{u\sin\theta}{g}$:
\begin{equation}
  H = y(t_{\text{apex}}) = u\sin\theta\left(\frac{u\sin\theta}{g}\right) - \frac{1}{2}g\left(\frac{u\sin\theta}{g}\right)^2 = \frac{u^2\sin^2\theta}{2g}
\end{equation}
\textbf{3. Horizontal Range ($R$):} Evaluating horizontal coordinate at $t = T$:
\begin{equation}
  R = x(T) = (u\cos\theta)\left(\frac{2u\sin\theta}{g}\right) = \frac{u^2(2\sin\theta\cos\theta)}{g} = \frac{u^2\sin 2\theta}{g}
\end{equation}
\textbf{4. Cartesian Equation of Trajectory:} Eliminating $t = \frac{x}{u\cos\theta}$:
\begin{equation}
  y(x) = u\sin\theta\left(\frac{x}{u\cos\theta}\right) - \frac{g}{2}\left(\frac{x}{u\cos\theta}\right)^2 = x\tan\theta - \frac{gx^2}{2u^2\cos^2\theta} = x\tan\theta\left(1 - \frac{x}{R}\right)
\end{equation}
\end{formulabox}

\begin{usecasebox}[High-Yield Projectile Properties \& Incline Formulas]
\begin{itemize}
  \item \textbf{Complementary Angles of Projection:} For a given initial velocity $u$, the range is identical for launch angles $\theta$ and $(90^\circ - \theta)$:
  \begin{equation}
    R(\theta) = R(90^\circ - \theta) = \frac{u^2\sin 2\theta}{g}
  \end{equation}
  The corresponding maximum heights $H_1, H_2$ satisfy the master identity:
  \begin{equation}
    R = 4\sqrt{H_1 H_2}, \qquad H_1 + H_2 = \frac{u^2}{2g}
  \end{equation}
  \item \textbf{Horizontal Projectile from Cliff of Height $h$:} Launched horizontally with speed $u$:
  \begin{equation}
    t_{\text{flight}} = \sqrt{\frac{2h}{g}}, \qquad R = u\sqrt{\frac{2h}{g}}, \qquad v_{\text{strike}} = \sqrt{u^2 + 2gh}
  \end{equation}
  \item \textbf{Projectile Motion on an Inclined Plane:} For a plane inclined at angle $\alpha$ to horizontal, with launch angle $\beta$ relative to the inclined plane:
  \begin{align}
    \text{Time of flight: } T &= \frac{2u\sin\beta}{g\cos\alpha} \\
    \text{Range up the incline: } R_{\text{up}} &= \frac{u^2}{g\cos^2\alpha}\left[\sin(2\beta + \alpha) - \sin\alpha\right] \\
    \text{Condition for Maximum Range: } \beta &= \frac{\pi}{4} - \frac{\alpha}{2} \implies R_{\max,\text{up}} = \frac{u^2}{g(1 + \sin\alpha)}
  \end{align}
\end{itemize}
\end{usecasebox}

\begin{formulabox}[Relative Velocity Formulations in 2D: River-Boat \& Rain-Man]
Relative velocity of body $A$ with respect to body $B$ is defined by:
\begin{equation}
  \mathbf{v}_{AB} = \mathbf{v}_A - \mathbf{v}_B
\end{equation}
\textbf{1. River-Boat Problem:} Let river width be $d$, river flow speed be $v_r\hat{\mathbf{i}}$, and boat velocity relative to water be $\mathbf{v}_{br}$:
\begin{itemize}
  \item \textbf{Shortest Time Crossing:} Steer directly across ($\theta = 90^\circ$ to river flow):
  \begin{equation}
    t_{\min} = \frac{d}{v_{br}}, \qquad \text{Drift } x = v_r \cdot t_{\min} = \frac{v_r d}{v_{br}}
  \end{equation}
  \item \textbf{Shortest Path Crossing (Zero Drift):} Feasible only when $v_{br} > v_r$. Steer at angle $\theta$ upstream such that horizontal component cancels river flow:
  \begin{equation}
    \sin\theta = \frac{v_r}{v_{br}}, \qquad v_{\text{resultant}} = \sqrt{v_{br}^2 - v_r^2}, \qquad t = \frac{d}{\sqrt{v_{br}^2 - v_r^2}}
  \end{equation}
\end{itemize}
\textbf{2. Rain-Man Problem:} Rain falling vertically with velocity $\mathbf{v}_r = -v_r\hat{\mathbf{j}}$, observer walking with $\mathbf{v}_m = v_m\hat{\mathbf{i}}$:
\begin{equation}
  \mathbf{v}_{r,m} = \mathbf{v}_r - \mathbf{v}_m = -v_m\hat{\mathbf{i}} - v_r\hat{\mathbf{j}}, \qquad \tan\theta = \frac{|\mathbf{v}_m|}{|\mathbf{v}_r|} = \frac{v_m}{v_r}
\end{equation}
The umbrella must be oriented tilted forward at angle $\theta = \tan^{-1}(v_m/v_r)$ relative to vertical.
\end{formulabox}

\begin{trapbox}[NEET Trap: Minimum Separation in 2D Interception Problems]
When two bodies $A$ and $B$ move with constant velocities $\mathbf{v}_A$ and $\mathbf{v}_B$ from initial positions $\mathbf{r}_A$ and $\mathbf{r}_B$:
\begin{itemize}
  \item \textbf{Wrong Approach:} Differentiating scalar distance formula algebraically leads to enormous computation time.
  \item \textbf{Correct Vector Method:} Freeze body $B$ at the origin of a relative frame. Body $A$ moves along a straight line with constant relative velocity $\mathbf{v}_{AB} = \mathbf{v}_A - \mathbf{v}_B$ starting from relative position $\mathbf{r}_{AB} = \mathbf{r}_A - \mathbf{r}_B$.
  \item The shortest distance is simply the length of the perpendicular from $B$ to the relative velocity line:
  \begin{equation}
    d_{\min} = \frac{|\mathbf{r}_{AB} \times \mathbf{v}_{AB}|}{|\mathbf{v}_{AB}|}, \qquad t_{\min} = -\frac{\mathbf{r}_{AB} \cdot \mathbf{v}_{AB}}{|\mathbf{v}_{AB}|^2}
  \end{equation}
\end{itemize}
\end{trapbox}

% =========================================================================
% SECTION 4: NEWTON'S LAWS OF MOTION \& CONSTRAINED SYSTEM DYNAMICS
% =========================================================================
\section{Newton's Laws of Motion, Constrained Dynamics \& Pseudo-Forces}

\begin{conceptbox}[Foundations of Newtonian Mechanics \& Frame Relativity]
\begin{enumerate}[leftmargin=*]
  \item \textbf{Newton's First Law (Law of Inertia):} In the absence of an unbalanced net external force, a body maintains its state of rest or uniform rectilinear motion. This establishes the existence of \textbf{Inertial Reference Frames}.
  \item \textbf{Newton's Second Law:} The rate of change of linear momentum $\mathbf{p} = m\mathbf{v}$ is proportional to and in the direction of the applied net force:
  \begin{equation}
    \mathbf{F}_{\text{net}} = \frac{d\mathbf{p}}{dt} = m\frac{d\mathbf{v}}{dt} + \mathbf{v}\frac{dm}{dt} = m\mathbf{a} \quad (\text{for constant mass system})
  \end{equation}
  \item \textbf{Newton's Third Law:} Forces occur strictly in pairs. When body $A$ exerts force $\mathbf{F}_{BA}$ on body $B$, body $B$ exerts an instantaneous, collinear, equal and opposite reaction force on $A$: $\mathbf{F}_{AB} = -\mathbf{F}_{BA}$. Action and reaction act on \textbf{different bodies} and never cancel each other.
  \item \textbf{Non-Inertial Reference Frames \& Pseudo-Force:} In a frame accelerating with acceleration $\mathbf{a}_f$, Newton's second law is restored by introducing the inertial pseudo-force:
  \begin{equation}
    \mathbf{F}_{\text{pseudo}} = -m\mathbf{a}_f
  \end{equation}
\end{enumerate}
\end{conceptbox}

\begin{formulabox}[Equilibrium, Lami's Theorem \& String/Wedge Constraints]
\textbf{1. Concurrent Force Equilibrium (Lami's Theorem):} If three coplanar forces $\mathbf{F}_1, \mathbf{F}_2, \mathbf{F}_3$ acting on a particle maintain equilibrium ($\mathbf{F}_1 + \mathbf{F}_2 + \mathbf{F}_3 = \mathbf{0}$):
\begin{equation}
  \frac{F_1}{\sin\alpha} = \frac{F_2}{\sin\beta} = \frac{F_3}{\sin\gamma}
\end{equation}
where $\alpha, \beta, \gamma$ are the angles opposite to forces $\mathbf{F}_1, \mathbf{F}_2, \mathbf{F}_3$ respectively.

\textbf{2. Virtual Work Principle for String Constraints:} For an inextensible, taut string passing over ideal frictionless massless pulleys, the total internal work done by string tension across the system is identically zero:
\begin{equation}
  \sum_{i} \mathbf{T}_i \cdot \mathbf{x}_i = 0 \implies \sum_{i} \mathbf{T}_i \cdot \mathbf{v}_i = 0 \implies \sum_{i} \mathbf{T}_i \cdot \mathbf{a}_i = 0
\end{equation}
\textbf{3. Wedge Constraints:} Two bodies maintaining continuous contact without penetration must have identical velocity and acceleration components along the \textbf{common normal} to their contact interface:
\begin{equation}
  (\mathbf{v}_1 \cdot \hat{\mathbf{n}}) = (\mathbf{v}_2 \cdot \hat{\mathbf{n}}), \qquad (\mathbf{a}_1 \cdot \hat{\mathbf{n}}) = (\mathbf{a}_2 \cdot \hat{\mathbf{n}})
\end{equation}
\end{formulabox}

\begin{usecasebox}[Atwood Machine Derivations \& High-Yield Accelerations]
\begin{itemize}
  \item \textbf{Standard Simple Atwood Machine:} Masses $m_1 > m_2$ connected over an ideal pulley:
  \begin{equation}
    a = \left(\frac{m_1 - m_2}{m_1 + m_2}\right)g, \qquad T = \left(\frac{2m_1 m_2}{m_1 + m_2}\right)g, \qquad \text{Force on Pulley Support } F_p = 2T = \frac{4m_1 m_2 g}{m_1 + m_2}
  \end{equation}
  \item \textbf{Movable Pulley System (Mechanical Advantage):} Mass $m_1$ on smooth horizontal table connected to movable pulley holding hanging mass $m_2$. Since $x_2 = x_1/2$, constraint yields $a_1 = 2a_2$:
  \begin{equation}
    a_2 = \frac{m_2 g}{4m_1 + m_2}, \qquad a_1 = \frac{2m_2 g}{4m_1 + m_2}, \qquad T = \frac{2m_1 m_2 g}{4m_1 + m_2}
  \end{equation}
  \item \textbf{Apparent Weight in an Accelerating Elevator:}
  \begin{align}
    \text{Accelerating upward with } a &: N = m(g + a) \\
    \text{Accelerating downward with } a &: N = m(g - a) \\
    \text{Free fall ($a = g$)} &: N = 0 \quad (\text{Weightlessness})
  \end{align}
\end{itemize}
\end{usecasebox}

\begin{trapbox}[NEET Trap: Spring Balance Reading During Transients \& Cut Strings]
\begin{itemize}
  \item \textbf{Spring vs String Cut Difference:} A light spring exerts force $F = kx$. The displacement $x$ cannot change discontinuously. Therefore, immediately after cutting an adjacent string/support at $t = 0^+$, the spring force \textbf{remains completely unchanged} ($F_{\text{spring}}(0^+) = F_{\text{spring}}(0^-)$).
  \item In contrast, string tension is an instantaneous constraint force that drops to zero or readjusts instantaneously ($T(0^+) \neq T(0^-)$).
\end{itemize}
\end{trapbox}

% =========================================================================
% SECTION 5: FRICTION DYNAMICS
% =========================================================================
\section{Friction: Microscopic Mechanism, Slopes \& Multi-Block Dynamics}

\begin{conceptbox}[Laws of Dry Friction (Amontons--Coulomb Model)]
Friction is an electromagnetic contact force arising from microscopic adhesive bonding (cold welding of microscopic asperities) and mechanical interlocking between contacting surfaces.
\begin{enumerate}[leftmargin=*]
  \item \textbf{Static Friction ($f_s$):} A self-adjusting reaction force that opposes the tendency of relative motion. It matches the net applied tangential force up to a limiting threshold:
  \begin{equation}
    0 \le f_s \le f_{s,\max} = \mu_s N
  \end{equation}
  where $\mu_s$ is the coefficient of static friction, and $N$ is the normal contact force.
  \item \textbf{Kinetic Friction ($f_k$):} Once relative slip begins, the opposing frictional force becomes nearly independent of relative sliding velocity and contact area:
  \begin{equation}
    f_k = \mu_k N \quad (\text{with } \mu_k < \mu_s)
  \end{equation}
  \item \textbf{Angle of Friction ($\lambda$) \& Angle of Repose ($\theta_r$):}
  \begin{equation}
    \tan\lambda = \frac{f_{s,\max}}{N} = \frac{\mu_s N}{N} = \mu_s, \qquad \tan\theta_r = \mu_s \implies \theta_r = \lambda = \tan^{-1}\mu_s
  \end{equation}
  A body resting on an inclined plane begins spontaneous sliding only when inclination $\theta > \theta_r$.
\end{enumerate}
\end{conceptbox}

\begin{formulabox}[Dynamics on Rough Inclined Planes \& Stacked Block Systems]
\textbf{1. Acceleration on Rough Incline:}
\begin{align}
  \text{Sliding down incline ($\theta > \theta_r$): } a_{\text{down}} &= g(\sin\theta - \mu_k\cos\theta) \\
  \text{Projected upward along incline: } a_{\text{up}} &= -g(\sin\theta + \mu_k\cos\theta)
\end{align}
Ratio of time taken to slide down rough incline vs smooth incline of identical length:
\begin{equation}
  t_{\text{rough}} = n \cdot t_{\text{smooth}} \implies \mu_k = \tan\theta\left(1 - \frac{1}{n^2}\right)
\end{equation}
\textbf{2. Minimum Force to Drag a Block on a Rough Horizontal Surface:}
Pulling at an upward angle $\theta$ to the horizontal reduces normal force $N = mg - F\sin\theta$:
\begin{equation}
  F\cos\theta = \mu_s(mg - F\sin\theta) \implies F(\theta) = \frac{\mu_s mg}{\cos\theta + \mu_s\sin\theta}
\end{equation}
Minimizing $F(\theta)$ by setting $\frac{d}{d\theta}(\cos\theta + \mu_s\sin\theta) = 0 \implies \tan\theta = \mu_s$:
\begin{equation}
  F_{\min} = \frac{\mu_s mg}{\sqrt{1 + \mu_s^2}} \quad \text{at angle } \theta = \tan^{-1}\mu_s
\end{equation}
\end{formulabox}

\begin{usecasebox}[Two-Block Stacked Friction Problem Archetype]
Consider block $A$ of mass $m_1$ resting atop block $B$ of mass $m_2$, situated on a frictionless horizontal floor. Friction coefficient between $A$ and $B$ is $\mu_s = \mu_k = \mu$.
\begin{itemize}
  \item \textbf{Case I: External Force $F$ Applied to Lower Block $B$ ($m_2$):}
  The only force driving block $A$ forward is static friction $f_s \le \mu m_1 g$.
  Maximum allowable acceleration without relative slip:
  \begin{equation}
    a_{\max} = \frac{f_{s,\max}}{m_1} = \frac{\mu m_1 g}{m_1} = \mu g
  \end{equation}
  Threshold force for no slipping:
  \begin{equation}
    F_{\text{threshold}} = (m_1 + m_2)a_{\max} = \mu(m_1 + m_2)g
  \end{equation}
  If $F \le F_{\text{threshold}}$, both blocks accelerate together at $a = \frac{F}{m_1 + m_2}$.
  If $F > F_{\text{threshold}}$, slipping occurs:
  \begin{equation}
    a_1 = \mu g, \qquad a_2 = \frac{F - \mu m_1 g}{m_2}
  \end{equation}
  \item \textbf{Case II: External Force $F$ Applied to Upper Block $A$ ($m_1$):}
  Here friction on $B$ drives $B$ forward: $f_s \le \mu m_1 g \implies a_{\max} = \frac{\mu m_1 g}{m_2}$.
  \begin{equation}
    F_{\text{threshold}} = (m_1 + m_2)a_{\max} = \mu m_1 g\left(1 + \frac{m_1}{m_2}\right)
  \end{equation}
\end{itemize}
\end{usecasebox}

\begin{trapbox}[NEET Trap: Direction of Static Friction on Walking \& Driving Wheels]
\begin{itemize}
  \item \textbf{Human Walking:} The foot pushes the ground backward. By Newton's third law, static friction exerted by the ground on the foot acts \textbf{forward}, providing propulsive acceleration!
  \item \textbf{Automobile / Bicycle Wheels:} For the rear (driving) wheel of a bicycle, engine/pedal torque pushes the ground backward, so ground friction acts \textbf{forward}. For the front (driven) wheel, ground friction acts \textbf{backward} to generate forward angular rotation.
\end{itemize}
\end{trapbox}

% =========================================================================
% SECTION 6: DYNAMICS OF CIRCULAR MOTION
% =========================================================================
\section{Dynamics of Circular Motion, Banking \& Conical Pendulum}

\begin{conceptbox}[Kinematics \& Force Mechanics in Curvilinear Trajectories]
In circular motion of radius $r$, the acceleration vector decomposes along radial (centripetal) and tangential directions:
\begin{equation}
  \mathbf{a} = \mathbf{a}_c + \mathbf{a}_t = -\left(\frac{v^2}{r}\right)\hat{\mathbf{r}} + \left(\frac{dv}{dt}\right)\hat{\boldsymbol{\theta}} = -(\omega^2 r)\hat{\mathbf{r}} + (\alpha r)\hat{\boldsymbol{\theta}}
\end{equation}
\begin{enumerate}[leftmargin=*]
  \item \textbf{Centripetal Force:} A net real force directed toward the instantaneous center of curvature required to continuously bend the velocity vector:
  \begin{equation}
    F_c = m a_c = \frac{m v^2}{r} = m \omega^2 r
  \end{equation}
  \item \textbf{Centrifugal Force:} In the rotating reference frame attached to the orbiting particle, a radially outward pseudo-force of magnitude $F_{\text{cf}} = m\omega^2 r$ is applied.
\end{enumerate}
\end{conceptbox}

\begin{formulabox}[Complete Derivation of Banking Dynamics on Curved Tracks]
Consider a vehicle negotiating a curve of radius $r$ on a road banked at angle $\theta$ with static friction coefficient $\mu_s$:
\begin{enumerate}[leftmargin=*]
  \item \textbf{Unbanked Level Curved Road:} Friction alone provides centripetal acceleration:
  \begin{equation}
    f_s \le \mu_s N = \mu_s mg \implies \frac{mv^2}{r} \le \mu_s mg \implies v_{\max} = \sqrt{\mu_s g r}
  \end{equation}
  \item \textbf{Frictionless Banked Road (Optimum Speed $v_0$):} Normal force component $N\sin\theta$ provides centripetal acceleration, while $N\cos\theta$ balances gravity:
  \begin{equation}
    N\sin\theta = \frac{mv_0^2}{r}, \quad N\cos\theta = mg \implies \tan\theta = \frac{v_0^2}{rg} \implies v_0 = \sqrt{rg\tan\theta}
  \end{equation}
  \item \textbf{Rough Banked Road (Maximum Safe Velocity $v_{\max}$):} Vehicle tends to skid outward up the incline; limiting friction $f_s = \mu_s N$ acts downward along the incline:
  \begin{align}
    N\cos\theta - f_s\sin\theta = mg &\implies N(\cos\theta - \mu_s\sin\theta) = mg \\
    N\sin\theta + f_s\cos\theta = \frac{mv^2}{r} &\implies N(\sin\theta + \mu_s\cos\theta) = \frac{mv^2}{r}
  \end{align}
  Dividing the two equations yields the exact NEET master formula:
  \begin{equation}
    v_{\max} = \sqrt{rg\left(\frac{\tan\theta + \mu_s}{1 - \mu_s\tan\theta}\right)}
  \end{equation}
  \item \textbf{Minimum Safe Velocity (Preventing Inward Slide $v_{\min}$):} Limiting friction acts upward along the incline:
  \begin{equation}
    v_{\min} = \sqrt{rg\left(\frac{\tan\theta - \mu_s}{1 + \mu_s\tan\theta}\right)} \quad (\text{if } \tan\theta > \mu_s)
  \end{equation}
\end{enumerate}
\end{formulabox}

\begin{usecasebox}[Conical Pendulum Dynamics \& Period Derivation]
A bob of mass $m$ attached to a string of length $L$ traces a horizontal circle of radius $r = L\sin\theta$ at constant angular velocity $\omega$, making semi-vertical angle $\theta$:
\begin{align}
  \text{Vertical force balance: } T\cos\theta &= mg \implies T = \frac{mg}{\cos\theta} \\
  \text{Radial centripetal equation: } T\sin\theta &= m\omega^2 r = m\omega^2(L\sin\theta) \implies T = m\omega^2 L
\end{align}
Equating expressions for string tension $T$:
\begin{equation}
  \frac{mg}{\cos\theta} = m\omega^2 L \implies \omega = \sqrt{\frac{g}{L\cos\theta}}
\end{equation}
The time period of the conical pendulum is:
\begin{equation}
  T_{\text{period}} = \frac{2\pi}{\omega} = 2\pi\sqrt{\frac{L\cos\theta}{g}} = 2\pi\sqrt{\frac{h}{g}}
\end{equation}
where $h = L\cos\theta$ is the vertical height of the apex above the circular plane of motion.
\end{usecasebox}

\begin{trapbox}[NEET Trap: Death Well / Rotor Minimum Angular Velocity]
In a cylindrical rotor (death well) of radius $R$, a person rests against the vertical wall without a floor.
\begin{itemize}
  \item The normal force $N = m\omega^2 R$ acts \textbf{horizontally} providing centripetal acceleration.
  \item Static friction acts \textbf{vertically upward} to balance weight: $f_s = mg$.
  \item Condition to prevent sliding down:
  \begin{equation}
    f_s \le \mu_s N \implies mg \le \mu_s(m\omega^2 R) \implies \omega_{\min} = \sqrt{\frac{g}{\mu_s R}}, \quad v_{\min} = \sqrt{\frac{gR}{\mu_s}}
  \end{equation}
  Notice that $\mu_s$ appears in the denominator, in stark contrast to level road turns!
\end{itemize}
\end{trapbox}

% =========================================================================
% SECTION 7: WORK, ENERGY, POWER \& VERTICAL CIRCULAR MOTION
% =========================================================================
\section{Work, Energy, Power \& Vertical Circular Dynamics}

\begin{conceptbox}[Calculus Formulations of Work \& Potential Energy Functions]
\begin{enumerate}[leftmargin=*]
  \item \textbf{Work Done by Variable Force:} For a particle moving along path $C$ from $\mathbf{r}_i$ to $\mathbf{r}_f$:
  \begin{equation}
    W = \int_{\mathbf{r}_i}^{\mathbf{r}_f} \mathbf{F} \cdot d\mathbf{r} = \int_{x_i}^{x_f} F_x\,dx + \int_{y_i}^{y_f} F_y\,dy + \int_{z_i}^{z_f} F_z\,dz
  \end{equation}
  \item \textbf{Conservative Force Criterion:} A force is conservative if the work done around any closed loop vanishes identically: $\oint \mathbf{F}\cdot d\mathbf{r} = 0 \iff \boldsymbol{\nabla}\times\mathbf{F} = \mathbf{0}$.
  \item \textbf{Potential Energy Field ($U$):} For conservative forces, potential energy is defined such that:
  \begin{equation}
    \Delta U = U_f - U_i = -W_{\text{int}} = -\int_{\mathbf{r}_i}^{\mathbf{r}_f} \mathbf{F}\cdot d\mathbf{r} \implies \mathbf{F} = -\boldsymbol{\nabla}U = -\left(\frac{\partial U}{\partial x}\hat{\mathbf{i}} + \frac{\partial U}{\partial y}\hat{\mathbf{j}} + \frac{\partial U}{\partial z}\hat{\mathbf{k}}\right)
  \end{equation}
  \item \textbf{Equilibrium States \& Stability Calculus:}
  \begin{align}
    \text{Equilibrium Condition: } &\frac{dU}{dx} = 0 \implies F = 0 \\
    \text{Stable Equilibrium: } &\frac{d^2 U}{dx^2} > 0 \quad (\text{Local Potential Minimum, restoring force}) \\
    \text{Unstable Equilibrium: } &\frac{d^2 U}{dx^2} < 0 \quad (\text{Local Potential Maximum, destabilizing}) \\
    \text{Neutral Equilibrium: } &\frac{d^2 U}{dx^2} = 0 \quad (\text{Flat potential valley})
  \end{align}
\end{enumerate}
\end{conceptbox}

\begin{formulabox}[Work-Energy Theorem Rigorous Proof \& Spring Potential Energy]
\textbf{1. Work-Energy Theorem for Variable Force in 1D:}
\begin{equation}
  W = \int_{x_i}^{x_f} F\,dx = \int_{x_i}^{x_f} m\left(\frac{dv}{dt}\right)dx = \int_{x_i}^{x_f} m v\left(\frac{dv}{dx}\right)dx = \int_{v_i}^{v_f} m v\,dv
\end{equation}
Integrating yields the fundamental Work-Energy Theorem:
\begin{equation}
  W_{\text{net}} = \left[\frac{1}{2}m v^2\right]_{v_i}^{v_f} = \frac{1}{2}m v_f^2 - \frac{1}{2}m v_i^2 = \Delta K
\end{equation}
\textbf{2. Elastic Potential Energy of Ideal Spring:}
According to Hooke's Law, restoring spring force is $F_s = -kx$. Work done by external agent stretching from $0$ to $x$:
\begin{equation}
  W_{\text{ext}} = \int_{0}^{x} kx\,dx = \frac{1}{2}kx^2 \implies U_s(x) = \frac{1}{2}kx^2
\end{equation}
\end{formulabox}

\begin{formulabox}[Complete Dynamics of Vertical Circular Motion (String vs Rod)]
Let a bob of mass $m$ on a string of length $l$ be projected horizontally at the lowest point (Point $L$) with velocity $u$.
At angle $\theta$ relative to the downward vertical:
\begin{align}
  \text{Conservation of Mechanical Energy: } &\frac{1}{2}mu^2 = \frac{1}{2}mv^2(\theta) + mgl(1 - \cos\theta) \nonumber \\
  &\implies v^2(\theta) = u^2 - 2gl(1 - \cos\theta) \\
  \text{Radial Equation for String Tension: } &T(\theta) - mg\cos\theta = \frac{mv^2(\theta)}{l} \nonumber \\
  &\implies T(\theta) = \frac{mu^2}{l} - mg(2 - 3\cos\theta)
\end{align}
\textbf{Key Critical Boundary Values for Complete Vertical Loop (String):}
\begin{enumerate}[leftmargin=*]
  \item \textbf{At Highest Point ($\theta = \pi$):} $\cos\pi = -1 \implies v_{\text{top}}^2 = u^2 - 4gl$.
  To maintain tautness, $T_{\text{top}} \ge 0$:
  \begin{equation}
    T_{\text{top}} = \frac{m v_{\text{top}}^2}{l} - mg \ge 0 \implies v_{\text{top}} \ge \sqrt{gl}
  \end{equation}
  \item \textbf{Critical Minimum Launch Speed at Bottom ($u_{\min}$):}
  \begin{equation}
    u^2 - 4gl \ge gl \implies u_{\min} = \sqrt{5gl}
  \end{equation}
  \item \textbf{Speed at Horizontal Position ($\theta = \pi/2$):}
  \begin{equation}
    v_{\text{horiz}} = \sqrt{u^2 - 2gl} \ge \sqrt{3gl}
  \end{equation}
  \item \textbf{Tension Difference Between Bottom and Top:}
  \begin{equation}
    T_{\text{bottom}} = \frac{mu^2}{l} + mg, \qquad T_{\text{top}} = \frac{mu^2}{l} - 5mg \implies T_{\text{bottom}} - T_{\text{top}} = 6mg
  \end{equation}
  \item \textbf{Looping in a Light Rigid Rod:} A rod can support compressive stress. Looping requires only $v_{\text{top}} \ge 0$:
  \begin{equation}
    v_{\text{top}} \ge 0 \implies u_{\min,\text{rod}} = \sqrt{4gl} = 2\sqrt{gl}
  \end{equation}
\end{enumerate}
\end{formulabox}

\begin{usecasebox}[Regimes of Motion in Vertical Circle Based on Launch Velocity $u$]
\begin{itemize}
  \item \textbf{Regime I ($0 < u \le \sqrt{2gl}$): Oscillation Regime.} Velocity vanishes ($v=0$) before or at horizontal ($\theta \le 90^\circ$). Tension remains strictly positive ($T > 0$). The particle oscillates as a simple pendulum.
  \item \textbf{Regime II ($\sqrt{2gl} < u < \sqrt{5gl}$): String Slackening Regime.} Tension vanishes ($T=0$) before velocity vanishes, at an angle $\theta$ in the upper hemisphere ($90^\circ < \theta < 180^\circ$):
  \begin{equation}
    \cos\theta = \frac{2gl - u^2}{3gl} \quad (\text{negative value})
  \end{equation}
  The particle leaves the circular path and executes parabolic projectile motion under gravity!
  \item \textbf{Regime III ($u \ge \sqrt{5gl}$): Complete Looping Regime.} String never slacks ($T \ge 0$ throughout).
\end{itemize}
\end{usecasebox}

\begin{trapbox}[NEET Trap: Instantaneous Power vs Average Power]
Power is the rate of doing work:
\begin{equation}
  P_{\text{avg}} = \frac{\Delta W}{\Delta t}, \qquad P_{\text{inst}} = \frac{dW}{dt} = \mathbf{F} \cdot \mathbf{v} = F v \cos\phi
\end{equation}
\begin{itemize}
  \item In uniform circular motion, $\mathbf{F} \perp \mathbf{v}$ at every instant ($\cos 90^\circ = 0$). Thus, instantaneous power delivered by centripetal force is \textbf{identically zero} ($P_{\text{inst}} = 0$), and net work done over any interval is zero.
  \item For an engine delivering constant power $P$, velocity scales as $v \propto t^{1/2}$ and displacement scales as $x \propto t^{3/2}$.
\end{itemize}
\end{trapbox}

% =========================================================================
% SECTION 8: MOMENTUM CONSERVATION, COLLISIONS \& VARIABLE MASS SYSTEMS
% =========================================================================
\section{Conservation of Linear Momentum, Collisions \& Variable Mass Systems}

\begin{conceptbox}[Conservation of Linear Momentum \& Coefficient of Restitution]
\begin{enumerate}[leftmargin=*]
  \item \textbf{Principle of Conservation of Linear Momentum:} For an isolated system of particles with zero net external force ($\sum \mathbf{F}_{\text{ext}} = \mathbf{0}$):
  \begin{equation}
    \frac{d\mathbf{P}_{\text{total}}}{dt} = \mathbf{0} \implies \mathbf{P}_{\text{total}} = \sum m_i \mathbf{v}_i = \text{constant}
  \end{equation}
  \item \textbf{Impulse-Momentum Theorem:} The time integral of net force equals the total change in linear momentum:
  \begin{equation}
    \mathbf{J} = \int_{t_1}^{t_2} \mathbf{F}\,dt = \Delta\mathbf{p} = \mathbf{p}_2 - \mathbf{p}_1
  \end{equation}
  \item \textbf{Newton's Law of Restitution:} The coefficient of restitution $e$ quantifies the elasticity along the line of impact (common normal):
  \begin{equation}
    e = \frac{\text{Velocity of Separation}}{\text{Velocity of Approach}} = \frac{v_2 - v_1}{u_1 - u_2}
  \end{equation}
  \begin{itemize}
    \item $e = 1$: Perfectly elastic collision (kinetic energy conserved).
    \item $0 < e < 1$: Inelastic collision (partial loss of kinetic energy to heat/deformation).
    \item $e = 0$: Perfectly inelastic collision (bodies stick together, maximum kinetic energy loss).
  \end{itemize}
\end{enumerate}
\end{conceptbox}

\begin{formulabox}[Complete Derivation of Post-Collision Velocities \& Kinetic Energy Loss]
Consider two masses $m_1, m_2$ moving with initial velocities $u_1, u_2$ undergoing a 1D head-on collision with coefficient of restitution $e$:
\begin{align}
  \text{Conservation of Momentum: } &m_1 u_1 + m_2 u_2 = m_1 v_1 + m_2 v_2 \\
  \text{Restitution Equation: } &v_2 - v_1 = e(u_1 - u_2) \implies v_2 = v_1 + e(u_1 - u_2)
\end{align}
Substituting $v_2$ into the momentum equation:
\begin{align}
  m_1 u_1 + m_2 u_2 &= m_1 v_1 + m_2[v_1 + e(u_1 - u_2)] = (m_1 + m_2)v_1 + e m_2(u_1 - u_2) \nonumber \\
  \implies (m_1 + m_2)v_1 &= m_1 u_1 + m_2 u_2 - e m_2 u_1 + e m_2 u_2 = (m_1 - e m_2)u_1 + (1+e)m_2 u_2
\end{align}
This yields the universal post-collision velocity formulas:
\begin{align}
  v_1 &= \left(\frac{m_1 - e m_2}{m_1 + m_2}\right)u_1 + \left(\frac{(1+e)m_2}{m_1 + m_2}\right)u_2 \\
  v_2 &= \left(\frac{(1+e)m_1}{m_1 + m_2}\right)u_1 + \left(\frac{m_2 - e m_1}{m_1 + m_2}\right)u_2
\end{align}
\textbf{Kinetic Energy Loss in Inelastic Collisions ($\Delta K_{\text{loss}}$):}
\begin{equation}
  \Delta K_{\text{loss}} = K_i - K_f = \frac{1}{2}\left(\frac{m_1 m_2}{m_1 + m_2}\right)(u_1 - u_2)^2(1 - e^2) = \frac{1}{2}\mu (u_{\text{rel}})^2(1 - e^2)
\end{equation}
where $\mu = \frac{m_1 m_2}{m_1 + m_2}$ is the reduced mass of the two-body colliding system.
\end{formulabox}

\begin{usecasebox}[High-Yield Collision Archetypes \& Variable Mass Systems]
\begin{itemize}
  \item \textbf{Equal Masses Elastic Collision ($m_1 = m_2, e = 1$):} Velocities swap completely:
  \begin{equation}
    v_1 = u_2, \qquad v_2 = u_1
  \end{equation}
  \item \textbf{Massive Target at Rest ($m_2 \gg m_1, u_2 = 0, e = 1$):}
  \begin{equation}
    v_1 = -u_1, \qquad v_2 \approx 0 \quad (\text{Projectile rebounds with inverted speed})
  \end{equation}
  \item \textbf{Massive Projectile on Light Target ($m_1 \gg m_2, u_2 = 0, e = 1$):}
  \begin{equation}
    v_1 \approx u_1, \qquad v_2 \approx 2u_1 \quad (\text{Target kicked at double projectile velocity})
  \end{equation}
  \item \textbf{Successive Bounces of Ball Dropped from Height $h_0$:}
  \begin{align}
    \text{Height after } n^{\text{th}}\text{ bounce: } h_n &= e^{2n}h_0 \\
    \text{Total distance traveled until rest: } H_{\text{total}} &= h_0 + 2\sum_{n=1}^\infty h_n = h_0\left(\frac{1+e^2}{1-e^2}\right) \\
    \text{Total time elapsed until rest: } T_{\text{total}} &= \sqrt{\frac{2h_0}{g}} + 2\sum_{n=1}^\infty \sqrt{\frac{2h_n}{g}} = \sqrt{\frac{2h_0}{g}}\left(\frac{1+e}{1-e}\right)
  \end{align}
  \item \textbf{Rocket Propulsion (Tsiolkovsky Equation):} For variable mass system with exhaust velocity $u_{\text{rel}}$:
  \begin{equation}
    v(t) = v_0 + u_{\text{rel}}\ln\left(\frac{m_0}{m(t)}\right) - gt
  \end{equation}
\end{itemize}
\end{usecasebox}

\begin{trapbox}[NEET Trap: Oblique Collisions \& Line of Impact Conservation]
In two-dimensional (oblique) collisions between smooth spheres:
\begin{itemize}
  \item Forces during collision act \textbf{solely along the line of impact} (the common normal connecting the sphere centers at the instant of contact).
  \item Consequently, momentum is conserved along the line of impact, and Newton's restitution formula $e = -\frac{v_{2n} - v_{1n}}{u_{2n} - u_{1n}}$ applies \textbf{only} to normal velocity components.
  \item Tangential velocity components parallel to the contact plane experience zero impulse and remain \textbf{strictly unchanged}:
  \begin{equation}
    v_{1t} = u_{1t}, \qquad v_{2t} = u_{2t}
  \end{equation}
  Applying restitution to total velocity vectors without normal projection is an immediate trap!
\end{itemize}
\end{trapbox}

\begin{ncertbox}[NCERT Critical Extraction: Collision Classifications \& Nuclear Moderation]
\textbf{NCERT Highlight:}
\begin{itemize}
  \item In a nuclear reactor, fast neutrons ($m \approx 1\,\mathrm{u}$) produced by uranium fission must be slowed down to thermal speeds to sustain a chain reaction. Moderators containing light nuclei such as protons or carbon ($m_2 \approx m_1$) are used.
  \item In an elastic head-on collision with a target of comparable mass ($m_1 \approx m_2$) at rest, the incoming neutron transfers virtually all of its kinetic energy ($100\%$) because $v_1 \approx 0$ and $v_2 \approx u_1$. Heavy stationary nuclei ($m_2 \gg m_1$) merely rebound the neutron elastically without absorbing its kinetic energy.
\end{itemize}
\end{ncertbox}

\chapter{System of Particles, Rotational Dynamics, Gravitation \& Fluid Mechanics}

In this chapter, we bridge translational physics with rotational mechanics, celestial gravitation, and the mechanics of deformable media. Topics include the calculus of continuous mass distributions, moment of inertia theorems, angular momentum conservation, pure rolling on rough inclines, Kepler's laws, gravitational field/potential energy formulations, elastic moduli, surface tension, capillarity, and Bernoulli's hydrodynamic theorem.

\section{Centre of Mass, Continuous Mass Distributions \& System Dynamics}

\begin{conceptbox}[Theoretical Postulates of Centre of Mass \& Invariant Trajectories]
The Centre of Mass ($\mathrm{CM}$) of a system of particles represents a unique point where the entire mass of the system can be considered concentrated for translational dynamic analysis.
\begin{enumerate}[leftmargin=*]
  \item \textbf{Discrete Particle Formulations:} For a collection of $N$ point masses $m_1, m_2, \dots, m_N$ with position vectors $\mathbf{r}_1, \mathbf{r}_2, \dots, \mathbf{r}_N$:
  \begin{equation}
    \mathbf{r}_{\text{cm}} = \frac{\sum_{i=1}^N m_i \mathbf{r}_i}{\sum_{i=1}^N m_i} = \frac{1}{M}\sum_{i=1}^N m_i \mathbf{r}_i
  \end{equation}
  \item \textbf{Continuous Mass Distributions (Calculus Formulation):}
  \begin{equation}
    \mathbf{r}_{\text{cm}} = \frac{1}{M}\int \mathbf{r}\,dm, \qquad x_{\text{cm}} = \frac{1}{M}\int x\,dm, \qquad y_{\text{cm}} = \frac{1}{M}\int y\,dm
  \end{equation}
  \item \textbf{System Equations of Motion \& Conservation of Linear Momentum:}
  \begin{equation}
    \mathbf{P}_{\text{total}} = M\mathbf{v}_{\text{cm}}, \qquad \mathbf{F}_{\text{ext}} = \frac{d\mathbf{P}_{\text{total}}}{dt} = M\mathbf{a}_{\text{cm}}
  \end{equation}
  Since internal mutual forces between system particles occur in collinear, equal and opposite action-reaction pairs by Newton's third law, they cancel identically in vector summation: $\sum \mathbf{F}_{\text{int}} = \mathbf{0}$. If $\mathbf{F}_{\text{ext}} = \mathbf{0}$, the centre of mass maintains constant velocity ($\mathbf{v}_{\text{cm}} = \text{constant}$), irrespective of internal explosions or collisions.
\end{enumerate}
\end{conceptbox}

\begin{formulabox}[Master Coordinates of Centre of Mass for Continuous Rigid Bodies]
Coordinates are defined from the geometric center or vertex as reference origin:
\begin{align}
  \text{Uniform Thin Rod of Length } L &: x_{\text{cm}} = \frac{L}{2} \\
  \text{Uniform Semicircular Ring of Radius } R &: y_{\text{cm}} = \frac{2R}{\pi} \approx 0.637 R \\
  \text{Uniform Semicircular Disc of Radius } R &: y_{\text{cm}} = \frac{4R}{3\pi} \approx 0.424 R \\
  \text{Uniform Hemispherical Shell of Radius } R &: y_{\text{cm}} = \frac{R}{2} = 0.500 R \\
  \text{Uniform Solid Hemisphere of Radius } R &: y_{\text{cm}} = \frac{3R}{8} = 0.375 R \\
  \text{Uniform Solid Cone of Height } h &: y_{\text{cm}} = \frac{h}{4} \quad (\text{from base}) \\
  \text{Uniform Hollow Cone of Height } h &: y_{\text{cm}} = \frac{h}{3} \quad (\text{from base})
\end{align}
\end{formulabox}

\begin{usecasebox}[Cavity Problems \& Mid-Air Projectile Explosions]
\begin{itemize}
  \item \textbf{Negative Mass Technique for Cavities:} When a circular or spherical cavity of mass $m_2$ with centroid $\mathbf{r}_2$ is excised from an intact parent body of mass $m_1$ with centroid $\mathbf{r}_1$:
  \begin{equation}
    \mathbf{r}_{\text{cm, remaining}} = \frac{m_1 \mathbf{r}_1 - m_2 \mathbf{r}_2}{m_1 - m_2} = \frac{A_1 \mathbf{r}_1 - A_2 \mathbf{r}_2}{A_1 - A_2} \quad (\text{for laminar 2D bodies})
  \end{equation}
  \item \textbf{Projectile Explosion Invariance:} If a projectile explodes in flight due to internal chemical forces, $\mathbf{F}_{\text{ext}} = M\mathbf{g}$ remains unchanged. Hence, the centre of mass of all fragments continues along the identical parabolic trajectory until the first fragment strikes the ground.
\end{itemize}
\end{usecasebox}

\begin{trapbox}[NEET Trap: Shift of Centre of Mass on Frictionless Floors]
Consider a man of mass $m$ walking a distance $x$ relative to a boat of mass $M$ floating on calm, frictionless water.
\begin{itemize}
  \item Since no external horizontal force acts on the (man + boat) system, the centre of mass cannot shift horizontally: $\Delta x_{\text{cm}} = 0$.
  \item If the boat displaces by distance $X_{\text{boat}}$ in the opposite direction:
  \begin{equation}
    m(x - X_{\text{boat}}) - M X_{\text{boat}} = 0 \implies X_{\text{boat}} = \left(\frac{m}{M + m}\right)x
  \end{equation}
  Assuming the boat moves distance $x$ is a common pitfall. The boat displacement is strictly proportional to reduced system fractions!
\end{itemize}
\end{trapbox}

% =========================================================================
% SECTION 2: RIGID BODY ROTATION \& MOMENT OF INERTIA DERIVATIONS
% =========================================================================
\section{Rigid Body Rotation, Moment of Inertia \& Parallel/Perpendicular Theorems}

\begin{conceptbox}[Rotational Kinematics \& Physical Meaning of Moment of Inertia]
A rigid body maintains constant mutual distances between all constituent particles.
\begin{enumerate}[leftmargin=*]
  \item \textbf{Rotational Kinematics (Constant Angular Acceleration $\alpha$):}
  \begin{equation}
    \omega = \omega_0 + \alpha t, \qquad \theta = \omega_0 t + \frac{1}{2}\alpha t^2, \qquad \omega^2 = \omega_0^2 + 2\alpha\theta, \qquad \theta_n = \omega_0 + \frac{\alpha}{2}(2n - 1)
  \end{equation}
  \item \textbf{Moment of Inertia ($I$) \& Radius of Gyration ($k$):} The quantitative measure of rotational inertia about a specified axis of rotation:
  \begin{equation}
    I = \sum_{i=1}^N m_i r_i^2 = \int r^2\,dm, \qquad I = M k^2 \implies k = \sqrt{\frac{I}{M}}
  \end{equation}
  The radius of gyration $k$ represents the effective radial distance from the rotational axis at which the total mass $M$ could be concentrated without altering its moment of inertia.
\end{enumerate}
\end{conceptbox}

\begin{formulabox}[Proofs of Fundamental Moment of Inertia Theorems]
\textbf{1. Parallel Axis Theorem (Steiner's Theorem):}
Let $I_{\text{cm}}$ be the moment of inertia about an axis passing through the centre of mass, and $I$ be the moment of inertia about a parallel axis located at perpendicular distance $d$.
Let the origin be at the centre of mass ($\int x\,dm = 0, \int y\,dm = 0$). Then:
\begin{align}
  I &= \int r^2\,dm = \int \left[(x + d)^2 + y^2\right] dm = \int (x^2 + y^2)dm + 2d\int x\,dm + d^2\int dm \nonumber \\
  &= I_{\text{cm}} + 0 + M d^2 \implies I = I_{\text{cm}} + M d^2
\end{align}
\textit{Domain of Validity:} Applies to any rigid body of arbitrary three-dimensional shape.

\textbf{2. Perpendicular Axis Theorem:}
For a flat two-dimensional planar lamina situated in the $x-y$ plane ($z = 0$):
\begin{equation}
  I_x = \int y^2\,dm, \qquad I_y = \int x^2\,dm, \qquad I_z = \int (x^2 + y^2)dm = I_x + I_y
\end{equation}
\textit{Domain of Validity:} Strictly restricted to 2D laminar planar geometries. Never apply to 3D bodies!
\end{formulabox}

\begin{formulabox}[Calculus Derivations of Moment of Inertia for Standard Geometries]
\textbf{1. Thin Uniform Rod of Mass $M$ and Length $L$:}
Linear mass density $\lambda = M/L$. Elemental mass $dm = \lambda\,dx$:
\begin{itemize}
  \item \textbf{Transverse axis through center ($x = 0$):}
  \begin{equation}
    I_{\text{cm}} = \int_{-L/2}^{L/2} x^2 \left(\frac{M}{L}\right)dx = \frac{M}{L}\left[\frac{x^3}{3}\right]_{-L/2}^{L/2} = \frac{M}{3L}\left[\frac{L^3}{8} - \left(-\frac{L^3}{8}\right)\right] = \frac{1}{12}ML^2
  \end{equation}
  \item \textbf{Transverse axis through end ($x = 0$ to $L$):} By parallel axis theorem:
  \begin{equation}
    I_{\text{end}} = I_{\text{cm}} + M\left(\frac{L}{2}\right)^2 = \frac{1}{12}ML^2 + \frac{1}{4}ML^2 = \frac{1}{3}ML^2
  \end{equation}
\end{itemize}
\textbf{2. Uniform Circular Disc of Mass $M$ and Radius $R$:}
Surface mass density $\sigma = \frac{M}{\pi R^2}$. Consider a thin concentric elemental ring of radius $r$ and radial width $dr$, mass $dm = \sigma(2\pi r dr) = \frac{2M}{R^2}r\,dr$:
\begin{equation}
  I_z = \int_0^R r^2\,dm = \frac{2M}{R^2}\int_0^R r^3\,dr = \frac{2M}{R^2}\left[\frac{R^4}{4}\right] = \frac{1}{2}MR^2
\end{equation}
By symmetry and perpendicular axis theorem ($I_z = I_x + I_y = 2I_{\text{dia}}$):
\begin{equation}
  I_{\text{diameter}} = \frac{1}{2}I_z = \frac{1}{4}MR^2
\end{equation}
\textbf{3. Uniform Solid Sphere of Mass $M$ and Radius $R$:}
Volume mass density $\rho = \frac{3M}{4\pi R^3}$. Slice the sphere into elemental discs perpendicular to the $x$-axis at position $x$, with radius $y = \sqrt{R^2 - x^2}$ and thickness $dx$.
Mass of elemental disc $dm = \rho(\pi y^2 dx) = \frac{3M}{4R^3}(R^2 - x^2)dx$.
The moment of inertia of this elemental disc about the $x$-axis is:
\begin{equation}
  dI = \frac{1}{2}y^2 dm = \frac{1}{2}(R^2 - x^2)\left[\frac{3M}{4R^3}(R^2 - x^2)dx\right] = \frac{3M}{8R^3}(R^2 - x^2)^2 dx
\end{equation}
Integrating from $x = -R$ to $+R$:
\begin{equation}
  I_{\text{dia}} = \frac{3M}{8R^3}\int_{-R}^{R} (R^4 - 2R^2 x^2 + x^4)dx = \frac{3M}{4R^3}\left[R^4 R - \frac{2R^2 R^3}{3} + \frac{R^5}{5}\right] = \frac{3M}{4R^3}\left[\frac{8}{15}R^5\right] = \frac{2}{5}MR^2
\end{equation}
About a tangential axis: $I_{\text{tangent}} = I_{\text{cm}} + MR^2 = \frac{2}{5}MR^2 + MR^2 = \frac{7}{5}MR^2$.

\textbf{4. Uniform Hollow Spherical Shell of Mass $M$ and Radius $R$:}
Surface mass density $\sigma = \frac{M}{4\pi R^2}$. Elemental ring at polar angle $\theta$ has radius $r = R\sin\theta$, width $R\,d\theta$, mass $dm = \frac{M}{4\pi R^2}(2\pi R\sin\theta \cdot R\,d\theta) = \frac{1}{2}M\sin\theta\,d\theta$:
\begin{equation}
  I_{\text{dia}} = \int_0^\pi (R\sin\theta)^2 dm = \frac{1}{2}MR^2\int_0^\pi \sin^3\theta\,d\theta = \frac{1}{2}MR^2\left(\frac{4}{3}\right) = \frac{2}{3}MR^2
\end{equation}
About a tangential axis: $I_{\text{tangent}} = \frac{2}{3}MR^2 + MR^2 = \frac{5}{3}MR^2$.
\end{formulabox}

\begin{usecasebox}[Moment of Inertia Master Reference Table]
\begin{center}
\begin{tabularx}{\linewidth}{|L|L|Y|Y|}
\toprule
\textbf{Rigid Body Geometry} & \textbf{Axis of Rotation} & \textbf{Moment of Inertia ($I$)} & \textbf{Radius of Gyration ($k^2/R^2$)} \\
\midrule
Thin Circular Ring & Central perpendicular & $MR^2$ & $1.0$ \\
Thin Circular Ring & Diametrical axis & $\frac{1}{2}MR^2$ & $0.5$ \\
Thin Circular Disc & Central perpendicular & $\frac{1}{2}MR^2$ & $0.5$ \\
Thin Circular Disc & Diametrical axis & $\frac{1}{4}MR^2$ & $0.25$ \\
Thin Hollow Cylinder & Longitudinal central & $MR^2$ & $1.0$ \\
Solid Cylinder & Longitudinal central & $\frac{1}{2}MR^2$ & $0.5$ \\
Solid Sphere & Diametrical axis & $\frac{2}{5}MR^2$ & $0.4$ \\
Solid Sphere & Tangential axis & $\frac{7}{5}MR^2$ & $1.4$ \\
Hollow Spherical Shell & Diametrical axis & $\frac{2}{3}MR^2$ & $0.67$ \\
Hollow Spherical Shell & Tangential axis & $\frac{5}{3}MR^2$ & $1.67$ \\
\bottomrule
\end{tabularx}
\end{center}
\end{usecasebox}

% =========================================================================
% SECTION 3: ROTATIONAL DYNAMICS \& PURE ROLLING MOTION ON INCLINES
% =========================================================================
\section{Rotational Dynamics, Angular Momentum \& Pure Rolling on Inclines}

\begin{conceptbox}[Torque, Angular Momentum \& Rolling Kinematics]
\begin{enumerate}[leftmargin=*]
  \item \textbf{Torque Vector \& Rotational Work:}
  \begin{equation}
    \boldsymbol{\tau} = \mathbf{r} \times \mathbf{F} = I\boldsymbol{\alpha}, \qquad W_{\text{rot}} = \int \tau\,d\theta, \qquad P_{\text{rot}} = \boldsymbol{\tau} \cdot \boldsymbol{\omega}
  \end{equation}
  \item \textbf{Angular Momentum Vector ($\mathbf{L}$):} For a particle $\mathbf{L} = \mathbf{r}\times\mathbf{p}$. For a rigid body rotating about a fixed principal symmetry axis:
  \begin{equation}
    \mathbf{L} = I\boldsymbol{\omega}, \qquad \boldsymbol{\tau}_{\text{ext}} = \frac{d\mathbf{L}}{dt} = I\frac{d\boldsymbol{\omega}}{dt} = I\boldsymbol{\alpha}
  \end{equation}
  If net external torque vanishes ($\sum \boldsymbol{\tau}_{\text{ext}} = \mathbf{0}$), total angular momentum is strictly conserved: $I_1 \omega_1 = I_2 \omega_2$.
  \item \textbf{Kinematics of Pure Rolling (Rolling Without Slipping):}
  A circular body of radius $R$ rolling on a stationary surface:
  \begin{equation}
    v_{\text{cm}} = \omega R, \qquad a_{\text{cm}} = \alpha R
  \end{equation}
  The contact point $P$ is instantaneously at rest: $\mathbf{v}_P = \mathbf{v}_{\text{cm}} - \omega R\hat{\mathbf{i}} = \mathbf{0}$. The highest point moves at $\mathbf{v}_{\text{top}} = 2\mathbf{v}_{\text{cm}}$.
  \item \textbf{Total Kinetic Energy in Pure Rolling:}
  \begin{equation}
    K_{\text{total}} = K_{\text{trans}} + K_{\text{rot}} = \frac{1}{2}M v_{\text{cm}}^2 + \frac{1}{2}I_{\text{cm}}\omega^2 = \frac{1}{2}M v_{\text{cm}}^2\left(1 + \frac{k^2}{R^2}\right)
  \end{equation}
\end{enumerate}
\end{conceptbox}

\begin{formulabox}[Rigorous Derivation of Acceleration, Velocity \& Friction on an Incline]
Consider a symmetrical rigid body of mass $M$, radius $R$, and radius of gyration $k$ rolling without slipping down an inclined plane of inclination angle $\theta$ and vertical height $h$:
\begin{align}
  \text{Translational Equation of Motion along Incline: } &Mg\sin\theta - f_s = M a_{\text{cm}} \\
  \text{Rotational Equation about Centre of Mass: } &\tau = f_s R = I_{\text{cm}}\alpha = (M k^2)\left(\frac{a_{\text{cm}}}{R}\right) \implies f_s = M\left(\frac{k^2}{R^2}\right)a_{\text{cm}}
\end{align}
Substituting the friction expression into the translational equation:
\begin{equation}
  Mg\sin\theta = M a_{\text{cm}} + M\left(\frac{k^2}{R^2}\right)a_{\text{cm}} = M a_{\text{cm}}\left(1 + \frac{k^2}{R^2}\right)
\end{equation}
This yields the foundational rolling acceleration formula:
\begin{equation}
  a_{\text{cm}} = \frac{g\sin\theta}{1 + \frac{k^2}{R^2}}
\end{equation}
\textbf{1. Linear Velocity at the Base of the Incline ($v$):} Applying conservation of mechanical energy ($Mgh = K_{\text{total}}$):
\begin{equation}
  Mgh = \frac{1}{2}M v^2\left(1 + \frac{k^2}{R^2}\right) \implies v = \sqrt{\frac{2gh}{1 + \frac{k^2}{R^2}}}
\end{equation}
\textbf{2. Necessary Static Friction Force ($f_s$):}
\begin{equation}
  f_s = M\left(\frac{k^2}{R^2}\right)\left[\frac{g\sin\theta}{1 + \frac{k^2}{R^2}}\right] = \frac{Mg\sin\theta}{1 + \frac{R^2}{k^2}}
\end{equation}
\textbf{3. Minimum Friction Coefficient to Prevent Slipping ($\mu_{\min}$):}
To sustain pure rolling without sliding: $f_s \le \mu_s N = \mu_s Mg\cos\theta$:
\begin{equation}
  \frac{Mg\sin\theta}{1 + \frac{R^2}{k^2}} \le \mu_s Mg\cos\theta \implies \mu_s \ge \frac{\tan\theta}{1 + \frac{R^2}{k^2}}
\end{equation}
\end{formulabox}

\begin{usecasebox}[Race Down an Incline: Comparison of Standard Bodies]
From the acceleration formula $a \propto \frac{1}{1 + k^2/R^2}$, bodies with \textbf{smaller} $k^2/R^2$ possess greater acceleration and reach the bottom first:
\begin{equation}
  \left(\frac{k^2}{R^2}\right)_{\text{Solid Sphere}} (0.4) < \left(\frac{k^2}{R^2}\right)_{\text{Disc}} (0.5) < \left(\frac{k^2}{R^2}\right)_{\text{Hollow Sphere}} (0.67) < \left(\frac{k^2}{R^2}\right)_{\text{Ring}} (1.0)
\end{equation}
\begin{itemize}
  \item \textbf{Order of Acceleration and Striking Velocity:}
  \begin{equation}
    a_{\text{Solid Sphere}} > a_{\text{Disc/Cylinder}} > a_{\text{Hollow Sphere}} > a_{\text{Ring/Hollow Cylinder}}
  \end{equation}
  \item \textbf{Order of Time of Descent ($t = \sqrt{2L/a}$):}
  \begin{equation}
    t_{\text{Solid Sphere}} < t_{\text{Disc/Cylinder}} < t_{\text{Hollow Sphere}} < t_{\text{Ring/Hollow Cylinder}}
  \end{equation}
  The solid sphere always wins the race, irrespective of mass or radius!
\end{itemize}
\end{usecasebox}

\begin{trapbox}[NEET Trap: Work Done by Static Friction in Pure Rolling]
In pure rolling motion on a stationary surface:
\begin{itemize}
  \item The instantaneous velocity of the contact point is zero ($\mathbf{v}_{\text{contact}} = \mathbf{0}$).
  \item Consequently, the power delivered by static friction is $P = \mathbf{f}_s \cdot \mathbf{v}_{\text{contact}} = 0$.
  \item \textbf{Static friction does zero net work} ($W_{f_s} = 0$) during pure rolling. It acts purely as a kinematic constraint that transfers energy reversibly between translational and rotational modes without dissipating mechanical energy as thermal heat.
\end{itemize}
\end{trapbox}

% =========================================================================
% SECTION 4: UNIVERSAL GRAVITATION \& KEPLER'S LAWS
% =========================================================================
\section{Universal Gravitation, Kepler's Laws \& Field/Potential Calculus}

\begin{conceptbox}[Newton's Law of Gravitation \& Kepler's Three Planetary Laws]
\begin{enumerate}[leftmargin=*]
  \item \textbf{Newton's Universal Gravitation Law:} Between two point masses $M$ and $m$ separated by distance $r$:
  \begin{equation}
    \mathbf{F}_{12} = -\frac{G M m}{r^2}\hat{\mathbf{r}}_{12}, \qquad G = 6.674\times 10^{-11}\,\si{N.m^2.kg^{-2}}
  \end{equation}
  \item \textbf{Kepler's First Law (Law of Orbits):} All planets move in elliptical orbits with the Sun located at one of the two foci.
  \item \textbf{Kepler's Second Law (Law of Areas):} The radius vector connecting the Sun to a planet sweeps out equal areas in equal intervals of time.
  \item \textbf{Kepler's Third Law (Law of Periods):} The square of the orbital period $T$ is proportional to the cube of the semi-major axis $a$: $T^2 \propto a^3$.
\end{enumerate}
\end{conceptbox}

\begin{formulabox}[Proofs of Kepler's Second and Third Laws]
\textbf{1. Proof of Kepler's Second Law from Angular Momentum Conservation:}
Gravitational force is purely central: $\mathbf{F} = F(r)\hat{\mathbf{r}}$.
Net external torque on the planet about the Sun:
\begin{equation}
  \boldsymbol{\tau} = \mathbf{r} \times \mathbf{F} = \mathbf{r} \times [F(r)\hat{\mathbf{r}}] = \mathbf{0} \implies \frac{d\mathbf{L}}{dt} = \mathbf{0} \implies \mathbf{L} = m(\mathbf{r}\times\mathbf{v}) = \text{constant}
\end{equation}
In time $dt$, the radius vector sweeps out triangular area $d\mathbf{A} = \frac{1}{2}(\mathbf{r}\times d\mathbf{r})$:
\begin{equation}
  \frac{d\mathbf{A}}{dt} = \frac{1}{2}\left(\mathbf{r}\times\frac{d\mathbf{r}}{dt}\right) = \frac{1}{2}(\mathbf{r}\times\mathbf{v}) = \frac{\mathbf{L}}{2m} = \text{constant}
\end{equation}
Areal velocity is constant because gravity produces zero torque!

\textbf{2. Proof of Kepler's Third Law for Circular Orbit of Radius $r$:}
Gravitational force provides the necessary centripetal acceleration:
\begin{equation}
  \frac{G M m}{r^2} = m\omega^2 r = m\left(\frac{2\pi}{T}\right)^2 r \implies T^2 = \left(\frac{4\pi^2}{G M}\right)r^3
\end{equation}
\end{formulabox}

\begin{formulabox}[Gravitational Field \& Potential Formulations (Solid Sphere vs Spherical Shell)]
Gravitational field $\mathbf{E} = -\boldsymbol{\nabla}V$. The potential $V(\mathbf{r}) = -\int_\infty^{\mathbf{r}} \mathbf{E}\cdot d\mathbf{r}'$.
\begin{enumerate}[leftmargin=*]
  \item \textbf{Uniform Solid Sphere of Mass $M$ and Radius $R$:}
  \begin{align}
    \text{Outside ($r \ge R$): } &E_{\text{out}}(r) = \frac{GM}{r^2}, \qquad V_{\text{out}}(r) = -\frac{GM}{r} \\
    \text{Inside ($r < R$): } &E_{\text{in}}(r) = \frac{GM r}{R^3}, \qquad V_{\text{in}}(r) = -\frac{GM}{2R^3}\left(3R^2 - r^2\right)
  \end{align}
  At the center of the sphere ($r = 0$):
  \begin{equation}
    V_{\text{center}} = -\frac{3GM}{2R} = \frac{3}{2}V_{\text{surface}}
  \end{equation}
  \item \textbf{Uniform Spherical Shell of Mass $M$ and Radius $R$:}
  \begin{align}
    \text{Outside ($r \ge R$): } &E_{\text{out}}(r) = \frac{GM}{r^2}, \qquad V_{\text{out}}(r) = -\frac{GM}{r} \\
    \text{Inside ($r < R$): } &E_{\text{in}}(r) = 0, \qquad V_{\text{in}}(r) = -\frac{GM}{R} = \text{constant}
  \end{align}
  The interior of a spherical shell is an equipotential field-free cavity.
\end{enumerate}
\end{formulabox}

\begin{usecasebox}[Variation of $g$, Escape Speed \& Satellite Orbital Dynamics]
\begin{itemize}
  \item \textbf{Altitude Variation:} At height $h$ above Earth's surface:
  \begin{equation}
    g_h = g\left(1 + \frac{h}{R}\right)^{-2} \approx g\left(1 - \frac{2h}{R}\right) \quad (\text{for } h \ll R)
  \end{equation}
  \item \textbf{Depth Variation:} At depth $d$ below Earth's surface:
  \begin{equation}
    g_d = g\left(1 - \frac{d}{R}\right) \quad (\text{Linear decrease to zero at Earth's center})
  \end{equation}
  \item \textbf{Latitude Variation Due to Earth's Rotation ($\omega = 7.29\times 10^{-5}\,\si{rad.s^{-1}}$):}
  \begin{equation}
    g' = g - \omega^2 R\cos^2\lambda \implies g_{\text{equator}} = g - \omega^2 R, \qquad g_{\text{pole}} = g
  \end{equation}
  \item \textbf{Escape Velocity ($v_e$):} Minimum projection speed required to escape Earth's gravitational well:
  \begin{equation}
    \frac{1}{2}m v_e^2 - \frac{GMm}{R} = 0 \implies v_e = \sqrt{\frac{2GM}{R}} = \sqrt{2gR} \approx 11.2\,\si{km.s^{-1}}
  \end{equation}
  \item \textbf{Satellite Dynamics \& Energy Partition:} For orbit radius $r = R + h$:
  \begin{align}
    \text{Orbital Velocity: } &v_o = \sqrt{\frac{GM}{r}} = \frac{v_e}{\sqrt{2}} \\
    \text{Total Mechanical Energy: } &E = K + U = \frac{GMm}{2r} - \frac{GMm}{r} = -\frac{GMm}{2r} = -K = \frac{U}{2}
  \end{align}
\end{itemize}
\end{usecasebox}

\begin{trapbox}[NEET Trap: Independence of Escape Velocity from Angle \& Drag Paradox]
\begin{itemize}
  \item \textbf{Projection Angle Independence:} Escape speed depends solely on gravitational potential at the launch site: $v_e = \sqrt{-2V_{\text{launch}}} = \sqrt{2GM/R}$. It is completely independent of launch direction $\theta$ (as long as it does not intercept Earth's surface) and independent of projectile mass $m$.
  \item \textbf{Satellite Atmospheric Drag Paradox:} When satellite orbits encounter atmospheric resistance, total energy decreases ($E$ becomes more negative $\implies r$ decreases). Because $v_o = \sqrt{GM/r}$, orbital velocity \textbf{increases} as the orbit spirals downward!
\end{itemize}
\end{trapbox}

% =========================================================================
% SECTION 5: MECHANICAL PROPERTIES OF SOLIDS (ELASTICITY)
% =========================================================================
\section{Mechanical Properties of Solids: Stress, Strain \& Elastic Moduli}

\begin{conceptbox}[Foundations of Elasticity \& Stress-Strain Curve]
Elasticity is the property by virtue of which a material body regains its original shape and dimensions when deforming external forces are removed.
\begin{enumerate}[leftmargin=*]
  \item \textbf{Stress ($\sigma = F_{\text{int}}/A$):} Internal restoring force per unit area. SI Unit: $\si{N.m^{-2}}$ or $\si{Pa}$.
  \item \textbf{Strain ($\epsilon = \Delta L/L$):} Fractional deformation. Strictly dimensionless.
  \item \textbf{Hooke's Law:} Within the proportional limit, stress is directly proportional to strain: $\sigma = E \cdot \epsilon$, where $E$ is the modulus of elasticity.
  \item \textbf{Characteristic Points on the Tensile Stress-Strain Diagram:}
  \begin{itemize}
    \item \textbf{Proportional Limit:} Region of linear Hookean behavior ($\sigma \propto \epsilon$).
    \item \textbf{Elastic Limit (Yield Point $\sigma_y$):} Maximum stress beyond which permanent plastic deformation (permanent set) occurs.
    \item \textbf{Ultimate Tensile Strength ($\sigma_u$):} Maximum stress the material can sustain before necking.
    \item \textbf{Fracture Point:} Separation point. If fracture occurs immediately past yield, material is \textbf{brittle} (e.g., glass, cast iron). If wide plastic region exists, material is \textbf{ductile} (e.g., copper, aluminium).
  \end{itemize}
\end{enumerate}
\end{conceptbox}

\begin{formulabox}[Elastic Moduli Formulations \& Potential Energy Density]
\begin{align}
  \text{Young's Modulus: } &Y = \frac{\text{Tensile Stress}}{\text{Longitudinal Strain}} = \frac{F/A}{\Delta L/L} = \frac{F L}{A \Delta L} \\
  \text{Bulk Modulus: } &B = -\frac{\Delta P}{\Delta V / V} = -V\frac{\Delta P}{\Delta V}, \qquad \text{Compressibility } K = \frac{1}{B} \\
  \text{Shear Modulus (Rigidity): } &\eta = \frac{\text{Shearing Stress}}{\text{Shear Strain}} = \frac{F/A}{\theta} \\
  \text{Poisson's Ratio: } &\sigma_p = -\frac{\text{Lateral Strain}}{\text{Longitudinal Strain}} = -\frac{\Delta d/d}{\Delta L/L} \quad (-1 \le \sigma_p \le 0.5)
\end{align}
\textbf{Elastic Potential Energy Stored in Stretched Wire:}
\begin{equation}
  W = \int_0^{\Delta L} F(x)\,dx = \int_0^{\Delta L} \left(\frac{YA}{L}x\right)dx = \frac{1}{2}\left(\frac{YA}{L}\right)(\Delta L)^2 = \frac{1}{2}F\Delta L
\end{equation}
Energy density (energy per unit volume $u$):
\begin{equation}
  u = \frac{W}{A L} = \frac{1}{2}\left(\frac{F}{A}\right)\left(\frac{\Delta L}{L}\right) = \frac{1}{2}\times\text{Stress}\times\text{Strain} = \frac{1}{2}Y\epsilon^2 = \frac{\sigma^2}{2Y}
\end{equation}
\end{formulabox}

% =========================================================================
% SECTION 6: FLUID STATICS \& SURFACE TENSION PHENOMENA
% =========================================================================
\section{Fluid Statics, Pascal's Principle, Surface Tension \& Capillary Rise}

\begin{conceptbox}[Hydrostatic Pressure, Pascal's Law \& Archimedes' Principle]
\begin{enumerate}[leftmargin=*]
  \item \textbf{Hydrostatic Pressure Variation with Depth:} In an incompressible fluid of uniform density $\rho$:
  \begin{equation}
    \frac{dP}{dz} = -\rho g \implies P(h) = P_0 + \rho g h
  \end{equation}
  \item \textbf{Pascal's Law:} Pressure exerted anywhere upon a confined incompressible fluid is transmitted undiminished equally in all directions throughout the fluid:
  \begin{equation}
    \frac{F_1}{A_1} = \frac{F_2}{A_2} \implies F_2 = F_1\left(\frac{A_2}{A_1}\right) \quad (\text{Hydraulic Lift Principle})
  \end{equation}
  \item \textbf{Archimedes' Principle \& Law of Floatation:} Upward buoyant force equals the weight of displaced fluid:
  \begin{equation}
    F_B = \rho_{\text{fluid}} V_{\text{submerged}} g, \qquad \frac{V_{\text{submerged}}}{V_{\text{total}}} = \frac{\rho_{\text{body}}}{\rho_{\text{fluid}}}
  \end{equation}
\end{enumerate}
\end{conceptbox}

\begin{formulabox}[Surface Tension \& Young-Laplace Excess Pressure Derivations]
Surface tension $T = \frac{F}{L} = \frac{dW}{dA}$ represents surface free energy per unit area.
\begin{enumerate}[leftmargin=*]
  \item \textbf{Excess Pressure Inside a Spherical Liquid Drop of Radius $R$:}
  A drop has a \textbf{single} liquid-gas interface. Increasing radius from $R$ to $R+dR$:
  \begin{equation}
    dW_{\text{surface}} = T \cdot d(4\pi R^2) = T(8\pi R\,dR)
  \end{equation}
  Work done by excess pressure $\Delta P$:
  \begin{equation}
    dW_{\text{pressure}} = \Delta P(4\pi R^2)dR = T(8\pi R\,dR) \implies \Delta P_{\text{drop}} = \frac{2T}{R}
  \end{equation}
  \item \textbf{Excess Pressure Inside a Spherical Soap Bubble of Radius $R$:}
  A bubble possesses \textbf{two} liquid-gas interfaces (interior and exterior surfaces):
  \begin{equation}
    dW_{\text{surface}} = 2 \times T(8\pi R\,dR) = 16\pi R T\,dR \implies \Delta P_{\text{bubble}} = \frac{4T}{R}
  \end{equation}
\end{enumerate}
\end{formulabox}

\begin{formulabox}[Rigorous Calculus Derivation of Capillary Rise (Jurin's Law)]
Consider a narrow capillary tube of radius $r$ dipped vertically into a wetting liquid of density $\rho$ and contact angle $\theta_c$.
The curved meniscus forms a spherical cap of radius of curvature $R_{\text{meniscus}} = \frac{r}{\cos\theta_c}$.
By the Young-Laplace relation, the pressure just below the concave meniscus is reduced relative to atmospheric pressure $P_0$ by:
\begin{equation}
  \Delta P = \frac{2T}{R_{\text{meniscus}}} = \frac{2T\cos\theta_c}{r}
\end{equation}
To equilibrate with hydrostatic pressure at the flat external liquid surface at depth $h$:
\begin{equation}
  \rho g h = \frac{2T\cos\theta_c}{r} \implies h = \frac{2T\cos\theta_c}{r\rho g}
\end{equation}
\end{formulabox}

\begin{usecasebox}[Insufficient Capillary Length \& Droplet Coalescence Thermics]
\begin{itemize}
  \item \textbf{Capillary Tube of Insufficient Length $l < h$:} Liquid rises to the top lip but \textbf{never overflows}. The meniscus flattens, increasing its radius of curvature $R'$ such that:
  \begin{equation}
    h r = l R' \cos\theta_c \implies h R = l R' = \text{constant}
  \end{equation}
  \item \textbf{Coalescence of $n$ Identical Droplets:}
  Volume conservation: $n\left(\frac{4}{3}\pi r^3\right) = \frac{4}{3}\pi R^3 \implies R = n^{1/3}r$.
  Change in surface area: $\Delta A = 4\pi R^2 - n(4\pi r^2) = 4\pi R^2(1 - n^{1/3}) < 0$.
  Surface energy released as heat:
  \begin{equation}
    \Delta E = T|\Delta A| = 4\pi T R^3\left(\frac{1}{r} - \frac{1}{R}\right) = 3 V T\left(\frac{1}{r} - \frac{1}{R}\right)
  \end{equation}
  Resultant temperature increase: $\Delta T = \frac{3T}{\rho s}\left(\frac{1}{r} - \frac{1}{R}\right)$, where $s$ is specific heat.
\end{itemize}
\end{usecasebox}

% =========================================================================
% SECTION 7: HYDRODYNAMICS, BERNOULLI'S PRINCIPLE \& VISCOSITY
% =========================================================================
\section{Fluid Dynamics, Bernoulli's Theorem, Viscosity \& Stokes' Drag}

\begin{conceptbox}[Continuity, Bernoulli Theorem \& Viscous Shear Dynamics]
\begin{enumerate}[leftmargin=*]
  \item \textbf{Equation of Continuity:} For steady, incompressible, irrotational fluid flow:
  \begin{equation}
    A_1 v_1 = A_2 v_2 \implies A v = \text{constant}
  \end{equation}
  \item \textbf{Bernoulli's Principle:} In streamline flow of an ideal (incompressible, non-viscous) fluid, total mechanical energy per unit volume is constant along any streamline:
  \begin{equation}
    P + \frac{1}{2}\rho v^2 + \rho g y = \text{constant}
  \end{equation}
  \item \textbf{Newton's Law of Viscosity:} Viscous drag force between adjacent fluid layers undergoing laminar shear:
  \begin{equation}
    F = -\eta A \left(\frac{dv}{dz}\right)
  \end{equation}
  where $\eta$ is coefficient of viscosity (SI Unit: $\si{Pa.s} = \si{N.s.m^{-2}} = 10\,\text{Poise}$).
  \item \textbf{Reynolds Number ($Re$):} Dimensionless criterion for flow turbulence:
  \begin{equation}
    Re = \frac{\rho v d}{\eta} \quad (Re < 1000 \implies \text{Laminar}; \quad Re > 2000 \implies \text{Turbulent})
  \end{equation}
\end{enumerate}
\end{conceptbox}

\begin{formulabox}[Work-Energy Derivation of Bernoulli's Theorem]
Consider a fluid element of volume $\Delta V$ and mass $\Delta m = \rho\Delta V$ moving along a streamline pipe between stations 1 and 2:
\begin{align}
  \text{Work Done by Pressure Forces: } &W_P = P_1 A_1 \Delta x_1 - P_2 A_2 \Delta x_2 = (P_1 - P_2)\Delta V \\
  \text{Change in Kinetic Energy: } &\Delta K = \frac{1}{2}\Delta m(v_2^2 - v_1^2) = \frac{1}{2}\rho\Delta V(v_2^2 - v_1^2) \\
  \text{Change in Gravitational Potential Energy: } &\Delta U = \Delta m g(y_2 - y_1) = \rho\Delta V g(y_2 - y_1)
\end{align}
Applying the Work-Energy Theorem ($W_P = \Delta K + \Delta U$):
\begin{equation}
  (P_1 - P_2)\Delta V = \frac{1}{2}\rho\Delta V(v_2^2 - v_1^2) + \rho\Delta V g(y_2 - y_1)
\end{equation}
Dividing by $\Delta V$ and rearranging terms:
\begin{equation}
  P_1 + \frac{1}{2}\rho v_1^2 + \rho g y_1 = P_2 + \frac{1}{2}\rho v_2^2 + \rho g y_2 \implies P + \frac{1}{2}\rho v^2 + \rho g y = \text{constant}
\end{equation}
\end{formulabox}

\begin{formulabox}[Venturimeter Flow Rate \& Torricelli's Law of Efflux]
\textbf{1. Venturimeter Flow Rate Derivation:}
A horizontal pipe ($y_1 = y_2$) with inlet area $A_1$ and throat area $A_2$. By continuity $v_2 = (A_1/A_2)v_1$:
\begin{equation}
  P_1 - P_2 = \frac{1}{2}\rho(v_2^2 - v_1^2) = \frac{1}{2}\rho v_1^2\left(\frac{A_1^2}{A_2^2} - 1\right) = \rho_m g h
\end{equation}
The volume flow rate $Q$ is given by:
\begin{equation}
  Q = A_1 v_1 = A_1 A_2 \sqrt{\frac{2(P_1 - P_2)}{\rho(A_1^2 - A_2^2)}} = A_1 A_2 \sqrt{\frac{2\rho_m g h}{\rho(A_1^2 - A_2^2)}}
\end{equation}
\textbf{2. Torricelli's Law of Efflux:}
For an open tank filled to total depth $H$ with an orifice at depth $h$ below the liquid surface:
\begin{align}
  \text{Velocity of Efflux: } &v = \sqrt{2gh} \\
  \text{Horizontal Range of Jet: } &R = v \cdot t = \sqrt{2gh}\sqrt{\frac{2(H-h)}{g}} = 2\sqrt{h(H-h)} \\
  \text{Maximum Horizontal Range: } &R_{\max} = H \quad \text{when orifice is at depth } h = \frac{H}{2}
\end{align}
Orifices at complementary depths $h$ and $(H-h)$ produce identical horizontal ranges $R$.
\end{formulabox}

\begin{formulabox}[Stokes' Law, Terminal Velocity \& Poiseuille's Pipe Flow Equation]
\textbf{1. Stokes' Law for Spherical Particle of Radius $r$:}
\begin{equation}
  F_v = 6\pi\eta r v
\end{equation}
\textbf{2. Rigorous Derivation of Terminal Velocity ($v_t$):}
For a sphere of density $\rho$ falling in fluid of density $\sigma$. At terminal equilibrium:
\begin{equation}
  \text{Gravity } F_g = \text{Buoyant Force } F_B + \text{Viscous Drag } F_v
\end{equation}
\begin{equation}
  \frac{4}{3}\pi r^3 \rho g = \frac{4}{3}\pi r^3 \sigma g + 6\pi\eta r v_t \implies 6\pi\eta r v_t = \frac{4}{3}\pi r^3(\rho - \sigma)g
\end{equation}
Solving for terminal velocity $v_t$:
\begin{equation}
  v_t = \frac{2}{9}\frac{r^2(\rho - \sigma)g}{\eta}
\end{equation}
\textbf{3. Poiseuille's Volume Flow Rate Equation for Cylindrical Pipe:}
For laminar flow of fluid through a capillary tube of internal radius $r$ and length $L$ under pressure differential $\Delta P$:
\begin{equation}
  Q = \frac{\pi \Delta P r^4}{8\eta L}, \qquad \text{Fluid Resistance } R_f = \frac{\Delta P}{Q} = \frac{8\eta L}{\pi r^4}
\end{equation}
\end{formulabox}

\begin{trapbox}[NEET Trap: Terminal Velocity Sign for Bubbles \& Poiseuille Radius Scaling]
\begin{itemize}
  \item \textbf{Air Bubble Rising in Liquid:} Since $\rho_{\text{air}} < \sigma_{\text{liquid}}$, the factor $(\rho - \sigma)$ is \textbf{negative}. Terminal velocity is negative, meaning the bubble ascends upward at $v_t = \frac{2}{9}\frac{r^2(\sigma - \rho)g}{\eta}$.
  \item \textbf{Sensitivity of Flow Rate to Radius (Poiseuille $r^4$ Law):} Flow rate scales as $Q \propto r^4$. If a blood vessel or pipe radius is reduced by a mere $10\%$ ($r' = 0.9r$), the flow rate drops by $1 - (0.9)^4 \approx 34.4\%$. To maintain the original flow rate, pressure gradient must increase by $(1/0.9)^4 \approx 1.52$ ($52\%$ increase)!
\end{itemize}
\end{trapbox}

\begin{ncertbox}[NCERT Critical Extraction: Dynamic Aerodynamic Lift \& Magnus Effect]
\textbf{NCERT Highlight:}
\begin{itemize}
  \item \textbf{Magnus Effect:} When a spinning cylinder or tennis ball travels through air, the boundary layer drags air around it. On the side where surface motion coincides with flow velocity, air speed is higher, resulting in lower pressure by Bernoulli's principle. On the opposite side, air speed is retarded, producing higher pressure. The resulting net transverse pressure differential generates dynamic curve deflection.
  \item \textbf{Aerofoil Design:} The upper surface of an airplane wing is contoured with greater curvature than the lower surface. Streamlines bunch together above the wing, increasing velocity ($v_{\text{top}} > v_{\text{bottom}}$), creating a vertical upward pressure gradient that generates the aerodynamic lift required for flight.
\end{itemize}
\end{ncertbox}

% =========================================================================
% PART II: HEAT, THERMODYNAMICS & WAVE MECHANICS
% =========================================================================
\part{Thermal Physics, Thermodynamics \& Wave Mechanics}

\chapter{Thermal Properties, Thermodynamics, Kinetic Theory \& Oscillations}

This chapter covers the thermal and vibrational behavior of matter. We explore thermal expansion, calorimetry, radiation laws, the first and second laws of thermodynamics with explicit NMC boundary exclusions, the molecular kinetic theory of gases with degrees of freedom and equipartition of energy, simple harmonic motion kinematics and spring-pendulum dynamics, and wave motion in strings and organ pipes.

\begin{mentorbox}[A Personal Note from Satya: Why Sleep is Non-Negotiable for Memory Consolidation]
\textbf{Hey, let's take a quick breath before we dive into Thermodynamics and Waves.}

If you've been grinding Mechanics for four hours straight and your eyes are burning, close this book right now. I'm serious. Go take a 20-minute nap, or if it's late at night, turn off the study lamp and sleep.

Here is the cold, biological truth that coaching institutes rarely tell you: \textit{your brain is an organic organ, not a silicon hard drive.} When you solve complex physics problems---free-body diagrams, torque balances, fluid pressure integrals---your brain builds temporary synaptic patterns in the hippocampus. But those memories are volatile. 

During Slow-Wave Sleep (SWS) and REM sleep, your brain literally replays those firing patterns, transferring the concepts from your fragile short-term memory into the stable neocortex for permanent consolidation. When you pull an all-nighter or cut your sleep down to four hours, two devastating things happen:
\begin{enumerate}[leftmargin=*]
  \item \textbf{Synaptic Fatigue \& Working Memory Collapse:} Your prefrontal processing speed drops by nearly 30\%. On paper, you think you're working hard; in reality, you stare at a simple Mayer's relation ($C_p - C_v = R$) or spring-pendulum phase constant for ten minutes, unable to connect two basic equations.
  \item \textbf{Fatal Sign Slips in Exam Conditions:} Sleep deprivation specifically impairs attention-to-detail circuits. You miss minus signs, you invert ratios, and you misread ``isobaric'' as ``isochoric.'' That's where you lose 5 marks per question.
\end{enumerate}

Aim for 7 to 8 hours of uninterrupted sleep every single night. Sleeping is not ``wasting time''---it is the active biological process that turns the physics you practiced today into permanent, instinctive recall on exam day. Rest well, hydrate, and when you wake up fresh, this chapter on thermodynamics will feel ten times clearer.
\end{mentorbox}

\section{Thermal Properties of Matter, Thermal Expansion \& Calorimetry}

\begin{conceptbox}[Foundations of Heat, Temperature \& Thermometry]
\textbf{1. Heat vs. Temperature:}
\begin{itemize}
  \item \textbf{Heat ($Q$):} Energy in transit across a boundary solely driven by a macroscopic temperature gradient. Heat is a path-dependent boundary phenomenon, not an intrinsic property stored within a system ($[Q] = \mathrm{J}$ or $\mathrm{cal}$, where $1\,\mathrm{cal} = 4.186\,\mathrm{J}$).
  \item \textbf{Temperature ($T$):} Macroscopic measure of the average translational kinetic energy per molecule of a substance. It dictates the direction of spontaneous thermal energy transfer (from higher $T$ to lower $T$).
\end{itemize}
\textbf{2. Thermometric Scales \& Interconversions:}
Any linear thermometric property $X(T)$ (e.g., liquid column length, platinum wire resistance, gas pressure at constant volume) satisfies:
\begin{equation}
  \frac{T - \mathrm{LFP}}{\mathrm{UFP} - \mathrm{LFP}} = \text{constant} \implies \frac{C - 0}{100} = \frac{F - 32}{180} = \frac{K - 273.15}{100} = \frac{R - 0}{80}
\end{equation}
where $\mathrm{LFP}$ is the Lower Fixed Point (ice point at $1\,\mathrm{atm}$) and $\mathrm{UFP}$ is the Upper Fixed Point (steam point at $1\,\mathrm{atm}$).
\begin{equation}
  \frac{C}{5} = \frac{F - 32}{9} = \frac{K - 273.15}{5}
\end{equation}
\textit{Key Coincidence Points:}
\begin{itemize}
  \item $C = F = -40^\circ \implies -40^\circ\mathrm{C} = -40^\circ\mathrm{F}$.
  \item $F = K = 574.25 \implies 574.25^\circ\mathrm{F} = 574.25\,\mathrm{K}$.
  \item Incremental temperature step equivalence: $\Delta T = 1\,\mathrm{K} \equiv 1^\circ\mathrm{C} \equiv 1.8^\circ\mathrm{F}$.
\end{itemize}
\end{conceptbox}

\begin{formulabox}[Thermal Expansion in Solids \& Liquids: Master Governing Equations]
\textbf{1. Anisotropic and Isotropic Thermal Expansion in Solids:}
\begin{itemize}
  \item \textbf{Linear Expansion:}
  \begin{equation}
    L(T) = L_0 (1 + \alpha \Delta T) \implies \Delta L = L_0 \alpha \Delta T, \quad \text{where } \alpha = \frac{1}{L_0}\frac{dL}{dT} \quad [\mathrm{K}^{-1} \text{ or } ^\circ\mathrm{C}^{-1}]
  \end{equation}
  \item \textbf{Superficial (Areal) Expansion:}
  \begin{equation}
    A(T) = A_0 (1 + \beta \Delta T) \implies \Delta A = A_0 \beta \Delta T, \quad \text{where } \beta \approx 2\alpha
  \end{equation}
  \item \textbf{Volumetric (Cubical) Expansion:}
  \begin{equation}
    V(T) = V_0 (1 + \gamma \Delta T) \implies \Delta V = V_0 \gamma \Delta T, \quad \text{where } \gamma \approx 3\alpha
  \end{equation}
  \item \textbf{Isotropic Ratio:} $\alpha : \beta : \gamma = 1 : 2 : 3 \iff \alpha = \frac{\beta}{2} = \frac{\gamma}{3}$.
\end{itemize}
\textbf{2. Proof of $\gamma = 3\alpha$ for Isotropic Solids:}
Consider a cube of initial edge length $L_0$ and initial volume $V_0 = L_0^3$. When temperature rises by $\Delta T$:
\begin{equation}
  V = L^3 = [L_0(1 + \alpha\Delta T)]^3 = L_0^3(1 + 3\alpha\Delta T + 3\alpha^2\Delta T^2 + \alpha^3\Delta T^3)
\end{equation}
Since $\alpha \sim 10^{-5}\,\mathrm{K}^{-1}$, higher-order terms $\alpha^2\Delta T^2 \ll 1$ and $\alpha^3\Delta T^3 \ll 1$:
\begin{equation}
  V \approx V_0(1 + 3\alpha\Delta T) \equiv V_0(1 + \gamma\Delta T) \implies \mathbf{\gamma = 3\alpha}
\end{equation}
\textbf{3. Apparent vs. Real Expansion of Liquids in Vessels:}
\begin{equation}
  \gamma_{\text{real}} = \gamma_{\text{apparent}} + \gamma_{\text{vessel}}
\end{equation}
\begin{itemize}
  \item If $\gamma_{\text{real}} > \gamma_{\text{vessel}}$, liquid level rises upon heating.
  \item If $\gamma_{\text{real}} < \gamma_{\text{vessel}}$, liquid level falls upon heating.
  \item If $\gamma_{\text{real}} = \gamma_{\text{vessel}}$, liquid level remains invariant regardless of temperature.
\end{itemize}
\textbf{4. Density Variation with Temperature:}
Let mass $M$ occupy volume $V_0$ at $T_0$ ($\rho_0 = M/V_0$). At temperature $T_0 + \Delta T$:
\begin{equation}
  \rho(T) = \frac{M}{V(T)} = \frac{M}{V_0(1 + \gamma\Delta T)} = \rho_0(1 + \gamma\Delta T)^{-1} \approx \mathbf{\rho_0(1 - \gamma\Delta T)}
\end{equation}
\end{formulabox}

\begin{usecasebox}[Thermal Stress in Rigidly Clamped Rods \& Bimetallic Strips]
\textbf{1. Thermal Stress Formulation:}
When a rod of length $L_0$, cross-sectional area $A$, and Young's modulus $Y$ is clamped rigidly between two immovable walls and heated by $\Delta T$:
\begin{itemize}
  \item Free thermal expansion prevented: $\Delta L = L_0 \alpha \Delta T$.
  \item Compressive thermal strain induced: $\varepsilon_{\text{thermal}} = \frac{\Delta L}{L_0} = \alpha \Delta T$.
  \item Compressive stress developed:
  \begin{equation}
    \mathbf{\sigma_{\text{th}} = \frac{F}{A} = Y \varepsilon_{\text{thermal}} = Y \alpha \Delta T}
  \end{equation}
  \item Thermal force exerted on clamps: $\mathbf{F = Y A \alpha \Delta T}$. Note that $\sigma_{\text{th}}$ is completely independent of the rod length $L_0$.
\end{itemize}
\textbf{2. Bimetallic Strip Curvature:}
Two strips of equal thickness $d$, made of metals with expansion coefficients $\alpha_1 > \alpha_2$, are riveted together. Upon heating, metal 1 expands more than metal 2, forcing the composite strip into an arc of radius $R$:
\begin{equation}
  R \approx \frac{d}{(\alpha_1 - \alpha_2)\Delta T}
\end{equation}
The strip bends with the higher-$\alpha$ metal on the convex outer side. Upon cooling ($\Delta T < 0$), it bends in the reverse direction (higher-$\alpha$ metal on the concave inner side).
\end{usecasebox}

\begin{formulabox}[Calorimetry, Specific Heat Capacity \& Latent Heat]
\textbf{1. Heat Capacities:}
\begin{itemize}
  \item \textbf{Specific Heat Capacity ($s$ or $c$):} Heat required per unit mass to raise temperature by $1\,\mathrm{K}$:
  \begin{equation}
    dQ = m s dT \implies Q = m \int_{T_1}^{T_2} s(T) dT = m s \Delta T \quad (\text{if } s = \text{const})
  \end{equation}
  For water: $s_w = 1.0\,\mathrm{cal}/(\mathrm{g}\cdot^\circ\mathrm{C}) = 4186\,\mathrm{J}/(\mathrm{kg}\cdot\mathrm{K})$. For ice: $s_{\text{ice}} \approx 0.5\,\mathrm{cal}/(\mathrm{g}\cdot^\circ\mathrm{C}) = 2100\,\mathrm{J}/(\mathrm{kg}\cdot\mathrm{K})$.
  \item \textbf{Molar Heat Capacity ($C$):} $C = \frac{1}{n}\frac{dQ}{dT} = M_{\text{molar}} s$.
  \item \textbf{Water Equivalent of Calorimeter ($W$):} Mass of water that absorbs/releases the same heat as the calorimeter for an identical temperature change:
  \begin{equation}
    W = m_{\text{cal}} \frac{s_{\text{cal}}}{s_w} \implies Q = (m_w + W) s_w \Delta T
  \end{equation}
\end{itemize}
\textbf{2. Latent Heat \& Phase Changes (Isothermal Heat Addition):}
\begin{equation}
  Q = m L
\end{equation}
\begin{itemize}
  \item \textbf{Latent Heat of Fusion of Ice ($L_f$):} $L_f \approx 80\,\mathrm{cal}/\mathrm{g} = 3.33 \times 10^5\,\mathrm{J}/\mathrm{kg}$ at $0^\circ\mathrm{C}$.
  \item \textbf{Latent Heat of Vaporization of Water ($L_v$):} $L_v \approx 540\,\mathrm{cal}/\mathrm{g} = 2.26 \times 10^6\,\mathrm{J}/\mathrm{kg}$ at $100^\circ\mathrm{C}$.
\end{itemize}
\textbf{3. Principle of Calorimetry (Conservation of Energy):}
In a thermally isolated calorimeter:
\begin{equation}
  \mathbf{\sum Q_{\text{lost by hot bodies}} = \sum Q_{\text{gained by cold bodies}}}
\end{equation}
\end{formulabox}

\begin{trapbox}[Calorimetry Pitfalls: Ice-Water-Steam Equilibrium Traps]
\textbf{Common Error:} Blindly applying $m_1 s_1 (T_1 - T_f) = m_2 s_2 (T_f - T_2)$ without verifying whether sufficient heat exists to melt all ice or condense all steam.
\begin{itemize}
  \item \textbf{Ice + Water Scenario:} Always calculate $Q_{\text{available}}$ from water cooling to $0^\circ\mathrm{C}$ vs. $Q_{\text{melt}}$ required to melt all ice at $0^\circ\mathrm{C}$. If $Q_{\text{available}} < m_{\text{ice}} L_f$, the final mixture temperature is \textbf{strictly $0^\circ\mathrm{C}$}, and only a fraction $m' = Q_{\text{available}}/L_f$ of ice melts.
  \item \textbf{Steam + Water Scenario:} Always verify if water can absorb all latent heat of steam ($540\,\mathrm{cal}/\mathrm{g}$) before heating to $100^\circ\mathrm{C}$. If not, final equilibrium temperature remains \textbf{strictly $100^\circ\mathrm{C}$}.
\end{itemize}
\end{trapbox}

\section{Modes of Heat Transfer \& Radiation Laws}

\begin{conceptbox}[Thermal Conduction: Fourier's Law \& Composite Slabs]
\textbf{1. Fourier's Law of Steady-State Conduction:}
Thermal energy flowing per unit time through a cross-sectional area $A$ normal to a temperature gradient $dT/dx$:
\begin{equation}
  \mathbf{H = \frac{dQ}{dt} = -k A \frac{dT}{dx}}
\end{equation}
where $k$ is thermal conductivity ($\mathrm{W}\cdot\mathrm{m}^{-1}\cdot\mathrm{K}^{-1}$). For a uniform rod of length $L$ with ends maintained at $T_H$ and $T_C$:
\begin{equation}
  H = \frac{k A (T_H - T_C)}{L} = \frac{T_H - T_C}{R_{\text{th}}}, \quad \text{where } \mathbf{R_{\text{th}} = \frac{L}{k A}} \quad [\mathrm{K}\cdot\mathrm{W}^{-1}]
\end{equation}
$R_{\text{th}}$ is the \textbf{thermal resistance}, directly analogous to electrical resistance $R_{\text{elec}} = \frac{L}{\sigma A}$.
\end{conceptbox}

\begin{formulabox}[Series \& Parallel Thermal Conductors: Derivations]
\textbf{1. Series Combination of Two Conducting Slabs:}
Consider two slabs of lengths $L_1, L_2$, areas $A$, and conductivities $k_1, k_2$ placed in end-to-end contact between temperatures $T_1$ and $T_2$ ($T_1 > T_2$). Let $T_0$ be the interface temperature. Under steady state, the heat current $H$ is identical through both:
\begin{equation}
  H = \frac{k_1 A (T_1 - T_0)}{L_1} = \frac{k_2 A (T_0 - T_2)}{L_2} \implies R_{\text{eq}} = R_1 + R_2 = \frac{L_1}{k_1 A} + \frac{L_2}{k_2 A}
\end{equation}
Equating expressions to solve for the interface temperature $T_0$:
\begin{equation}
  \mathbf{T_0 = \frac{\frac{k_1}{L_1} T_1 + \frac{k_2}{L_2} T_2}{\frac{k_1}{L_1} + \frac{k_2}{L_2}} = \frac{k_1 L_2 T_1 + k_2 L_1 T_2}{k_1 L_2 + k_2 L_1}}
\end{equation}
Equivalent thermal conductivity $k_{\text{eq}}$ for total length $L_1 + L_2$:
\begin{equation}
  \frac{L_1 + L_2}{k_{\text{eq}} A} = \frac{L_1}{k_1 A} + \frac{L_2}{k_2 A} \implies \mathbf{k_{\text{eq}} = \frac{L_1 + L_2}{\frac{L_1}{k_1} + \frac{L_2}{k_2}}} \quad \left(\text{If } L_1 = L_2: \; k_{\text{eq}} = \frac{2k_1 k_2}{k_1 + k_2}\right)
\end{equation}
\textbf{2. Parallel Combination of Two Conducting Slabs:}
Plates of areas $A_1, A_2$ and common length $L$ connected between reservoirs at $T_1$ and $T_2$:
\begin{equation}
  H_{\text{net}} = H_1 + H_2 \implies \frac{1}{R_{\text{eq}}} = \frac{1}{R_1} + \frac{1}{R_2} \implies \mathbf{k_{\text{eq}} = \frac{k_1 A_1 + k_2 A_2}{A_1 + A_2}}
\end{equation}
\end{formulabox}

\begin{conceptbox}[Thermal Radiation: Stefan-Boltzmann \& Wien's Laws]
\textbf{1. Blackbody Radiation \& Prevost's Theory:}
Prevost's theory asserts that all bodies at temperatures $T > 0\,\mathrm{K}$ emit thermal radiation continuously to, and absorb radiation from, their surroundings. Thermal equilibrium represents dynamic parity between emission and absorption.
\begin{itemize}
  \item \textbf{Stefan-Boltzmann Law:} Total radiant energy emitted per unit surface area per second by a blackbody:
  \begin{equation}
    E_b = \sigma T^4
  \end{equation}
  where $\sigma = 5.670 \times 10^{-8}\,\mathrm{W}\cdot\mathrm{m}^{-2}\cdot\mathrm{K}^{-4}$ is the Stefan-Boltzmann constant.
  \item For a non-blackbody of surface area $A$ and emissivity $e$ ($0 < e < 1$):
  \begin{equation}
    P_{\text{emit}} = e \sigma A T^4
  \end{equation}
  \item \textbf{Net Radiative Heat Transfer with Surroundings at $T_0$:}
  \begin{equation}
    \mathbf{P_{\text{net}} = \frac{dQ}{dt} = e \sigma A (T^4 - T_0^4)}
  \end{equation}
  \item \textbf{Wien's Displacement Law:} The wavelength $\lambda_m$ corresponding to maximum spectral emissive power is inversely proportional to absolute temperature:
  \begin{equation}
    \mathbf{\lambda_m T = b = 2.898 \times 10^{-3}\,\mathrm{m}\cdot\mathrm{K}}
  \end{equation}
  As temperature increases, peak emission shifts towards shorter wavelengths (higher frequencies).
\end{itemize}
\end{conceptbox}

\begin{formulabox}[Newton's Law of Cooling: Derivation \& Analytical Form]
\textbf{1. Calculus Derivation from Stefan-Boltzmann Law:}
Consider a body of mass $m$, specific heat capacity $s$, and surface area $A$ at temperature $T$ cooling in surroundings at temperature $T_0$, where $\Delta T = T - T_0 \ll T_0$.
\begin{equation}
  T^4 - T_0^4 = (T_0 + \Delta T)^4 - T_0^4 = T_0^4 \left[\left(1 + \frac{\Delta T}{T_0}\right)^4 - 1\right]
\end{equation}
Using the binomial expansion $(1 + x)^4 \approx 1 + 4x$ for $x = \Delta T / T_0 \ll 1$:
\begin{equation}
  T^4 - T_0^4 \approx T_0^4 \left(1 + 4\frac{\Delta T}{T_0} - 1\right) = 4 T_0^3 (T - T_0)
\end{equation}
Substituting this into the net radiant heat loss rate:
\begin{equation}
  \frac{dQ}{dt} = -m s \frac{dT}{dt} = e \sigma A \cdot 4 T_0^3 (T - T_0) \implies \mathbf{\frac{dT}{dt} = -k (T - T_0)}
\end{equation}
where cooling constant $k = \frac{4 e \sigma A T_0^3}{m s}$.
\textbf{2. Exact Exponential Solution:}
\begin{equation}
  \int_{T_i}^{T(t)} \frac{dT}{T - T_0} = -k \int_0^t dt \implies \mathbf{T(t) - T_0 = (T_i - T_0) e^{-k t}}
\end{equation}
\textbf{3. Practical Average Approximation Formula for Problems:}
If a body cools from $T_1$ to $T_2$ in time $t$:
\begin{equation}
  \mathbf{\frac{T_1 - T_2}{t} = K \left(\frac{T_1 + T_2}{2} - T_0\right)}
\end{equation}
\end{formulabox}

\begin{ncertbox}[NCERT Extraction: Anomalous Expansion of Water]
Water exhibits anomalous expansion between $0^\circ\mathrm{C}$ and $4^\circ\mathrm{C}$. When water is heated from $0^\circ\mathrm{C}$, its volume \textbf{decreases} and its density \textbf{increases}, reaching a maximum density of $\rho_{\max} \approx 1.000\,\mathrm{g/cm}^3$ at $3.98^\circ\mathrm{C} \approx 4^\circ\mathrm{C}$. Above $4^\circ\mathrm{C}$, water expands normally with temperature.
\begin{itemize}
  \item \textit{Ecological Consequence:} In freezing winters, surface water cools to $4^\circ\mathrm{C}$ and sinks. Once the entire body reaches $4^\circ\mathrm{C}$, further cooling drops surface water to $0^\circ\mathrm{C}$ where it freezes into ice ($\rho_{\text{ice}} < \rho_{\text{water}}$). The ice layer floats on top, acting as a thermal insulator ($k_{\text{ice}} \approx 2.2\,\mathrm{W/m\cdot K}$), preserving liquid water at $4^\circ\mathrm{C}$ beneath and sustaining aquatic life.
\end{itemize}
\end{ncertbox}

\section{Thermodynamics: First Law, Processes \& Mayer's Relation}

\begin{conceptbox}[Thermodynamic Systems, State Functions \& First Law]
\textbf{1. Thermodynamic Equilibrium \& Zeroth Law:}
Two systems in mutual thermal equilibrium with a third system are in thermal equilibrium with each other. This establishes temperature as the unique scalar state property governing thermal equilibrium.
\textbf{2. Internal Energy ($U$):}
Internal energy is a unique thermodynamic \textbf{state function} reflecting the sum of molecular kinetic and intermolecular potential energies. For an ideal gas (zero intermolecular potential):
\begin{equation}
  U = U(T) = n C_v T \implies \mathbf{dU = n C_v dT}
\end{equation}
Across any cyclic process, $\Delta U_{\text{cycle}} = \oint dU = 0$.
\textbf{3. First Law of Thermodynamics (Energy Conservation):}
\begin{equation}
  \mathbf{\Delta Q = \Delta U + W} \iff \mathbf{dQ = dU + dW} = n C_v dT + P dV
\end{equation}
\textit{Physics Sign Conventions:}
\begin{itemize}
  \item $Q > 0$ when heat is absorbed by the system; $Q < 0$ when heat is released.
  \item $W = \int P dV > 0$ when work is done \textbf{by} the system (expansion, $dV > 0$); $W < 0$ when work is done \textbf{on} the system (compression, $dV < 0$).
\end{itemize}
\end{conceptbox}

\begin{formulabox}[Work Done in Master Thermodynamic Processes: Step-by-Step Derivations]
\textbf{1. Isochoric Process (Constant Volume, $V = \text{const}$):}
\begin{equation}
  dV = 0 \implies \mathbf{W = 0}, \quad \mathbf{\Delta Q = \Delta U = n C_v \Delta T}
\end{equation}
\textbf{2. Isobaric Process (Constant Pressure, $P = \text{const}$):}
\begin{equation}
  \mathbf{W = P \Delta V = P(V_2 - V_1) = n R (T_2 - T_1)}
\end{equation}
\begin{equation}
  \Delta Q = n C_p \Delta T, \quad \Delta U = n C_v \Delta T \implies \mathbf{\Delta U : W : \Delta Q = C_v : R : C_p = 1 : (\gamma - 1) : \gamma}
\end{equation}
\textbf{3. Isothermal Process (Constant Temperature, $T = \text{const}$):}
Since $T = \text{const}$, $\Delta U = 0 \implies \Delta Q = W$. Using $P = \frac{nRT}{V}$:
\begin{equation}
  \mathbf{W = \int_{V_1}^{V_2} P dV = n R T \int_{V_1}^{V_2} \frac{dV}{V} = n R T \ln\left(\frac{V_2}{V_1}\right) = n R T \ln\left(\frac{P_1}{P_2}\right) = 2.303 n R T \log_{10}\left(\frac{V_2}{V_1}\right)}
\end{equation}
\textbf{4. Adiabatic Process (Zero Heat Exchange, $\Delta Q = 0$):}
Governing equations: $P V^\gamma = \text{const}$, $T V^{\gamma-1} = \text{const}$, $P^{1-\gamma} T^\gamma = \text{const}$.
From First Law: $dQ = 0 \implies dW = -dU = -n C_v dT$.
\begin{equation}
  W = -\int_{T_1}^{T_2} n C_v dT = -n C_v (T_2 - T_1) = n \left(\frac{R}{\gamma - 1}\right)(T_1 - T_2)
\end{equation}
Using $nRT_1 = P_1 V_1$ and $nRT_2 = P_2 V_2$:
\begin{equation}
  \mathbf{W = \frac{P_1 V_1 - P_2 V_2}{\gamma - 1} = \frac{n R (T_1 - T_2)}{\gamma - 1} = -\Delta U}
\end{equation}
\textbf{5. General Polytropic Process ($P V^m = \text{const}$):}
\begin{equation}
  W = \frac{P_1 V_1 - P_2 V_2}{m - 1}, \qquad \mathbf{C = C_v + \frac{R}{1 - m} = \frac{R}{\gamma - 1} + \frac{R}{1 - m}}
\end{equation}
\end{formulabox}

\begin{formulabox}[Mayer's Relation Derivation: $C_p - C_v = R$]
Consider 1 mole of an ideal gas ($n=1$).
\begin{enumerate}
  \item Heat added at constant volume: $dQ_v = C_v dT$. Since $dV = 0$, $dW = 0$, so by First Law:
  \begin{equation}
    dU = dQ_v = C_v dT
  \end{equation}
  Because internal energy of an ideal gas depends solely on temperature, $dU = C_v dT$ holds across \textbf{any} process.
  \item Heat added at constant pressure: $dQ_p = C_p dT$. By First Law:
  \begin{equation}
    dQ_p = dU + dW = C_v dT + P dV
  \end{equation}
  \item Differentiating the ideal gas equation of state $P V = R T$ at constant pressure $P$:
  \begin{equation}
    P dV = R dT
  \end{equation}
  \item Substituting $P dV = R dT$ into the First Law expression:
  \begin{equation}
    C_p dT = C_v dT + R dT \implies (C_p - C_v) dT = R dT \implies \mathbf{C_p - C_v = R}
  \end{equation}
\end{enumerate}
\textit{In terms of specific heats (per unit mass):} $c_p - c_v = \frac{R}{M_{\text{molar}}}$.
\end{formulabox}

\begin{usecasebox}[P-V Indicator Diagram Slopes \& Cyclic Process Work]
\textbf{1. Comparison of Adiabatic vs. Isothermal Slopes:}
\begin{itemize}
  \item Isothermal curve: $P V = \text{const} \implies P dV + V dP = 0 \implies \mathbf{\left(\frac{dP}{dV}\right)_{\text{iso}} = -\frac{P}{V}}$.
  \item Adiabatic curve: $P V^\gamma = \text{const} \implies V^\gamma dP + \gamma P V^{\gamma-1} dV = 0 \implies \mathbf{\left(\frac{dP}{dV}\right)_{\text{ad}} = -\gamma \frac{P}{V}}$.
  \item \textbf{Fundamental Slope Theorem:}
  \begin{equation}
    \mathbf{\left(\frac{dP}{dV}\right)_{\text{adiabatic}} = \gamma \left(\frac{dP}{dV}\right)_{\text{isothermal}}}
  \end{equation}
  Since $\gamma > 1$, the adiabatic curve is \textbf{$\gamma$ times steeper} than the isothermal curve at any point of intersection.
\end{itemize}
\textbf{2. Expansion from Common Initial State $(P_0, V_0)$ to Final Volume $V_f > V_0$:}
\begin{equation}
  P_{\text{isobaric}} > P_{\text{isothermal}} > P_{\text{adiabatic}} \implies \mathbf{W_{\text{isobaric}} > W_{\text{isothermal}} > W_{\text{adiabatic}}}
\end{equation}
\textbf{3. Compression from Common Initial State $(P_0, V_0)$ to Final Volume $V_f < V_0$:}
Work done \textbf{on} the gas is largest for adiabatic compression:
\begin{equation}
  |W_{\text{adiabatic}}| > |W_{\text{isothermal}}| > |W_{\text{isobaric}}|
\end{equation}
\textbf{4. Net Work in Cyclic Processes:}
The net work in any closed cycle on a $P-V$ diagram equals the enclosed loop area:
\begin{equation}
  W_{\text{net}} = \oint P dV = \begin{cases} > 0 & \text{Clockwise cycle (heat engine, net work output)} \\ < 0 & \text{Counter-clockwise cycle (work absorbed)} \end{cases}
\end{equation}
\end{usecasebox}

\section{Second Law of Thermodynamics \& Syllabus Boundaries}

\begin{conceptbox}[Second Law of Thermodynamics: Clausius \& Kelvin-Planck Formulations]
The First Law dictates energy conservation but fails to predict the spontaneous direction of natural processes. The Second Law establishes irreversibility:
\begin{itemize}
  \item \textbf{Kelvin-Planck Statement:} It is impossible to construct a heat engine operating in a thermodynamic cycle that absorbs thermal energy from a single thermal reservoir and converts $100\%$ of this energy into mechanical work without rejecting any heat to a colder reservoir. (No engine can operate at $100\%$ thermal efficiency; perpetual motion machines of the second kind are impossible).
  \item \textbf{Clausius Statement:} It is impossible to construct a cyclic device that transfers heat spontaneously from a cooler body to a hotter body without requiring external mechanical work input.
  \item \textbf{Equivalence:} A violation of the Kelvin-Planck statement directly entails a violation of the Clausius statement, and vice versa.
\end{itemize}
\end{conceptbox}

\begin{trapbox}[Syllabus Boundary Alert: Carnot Engine \& Refrigerator Deletion]
\textbf{NMC REVISION WARNING FOR NEET UG:}
Under the updated NEET UG syllabus issued by the National Medical Commission (NMC):
\begin{itemize}
  \item \textbf{DELETED TOPICS:} The theoretical formulation of the \textbf{Carnot Engine}, Carnot cycle efficiency ($\eta = 1 - T_C/T_H$), Carnot theorem proofs, and \textbf{Refrigerators / Heat Pumps} (coefficient of performance $\beta = Q_C/W = T_C/(T_H - T_C)$) have been \textbf{formally removed}.
  \item \textbf{EXAM ORIENTATION:} Questions will \textbf{not} test Carnot cycle calculations. Candidates must focus strictly on General Thermodynamic Cycles (triangular, rectangular $P-V$ loops), First Law applications ($\Delta Q = \Delta U + W$), work in isothermal/adiabatic/isobaric expansions, and the conceptual statements of the Second Law.
\end{itemize}
\end{trapbox}

\section{Kinetic Theory of Gases (KTG): Molecular Pressure \& Equipartition}

\begin{conceptbox}[Postulates of the Kinetic Theory of an Ideal Gas]
An ideal gas is modeled micro-mechanically under the following foundational postulates:
\begin{enumerate}
  \item The gas comprises an enormous number of identical, tiny, hard spherical molecules in ceaseless, chaotic random motion across all directions with all possible speeds.
  \item The physical volume occupied by the gas molecules themselves is negligibly small compared to the total macroscopic volume of the container ($V_{\text{molecules}} \ll V_{\text{container}}$).
  \item No intermolecular forces (attractive or repulsive) act between molecules except during collisions. Hence, potential energy is strictly zero; all internal energy is purely translational/rotational kinetic energy.
  \item All intermolecular collisions and collisions with container walls are perfectly elastic ($\Delta K = 0$). Collision durations are negligible compared to transit times between collisions.
  \item Gravity has negligible effect on molecular trajectory between collisions due to high speeds ($\sim 10^2 - 10^3\,\mathrm{m/s}$).
\end{enumerate}
\end{conceptbox}

\begin{formulabox}[First-Principles Derivation of Ideal Gas Pressure: $P = \frac{1}{3}\rho v_{\text{rms}}^2$]
Consider $N$ identical gas molecules, each of mass $m$, enclosed in a cubical box of side $L$ and volume $V = L^3$.
\begin{enumerate}
  \item Let a molecule possess velocity vector $\mathbf{v} = v_x \hat{\mathbf{i}} + v_y \hat{\mathbf{j}} + v_z \hat{\mathbf{k}}$ with speed squared $v^2 = v_x^2 + v_y^2 + v_z^2$.
  \item When this molecule collides elastically with the wall perpendicular to the $x$-axis at $x = L$:
  \begin{equation}
    p_{x,\text{initial}} = m v_x, \qquad p_{x,\text{final}} = -m v_x \implies \Delta p_x = -2 m v_x
  \end{equation}
  Momentum imparted to the wall per collision is $+2 m v_x$.
  \item Time elapsed between consecutive collisions of this molecule with the same wall:
  \begin{equation}
    \Delta t = \frac{2 L}{v_x}
  \end{equation}
  \item Average force exerted by this single molecule on the wall:
  \begin{equation}
    F_1 = \frac{\Delta p_{\text{wall}}}{\Delta t} = \frac{2 m v_x}{2 L / v_x} = \frac{m v_x^2}{L}
  \end{equation}
  \item Total force exerted by all $N$ molecules on the wall of area $A = L^2$:
  \begin{equation}
    F_{\text{total}} = \sum_{i=1}^N \frac{m v_{xi}^2}{L} = \frac{m}{L} N \overline{v_x^2}
  \end{equation}
  \item Since molecular motion is isotropic, $\overline{v_x^2} = \overline{v_y^2} = \overline{v_z^2} = \frac{1}{3}\overline{v^2} = \frac{1}{3} v_{\text{rms}}^2$.
  \item Pressure $P = \frac{F_{\text{total}}}{A} = \frac{m N \overline{v_x^2}}{L \cdot L^2} = \frac{1}{3}\frac{N m}{V} v_{\text{rms}}^2$:
  \begin{equation}
    \mathbf{P = \frac{1}{3}\rho v_{\text{rms}}^2 = \frac{1}{3}\frac{M_{\text{total}}}{V} v_{\text{rms}}^2}
  \end{equation}
\end{enumerate}
\end{formulabox}

\begin{formulabox}[Kinetic Interpretation of Temperature \& Molecular Speed Spectrum]
\textbf{1. Average Kinetic Energy per Molecule:}
Rewriting the pressure equation as $P V = \frac{2}{3} \left(N \cdot \frac{1}{2}m v_{\text{rms}}^2\right) = \frac{2}{3} E_{\text{trans}}$.
Equating to the ideal gas law $P V = N k_B T$ (where $k_B = R/N_A = 1.381 \times 10^{-23}\,\mathrm{J/K}$):
\begin{equation}
  \frac{2}{3} E_{\text{trans}} = N k_B T \implies \mathbf{\overline{K}_{\text{trans}} = \frac{1}{2}m v_{\text{rms}}^2 = \frac{3}{2}k_B T}
\end{equation}
Average translational kinetic energy depends \textbf{solely on absolute temperature $T$}, invariant of molecular mass or gas species!
\textbf{2. Molecular Speed Hierarchy:}
\begin{itemize}
  \item \textbf{Root-Mean-Square Speed ($v_{\text{rms}}$):}
  \begin{equation}
    \mathbf{v_{\text{rms}} = \sqrt{\frac{3 k_B T}{m}} = \sqrt{\frac{3 R T}{M_{\text{molar}}}} = \sqrt{\frac{3 P}{\rho}}}
  \end{equation}
  \item \textbf{Average Speed ($v_{\text{avg}}$):}
  \begin{equation}
    \mathbf{v_{\text{avg}} = \sqrt{\frac{8 k_B T}{\pi m}} = \sqrt{\frac{8 R T}{\pi M_{\text{molar}}}}} \approx 0.921 v_{\text{rms}}
  \end{equation}
  \item \textbf{Most Probable Speed ($v_{\text{mp}}$):}
  \begin{equation}
    \mathbf{v_{\text{mp}} = \sqrt{\frac{2 k_B T}{m}} = \sqrt{\frac{2 R T}{M_{\text{molar}}}}} \approx 0.816 v_{\text{rms}}
  \end{equation}
  \item \textbf{Golden Speed Ratio:}
  \begin{equation}
    \mathbf{v_{\text{mp}} : v_{\text{avg}} : v_{\text{rms}} = \sqrt{2} : \sqrt{\frac{8}{\pi}} : \sqrt{3} \approx 1 : 1.128 : 1.224}
  \end{equation}
\end{itemize}
\end{formulabox}

\begin{formulabox}[Law of Equipartition of Energy, Degrees of Freedom \& Specific Heat Ratios]
\textbf{1. Equipartition Principle:}
In thermal equilibrium at temperature $T$, the total internal energy is partitioned equally among all quadratic independent degrees of freedom ($f$), each contributing an average energy of $\mathbf{\frac{1}{2}k_B T}$ per molecule:
\begin{equation}
  \overline{E}_{\text{molecule}} = \frac{f}{2}k_B T, \quad U_{\text{molar}} = \frac{f}{2} R T \implies \mathbf{C_v = \frac{f}{2}R}, \quad \mathbf{C_p = C_v + R = \left(\frac{f}{2} + 1\right)R}, \quad \mathbf{\gamma = \frac{C_p}{C_v} = 1 + \frac{2}{f}}
\end{equation}
\textbf{2. Comprehensive Gas Classification Table:}
\begin{center}
\begin{tabularx}{\linewidth}{|l|c|c|c|c|Y|}
\hline
\textbf{Molecular Atomicity} & \textbf{DoF ($f$)} & $C_v$ & $C_p$ & $\gamma = C_p / C_v$ & \textbf{Representative Gases} \\
\hline
\textbf{Monoatomic} & $3$ trans & $\frac{3}{2}R$ & $\frac{5}{2}R$ & $\frac{5}{3} \approx 1.67$ & $\mathrm{He, Ne, Ar, Kr}$ \\
\hline
\textbf{Diatomic (Rigid, Room $T$)} & $3\text{ trans} + 2\text{ rot} = 5$ & $\frac{5}{2}R$ & $\frac{7}{2}R$ & $\frac{7}{5} = 1.40$ & $\mathrm{H_2, N_2, O_2, CO}$ \\
\hline
\textbf{Diatomic (Vibrating, High $T$)} & $5 + 2\text{ vib} = 7$ & $\frac{7}{2}R$ & $\frac{9}{2}R$ & $\frac{9}{7} \approx 1.29$ & $\mathrm{Cl_2, Br_2}$ at high $T$ \\
\hline
\textbf{Non-linear Polyatomic} & $3\text{ trans} + 3\text{ rot} = 6$ & $3R$ & $4R$ & $\frac{4}{3} \approx 1.33$ & $\mathrm{H_2O, NH_3, CH_4}$ \\
\hline
\end{tabularx}
\end{center}
\textbf{3. Non-Reacting Mixtures of Gases:}
For a mixture of $n_1$ moles of gas 1 ($C_{v1}, \gamma_1$) and $n_2$ moles of gas 2 ($C_{v2}, \gamma_2$):
\begin{equation}
  \mathbf{C_{v,\text{mix}} = \frac{n_1 C_{v1} + n_2 C_{v2}}{n_1 + n_2}}, \qquad \mathbf{C_{p,\text{mix}} = \frac{n_1 C_{p1} + n_2 C_{p2}}{n_1 + n_2}}, \qquad \mathbf{\gamma_{\text{mix}} = \frac{C_{p,\text{mix}}}{C_{v,\text{mix}}} = \frac{n_1 C_{p1} + n_2 C_{p2}}{n_1 C_{v1} + n_2 C_{v2}}}
\end{equation}
\textbf{4. Mean Free Path ($\lambda$):}
Average distance traversed by a gas molecule between successive collisions:
\begin{equation}
  \mathbf{\lambda = \frac{1}{\sqrt{2} \pi n_v d^2} = \frac{k_B T}{\sqrt{2} \pi d^2 P}}
\end{equation}
where $n_v = N/V$ is number density, $d$ is molecular diameter. Notice that $\lambda \propto T$ at constant pressure, and $\lambda \propto \frac{1}{P}$ at constant temperature.
\end{formulabox}

\section{Simple Harmonic Motion (SHM): Kinematics, Energy \& Oscillators}

\begin{conceptbox}[Mathematical Foundations of Simple Harmonic Motion]
SHM is bounded periodic motion wherein the restoring force/torque is directly proportional to displacement from the stable equilibrium position and directed towards that equilibrium:
\begin{equation}
  F(x) = -k x \implies m \frac{d^2 x}{dt^2} + k x = 0 \implies \mathbf{\frac{d^2 x}{dt^2} + \omega^2 x = 0} \quad \left(\text{where } \omega = \sqrt{\frac{k}{m}}\right)
\end{equation}
\textbf{General Solution \& Kinematic Equations:}
\begin{itemize}
  \item \textbf{Displacement:} $x(t) = A \sin(\omega t + \phi_0)$, where $A$ is amplitude, $\omega = 2\pi/T = 2\pi\nu$ is angular frequency, and $\phi_0$ is initial epoch.
  \item \textbf{Velocity:}
  \begin{equation}
    v(t) = \frac{dx}{dt} = A\omega \cos(\omega t + \phi_0) = \mathbf{\pm \omega \sqrt{A^2 - x^2}}
  \end{equation}
  $v_{\max} = \omega A$ occurs at mean position ($x=0$); $v=0$ at extreme positions ($x = \pm A$).
  \item \textbf{Acceleration:}
  \begin{equation}
    a(t) = \frac{dv}{dt} = -\omega^2 A \sin(\omega t + \phi_0) = \mathbf{-\omega^2 x}
  \end{equation}
  $a_{\max} = \omega^2 A$ occurs at extremes ($x = \pm A$); $a=0$ at mean position ($x=0$).
  \item \textbf{Phase Relationships:} Velocity leads displacement by $\pi/2$ radians ($90^\circ$); acceleration leads displacement by $\pi$ radians ($180^\circ$).
\end{itemize}
\end{conceptbox}

\begin{formulabox}[Energy Conservation in SHM: Space \& Time Averages]
\textbf{1. Instantaneous Energies as Functions of Position ($x$):}
\begin{itemize}
  \item \textbf{Potential Energy ($U$):} $U(x) = \int_0^x kx dx = \mathbf{\frac{1}{2} k x^2 = \frac{1}{2} m \omega^2 x^2}$ (assuming $U(0) = 0$).
  \item \textbf{Kinetic Energy ($K$):} $K(x) = \frac{1}{2} m v^2 = \mathbf{\frac{1}{2} m \omega^2 (A^2 - x^2)}$.
  \item \textbf{Total Mechanical Energy ($E$):}
  \begin{equation}
    \mathbf{E = K(x) + U(x) = \frac{1}{2} m \omega^2 A^2 = 2 \pi^2 m \nu^2 A^2 = \text{constant}}
  \end{equation}
  Total energy is completely independent of time $t$ and position $x$.
\end{itemize}
\textbf{2. Position where $K = U$:}
\begin{equation}
  \frac{1}{2}m\omega^2(A^2 - x^2) = \frac{1}{2}m\omega^2 x^2 \implies 2x^2 = A^2 \implies \mathbf{x = \pm \frac{A}{\sqrt{2}}}
\end{equation}
\textbf{3. Space Averages vs. Time Averages over One Cycle:}
\begin{itemize}
  \item \textbf{Time Averages ($\langle \dots \rangle_t = \frac{1}{T}\int_0^T \dots dt$):}
  \begin{equation}
    \mathbf{\langle K \rangle_t = \frac{1}{4} m \omega^2 A^2 = \frac{1}{2}E}, \qquad \mathbf{\langle U \rangle_t = \frac{1}{4} m \omega^2 A^2 = \frac{1}{2}E}
  \end{equation}
  \item \textbf{Space Averages ($\langle \dots \rangle_x = \frac{1}{2A}\int_{-A}^A \dots dx$):}
  \begin{equation}
    \mathbf{\langle K \rangle_x = \frac{2}{3}E}, \qquad \mathbf{\langle U \rangle_x = \frac{1}{3}E}
  \end{equation}
  \item \textbf{Frequency of Oscillations:} If displacement oscillates with frequency $\nu$, both kinetic energy $K(t)$ and potential energy $U(t)$ oscillate with \textbf{double frequency $2\nu$}.
\end{itemize}
\end{formulabox}

\begin{formulabox}[Mechanical Oscillators: Spring-Mass Systems \& Simple Pendulums]
\textbf{1. Spring-Block Oscillators:}
\begin{equation}
  T = 2\pi \sqrt{\frac{m}{k}}
\end{equation}
\begin{itemize}
  \item \textbf{Series Combination:} $\frac{1}{k_{\text{eq}}} = \frac{1}{k_1} + \frac{1}{k_2} \implies k_{\text{eq}} = \frac{k_1 k_2}{k_1 + k_2} \implies T = 2\pi\sqrt{\frac{m(k_1+k_2)}{k_1 k_2}}$.
  \item \textbf{Parallel Combination:} $k_{\text{eq}} = k_1 + k_2 \implies T = 2\pi\sqrt{\frac{m}{k_1 + k_2}}$.
  \item \textbf{Spring Cutting:} Spring constant is inversely proportional to length: $k \cdot L = \text{const}$. If cut into length ratio $m : n$, spring constants become $k_1 = \frac{m+n}{m}k$ and $k_2 = \frac{m+n}{n}k$.
  \item \textbf{Two-Body Oscillator Connected by Spring:} Reduced mass $\mu = \frac{m_1 m_2}{m_1 + m_2} \implies \mathbf{T = 2\pi\sqrt{\frac{\mu}{k}}}$.
\end{itemize}
\textbf{2. Simple Pendulum Derivation:}
A bob of mass $m$ suspended by a light string of length $L$. For angular displacement $\theta$:
\begin{equation}
  \tau_{\text{restoring}} = -m g L \sin\theta \approx -m g L \theta \quad (\text{for paraxial small angles } \theta \ll 1)
\end{equation}
Using $\tau = I \alpha = (m L^2) \frac{d^2\theta}{dt^2}$:
\begin{equation}
  m L^2 \frac{d^2\theta}{dt^2} + m g L \theta = 0 \implies \frac{d^2\theta}{dt^2} + \frac{g}{L}\theta = 0 \implies \mathbf{\omega = \sqrt{\frac{g}{L}}}, \quad \mathbf{T = 2\pi\sqrt{\frac{L}{g}}}
\end{equation}
\textit{Accelerating Frames \& Effective Gravity:}
\begin{itemize}
  \item Elevator accelerating upwards with $a$: $g_{\text{eff}} = g + a \implies T = 2\pi\sqrt{\frac{L}{g+a}}$ (period decreases).
  \item Elevator accelerating downwards with $a$: $g_{\text{eff}} = g - a \implies T = 2\pi\sqrt{\frac{L}{g-a}}$ (period increases; free fall $a=g \implies T \to \infty$).
  \item Horizontally accelerating vehicle: $g_{\text{eff}} = \sqrt{g^2 + a^2} \implies T = 2\pi\sqrt{\frac{L}{\sqrt{g^2 + a^2}}}$.
  \item Pendulum in liquid of density $\sigma$ (bob density $\rho$): $g_{\text{eff}} = g\left(1 - \frac{\sigma}{\rho}\right) \implies T = 2\pi\sqrt{\frac{L}{g(1 - \sigma/\rho)}}$.
\end{itemize}
\textbf{3. Seconds Pendulum:} A pendulum whose time period is exactly $T = 2.0\,\mathrm{s}$ ($1\,\mathrm{s}$ tick each half-swing):
\begin{equation}
  L = \frac{g T^2}{4\pi^2} = \frac{9.8 \times 4}{4\pi^2} \approx \mathbf{0.993\,\mathrm{m} \approx 1.0\,\mathrm{m}}
\end{equation}
\end{formulabox}

\begin{usecasebox}[Phase Circle (Phasor) Technique for Instantaneous Transit Times]
To compute the transit time between any two positions $x_1$ and $x_2$ in SHM, project the motion onto a reference circle of radius $A$ rotating counterclockwise at angular velocity $\omega$:
\begin{equation}
  x = A \sin\theta, \quad \theta = \omega t \implies t = \frac{\Delta\theta}{\omega} = \frac{\Delta\theta}{2\pi} T
\end{equation}
\textit{High-Yield Benchmarks (Starting from Mean Position $x=0$):}
\begin{itemize}
  \item From $x = 0$ to $x = \frac{A}{2}$: $\theta_1 = 0 \to \theta_2 = 30^\circ = \frac{\pi}{6} \implies \mathbf{t = \frac{T}{12}}$.
  \item From $x = \frac{A}{2}$ to $x = A$: $\Delta\theta = 90^\circ - 30^\circ = 60^\circ = \frac{\pi}{3} \implies \mathbf{t = \frac{T}{6}}$.
  \item Notice: Traveling the outer half of amplitude takes \textbf{twice as long} as the inner half ($T/6$ vs $T/12$) due to deceleration.
  \item From $x = 0$ to $x = \frac{A}{\sqrt{2}}$: $\theta = 45^\circ = \frac{\pi}{4} \implies \mathbf{t = \frac{T}{8}}$.
\end{itemize}
\end{usecasebox}

\section{Wave Motion, Stretched Strings \& Organ Pipes}

\begin{conceptbox}[Kinematics of 1D Harmonic Progressive Waves]
A 1D transverse or longitudinal sinusoidal traveling wave propagating along the $+x$ direction:
\begin{equation}
  \mathbf{y(x,t) = A \sin(k x - \omega t + \phi) = A \sin\left[2\pi\left(\frac{x}{\lambda} - \nu t\right) + \phi\right]}
\end{equation}
where $k = \frac{2\pi}{\lambda}$ is the angular wave number, $\omega = 2\pi\nu$ is angular frequency.
\begin{itemize}
  \item \textbf{Wave Speed ($v$):} Speed of propagation of constant phase profile:
  \begin{equation}
    k x - \omega t = \text{const} \implies k \frac{dx}{dt} - \omega = 0 \implies \mathbf{v = \frac{\omega}{k} = \nu \lambda}
  \end{equation}
  \item \textbf{Particle Velocity ($v_p$) vs. Wave Velocity ($v$):}
  \begin{equation}
    v_p = \frac{\partial y}{\partial t} = -\omega A \cos(k x - \omega t), \qquad \frac{\partial y}{\partial x} = k A \cos(k x - \omega t)
  \end{equation}
  \begin{equation}
    \mathbf{v_p = -v \left(\frac{\partial y}{\partial x}\right) = -v \times (\text{Wave profile slope})}
  \end{equation}
  If wave moves right ($v>0$) and slope is positive ($\partial y/\partial x > 0$), particle moves downwards ($v_p < 0$).
  \item \textbf{Particle Acceleration:} $a_p = \frac{\partial^2 y}{\partial t^2} = -\omega^2 y$.
\end{itemize}
\end{conceptbox}

\begin{formulabox}[Speeds of Transverse Waves on Strings \& Sound in Gases]
\textbf{1. Transverse Wave on Stretched String:}
\begin{equation}
  \mathbf{v = \sqrt{\frac{T}{\mu}}}
\end{equation}
where $T$ is tension ($\mathrm{N}$), and $\mu = \frac{M}{L} = \rho A$ is linear mass density ($\mathrm{kg/m}$).
\textit{Derivation via Centripetal Equilibrium of Arc Element:}
Consider a small curved element of string of length $\Delta l = R \Delta\theta$ moving with wave speed $v$ along a circular crest of radius $R$. Mass of element $\Delta m = \mu \Delta l = \mu R \Delta\theta$.
The inward restoring component of tension:
\begin{equation}
  F_{\text{inward}} = 2 T \sin\left(\frac{\Delta\theta}{2}\right) \approx T \Delta\theta
\end{equation}
Equating to centripetal force:
\begin{equation}
  T \Delta\theta = \frac{\Delta m \cdot v^2}{R} = \frac{(\mu R \Delta\theta) v^2}{R} = \mu v^2 \Delta\theta \implies \mathbf{v = \sqrt{\frac{T}{\mu}}}
\end{equation}
\textbf{2. Speed of Longitudinal Waves (Sound) in Fluids \& Gases:}
\begin{itemize}
  \item General relation: $v = \sqrt{\frac{B}{\rho}}$, where $B$ is bulk modulus of elasticity.
  \item \textbf{Newton's Isothermal Formula (Erroneous):} Newton assumed sound compressions/rarefactions occur isothermally ($B = P$).
  \begin{equation}
    v_{\text{Newton}} = \sqrt{\frac{P}{\rho}} \approx \sqrt{\frac{1.013 \times 10^5}{1.293}} \approx 280\,\mathrm{m/s} \quad (\approx 15\% \text{ below experimental } 332\,\mathrm{m/s})
  \end{equation}
  \item \textbf{Laplace's Adiabatic Correction (Accurate):} Compressions and rarefactions occur so rapidly that heat exchange is zero ($dQ=0 \implies B_{\text{adiabatic}} = \gamma P$):
  \begin{equation}
    \mathbf{v = \sqrt{\frac{\gamma P}{\rho}} = \sqrt{\frac{\gamma R T}{M_{\text{molar}}}}}
  \end{equation}
  For air ($\gamma \approx 1.4, M \approx 28.8 \times 10^{-3}\,\mathrm{kg/mol}$ at $273\,\mathrm{K}$): $v = 332\,\mathrm{m/s}$ (perfect agreement).
  \item \textbf{Dependence on Environmental Factors:}
  \begin{itemize}
    \item \textit{Pressure:} Invariant of pressure at constant temperature (since $P/\rho = RT/M = \text{const}$).
    \item \textit{Temperature:} $v \propto \sqrt{T} \implies v_t \approx v_0 + 0.61 t$ (for small temperatures $t^\circ\mathrm{C}$).
    \item \textit{Humidity:} Moist air is less dense than dry air ($\rho_{\text{moist}} < \rho_{\text{dry}}$), so speed of sound \textbf{increases} with humidity.
  \end{itemize}
\end{itemize}
\end{formulabox}

\begin{formulabox}[Standing Waves: Strings \& Organ Pipes with End Corrections]
\textbf{1. Superposition \& Standing Wave Equation on String Fixed at Both Ends:}
Two identical waves traveling in opposite directions: $y_1 = A\sin(kx - \omega t)$ and $y_2 = A\sin(kx + \omega t)$.
\begin{equation}
  y_{\text{net}} = y_1 + y_2 = \mathbf{[2A \sin(kx)] \cos(\omega t)}
\end{equation}
\begin{itemize}
  \item \textbf{Nodes ($y=0$):} $\sin(kx) = 0 \implies kx = n\pi \implies \mathbf{x = \frac{n\lambda}{2}}$ ($n = 0, 1, 2, \dots$).
  \item \textbf{Antinodes ($|y|=2A$):} $|\sin(kx)| = 1 \implies \mathbf{x = (2n-1)\frac{\lambda}{4}}$.
  \item Boundary condition: nodes at $x = 0$ and $x = L \implies L = n\frac{\lambda}{2} \implies \mathbf{\nu_n = n \left(\frac{v}{2L}\right) = n \nu_1}$ (\textbf{All harmonics present}: $1 : 2 : 3 : \dots$).
\end{itemize}
\textbf{2. Closed Organ Pipe (One End Closed, One End Open):}
Boundary condition: Node at closed end, Antinode at open end. Length $L = (2n-1)\frac{\lambda}{4}$:
\begin{equation}
  \mathbf{\nu_n = (2n-1)\frac{v}{4L} = (2n-1)\nu_1} \qquad \text{where fundamental } \nu_1 = \frac{v}{4L}
\end{equation}
\textbf{Only odd harmonics exist} in a closed organ pipe ($\nu_1 : \nu_3 : \nu_5 = 1 : 3 : 5 : \dots$). No even harmonics.
\textbf{3. Open Organ Pipe (Both Ends Open):}
Boundary condition: Antinodes at both open ends. Length $L = n\frac{\lambda}{2}$:
\begin{equation}
  \mathbf{\nu_n = n \frac{v}{2L} = n \nu_1} \qquad \text{where fundamental } \nu_1 = \frac{v}{2L}
\end{equation}
\textbf{All harmonics exist} ($\nu_1 : \nu_2 : \nu_3 = 1 : 2 : 3 : \dots$). Tone is richer/sweeter than closed pipe.
\textbf{4. End Correction ($e$):}
Antinode forms slightly outside the physical open tube boundary due to finite acoustic inertia of moving air:
\begin{equation}
  e \approx 0.6 r \quad (\text{where } r \text{ is tube inner radius})
\end{equation}
\begin{itemize}
  \item For Closed Pipe: $L_{\text{eff}} = L + e = L + 0.6r$.
  \item For Open Pipe: $L_{\text{eff}} = L + 2e = L + 1.2r$.
  \item \textbf{Resonance Tube Method to eliminate $e$:}
  Using 1st resonance length $l_1$ and 2nd resonance length $l_2$ for a tuning fork of frequency $\nu$:
  \begin{equation}
    l_1 + e = \frac{\lambda}{4}, \qquad l_2 + e = \frac{3\lambda}{4} \implies l_2 - l_1 = \frac{\lambda}{2} \implies \mathbf{\lambda = 2(l_2 - l_1)}
  \end{equation}
  \begin{equation}
    \mathbf{v = 2\nu(l_2 - l_1)}, \qquad \mathbf{e = \frac{l_2 - 3l_1}{2}}
  \end{equation}
\end{itemize}
\end{formulabox}

\begin{formulabox}[Acoustic Beats: Analytical Derivation]
Consider superposition of two harmonic sound waves of equal amplitude $A$ and slightly differing frequencies $\nu_1, \nu_2$ ($\omega_1 \approx \omega_2$):
\begin{equation}
  y_1 = A\sin(\omega_1 t), \qquad y_2 = A\sin(\omega_2 t)
\end{equation}
\begin{equation}
  y_{\text{net}} = y_1 + y_2 = 2A \cos\left(\frac{\omega_1 - \omega_2}{2}t\right) \sin\left(\frac{\omega_1 + \omega_2}{2}t\right) = A_{\text{mod}}(t) \sin(\omega_{\text{avg}} t)
\end{equation}
The modulated amplitude is $A_{\text{mod}}(t) = 2A \cos(\pi(\nu_1 - \nu_2)t)$.
Intensity $I \propto A_{\text{mod}}^2 \propto \cos^2(\pi(\nu_1 - \nu_2)t) = \frac{1 + \cos(2\pi(\nu_1 - \nu_2)t)}{2}$.
Maximum intensity occurs whenever $2\pi(\nu_1 - \nu_2)t = n\pi \implies t = \frac{n}{|\nu_1 - \nu_2|}$.
Time interval between consecutive waxing/waning: $\Delta t = \frac{1}{|\nu_1 - \nu_2|}$.
\begin{equation}
  \mathbf{\nu_{\text{beat}} = |\nu_1 - \nu_2|}
\end{equation}
\textit{Tuning Fork Rules:}
\begin{itemize}
  \item Loading fork with wax increases inertia $\implies \nu$ decreases.
  \item Filing fork prongs decreases mass $\implies \nu$ increases.
\end{itemize}
\end{formulabox}

\begin{trapbox}[Syllabus Boundary Alert: Complete Deletion of Doppler Effect]
\textbf{NMC REVISION WARNING FOR NEET UG:}
Under the revised NMC curriculum for NEET UG 2027:
\begin{itemize}
  \item \textbf{STRICTLY DELETED:} The \textbf{Doppler Effect in Sound Waves} (apparent frequency shifts $f' = f\frac{v \pm v_0}{v \mp v_s}$, moving observers, moving sources, echo reflection from cliffs, acoustic shockwaves) has been \textbf{completely eliminated} from the syllabus.
  \item \textbf{EXAM ORIENTATION:} Candidates must not solve Doppler frequency formulas. Instead, master Standing Waves, Harmonics in Organ Pipes, Resonance Tube end corrections, and Beats.
\end{itemize}
\end{trapbox}

\section{Master High-Yield Problem Archetypes \& Worked Solutions}

\begin{usecasebox}[Archetype 1: Ice-Water-Steam Calorimetry Mixing Equilibrium]
\textbf{Problem:} $10\,\mathrm{g}$ of ice at $-10^\circ\mathrm{C}$ is mixed with $10\,\mathrm{g}$ of steam at $100^\circ\mathrm{C}$ in an insulated calorimeter of negligible heat capacity. Determine the final equilibrium temperature and the constitution of the mixture.
\textbf{Step-by-Step Solution:}
\begin{enumerate}
  \item Heat required to bring ice from $-10^\circ\mathrm{C}$ to liquid water at $100^\circ\mathrm{C}$:
  \begin{itemize}
    \item Warm ice to $0^\circ\mathrm{C}$: $Q_1 = m_{\text{ice}} s_{\text{ice}} \Delta T = 10 \times 0.5 \times 10 = 50\,\mathrm{cal}$.
    \item Melt ice at $0^\circ\mathrm{C}$: $Q_2 = m_{\text{ice}} L_f = 10 \times 80 = 800\,\mathrm{cal}$.
    \item Warm water to $100^\circ\mathrm{C}$: $Q_3 = m_{\text{ice}} s_w \Delta T = 10 \times 1.0 \times 100 = 1000\,\mathrm{cal}$.
    \item Total heat absorbed to reach $100^\circ\mathrm{C}$ water: $Q_{\text{absorb}} = 50 + 800 + 1000 = \mathbf{1850\,\mathrm{cal}}$.
  \end{itemize}
  \item Total latent heat available from steam condensing at $100^\circ\mathrm{C}$:
  \begin{equation}
    Q_{\text{condense, total}} = m_{\text{steam}} L_v = 10 \times 540 = \mathbf{5400\,\mathrm{cal}}
  \end{equation}
  \item Notice that $Q_{\text{condense, total}} = 5400\,\mathrm{cal} > Q_{\text{absorb}} = 1850\,\mathrm{cal}$.
  Therefore, not all steam condenses! The mixture reaches dynamic equilibrium at \textbf{$T_{\text{eq}} = 100^\circ\mathrm{C}$}.
  \item Mass of steam condensed to supply $1850\,\mathrm{cal}$:
  \begin{equation}
    m_{\text{condensed}} = \frac{1850\,\mathrm{cal}}{540\,\mathrm{cal/g}} \approx 3.43\,\mathrm{g}
  \end{equation}
  \item Surviving steam: $m_{\text{steam, remaining}} = 10 - 3.43 = \mathbf{6.57\,\mathrm{g}}$.
  \item Total water at $100^\circ\mathrm{C}$: $m_{\text{water}} = 10 + 3.43 = \mathbf{13.43\,\mathrm{g}}$.
\end{enumerate}
\end{usecasebox}

\begin{usecasebox}[Archetype 2: Mixture of Distinct Polyatomic \& Monoatomic Gases]
\textbf{Problem:} $2$ moles of Helium (monoatomic, $\gamma_1 = 5/3$) are mixed with $3$ moles of Oxygen (diatomic, rigid, $\gamma_2 = 7/5$). Calculate $C_{v,\text{mix}}$, $C_{p,\text{mix}}$, and the adiabatic exponent $\gamma_{\text{mix}}$.
\textbf{Step-by-Step Solution:}
\begin{enumerate}
  \item Molar heat capacities of individual gases:
  \begin{itemize}
    \item Helium (monoatomic): $C_{v1} = \frac{3}{2}R$, $C_{p1} = \frac{5}{2}R$.
    \item Oxygen (diatomic): $C_{v2} = \frac{5}{2}R$, $C_{p2} = \frac{7}{2}R$.
  \end{itemize}
  \item Total moles: $n_{\text{total}} = 2 + 3 = 5\,\mathrm{mol}$.
  \item Equivalent $C_{v,\text{mix}}$:
  \begin{equation}
    C_{v,\text{mix}} = \frac{n_1 C_{v1} + n_2 C_{v2}}{n_1 + n_2} = \frac{2 \left(\frac{3}{2}R\right) + 3 \left(\frac{5}{2}R\right)}{5} = \frac{3R + 7.5R}{5} = \frac{10.5R}{5} = \mathbf{2.1 R} = \frac{21}{10}R
  \end{equation}
  \item Equivalent $C_{p,\text{mix}}$:
  \begin{equation}
    C_{p,\text{mix}} = C_{v,\text{mix}} + R = 2.1 R + R = \mathbf{3.1 R} = \frac{31}{10}R
  \end{equation}
  \item Equivalent adiabatic index $\gamma_{\text{mix}}$:
  \begin{equation}
    \gamma_{\text{mix}} = \frac{C_{p,\text{mix}}}{C_{v,\text{mix}}} = \frac{3.1 R}{2.1 R} = \frac{31}{21} \approx \mathbf{1.476}
  \end{equation}
\end{enumerate}
\end{usecasebox}

\begin{usecasebox}[Archetype 3: Resonance Tube End Correction \& Speed of Sound]
\textbf{Problem:} A tuning fork of frequency $\nu = 512\,\mathrm{Hz}$ excites the first two resonances in a resonance water tube at lengths $l_1 = 15.2\,\mathrm{cm}$ and $l_2 = 48.0\,\mathrm{cm}$. Compute the speed of sound in air and the inner pipe radius $r$.
\textbf{Step-by-Step Solution:}
\begin{enumerate}
  \item The fundamental difference between resonances isolates the half-wavelength:
  \begin{equation}
    \frac{\lambda}{2} = l_2 - l_1 = 48.0\,\mathrm{cm} - 15.2\,\mathrm{cm} = 32.8\,\mathrm{cm} \implies \lambda = 65.6\,\mathrm{cm} = 0.656\,\mathrm{m}
  \end{equation}
  \item Speed of sound in air:
  \begin{equation}
    v = \nu \lambda = 512\,\mathrm{s}^{-1} \times 0.656\,\mathrm{m} = \mathbf{335.87\,\mathrm{m/s}}
  \end{equation}
  \item End correction $e$:
  \begin{equation}
    e = \frac{l_2 - 3l_1}{2} = \frac{48.0 - 3(15.2)}{2} = \frac{48.0 - 45.6}{2} = \frac{2.4}{2} = \mathbf{1.2\,\mathrm{cm}}
  \end{equation}
  \item Inner radius $r$ of the tube using $e \approx 0.6 r$:
  \begin{equation}
    r = \frac{e}{0.6} = \frac{1.2\,\mathrm{cm}}{0.6} = \mathbf{2.0\,\mathrm{cm}}
  \end{equation}
\end{enumerate}
\end{usecasebox}

% =========================================================================
% PART III: ELECTRODYNAMICS, MAGNETISM & ALTERNATING CURRENTS
% =========================================================================
\part{Electrodynamics, Magnetism \& Alternating Currents}

\chapter{Electrostatics, Capacitance, Current Electricity \& Electromagnetism}

Electrodynamics represents over 25\% of the NEET Physics paper. In this chapter, we build from Coulomb's law and continuous charge distributions to Gauss's flux theorem, dielectric physics and capacitance networks, Ohm's drift dynamics and DC circuits, the magnetic effects of moving charges via Biot-Savart and Ampère laws, electromagnetic induction, LCR resonance in AC circuits, and Maxwell's electromagnetic wave spectrum.

\begin{mentorbox}[A Personal Note from Satya: The Art of Knowing When to Stop (The 45-Second Rule)]
\textbf{Before we enter the world of circuits, flux, and resonance, let me share a crucial truth about exam psychology.}

NEET UG is not an IQ contest; it is a rapid-fire elimination game played under intense time pressure (200 questions in 200 minutes). The single most common reason why dedicated students score 520 instead of 650+ is not that they didn't know the theory---it's that \textit{they didn't know when to put their pen down.}

Here is the cognitive trap: You encounter a tricky question---say, an asymmetrical resistor cube or a magnetic force on an oblique helical trajectory. You spend 60 seconds setting up the equations. The algebra gets tangled. You feel stuck. But your ego whispers: \textit{``I've already spent two minutes on this! If I skip now, those two minutes are wasted!''} So you stay for another two minutes. You end up calculating in panic, and you mark an option. 

What just happened?
\begin{itemize}[leftmargin=*]
  \item You probably marked a wrong option under panic and got $-1$ instead of $+4$ (a devastating 5-mark swing).
  \item Even worse: you consumed 4 full minutes on one question. In that exact time, you could have read and solved \textbf{six} direct, guaranteed NCERT questions in Biology or Inorganic Chemistry!
\end{itemize}

\textbf{The 45-Second Golden Rule:}
Read the problem. If you cannot visualize the complete, crystal-clear solution path within 45 seconds, \textbf{circle the question number on your test booklet, put your pen down, and move on immediately.} Do not hesitate. Do not feel guilty. 

NEET top ranks belong to students who sweep through the paper collecting all the easy 4-pointers first in Round 1, and return to battle the complex derivations in Round 2. Knowing when to stop is not surrender; it is tactical mastery.
\end{mentorbox}

\section{Electrostatics, Coulomb's Law \& Continuous Charge Distributions}

\begin{conceptbox}[Foundations of Electric Charge \& Coulomb's Law in Vector Form]
\textbf{1. Fundamental Charge Properties:}
\begin{itemize}
  \item \textbf{Quantization of Charge:} Charge exists strictly as integer multiples of the elementary charge $e \approx 1.602 \times 10^{-19}\,\mathrm{C}$:
  \begin{equation}
    q = \pm n e \quad (n \in \mathbb{N})
  \end{equation}
  \item \textbf{Conservation of Charge:} The algebraic sum of fundamental charges within an isolated physical system remains invariant over time.
  \item \textbf{Invariance of Charge:} The magnitude of electric charge is independent of the reference frame speed ($q_{\text{rest}} = q_{\text{motion}}$), unlike relativistic mass.
\end{itemize}
\textbf{2. Coulomb's Inverse-Square Law in Vector Form:}
The electrostatic force exerted by point charge $q_1$ at position vector $\mathbf{r}_1$ on charge $q_2$ at $\mathbf{r}_2$:
\begin{equation}
  \mathbf{F}_{12} = \frac{1}{4\pi\varepsilon_0 \varepsilon_r} \frac{q_1 q_2}{|\mathbf{r}_2 - \mathbf{r}_1|^3} (\mathbf{r}_2 - \mathbf{r}_1) = \frac{1}{4\pi\varepsilon_0 \varepsilon_r} \frac{q_1 q_2}{r_{12}^2} \hat{\mathbf{r}}_{12}
\end{equation}
where $\varepsilon_0 = 8.854 \times 10^{-12}\,\mathrm{C}^2\cdot\mathrm{N}^{-1}\cdot\mathrm{m}^{-2}$ is the permittivity of free space, and $\frac{1}{4\pi\varepsilon_0} \approx 8.988 \times 10^9\,\mathrm{N}\cdot\mathrm{m}^2\cdot\mathrm{C}^{-2}$.
In a dielectric medium of relative permittivity (dielectric constant) $\kappa = \varepsilon_r > 1$:
\begin{equation}
  F_{\text{medium}} = \frac{F_{\text{vacuum}}}{\kappa}
\end{equation}
\textbf{3. Superposition Principle:}
The resultant force acting on any test charge $q_0$ due to a collection of discrete charges $q_1, q_2, \dots, q_n$ is the vector sum:
\begin{equation}
  \mathbf{F}_{\text{net}} = \sum_{i=1}^n \mathbf{F}_{i0} = \frac{q_0}{4\pi\varepsilon_0} \sum_{i=1}^n \frac{q_i}{|\mathbf{r}_0 - \mathbf{r}_i|^3} (\mathbf{r}_0 - \mathbf{r}_i)
\end{equation}
\end{conceptbox}

\begin{formulabox}[Electric Field \& Potential of Continuous Charge Geometries: Ring \& Disk]
\textbf{1. Electric Field Definition:}
$\mathbf{E} = \lim_{q_0 \to 0} \frac{\mathbf{F}}{q_0} = -\boldsymbol{\nabla} V$.
\textbf{2. Uniformly Charged Thin Circular Ring on its Symmetry Axis:}
Consider a ring of total charge $Q$ and radius $R$ positioned in the $y-z$ plane with symmetry axis along $x$.
\begin{itemize}
  \item By radial symmetry, perpendicular field components cancel: $dE_\perp = 0$.
  \item Axial field element: $dE_x = dE \cos\theta = \frac{1}{4\pi\varepsilon_0} \frac{dq}{x^2 + R^2} \frac{x}{\sqrt{x^2 + R^2}}$.
  \item Total integrated axial electric field:
  \begin{equation}
    \mathbf{E(x) = \frac{1}{4\pi\varepsilon_0} \frac{Q x}{(x^2 + R^2)^{3/2}} \hat{\mathbf{i}}}
  \end{equation}
  \item \textbf{At Ring Center ($x=0$):} $E(0) = 0$.
  \item \textbf{Far-Field Regime ($x \gg R$):} $E(x) \approx \frac{1}{4\pi\varepsilon_0}\frac{Q}{x^2}$ (acts as point charge).
  \item \textbf{Location of Maximum Electric Field:}
  \begin{equation}
    \frac{dE}{dx} = \frac{Q}{4\pi\varepsilon_0} \frac{(x^2 + R^2)^{3/2} - x \cdot \frac{3}{2}(x^2 + R^2)^{1/2}(2x)}{(x^2 + R^2)^3} = 0 \implies x^2 + R^2 - 3x^2 = 0 \implies \mathbf{x = \pm \frac{R}{\sqrt{2}}}
  \end{equation}
  \begin{equation}
    \mathbf{E_{\max} = \frac{1}{4\pi\varepsilon_0} \frac{2 Q}{3\sqrt{3} R^2}}
  \end{equation}
  \item \textbf{Electrostatic Potential on Ring Axis:}
  \begin{equation}
    \mathbf{V(x) = \frac{1}{4\pi\varepsilon_0} \frac{Q}{\sqrt{x^2 + R^2}}} \implies V_{\text{center}} = \frac{Q}{4\pi\varepsilon_0 R}
  \end{equation}
\end{itemize}
\textbf{3. Uniformly Charged Flat Circular Disk on its Axis:}
Integrating concentric rings of radius $r$ and charge $dq = \sigma (2\pi r dr)$ from $r=0$ to $R$:
\begin{equation}
  \mathbf{E(x) = \frac{\sigma}{2\varepsilon_0} \left(1 - \frac{x}{\sqrt{x^2 + R^2}}\right) \hat{\mathbf{i}}}, \qquad \mathbf{V(x) = \frac{\sigma}{2\varepsilon_0} \left(\sqrt{x^2 + R^2} - x\right)}
\end{equation}
For $R \to \infty$ (infinite sheet), $E = \frac{\sigma}{2\varepsilon_0}$ (constant, independent of $x$).
\end{formulabox}

\begin{formulabox}[Electric Dipole Dynamics: Field, Torque \& Potential Energy]
An electric dipole consists of two equal and opposite point charges $+q$ and $-q$ separated by a displacement vector $2\mathbf{a}$ pointing from $-q$ to $+q$:
\begin{equation}
  \mathbf{p} = q(2\mathbf{a}) \quad [\mathrm{C}\cdot\mathrm{m}]
\end{equation}
\textbf{1. Electric Field of Dipole in Coordinate Space ($r \gg a$):}
\begin{itemize}
  \item \textbf{Axial Point (End-on Position, $\theta = 0^\circ$):}
  \begin{equation}
    \mathbf{E_{\text{axial}} = \frac{1}{4\pi\varepsilon_0} \frac{2\mathbf{p}}{r^3}} \quad (\text{Collinear and parallel to } \mathbf{p})
  \end{equation}
  \item \textbf{Equatorial Point (Broadside-on Position, $\theta = 90^\circ$):}
  \begin{equation}
    \mathbf{E_{\text{eq}} = -\frac{1}{4\pi\varepsilon_0} \frac{\mathbf{p}}{r^3}} \quad (\text{Antiparallel to } \mathbf{p})
  \end{equation}
  \item \textbf{General Coordinate Point $(r, \theta)$:}
  \begin{equation}
    \mathbf{|\mathbf{E}| = \frac{1}{4\pi\varepsilon_0} \frac{p}{r^3} \sqrt{1 + 3\cos^2\theta}}, \qquad \tan\alpha = \frac{1}{2}\tan\theta
  \end{equation}
  where $\alpha$ is the angle between $\mathbf{E}$ and position vector $\mathbf{r}$.
  \item \textbf{Electrostatic Potential at $(r, \theta)$:}
  \begin{equation}
    \mathbf{V(r, \theta) = \frac{1}{4\pi\varepsilon_0} \frac{\mathbf{p}\cdot\hat{\mathbf{r}}}{r^2} = \frac{1}{4\pi\varepsilon_0} \frac{p\cos\theta}{r^2}}
  \end{equation}
  Notice: $V_{\text{eq}} = V(r, 90^\circ) = 0$ everywhere along the equatorial bisector plane.
\end{itemize}
\textbf{2. External Uniform Electric Field Interactions:}
\begin{itemize}
  \item Net Force: $\mathbf{F_{\text{net}} = q\mathbf{E} + (-q)\mathbf{E} = 0}$ (zero translational acceleration).
  \item Torque:
  \begin{equation}
    \mathbf{\boldsymbol{\tau} = \mathbf{p} \times \mathbf{E}} \implies |\boldsymbol{\tau}| = p E \sin\theta
  \end{equation}
  Maximum torque occurs at $\theta = 90^\circ$ ($\tau_{\max} = pE$).
  \item Potential Energy of Dipole in External Field:
  \begin{equation}
    \mathbf{U = -\mathbf{p}\cdot\mathbf{E} = -p E \cos\theta}
  \end{equation}
  \begin{itemize}
    \item \textbf{Stable Equilibrium ($\theta = 0^\circ$):} $\mathbf{p} \parallel \mathbf{E} \implies \boldsymbol{\tau} = 0$, $U_{\min} = -pE$.
    \item \textbf{Unstable Equilibrium ($\theta = 180^\circ$):} $\mathbf{p} \parallel -\mathbf{E} \implies \boldsymbol{\tau} = 0$, $U_{\max} = +pE$.
  \end{itemize}
  \item Work done by external agent in rotating dipole from $\theta_1$ to $\theta_2$:
  \begin{equation}
    \mathbf{W_{\text{ext}} = \Delta U = p E (\cos\theta_1 - \cos\theta_2)}
  \end{equation}
\end{itemize}
\end{formulabox}

\section{Gauss's Law, Electric Flux \& Boundary Value Applications}

\begin{conceptbox}[Gauss's Law Formulation \& Electric Flux]
\textbf{1. Electric Flux ($\Phi_E$):}
Surface integral of the electric field vector through an oriented surface element $d\mathbf{A} = \hat{\mathbf{n}} dA$:
\begin{equation}
  \Phi_E = \iint_S \mathbf{E}\cdot d\mathbf{A} \quad [\mathrm{N}\cdot\mathrm{m}^2\cdot\mathrm{C}^{-1} \text{ or } \mathrm{V}\cdot\mathrm{m}]
\end{equation}
\textbf{2. Gauss's Law (Maxwell's First Equation):}
The net outward electric flux through any closed Gaussian surface enclosing net charge $q_{\text{enclosed}}$:
\begin{equation}
  \mathbf{\oint_{\partial V} \mathbf{E}\cdot d\mathbf{A} = \frac{q_{\text{enclosed}}}{\varepsilon_0}}
\end{equation}
\textit{Key Physical Insights:}
\begin{itemize}
  \item Net flux is completely independent of the shape, size, or geometry of the Gaussian boundary.
  \item Charges residing outside the closed surface contribute zero net flux through the surface (entering flux equals exiting flux).
  \item Electric field $\mathbf{E}$ on the left-hand side is the \textbf{net resultant field} produced by all charges (both inside and outside).
\end{itemize}
\end{conceptbox}

\begin{formulabox}[Master Applications of Gauss's Law: Exact Field Derivations]
\textbf{1. Infinitely Long Line Charge of Uniform Density $\lambda$ ($\mathrm{C/m}$):}
Using a coaxial Gaussian cylinder of radius $r$ and length $L$:
\begin{equation}
  \oint \mathbf{E}\cdot d\mathbf{A} = E(2\pi r L) = \frac{\lambda L}{\varepsilon_0} \implies \mathbf{E(r) = \frac{\lambda}{2\pi\varepsilon_0 r}\hat{\mathbf{r}}}
\end{equation}
\textbf{2. Uniformly Charged Infinite Thin Flat Sheet of Density $\sigma$ ($\mathrm{C/m^2}$):}
Using a Gaussian pillbox of cross-sectional area $A$ piercing normal to the sheet:
\begin{equation}
  \oint \mathbf{E}\cdot d\mathbf{A} = 2 E A = \frac{\sigma A}{\varepsilon_0} \implies \mathbf{E = \frac{\sigma}{2\varepsilon_0}\hat{\mathbf{n}}} \quad (\text{Independent of distance})
\end{equation}
\textbf{3. Surface of a Charged Conductor:}
Inside the conductor, $E=0$. Flux exits only from the outer face:
\begin{equation}
  E A = \frac{\sigma A}{\varepsilon_0} \implies \mathbf{E_{\text{conductor}} = \frac{\sigma}{\varepsilon_0}\hat{\mathbf{n}}} \quad (\text{Twice that of a non-conducting sheet})
\end{equation}
\textbf{4. Thin Spherical Shell of Radius $R$ and Total Charge $Q$:}
\begin{itemize}
  \item \textbf{Interior Region ($r < R$):} Gaussian surface encloses $q_{\text{encl}} = 0 \implies \mathbf{E_{\text{in}} = 0}$.
  \item Potential is constant throughout interior: $\mathbf{V_{\text{in}} = \frac{1}{4\pi\varepsilon_0}\frac{Q}{R}}$.
  \item \textbf{Exterior Region ($r \ge R$):} Gaussian sphere encloses $Q \implies \mathbf{E_{\text{out}}(r) = \frac{1}{4\pi\varepsilon_0}\frac{Q}{r^2}\hat{\mathbf{r}}}$.
  \item Exterior potential: $\mathbf{V_{\text{out}}(r) = \frac{1}{4\pi\varepsilon_0}\frac{Q}{r}}$.
\end{itemize}
\textbf{5. Uniformly Charged Non-Conducting Solid Sphere of Radius $R$ and Total Charge $Q$:}
\begin{itemize}
  \item Volume charge density: $\rho = \frac{Q}{\frac{4}{3}\pi R^3}$.
  \item \textbf{Interior Region ($r \le R$):} $q_{\text{encl}} = \rho \left(\frac{4}{3}\pi r^3\right) = Q \frac{r^3}{R^3}$:
  \begin{equation}
    E(r) \cdot 4\pi r^2 = \frac{Q r^3}{\varepsilon_0 R^3} \implies \mathbf{E_{\text{in}}(r) = \frac{1}{4\pi\varepsilon_0}\frac{Q r}{R^3}\hat{\mathbf{r}} = \frac{\rho r}{3\varepsilon_0}\hat{\mathbf{r}}}
  \end{equation}
  Electric field increases linearly with radial distance $r$ from center.
  \item \textbf{Interior Potential:} $V_{\text{in}}(r) = -\int_\infty^R E_{\text{out}} dr - \int_R^r E_{\text{in}} dr$:
  \begin{equation}
    \mathbf{V_{\text{in}}(r) = \frac{Q}{4\pi\varepsilon_0}\left[\frac{3R^2 - r^2}{2R^3}\right]} \implies \mathbf{V_{\text{center}} = \frac{3}{2} V_{\text{surface}} = 1.5 V_{\text{surface}}}
  \end{equation}
  \item \textbf{Exterior Region ($r \ge R$):} Identical to point charge:
  \begin{equation}
    \mathbf{E_{\text{out}}(r) = \frac{1}{4\pi\varepsilon_0}\frac{Q}{r^2}\hat{\mathbf{r}}}, \qquad \mathbf{V_{\text{out}}(r) = \frac{1}{4\pi\varepsilon_0}\frac{Q}{r}}
  \end{equation}
\end{itemize}
\end{formulabox}

\begin{ncertbox}[NCERT Extraction: Electrostatic Shielding \& Cavity Boundary Properties]
Inside any hollow conducting cavity of arbitrary shape, the electric field is \textbf{identically zero} ($\mathbf{E} = 0$), provided no charge is enclosed within the cavity. The potential throughout the cavity remains constant and strictly equals the potential on the conducting surface ($V_{\text{cavity}} = V_{\text{surface}}$).
\begin{itemize}
  \item \textit{Faraday Cage Principle:} Any sensitive electrical apparatus enclosed in a hollow metallic enclosure is completely shielded from external electrical disturbances.
  \item \textit{Practical Safety:} During lightning storms, it is significantly safer to remain inside a closed metallic car or bus than under a tree or in the open, as lightning charges flow along the outer conducting body directly to the ground without penetrating the passenger interior.
\end{itemize}
\end{ncertbox}

\begin{trapbox}[Syllabus Boundary Alert: Van de Graaff Generator Deletion]
\textbf{NMC REVISION WARNING FOR NEET UG:}
Under the revised NEET UG syllabus issued by the National Medical Commission:
\begin{itemize}
  \item \textbf{STRICTLY DELETED:} The \textbf{Van de Graaff Generator} (high voltage production via belt charge transport, corona discharge, insulating column, spherical dome collection) has been \textbf{completely eliminated}.
  \item \textbf{EXAM ORIENTATION:} Candidates must master Gauss's law boundary applications (spherical shells, solid spheres, line charges, and conductor electrostatic shielding) without studying the obsolete Van de Graaff machine.
\end{itemize}
\end{trapbox}

\section{Capacitance, Dielectric Physics \& Energy Storage}

\begin{conceptbox}[Dielectric Polarization \& Capacitance Principles]
\textbf{1. Capacitance Definition:}
The capacity of a conductor to store electric charge per unit potential rise:
\begin{equation}
  C = \frac{Q}{V} \quad [\mathrm{Farad}, \mathrm{F} = \mathrm{C}\cdot\mathrm{V}^{-1}]
\end{equation}
For an isolated conducting sphere of radius $R$: $C = 4\pi\varepsilon_0 R$.
\textbf{2. Parallel Plate Capacitor in Vacuum:}
Plates of area $A$ separated by vacuum gap $d$ with uniform field $E = \frac{\sigma}{\varepsilon_0} = \frac{Q}{\varepsilon_0 A}$ and $V = Ed$:
\begin{equation}
  \mathbf{C_0 = \frac{\varepsilon_0 A}{d}}
\end{equation}
\textbf{3. Dielectric Material Polarization:}
When an insulating dielectric is inserted into an external electric field $\mathbf{E}_0$, bound molecular dipoles align, generating a macroscopic polarization vector $\mathbf{P} = \varepsilon_0 \chi_e \mathbf{E}$.
This induces surface bound charge density $\sigma_b = \mathbf{P}\cdot\hat{\mathbf{n}}$ producing an internal opposing depolarizing field $\mathbf{E}_{\text{induced}} = \frac{\sigma_b}{\varepsilon_0}$.
\begin{equation}
  \mathbf{E_{\text{net}} = \mathbf{E}_0 - \mathbf{E}_{\text{induced}} = \frac{\mathbf{E}_0}{\kappa}} \quad (\text{Field is attenuated by factor } \kappa)
\end{equation}
where $\kappa = 1 + \chi_e = \varepsilon_r$ is the dielectric constant.
\textbf{4. Bound Induced Surface Charge:}
\begin{equation}
  E_{\text{net}} = \frac{\sigma - \sigma_b}{\varepsilon_0} = \frac{\sigma}{\kappa\varepsilon_0} \implies \sigma_b = \sigma\left(1 - \frac{1}{\kappa}\right) \implies \mathbf{q_b = q\left(1 - \frac{1}{\kappa}\right)}
\end{equation}
\end{conceptbox}

\begin{formulabox}[Capacitor with Partial Dielectric Slab \& Network Combinations]
\textbf{1. Partial Dielectric Slab Insertion (Thickness $t < d$):}
Across the vacuum gap $(d-t)$, electric field is $E_0 = \frac{Q}{\varepsilon_0 A}$. Inside the slab of thickness $t$, field is $E = \frac{E_0}{\kappa}$.
Total potential difference between plates:
\begin{equation}
  V = E_0 (d - t) + E t = E_0 \left[(d - t) + \frac{t}{\kappa}\right] = \frac{Q}{\varepsilon_0 A} \left[d - t\left(1 - \frac{1}{\kappa}\right)\right]
\end{equation}
\begin{equation}
  \mathbf{C = \frac{Q}{V} = \frac{\varepsilon_0 A}{d - t\left(1 - \frac{1}{\kappa}\right)}}
\end{equation}
\textit{Special Cases:}
\begin{itemize}
  \item Complete dielectric filling ($t=d$): $C = \kappa \frac{\varepsilon_0 A}{d} = \kappa C_0$.
  \item Conducting slab inserted ($t < d$, $\kappa \to \infty$): $C = \frac{\varepsilon_0 A}{d - t}$.
\end{itemize}
\textbf{2. Grouping of Capacitors:}
\begin{itemize}
  \item \textbf{Series Combination (Charge $Q$ is identical):}
  \begin{equation}
    V = V_1 + V_2 + \dots \implies \mathbf{\frac{1}{C_{\text{eq}}} = \sum_{i=1}^n \frac{1}{C_i}} \quad \left(\text{For two: } C_{\text{eq}} = \frac{C_1 C_2}{C_1 + C_2}\right)
  \end{equation}
  \item \textbf{Parallel Combination (Potential difference $V$ is identical):}
  \begin{equation}
    Q = Q_1 + Q_2 + \dots \implies \mathbf{C_{\text{eq}} = \sum_{i=1}^n C_i}
  \end{equation}
\end{itemize}
\textbf{3. Electrostatic Energy Stored \& Field Energy Density:}
Work required to assemble charge $Q$ on plates against potential $V$:
\begin{equation}
  \mathbf{U = \int_0^Q V(q) dq = \int_0^Q \frac{q}{C} dq = \frac{Q^2}{2C} = \frac{1}{2} C V^2 = \frac{1}{2} Q V}
\end{equation}
\textbf{4. Energy Density in Electric Field:}
Dividing energy $U$ by the volume of the dielectric space ($A \cdot d$):
\begin{equation}
  u_E = \frac{U}{Ad} = \frac{\frac{1}{2}(\kappa \varepsilon_0 A/d)(Ed)^2}{Ad} \implies \mathbf{u_E = \frac{1}{2} \kappa \varepsilon_0 E^2}
\end{equation}
\textbf{5. Attractive Mechanical Force Between Plates:}
\begin{equation}
  \mathbf{F = -\frac{dU}{dx} = \frac{Q^2}{2\varepsilon_0 A} = \frac{1}{2} Q E_0 = \frac{1}{2} \varepsilon_0 A E_0^2}
\end{equation}
\end{formulabox}

\begin{usecasebox}[Dielectric Insertion: Battery Connected vs. Battery Disconnected]
\begin{center}
\begin{tabularx}{\linewidth}{|l|Y|Y|}
\hline
\textbf{Physical Parameter} & \textbf{Case A: Battery Disconnected ($Q = \text{const}$)} & \textbf{Case B: Battery Remains Connected ($V = \text{const}$)} \\
\hline
\textbf{Charge ($Q$)} & Unchanged: $Q = Q_0$ & Increases: $Q = \kappa Q_0$ \\
\hline
\textbf{Capacitance ($C$)} & Increases: $C = \kappa C_0$ & Increases: $C = \kappa C_0$ \\
\hline
\textbf{Potential Difference ($V$)} & Decreases: $V = V_0 / \kappa$ & Unchanged: $V = V_0$ \\
\hline
\textbf{Electric Field ($E$)} & Decreases: $E = E_0 / \kappa$ & Unchanged: $E = E_0$ \\
\hline
\textbf{Stored Energy ($U$)} & Decreases: $U = U_0 / \kappa$ & Increases: $U = \kappa U_0$ \\
\hline
\textbf{Battery Work Done} & Zero (battery isolated) & $W_{\text{battery}} = \Delta Q \cdot V_0 = (\kappa - 1) C_0 V_0^2$ \\
\hline
\end{tabularx}
\end{center}
\end{usecasebox}

\section{Current Electricity, Drift Dynamics \& DC Resistive Networks}

\begin{conceptbox}[Drift Velocity, Microscopic Ohm's Law \& Mobility]
\textbf{1. Microscopic Mechanism of Conduction in Metals:}
Under zero external field, conduction electrons execute chaotic thermal motion at speeds $v_{\text{th}} \sim 10^5\,\mathrm{m/s}$ with random directions ($\langle \mathbf{v}_{\text{th}} \rangle = 0$).
When an external electric field $\mathbf{E}$ is applied, each electron experiences an electrostatic force $\mathbf{F} = -e\mathbf{E}$, imparting constant acceleration $\mathbf{a} = -\frac{e\mathbf{E}}{m_e}$.
Electrons repeatedly scatter elastically from vibrating crystal lattice ions. The average time elapsed between consecutive scattering collisions is the \textbf{relaxation time} $\tau \sim 10^{-14}\,\mathrm{s}$.
\textbf{2. Drift Velocity ($\mathbf{v}_d$) Derivation:}
\begin{equation}
  \mathbf{v}_d = \langle \mathbf{v}_{\text{initial}} + \mathbf{a} t \rangle = 0 + \mathbf{a}\tau \implies \mathbf{\mathbf{v}_d = -\frac{e \mathbf{E} \tau}{m_e}}
\end{equation}
Typically, drift speed is remarkably slow: $v_d \sim 10^{-4}\,\mathrm{m/s} = 0.1\,\mathrm{mm/s}$.
\textbf{3. Current Density \& Microscopic Ohm's Law:}
Consider a conductor of cross-sectional area $A$ with free electron density $n$. In time $\Delta t$, charge passing through is $\Delta Q = n e A v_d \Delta t$.
\begin{equation}
  I = \frac{\Delta Q}{\Delta t} = \mathbf{n e A v_d} \implies \mathbf{J = \frac{I}{A} = n e v_d = \left(\frac{n e^2 \tau}{m_e}\right) E}
\end{equation}
\begin{equation}
  \mathbf{\mathbf{J} = \sigma \mathbf{E} = \frac{1}{\rho}\mathbf{E}} \quad (\text{Vector / Microscopic Ohm's Law})
\end{equation}
where electrical conductivity $\sigma$ and resistivity $\rho$:
\begin{equation}
  \mathbf{\sigma = \frac{n e^2 \tau}{m_e}}, \qquad \mathbf{\rho = \frac{m_e}{n e^2 \tau}}
\end{equation}
\textbf{4. Electron Mobility ($\mu$):}
Drift speed per unit electric field:
\begin{equation}
  \mathbf{\mu = \frac{v_d}{E} = \frac{e \tau}{m_e}} \quad [\mathrm{m}^2\cdot\mathrm{V}^{-1}\cdot\mathrm{s}^{-1}]
\end{equation}
\end{conceptbox}

\begin{formulabox}[Temperature Dependence, Cell Grouping \& Power Transfer Theorem]
\textbf{1. Temperature Dependence of Resistance:}
For metallic conductors, increasing temperature intensifies lattice vibrations, reducing relaxation time $\tau$:
\begin{equation}
  R(T) = R_0 [1 + \alpha (T - T_0)] \implies \mathbf{\alpha = \frac{R_T - R_0}{R_0 \Delta T}} \quad [^\circ\mathrm{C}^{-1} \text{ or } \mathrm{K}^{-1}]
\end{equation}
For metals, $\alpha > 0$ (positive). For semiconductors and electrolytes, charge carrier density surges exponentially with temperature ($n(T) \propto e^{-E_g / 2k_B T}$), overriding $\tau$ drops, giving $\mathbf{\alpha < 0}$ (negative).
\textbf{2. Terminal Potential Difference of a Real Cell ($\mathcal{E}, r$):}
\begin{itemize}
  \item \textbf{Discharging (Supplying current $I$ to load $R$):} $\mathbf{V = \mathcal{E} - I r} = \frac{\mathcal{E} R}{R + r}$.
  \item \textbf{Charging (External charger driving current into cell):} $\mathbf{V = \mathcal{E} + I r}$.
  \item \textbf{Open Circuit ($I=0$):} $V = \mathcal{E}$. Short circuit ($R=0$): $I_{\text{sc}} = \frac{\mathcal{E}}{r}, V = 0$.
\end{itemize}
\textbf{3. Mixed Grouping of Cells ($m$ rows of $n$ series cells each):}
Total cells $N = n \times m$. Total EMF $= n\mathcal{E}$; total internal resistance $= \frac{n r}{m}$.
Current through load $R$:
\begin{equation}
  I = \frac{n\mathcal{E}}{R + \frac{n r}{m}} = \frac{m n \mathcal{E}}{m R + n r}
\end{equation}
To maximize current, minimize the denominator $(mR + nr) = (\sqrt{mR} - \sqrt{nr})^2 + 2\sqrt{mnRr}$.
The minimum occurs when $\sqrt{mR} = \sqrt{nr}$:
\begin{equation}
  \mathbf{R = \frac{n r}{m}} \quad (\text{Load resistance equals total internal resistance})
\end{equation}
\textbf{4. Maximum Power Transfer Theorem:}
Power delivered to load: $P(R) = I^2 R = \frac{\mathcal{E}^2 R}{(R + r)^2}$. Setting $\frac{dP}{dR} = 0$:
\begin{equation}
  \frac{dP}{dR} = \mathcal{E}^2 \frac{(R + r)^2 - 2R(R + r)}{(R + r)^4} = 0 \implies R + r - 2R = 0 \implies \mathbf{R = r}
\end{equation}
\begin{equation}
  \mathbf{P_{\max} = \frac{\mathcal{E}^2}{4r}}
\end{equation}
\end{formulabox}

\begin{formulabox}[Kirchhoff's Laws \& Measurement Bridges]
\textbf{1. Kirchhoff's Laws:}
\begin{itemize}
  \item \textbf{Junction Rule (KCL):} At any circuit node: $\sum I_{\text{in}} = \sum I_{\text{out}}$ (Conservation of Charge).
  \item \textbf{Loop Rule (KVL):} Around any closed loop: $\sum \Delta V = \sum \mathcal{E} - \sum I R = 0$ (Conservation of Energy).
\end{itemize}
\textbf{2. Wheatstone Bridge Balance Derivation:}
A quadrilateral network of four resistors $P, Q, R, S$ with a galvanometer connected between junctions $B$ and $D$.
At balance, no current traverses the galvanometer ($I_g = 0 \implies V_B = V_D$).
Applying KVL to loop $ABDA$: $I_1 P - I_2 R = 0 \implies I_1 P = I_2 R$.
Applying KVL to loop $BCDB$: $I_1 Q - I_2 S = 0 \implies I_1 Q = I_2 S$.
Dividing the two equations yields the classical balance condition:
\begin{equation}
  \mathbf{\frac{P}{Q} = \frac{R}{S}}
\end{equation}
\textbf{3. Metre Bridge (Slide-Wire Bridge):}
A uniform wire of length $100\,\mathrm{cm}$ acts as ratio arms $P = \rho_w l$ and $Q = \rho_w (100 - l)$.
\begin{equation}
  \frac{R}{X} = \frac{l}{100 - l} \implies \mathbf{X = R \left(\frac{100 - l}{l}\right)}
\end{equation}
Maximum experimental precision occurs at the central null point ($l \approx 50\,\mathrm{cm}$).
\end{formulabox}

\begin{trapbox}[Syllabus Boundary Alert: Potentiometer \& Carbon Resistor Color Codes Deletion]
\textbf{NMC REVISION WARNING FOR NEET UG:}
Under the revised NEET UG syllabus issued by the National Medical Commission:
\begin{itemize}
  \item \textbf{STRICTLY DELETED:} The \textbf{Potentiometer} (principle of fall of potential $V \propto l$, comparison of EMFs of two cells $\mathcal{E}_1/\mathcal{E}_2 = l_1/l_2$, and determination of cell internal resistance $r = R(l_1/l_2 - 1)$) has been \textbf{completely deleted}.
  \item \textbf{STRICTLY DELETED:} \textbf{Color Coding of Carbon Resistors} (BBROYGBVGW mnemonic bands, tolerance multipliers) has been \textbf{completely deleted}.
  \item \textbf{RETAINED:} The \textbf{Metre Bridge} and \textbf{Wheatstone Bridge} are fully retained and form the exclusive focus for bridge measurements.
\end{itemize}
\end{trapbox}

\section{Magnetic Effects of Electric Current \& Moving Charges}

\begin{conceptbox}[Biot-Savart Law \& Magnetic Lorentz Force]
\textbf{1. Biot-Savart Law:}
The infinitesimal magnetic field $d\mathbf{B}$ generated by an elementary current length vector $I d\mathbf{l}$ at displacement $\mathbf{r}$:
\begin{equation}
  \mathbf{d\mathbf{B} = \frac{\mu_0}{4\pi} \frac{I (d\mathbf{l} \times \hat{\mathbf{r}})}{r^2} = \frac{\mu_0}{4\pi} \frac{I (d\mathbf{l} \times \mathbf{r})}{r^3}}
\end{equation}
where $\mu_0 = 4\pi \times 10^{-7}\,\mathrm{T}\cdot\mathrm{m}\cdot\mathrm{A}^{-1}$ is vacuum permeability.
\textbf{2. Magnetic Lorentz Force on a Moving Charge:}
\begin{equation}
  \mathbf{\mathbf{F} = q(\mathbf{E} + \mathbf{v} \times \mathbf{B})}
\end{equation}
\begin{itemize}
  \item Pure magnetic force: $\mathbf{F}_B = q(\mathbf{v} \times \mathbf{B})$. Since $\mathbf{F}_B \perp \mathbf{v}$, the instantaneous power is $P = \mathbf{F}_B \cdot \mathbf{v} = 0$.
  \item \textbf{Work Done by Magnetic Force is Identically Zero:} Magnetic forces change the direction of motion, never particle kinetic energy or speed.
  \item \textbf{Helical Trajectory in Uniform Field $\mathbf{B}$:}
  When particle velocity $\mathbf{v}$ makes angle $\theta$ with $\mathbf{B}$ ($v_\perp = v\sin\theta, v_\parallel = v\cos\theta$):
  \begin{equation}
    \mathbf{r = \frac{m v_\perp}{q B} = \frac{m v \sin\theta}{q B}}, \qquad \mathbf{T = \frac{2\pi m}{q B}}, \qquad \mathbf{\text{Pitch } p = v_\parallel T = \frac{2\pi m v \cos\theta}{q B}}
  \end{equation}
\end{itemize}
\end{conceptbox}

\begin{formulabox}[Master Magnetic Field Formulas: Straight Wires, Loops, Solenoids \& Toroids]
\textbf{1. Finite Straight Wire Subtending Angles $\theta_1, \theta_2$ at Distance $d$:}
\begin{equation}
  \mathbf{B = \frac{\mu_0 I}{4\pi d}(\sin\theta_1 + \sin\theta_2)}
\end{equation}
\begin{itemize}
  \item \textbf{Infinitely Long Straight Wire ($\theta_1 = \theta_2 = 90^\circ$):} $\mathbf{B = \frac{\mu_0 I}{2\pi d}}$.
  \item \textbf{Semi-Infinite Wire End ($\theta_1 = 90^\circ, \theta_2 = 0^\circ$):} $\mathbf{B = \frac{\mu_0 I}{4\pi d}}$.
\end{itemize}
\textbf{2. Circular Current Loop of Radius $R$ on Symmetry Axis ($x$):}
\begin{equation}
  \mathbf{B(x) = \frac{\mu_0 I R^2}{2(R^2 + x^2)^{3/2}} \hat{\mathbf{i}}} \implies \mathbf{B_{\text{center}} = \frac{\mu_0 I}{2R}}
\end{equation}
\textbf{3. Ampere's Circuital Law:}
\begin{equation}
  \mathbf{\oint \mathbf{B}\cdot d\mathbf{l} = \mu_0 I_{\text{enclosed}}}
\end{equation}
\begin{itemize}
  \item \textbf{Ideal Long Solenoid ($n$ turns/meter):} Interior field is axial and uniform: $\mathbf{B = \mu_0 n I}$. Near ends: $B_{\text{end}} = \frac{1}{2}\mu_0 n I$.
  \item \textbf{Toroid of Mean Radius $R$ and $N$ Total Turns:} $n = \frac{N}{2\pi R} \implies \mathbf{B = \mu_0 n I = \frac{\mu_0 N I}{2\pi R}}$.
\end{itemize}
\textbf{4. Force Between Parallel Currents (Definition of Ampere):}
\begin{equation}
  \mathbf{\frac{F}{L} = \frac{\mu_0 I_1 I_2}{2\pi d}}
\end{equation}
Parallel currents \textbf{attract}; antiparallel currents \textbf{repel}.
\end{formulabox}

\begin{formulabox}[Moving Coil Galvanometer Conversion \& Magnetic Materials]
\textbf{1. Moving Coil Galvanometer (MCG):}
Deflecting torque balances radial restoring suspension torque: $\tau = N I A B = C \theta \implies \theta = \left(\frac{NAB}{C}\right)I$.
\begin{itemize}
  \item \textbf{Conversion to Ammeter (Extending range to $I > I_g$):} Connect small shunt resistance $S$ in parallel:
  \begin{equation}
    (I - I_g) S = I_g R_g \implies \mathbf{S = \frac{I_g R_g}{I - I_g}}
  \end{equation}
  \item \textbf{Conversion to Voltmeter (Extending range to $V$):} Connect high multiplier resistance $R$ in series:
  \begin{equation}
    V = I_g (R_g + R) \implies \mathbf{R = \frac{V}{I_g} - R_g}
  \end{equation}
\end{itemize}
\textbf{2. Classification of Magnetic Materials:}
\begin{itemize}
  \item \textbf{Diamagnetic:} $\chi_m < 0$ (small, negative, independent of temperature). $\mu_r < 1$. Repelled by magnetic fields. Examples: $\mathrm{Bi, Cu, H_2O, NaCl}$.
  \item \textbf{Paramagnetic:} $\chi_m > 0$ (small, positive). Follows \textbf{Curie's Law}: $\chi_m = \frac{C}{T}$. $\mu_r > 1$. Weakly attracted. Examples: $\mathrm{Al, Na, O_2}$.
  \item \textbf{Ferromagnetic:} $\chi_m \gg 1$ (enormous, positive). Follows \textbf{Curie-Weiss Law} above Curie temperature $T_c$: $\chi_m = \frac{C}{T - T_c}$. Exhibits domains and hysteresis. Examples: $\mathrm{Fe, Co, Ni}$.
\end{itemize}
\end{formulabox}

\begin{trapbox}[Syllabus Boundary Alert: Cyclotron \& Earth's Magnetism Deletion]
\textbf{NMC REVISION WARNING FOR NEET UG:}
Under the revised NEET UG syllabus issued by the National Medical Commission:
\begin{itemize}
  \item \textbf{STRICTLY DELETED:} The \textbf{Cyclotron} (resonating RF frequency $\nu_c = \frac{qB}{2\pi m}$, kinetic energy $K_{\max} = \frac{q^2 B^2 R^2}{2m}$, relativistic limits) has been \textbf{completely removed}.
  \item \textbf{STRICTLY DELETED:} \textbf{Earth's Magnetic Elements} (Magnetic Declination $\theta$, Angle of Dip / Inclination $\delta$, Horizontal Component $B_H$, relation $B = \sqrt{B_H^2 + B_V^2}$) have been \textbf{completely removed}.
\end{itemize}
\end{trapbox}

\section{Electromagnetic Induction (EMI) \& Inductive Circuits}

\begin{conceptbox}[Faraday's Law, Lenz's Law \& Motional EMF]
\textbf{1. Magnetic Flux ($\Phi_B$):}
\begin{equation}
  \Phi_B = \iint \mathbf{B}\cdot d\mathbf{A} = B A \cos\theta \quad [\mathrm{Weber}, \mathrm{Wb} = \mathrm{T}\cdot\mathrm{m}^2]
\end{equation}
\textbf{2. Faraday's Law of Electromagnetic Induction:}
Induced EMF in any closed circuit equals the negative time rate of change of magnetic flux:
\begin{equation}
  \mathbf{\mathcal{E} = -\frac{d\Phi_B}{dt} = -N \frac{d\Phi_B}{dt}}
\end{equation}
\textbf{3. Lenz's Law (Conservation of Energy):}
The polarity of the induced EMF is such that it produces a current whose magnetic field opposes the original flux change that created it.
\textbf{4. Motional EMF Derivations:}
\begin{itemize}
  \item \textbf{Straight Conducting Rod of Length $L$ Moving with Velocity $\mathbf{v} \perp \mathbf{B}$:}
  Lorentz force on free electrons $\mathbf{F} = -e(\mathbf{v} \times \mathbf{B})$ drives charge accumulation until internal electrostatic field cancels magnetic force: $eE = evB \implies \mathbf{\mathcal{E} = B v L}$.
  \item \textbf{Conducting Rod of Length $L$ Rotating at Angular Velocity $\omega$ in Uniform $\mathbf{B}$:}
  Integrating elemental motional EMF $d\mathcal{E} = B v(r) dr = B (\omega r) dr$:
  \begin{equation}
    \mathbf{\mathcal{E} = \int_0^L B \omega r dr = \frac{1}{2} B \omega L^2}
  \end{equation}
\end{itemize}
\end{conceptbox}

\begin{formulabox}[Self-Inductance, Mutual Inductance \& Magnetic Energy Density]
\textbf{1. Self-Inductance ($L$):}
Total flux linkage is proportional to circulating current: $N \Phi_B = L I$.
\begin{equation}
  \mathbf{\mathcal{E}_L = -L \frac{dI}{dt}} \quad [\mathrm{Henry}, \mathrm{H} = \mathrm{V}\cdot\mathrm{s}\cdot\mathrm{A}^{-1}]
\end{equation}
For a long air-core solenoid of length $l$, area $A$, and turn density $n$:
\begin{equation}
  \Phi = B A = (\mu_0 n I) A \implies N \Phi = (n l)(\mu_0 n I A) = \mu_0 n^2 A l I \implies \mathbf{L = \mu_0 n^2 A l = \mu_0 \frac{N^2}{l} A}
\end{equation}
\textbf{2. Mutual Inductance ($M$):}
$N_2 \Phi_2 = M_{21} I_1 \implies \mathcal{E}_2 = -M \frac{dI_1}{dt}$. For coaxial solenoids:
\begin{equation}
  \mathbf{M = \mu_0 n_1 n_2 A l = \frac{\mu_0 N_1 N_2 A}{l}}
\end{equation}
Coupling coefficient: $\mathbf{k = \frac{M}{\sqrt{L_1 L_2}}}$ ($0 \le k \le 1$).
\textbf{3. Energy Stored in an Inductor:}
\begin{equation}
  \mathbf{U = \int_0^I L i di = \frac{1}{2} L I^2}
\end{equation}
\textbf{4. Magnetic Energy Density in a Magnetic Field:}
\begin{equation}
  \mathbf{u_B = \frac{U}{\text{Volume}} = \frac{B^2}{2\mu_0}}
\end{equation}
\end{formulabox}

\begin{ncertbox}[NCERT Extraction: Eddy Currents \& Practical Technological Applications]
When bulk conductors are subjected to changing magnetic flux, circulating closed induced current loops (\textbf{Eddy Currents}) are established in planes perpendicular to the magnetic lines of force.
\begin{itemize}
  \item \textbf{Electromagnetic Damping:} The metallic non-magnetic frame of a moving coil galvanometer develops opposing eddy currents as it deflects, rapidly bringing the pointer to rest at its steady reading without oscillatory hunting (dead-beat galvanometer).
  \item \textbf{Induction Furnaces:} High-frequency alternating magnetic fields generate massive eddy currents within metallic charges, producing Joule heat sufficient to melt metals in vacuum.
  \item \textbf{Magnetic Braking in Trains:} Electromagnets activated above rails induce opposing eddy currents in the rotating metallic wheels/rails, producing smooth, wear-free deceleration.
  \item \textbf{Mitigation in Transformers:} To minimize eddy current heating losses in transformers and motors, solid iron cores are replaced with thin, individually varnished/insulated metallic laminations stacked parallel to magnetic flux lines.
\end{itemize}
\end{ncertbox}

\section{Alternating Currents (AC) \& Resonant Circuits}

\begin{conceptbox}[Alternating Current Fundamentals: Phasors \& RMS Values]
A sinusoidal alternating EMF and current:
\begin{equation}
  V(t) = V_0 \sin(\omega t), \qquad I(t) = I_0 \sin(\omega t - \phi)
\end{equation}
\textbf{1. Mean (Average) Value over Half-Cycle:}
\begin{equation}
  I_{\text{avg}} = \frac{1}{\pi/\omega} \int_0^{\pi/\omega} I_0 \sin(\omega t) dt = \mathbf{\frac{2}{\pi} I_0 \approx 0.637 I_0}
\end{equation}
Full cycle average is strictly zero: $\langle I \rangle_{\text{full}} = 0$.
\textbf{2. Root-Mean-Square (RMS / Virtual) Value:}
Equating Joule heating over a full cycle $H = \int_0^T I^2 R dt = I_{\text{rms}}^2 R T$:
\begin{equation}
  I_{\text{rms}} = \sqrt{\frac{1}{T} \int_0^T I_0^2 \sin^2(\omega t) dt} = \mathbf{\frac{I_0}{\sqrt{2}} \approx 0.707 I_0}, \qquad \mathbf{V_{\text{rms}} = \frac{V_0}{\sqrt{2}}}
\end{equation}
AC meters (hot-wire ammeters and voltmeters) always read \textbf{RMS values}.
\end{conceptbox}

\begin{formulabox}[Series LCR Circuit: Impedance, Resonance \& Quality Factor]
In a series LCR circuit with supply voltage $V(t) = V_0 \sin(\omega t)$:
\begin{equation}
  V_R = I R, \qquad V_L = I X_L = I (\omega L), \qquad V_C = I X_C = I \left(\frac{1}{\omega C}\right)
\end{equation}
Since $V_L$ leads current by $\pi/2$ and $V_C$ lags current by $\pi/2$:
\begin{equation}
  V_0 = \sqrt{V_R^2 + (V_L - V_C)^2} = I_0 \sqrt{R^2 + \left(\omega L - \frac{1}{\omega C}\right)^2}
\end{equation}
\begin{equation}
  \mathbf{Z = \sqrt{R^2 + (X_L - X_C)^2} = \sqrt{R^2 + \left(\omega L - \frac{1}{\omega C}\right)^2}}, \qquad \mathbf{\tan\phi = \frac{X_L - X_C}{R}}
\end{equation}
\textbf{1. Electrical Resonance Condition:}
When inductive reactance equals capacitive reactance:
\begin{equation}
  X_L = X_C \implies \omega_0 L = \frac{1}{\omega_0 C} \implies \mathbf{\omega_0 = \frac{1}{\sqrt{LC}}}, \quad \mathbf{\nu_0 = \frac{1}{2\pi\sqrt{LC}}}
\end{equation}
At resonance:
\begin{itemize}
  \item Impedance is minimum and purely resistive: $\mathbf{Z_{\min} = R}$.
  \item Current amplitude is maximum and in phase with voltage: $\mathbf{I_{0,\max} = \frac{V_0}{R}}$, $\phi = 0$.
  \item Power factor is unity: $\cos\phi = 1$.
\end{itemize}
\textbf{2. Sharpness of Resonance \& Quality Factor ($Q$):}
\begin{equation}
  \mathbf{Q = \frac{\omega_0 L}{R} = \frac{1}{\omega_0 C R} = \frac{1}{R}\sqrt{\frac{L}{C}} = \frac{\omega_0}{2\Delta\omega} = \frac{\text{Resonant Frequency}}{\text{Bandwidth}}}
\end{equation}
High $Q$-factor indicates sharp selectivity and low damping.
\textbf{3. Average Power in AC Circuit \& Power Factor:}
\begin{equation}
  \mathbf{P_{\text{avg}} = V_{\text{rms}} I_{\text{rms}} \cos\phi}
\end{equation}
where $\cos\phi = \frac{R}{Z}$ is the \textbf{power factor}.
\begin{itemize}
  \item Pure $L$ or pure $C$: $\phi = \pm \pi/2 \implies \cos\phi = 0 \implies P_{\text{avg}} = 0$ (\textbf{Wattless Current} $I_w = I_{\text{rms}} \sin\phi$).
  \item Pure $R$ or LCR at resonance: $\phi = 0 \implies \cos\phi = 1 \implies P_{\text{avg}} = V_{\text{rms}} I_{\text{rms}}$.
\end{itemize}
\textbf{4. Transformers:}
\begin{equation}
  \mathbf{\frac{V_s}{V_p} = \frac{N_s}{N_p} = \frac{I_p}{I_s} = K} \quad (\text{Transformation Ratio})
\end{equation}
Efficiency: $\eta = \frac{P_{\text{output}}}{P_{\text{input}}} = \frac{V_s I_s \cos\phi_s}{V_p I_p \cos\phi_p}$.
\end{formulabox}

\section{Electromagnetic Waves \& Maxwell's Equations}

\begin{conceptbox}[Maxwell's Displacement Current \& Unified Field Equations]
\textbf{1. Maxwell's Displacement Current ($I_d$):}
Ampere's Circuital Law ($\oint \mathbf{B}\cdot d\mathbf{l} = \mu_0 I$) failed when applied to a charging capacitor boundary. Maxwell resolved this inconsistency by postulating that a time-varying electric field produces an equivalent displacement current:
\begin{equation}
  \Phi_E = E A = \left(\frac{Q}{\varepsilon_0 A}\right) A = \frac{Q}{\varepsilon_0} \implies I_d = \varepsilon_0 \frac{d\Phi_E}{dt} = \frac{dQ}{dt} = I_c
\end{equation}
\begin{equation}
  \mathbf{\oint \mathbf{B}\cdot d\mathbf{l} = \mu_0 (I_c + I_d) = \mu_0 I_c + \mu_0 \varepsilon_0 \frac{d\Phi_E}{dt}} \quad (\text{Ampere-Maxwell Law})
\end{equation}
\textbf{2. Transverse Nature \& Speed of EM Waves:}
Electric and magnetic vectors oscillate perpendicular to each other and perpendicular to the wave propagation direction:
\begin{equation}
  \mathbf{c = \frac{1}{\sqrt{\mu_0 \varepsilon_0}} = \frac{E_0}{B_0} \approx 3.00 \times 10^8\,\mathrm{m/s}}
\end{equation}
\textbf{3. Poynting Vector \& Radiation Pressure:}
Energy flux vector: $\mathbf{S = \frac{1}{\mu_0}(\mathbf{E} \times \mathbf{B})}$. Average intensity: $I = \frac{1}{2}\varepsilon_0 c E_0^2 = \frac{E_0 B_0}{2\mu_0}$.
\begin{equation}
  \text{Radiation Pressure: } P_{\text{rad}} = \begin{cases} \frac{I}{c} & \text{Perfectly Absorbing Surface} \\ \frac{2I}{c} & \text{Perfectly Reflecting Surface} \end{cases}
\end{equation}
\end{conceptbox}

\begin{formulabox}[The Electromagnetic Spectrum: Master Classification]
\begin{center}
\begin{tabularx}{\linewidth}{|l|c|c|Y|}
\hline
\textbf{Wave Type} & \textbf{Wavelength ($\lambda$)} & \textbf{Frequency ($\nu$)} & \textbf{Production Mechanism \& Major Use} \\
\hline
\textbf{Radio Waves} & $> 0.1\,\mathrm{m}$ & $< 3 \times 10^9\,\mathrm{Hz}$ & Rapid charge acceleration in LC antennas; Radio/TV communication \\
\hline
\textbf{Microwaves} & $1\,\mathrm{mm} - 0.1\,\mathrm{m}$ & $10^9 - 10^{11}\,\mathrm{Hz}$ & Klystron/Magnetron tubes; Radar navigation, Microwave ovens \\
\hline
\textbf{Infrared} & $700\,\mathrm{nm} - 1\,\mathrm{mm}$ & $10^{11} - 4 \times 10^{14}\,\mathrm{Hz}$ & Thermal molecular vibrations; Greenhouse effect, Remote controls \\
\hline
\textbf{Visible Light} & $400\,\mathrm{nm} - 700\,\mathrm{nm}$ & $4 - 7.5 \times 10^{14}\,\mathrm{Hz}$ & Electron transitions in outer atomic shells; Human vision \\
\hline
\textbf{Ultraviolet} & $1\,\mathrm{nm} - 400\,\mathrm{nm}$ & $10^{15} - 10^{17}\,\mathrm{Hz}$ & High-temperature inner shell transitions, sun; Sterilization, LASIK \\
\hline
\textbf{X-Rays} & $0.01\,\mathrm{nm} - 10\,\mathrm{nm}$ & $10^{16} - 10^{20}\,\mathrm{Hz}$ & Fast electrons decelerated on heavy metal target; Medical radiography \\
\hline
\textbf{Gamma Rays} & $< 10^{-3}\,\mathrm{nm}$ & $> 10^{19}\,\mathrm{Hz}$ & Nuclear de-excitation / nuclear isomeric transitions; Cancer radiotherapy \\
\hline
\end{tabularx}
\end{center}
\end{formulabox}

\section{Master High-Yield Problem Archetypes \& Worked Solutions}

\begin{usecasebox}[Archetype 1: Cube Resistor Network Equivalent Resistances]
\textbf{Problem:} Twelve identical resistors, each of resistance $R$, form the edges of a skeleton cube. Compute equivalent resistance between: (i) diagonally opposite body corners, (ii) opposite face corners, and (iii) adjacent edge corners.
\textbf{Step-by-Step Solution:}
\begin{enumerate}
  \item \textbf{Body Diagonal (Points $A$ and $G$):}
  Inject current $I$ at $A$ and extract at $G$. By threefold rotational symmetry around the body diagonal:
  \begin{itemize}
    \item 3 edges leave $A$, each carrying $I/3$.
    \item Each splits into 2 edges, each carrying $I/6$ across the middle 6 edges.
    \item These recombine into 3 edges entering $G$, each carrying $I/3$.
  \end{itemize}
  Total potential drop: $V_{AG} = \left(\frac{I}{3}\right)R + \left(\frac{I}{6}\right)R + \left(\frac{I}{3}\right)R = I R \left(\frac{1}{3} + \frac{1}{6} + \frac{1}{3}\right) = \frac{5}{6} I R$.
  \begin{equation}
    \mathbf{R_{\text{body}} = \frac{5}{6} R}
  \end{equation}
  \item \textbf{Face Diagonal (Points $A$ and $C$):}
  By mirror symmetry across the plane bisecting the face:
  \begin{equation}
    \mathbf{R_{\text{face}} = \frac{3}{4} R}
  \end{equation}
  \item \textbf{Adjacent Edge (Points $A$ and $B$):}
  \begin{equation}
    \mathbf{R_{\text{edge}} = \frac{7}{12} R}
  \end{equation}
\end{enumerate}
\end{usecasebox}

\begin{usecasebox}[Archetype 2: Series LCR Circuit Resonance \& Power Factor Calculation]
\textbf{Problem:} A series LCR circuit has $R = 10\,\Omega$, $L = 0.1\,\mathrm{H}$, and $C = 100\,\mu\mathrm{F}$ connected to an AC source $V(t) = 200 \sin(100 t)\,\mathrm{V}$. Calculate: (i) Inductive reactance $X_L$, (ii) Capacitive reactance $X_C$, (iii) Total impedance $Z$, (iv) Peak current $I_0$, (v) Resonant frequency $\omega_0$, and (vi) Quality factor $Q$.
\textbf{Step-by-Step Solution:}
\begin{enumerate}
  \item $\omega = 100\,\mathrm{rad/s}$, $V_0 = 200\,\mathrm{V}$.
  \item $X_L = \omega L = 100 \times 0.1 = \mathbf{10\,\Omega}$.
  \item $X_C = \frac{1}{\omega C} = \frac{1}{100 \times 100 \times 10^{-6}} = \frac{1}{10^{-2}} = \mathbf{100\,\Omega}$.
  \item Total impedance $Z$:
  \begin{equation}
    Z = \sqrt{R^2 + (X_L - X_C)^2} = \sqrt{10^2 + (10 - 100)^2} = \sqrt{100 + (-90)^2} = \sqrt{100 + 8100} = \sqrt{8200} \approx \mathbf{90.55\,\Omega}
  \end{equation}
  Circuit is strongly capacitive ($X_C > X_L$), current leads voltage.
  \item Peak current: $I_0 = \frac{V_0}{Z} = \frac{200}{90.55} \approx \mathbf{2.21\,\mathrm{A}}$.
  \item Resonant angular frequency $\omega_0$:
  \begin{equation}
    \omega_0 = \frac{1}{\sqrt{LC}} = \frac{1}{\sqrt{0.1 \times 10^{-4}}} = \frac{1}{\sqrt{10^{-5}}} = \frac{1000}{\sqrt{10}} \approx \mathbf{316.2\,\mathrm{rad/s}}
  \end{equation}
  \item Quality factor $Q$:
  \begin{equation}
    Q = \frac{1}{R}\sqrt{\frac{L}{C}} = \frac{1}{10}\sqrt{\frac{0.1}{10^{-4}}} = \frac{1}{10}\sqrt{1000} = \frac{31.62}{10} \approx \mathbf{3.16}
  \end{equation}
\end{enumerate}
\end{usecasebox}

% =========================================================================
% PART IV: OPTICS, MODERN PHYSICS & SEMICONDUCTORS
% =========================================================================
\part{Optics, Modern Physics \& Semiconductor Electronics}

\chapter{Ray Optics, Wave Optics, Quantum Radiation, Atomic Physics \& Semiconductors}

This chapter covers geometric and wave optics alongside modern quantum and solid-state physics. We investigate Snell's law, total internal reflection, the lens-maker's equation, optical instruments, Huygens' wavelets, Young's double-slit interference, single-slit diffraction, Einstein's photoelectric effect, de Broglie matter waves, Bohr's atomic postulates and spectral emissions, nuclear binding energy, and p-n junction diode electronics with digital logic gates.

\begin{mentorbox}[A Personal Note from Satya: The 100-Mistakes Principle (Why Trying Again Actually Rewires Your Brain)]
\textbf{You are more than halfway through this book. Let's talk frankly about the days when self-doubt creeps in.}

Every single medical student currently walking the corridors of AIIMS, MAMC, or KGMU had days where they scored 480 in a Sunday grand test, stared at their marks sheet in despair, and wondered if they simply weren't smart enough.

Here is what neuroscience tells us about high-stakes learning: \textit{your brain does not grow when you get a question right.} Getting a question right simply exercises an existing, pre-built neural circuit. But when you tackle a difficult problem, get it wrong, feel that sting of frustration, and then sit down to trace *exactly* where your reasoning derailed---that is the precise biological moment where your brain releases acetylcholine and dopamine, signaling your synapses to rebuild and strengthen.

\textbf{How to Turn Mistakes into a 700+ Score:}
\begin{enumerate}[leftmargin=*]
  \item \textbf{Maintain a Ruthless Error Logbook:} Never just say ``Oh, that was a silly calculation mistake.'' Categorize every single error into one of three buckets:
  \begin{itemize}
    \item \textit{Conceptual Blind Spot:} You didn't know the law or rule (e.g., forgetting that de Broglie wavelength for a charged particle depends on $qV$).
    \item \textit{Question Trap Slip:} You knew the concept, but NTA phrased it with ``incorrect statement'' or switched units from $\mu\text{m}$ to $\text{nm}$.
    \item \textit{Execution Panic:} You rushed through basic arithmetic under clock pressure.
  \end{itemize}
  \item \textbf{Re-solve Every Mistake 48 Hours Later:} An error is only truly eliminated when you can re-solve the question on a blank sheet of paper two days later without looking at the answer key.
\end{enumerate}

Every single mistake you catch and correct in this book is one fewer trap waiting for you on the actual NEET exam day. Treat every wrong answer not as an indictment of your talent, but as a solved mystery. Keep your head high, turn the page, and let's master Modern Physics.
\end{mentorbox}

\section{Ray Optics, Reflection, Refraction \& Total Internal Reflection}

\begin{conceptbox}[Foundations of Geometric Optics \& Sign Conventions]
\textbf{1. Cartesian Sign Convention (Strict Standard):}
\begin{itemize}
  \item All optical distances are measured along the principal axis starting from the optical pole ($P$) or optical centre ($O$).
  \item Distances measured in the direction of incident light are \textbf{positive} ($+$); distances opposite to incident light are \textbf{negative} ($-$).
  \item Heights measured perpendicularly upwards above the principal axis are \textbf{positive} ($+$); heights downwards are \textbf{negative} ($-$).
\end{itemize}
\textbf{2. Reflection at Spherical Mirrors:}
\begin{equation}
  \mathbf{\frac{1}{v} + \frac{1}{u} = \frac{1}{f} = \frac{2}{R}}
\end{equation}
\begin{itemize}
  \item \textbf{Concave Mirror:} $f < 0, R < 0$ (converging). Produces real inverted images for $u > |f|$, and virtual erect magnified image for $u < |f|$.
  \item \textbf{Convex Mirror:} $f > 0, R > 0$ (diverging). Produces virtual, erect, diminished images for all real object locations.
  \item \textbf{Transverse (Linear) Magnification:}
  \begin{equation}
    \mathbf{m = \frac{h_i}{h_o} = -\frac{v}{u} = \frac{f}{f - u} = \frac{f - v}{f}}
  \end{equation}
  $m < 0 \implies$ real, inverted; $m > 0 \implies$ virtual, erect.
  \item \textbf{Longitudinal Magnification for Small Objects ($du$):} $m_L = \frac{dv}{du} = -m^2 = -\left(\frac{v}{u}\right)^2$.
\end{itemize}
\end{conceptbox}

\begin{formulabox}[Refraction at Plane Interfaces, Apparent Depth \& Total Internal Reflection]
\textbf{1. Snell's Law of Refraction:}
\begin{equation}
  \mathbf{n_1 \sin i = n_2 \sin r \implies \frac{\sin i}{\sin r} = \frac{n_2}{n_1} = \,^1n_2 = \frac{v_1}{v_2} = \frac{\lambda_1}{\lambda_2}}
\end{equation}
Frequency $\nu$ remains invariant across any refractive boundary: $\nu_1 = \nu_2 = \nu$.
\textbf{2. Apparent Depth \& Normal Shift in Dense Media:}
An object placed at real depth $d$ in a medium of refractive index $\mu$ viewed normally from air:
\begin{equation}
  \mathbf{d_{\text{apparent}} = \frac{d}{\mu}}, \qquad \mathbf{\text{Normal Shift } \Delta t = d - d_{\text{app}} = d\left(1 - \frac{1}{\mu}\right)}
\end{equation}
For composite multiple horizontal layers: $d_{\text{apparent, total}} = \sum_{i=1}^n \frac{d_i}{\mu_i}$.
\textbf{3. Total Internal Reflection (TIR):}
When light propagates from an optically denser medium ($n_1$) toward a rarer medium ($n_2 < n_1$):
\begin{itemize}
  \item Critical angle $i_c$: angle of incidence where angle of refraction $r = 90^\circ$:
  \begin{equation}
    n_1 \sin i_c = n_2 \sin 90^\circ \implies \mathbf{\sin i_c = \frac{n_2}{n_1} = \frac{1}{\mu}} \quad (\text{if rarer medium is air, } n_2 = 1)
  \end{equation}
  \item For $i > i_c$, $100\%$ of incident light energy reflects back into the denser medium with zero transmission loss.
  \item \textbf{Optical Fibres:} Comprise a central cylindrical core ($\mu_1$) surrounded by a cladding ($\mu_2 < \mu_1$). Light admitted within the cone of acceptance undergoes repeated total internal reflections.
  \begin{equation}
    \mathbf{\text{Numerical Aperture (NA)} = \sin\theta_a = \sqrt{\mu_1^2 - \mu_2^2}}
  \end{equation}
\end{itemize}
\end{formulabox}

\begin{formulabox}[Refraction at a Single Spherical Surface: First-Principles Derivation]
Consider a convex spherical refracting interface of radius of curvature $R$ separating medium 1 ($\mu_1$) from medium 2 ($\mu_2$). A point object $O$ lies on the principal axis in medium 1 at distance $-u$. Light refracts to form an image $I$ at $+v$ in medium 2.
\begin{enumerate}
  \item Let ray from $O$ strike surface at $M$ near pole $P$. Let angles with principal axis be $\alpha$ (at $O$), $\beta$ (at $I$), and $\theta$ (at center of curvature $C$).
  \item In triangle $OMC$, exterior angle is incidence angle: $i = \alpha + \theta$.
  \item In triangle $IMC$, exterior angle is $\theta$: $\theta = r + \beta \implies r = \theta - \beta$.
  \item Under paraxial approximations ($\tan\phi \approx \phi \approx \sin\phi$ and $M \approx P$):
  \begin{equation}
    \alpha \approx \frac{PM}{u} = -\frac{h}{u}, \qquad \beta \approx \frac{PM}{v} = \frac{h}{v}, \qquad \theta \approx \frac{PM}{R} = \frac{h}{R}
  \end{equation}
  \item By Snell's Law for small angles: $\mu_1 i = \mu_2 r$.
  \begin{equation}
    \mu_1 (\alpha + \theta) = \mu_2 (\theta - \beta) \implies \mu_1 \alpha + \mu_2 \beta = (\mu_2 - \mu_1)\theta
  \end{equation}
  \item Substituting the paraxial arc angles:
  \begin{equation}
    \mu_1 \left(-\frac{h}{u}\right) + \mu_2 \left(\frac{h}{v}\right) = (\mu_2 - \mu_1)\frac{h}{R}
  \end{equation}
  Cancelling $h$:
  \begin{equation}
    \mathbf{\frac{\mu_2}{v} - \frac{\mu_1}{u} = \frac{\mu_2 - \mu_1}{R}}
  \end{equation}
\end{enumerate}
This master governing formula holds universally for both convex and concave surfaces, whether light travels from rarer to denser or denser to rarer, provided proper Cartesian signs are used for $u, v, R$.
\end{formulabox}

\section{Thin Lenses, Lens-Maker's Formula, Prisms \& Optical Instruments}

\begin{formulabox}[Lens-Maker's Formula \& Thin Lens Equation: Step-by-Step Proof]
Consider a thin lens of refractive index $\mu_2$ bounded by spherical surfaces of radii $R_1$ and $R_2$ immersed in a surrounding medium of refractive index $\mu_1$.
\begin{enumerate}
  \item \textbf{Refraction at First Surface ($R_1$):} Object at $u$, forms intermediate image at $v'$:
  \begin{equation}
    \frac{\mu_2}{v'} - \frac{\mu_1}{u} = \frac{\mu_2 - \mu_1}{R_1} \tag{Eq. 1}
  \end{equation}
  \item \textbf{Refraction at Second Surface ($R_2$):} The intermediate image $v'$ serves as a virtual object for the second interface, forming final image at $v$:
  \begin{equation}
    \frac{\mu_1}{v} - \frac{\mu_2}{v'} = \frac{\mu_1 - \mu_2}{R_2} = -\frac{\mu_2 - \mu_1}{R_2} \tag{Eq. 2}
  \end{equation}
  \item \textbf{Adding Equations (1) and (2):} The intermediate term $\frac{\mu_2}{v'}$ cancels completely:
  \begin{equation}
    \mu_1 \left(\frac{1}{v} - \frac{1}{u}\right) = (\mu_2 - \mu_1)\left(\frac{1}{R_1} - \frac{1}{R_2}\right)
  \end{equation}
  Dividing through by $\mu_1$:
  \begin{equation}
    \frac{1}{v} - \frac{1}{u} = \left(\frac{\mu_2}{\mu_1} - 1\right)\left(\frac{1}{R_1} - \frac{1}{R_2}\right)
  \end{equation}
  \item Setting $u \to -\infty \implies v = f$:
  \begin{equation}
    \mathbf{\frac{1}{f} = \left(\frac{\mu_{\text{lens}}}{\mu_{\text{medium}}} - 1\right)\left(\frac{1}{R_1} - \frac{1}{R_2}\right)} \quad (\text{Lens-Maker's Formula})
  \end{equation}
  \begin{equation}
    \mathbf{\frac{1}{v} - \frac{1}{u} = \frac{1}{f}} \quad (\text{Gaussian Thin Lens Formula})
  \end{equation}
\end{enumerate}
\textbf{3. Optical Power ($P$):}
\begin{equation}
  \mathbf{P = \frac{1}{f\,[\mathrm{m}]}} \quad [\mathrm{Dioptres}, \mathrm{D} = \mathrm{m}^{-1}]
\end{equation}
\textbf{4. Combination of Thin Lenses:}
\begin{itemize}
  \item \textbf{Lenses in Contact:} $\mathbf{P_{\text{eq}} = P_1 + P_2 \implies \frac{1}{F_{\text{eq}}} = \frac{1}{f_1} + \frac{1}{f_2}}$, magnification $m = m_1 \times m_2$.
  \item \textbf{Separated by Distance $d$:} $\mathbf{P_{\text{eq}} = P_1 + P_2 - d P_1 P_2 \implies \frac{1}{F_{\text{eq}}} = \frac{1}{f_1} + \frac{1}{f_2} - \frac{d}{f_1 f_2}}$.
\end{itemize}
\end{formulabox}

\begin{formulabox}[Refraction Through Prisms, Minimum Deviation \& Dispersion]
A triangular prism of refracting angle $A$ and refractive index $\mu$:
\begin{enumerate}
  \item \textbf{Geometric Relations:}
  \begin{equation}
    r_1 + r_2 = A, \qquad \mathbf{\delta = i + e - A}
  \end{equation}
  where $i$ is incidence angle, $e$ is emergence angle, and $\delta$ is angular deviation.
  \item \textbf{Condition for Minimum Deviation ($\delta = \delta_m$):}
  By symmetry and reversibility of light paths: $i = e$ and $r_1 = r_2 = \frac{A}{2}$.
  \begin{equation}
    \delta_m = 2i - A \implies i = \frac{A + \delta_m}{2}
  \end{equation}
  Applying Snell's law at first refracting surface ($\sin i = \mu \sin r_1$):
  \begin{equation}
    \mathbf{\mu = \frac{\sin\left(\frac{A + \delta_m}{2}\right)}{\sin\left(\frac{A}{2}\right)}} \quad (\text{Prism Formula})
  \end{equation}
  \item \textbf{Thin Prism Approximation ($A \le 10^\circ$):}
  \begin{equation}
    \sin\theta \approx \theta \implies \frac{A + \delta}{2} \approx \mu \frac{A}{2} \implies \mathbf{\delta = (\mu - 1) A}
  \end{equation}
  \item \textbf{Dispersion \& Dispersive Power ($\omega$):}
  \begin{itemize}
    \item Angular Dispersion: $\mathbf{\theta = \delta_V - \delta_R = (\mu_V - \mu_R) A}$.
    \item Mean Deviation (Yellow light): $\delta_Y = (\mu_Y - 1) A$.
    \item Dispersive Power: $\mathbf{\omega = \frac{\delta_V - \delta_R}{\delta_Y} = \frac{\mu_V - \mu_R}{\mu_Y - 1}}$ (Independent of prism angle $A$).
  \end{itemize}
\end{enumerate}
\end{formulabox}

\begin{formulabox}[Magnifying Power of Optical Instruments: Microscopes \& Telescopes]
\textbf{1. Simple Microscope (Magnifying Glass):}
\begin{itemize}
  \item Image formed at Near Point ($D = 25\,\mathrm{cm}$): $\mathbf{m = 1 + \frac{D}{f}}$ (Maximum magnification, strained eye).
  \item Image formed at Infinity (Normal Adjustment): $\mathbf{m = \frac{D}{f}}$ (Relaxed eye).
\end{itemize}
\textbf{2. Compound Microscope:}
Objective lens ($f_o$, small aperture) + Eyepiece ($f_e$, larger aperture).
\begin{itemize}
  \item Total Magnifying Power: $m = m_o \times m_e = \left(-\frac{v_o}{u_o}\right) m_e$.
  \item \textbf{Near Point Adjustment ($D$):} $\mathbf{m = -\frac{v_o}{u_o}\left(1 + \frac{D}{f_e}\right)}$, Tube length $L = v_o + |u_e|$.
  \item \textbf{Normal Adjustment (Infinity):} $\mathbf{m \approx -\frac{L}{f_o}\frac{D}{f_e}}$, Tube length $L = v_o + f_e$.
\end{itemize}
\textbf{3. Astronomical Refracting Telescope:}
Objective ($f_o$, large aperture) + Eyepiece ($f_e$, small aperture).
\begin{itemize}
  \item \textbf{Normal Adjustment (Image at Infinity):} $\mathbf{m = -\frac{f_o}{f_e}}$, Tube length $\mathbf{L = f_o + f_e}$.
  \item \textbf{Near Point Adjustment ($D$):} $\mathbf{m = -\frac{f_o}{f_e}\left(1 + \frac{f_e}{D}\right)}$, Tube length $L = f_o + \frac{f_e D}{f_e + D}$.
\end{itemize}
\end{formulabox}

\begin{trapbox}[Syllabus Boundary Alert: Deletion of Resolving Power of Optical Instruments]
\textbf{NMC REVISION WARNING FOR NEET UG:}
Under the revised NEET UG curriculum:
\begin{itemize}
  \item \textbf{STRICTLY DELETED:} The theoretical formulas for \textbf{Resolving Power of a Microscope} ($\text{RP} = \frac{2\mu\sin\theta}{1.22\lambda}$) and \textbf{Resolving Power of an Astronomical Telescope} ($\text{RP} = \frac{D}{1.22\lambda}$), along with Rayleigh's criterion angular limit $\theta = \frac{1.22\lambda}{D}$, have been \textbf{completely deleted}.
  \item \textbf{EXAM ORIENTATION:} Candidates must focus on Magnifying Power derivations, tube lengths, and ray diagrams of microscopes/telescopes without computing resolving powers.
\end{itemize}
\end{trapbox}

\section{Wave Optics: Huygens' Principle, Interference \& Diffraction}

\begin{conceptbox}[Wavefronts, Huygens' Secondary Wavelets \& Interference Physics]
\textbf{1. Wavefront Classification:}
A wavefront is the locus of all adjacent points in a medium vibrating with identical instantaneous phase:
\begin{itemize}
  \item \textbf{Spherical Wavefront:} Emitted by an isotropic point source; intensity falls as $I \propto \frac{1}{r^2}$, amplitude $A \propto \frac{1}{r}$.
  \item \textbf{Cylindrical Wavefront:} Emitted by a linear slit source; intensity $I \propto \frac{1}{r}$, amplitude $A \propto \frac{1}{\sqrt{r}}$.
  \item \textbf{Plane Wavefront:} Formed at infinite distance from any source; intensity and amplitude remain constant.
  \item Light rays are everywhere perpendicular (normal) to the advancing wavefront.
\end{itemize}
\textbf{2. Huygens' Principle:}
\begin{enumerate}
  \item Every point on a primary wavefront acts as a source of secondary spherical wavelets spreading in all forward directions at the speed of the wave.
  \item The forward envelope (tangential surface) touching all these secondary wavelets at any subsequent instant defines the new position of the wavefront.
\end{enumerate}
\textbf{3. Superposition of Two Coherent Waves:}
Let two waves of identical frequency have amplitudes $a_1, a_2$ and relative phase difference $\phi$:
\begin{equation}
  y_1 = a_1 \sin(\omega t), \qquad y_2 = a_2 \sin(\omega t + \phi)
\end{equation}
Resultant amplitude $A_R$ and resultant intensity $I_R$:
\begin{equation}
  \mathbf{A_R = \sqrt{a_1^2 + a_2^2 + 2a_1 a_2 \cos\phi}}, \qquad \mathbf{I_R = I_1 + I_2 + 2\sqrt{I_1 I_2}\cos\phi}
\end{equation}
If two sources are identical ($I_1 = I_2 = I_0$):
\begin{equation}
  \mathbf{I_R = 2I_0(1 + \cos\phi) = 4I_0 \cos^2\left(\frac{\phi}{2}\right)}
\end{equation}
\begin{itemize}
  \item \textbf{Constructive Interference (Bright Fringes):} $\phi = 2n\pi \iff \Delta x = n\lambda \implies \mathbf{I_{\max} = (\sqrt{I_1} + \sqrt{I_2})^2 = 4I_0}$.
  \item \textbf{Destructive Interference (Dark Fringes):} $\phi = (2n-1)\pi \iff \Delta x = (2n-1)\frac{\lambda}{2} \implies \mathbf{I_{\min} = (\sqrt{I_1} - \sqrt{I_2})^2 = 0}$.
  \item \textbf{Fringe Contrast / Visibility:} $V = \frac{I_{\max} - I_{\min}}{I_{\max} + I_{\min}} = \frac{2\sqrt{I_1 I_2}}{I_1 + I_2}$.
\end{itemize}
\end{conceptbox}

\begin{formulabox}[Young's Double Slit Experiment (YDSE): Path Difference \& Fringe Width]
Two narrow slits $S_1, S_2$ separated by distance $d$ illuminate a screen at distance $D \gg d$. Let $P$ be a point at height $y$ from the central line.
\begin{enumerate}
  \item Path difference $\Delta x = S_2 P - S_1 P$:
  \begin{equation}
    S_2 P^2 - S_1 P^2 = \left[D^2 + \left(y + \frac{d}{2}\right)^2\right] - \left[D^2 + \left(y - \frac{d}{2}\right)^2\right] = 2yd
  \end{equation}
  \begin{equation}
    (S_2 P - S_1 P)(S_2 P + S_1 P) = 2yd \implies \Delta x \approx \frac{2yd}{2D} \implies \mathbf{\Delta x = \frac{yd}{D} = d\sin\theta}
  \end{equation}
  \item \textbf{Bright Fringe Positions (Maxima):}
  \begin{equation}
    \Delta x = n\lambda \implies \frac{y_n d}{D} = n\lambda \implies \mathbf{y_n = \frac{n\lambda D}{d}} \quad (n = 0, \pm 1, \pm 2, \dots)
  \end{equation}
  \item \textbf{Dark Fringe Positions (Minima):}
  \begin{equation}
    \Delta x = (2n-1)\frac{\lambda}{2} \implies \mathbf{y_n' = (2n-1)\frac{\lambda D}{2d}} \quad (n = 1, 2, 3, \dots)
  \end{equation}
  \item \textbf{Fringe Width ($\beta$):} Distance separating any two consecutive bright or dark fringes:
  \begin{equation}
    \mathbf{\beta = y_{n+1} - y_n = \frac{\lambda D}{d}}
  \end{equation}
  $\beta$ is completely uniform across all fringe orders on the screen.
  \item \textbf{Angular Fringe Width ($\theta_0$):}
  \begin{equation}
    \mathbf{\theta_0 = \frac{\beta}{D} = \frac{\lambda}{d}} \quad (\text{Independent of screen distance } D)
  \end{equation}
  \item \textbf{Fringe Shift due to Transparent Thin Sheet of Thickness $t$ and Index $\mu$:}
  Introducing the sheet in front of one slit introduces additional optical path $(\mu - 1)t$:
  \begin{equation}
    \mathbf{\Delta y_0 = \frac{(\mu - 1)t D}{d} = \frac{(\mu - 1)t}{\lambda}\beta}
  \end{equation}
  The entire fringe pattern shifts toward the covered slit without changing fringe width $\beta$.
\end{enumerate}
\end{formulabox}

\begin{formulabox}[Single-Slit Fraunhofer Diffraction: Minima \& Central Maximum Width]
Light of wavelength $\lambda$ illuminates a single rectangular slit of width $a$.
\begin{enumerate}
  \item \textbf{Diffraction Minima Condition:}
  The slit aperture is divided into $2n$ equal zones that cancel pairwise:
  \begin{equation}
    \mathbf{a \sin\theta = n \lambda} \quad (n = \pm 1, \pm 2, \pm 3, \dots)
  \end{equation}
  For small angles: $\theta_n \approx \frac{n\lambda}{a}$.
  \item \textbf{Secondary Maxima Condition:}
  \begin{equation}
    \mathbf{a \sin\theta \approx (2n+1)\frac{\lambda}{2}} \quad (n = 1, 2, 3, \dots)
  \end{equation}
  Peak intensities drop precipitously: $I_0 : \frac{I_0}{22} : \frac{I_0}{61} : \dots$
  \item \textbf{Angular Width of Central Maximum ($2\theta_0$):}
  Bounded by the first diffraction minima on either side ($n = \pm 1$):
  \begin{equation}
    \mathbf{2\theta_0 = \frac{2\lambda}{a}}
  \end{equation}
  \item \textbf{Linear Width of Central Diffraction Maximum ($W$):}
  \begin{equation}
    \mathbf{W = 2 y_1 = 2 \left(\frac{\lambda D}{a}\right) = \frac{2\lambda D}{a}}
  \end{equation}
  Notice: The central diffraction peak is \textbf{twice as wide} as secondary maxima ($W = 2\beta_{\text{diff}}$).
\end{enumerate}
\end{formulabox}

\begin{trapbox}[Syllabus Boundary Alert: Complete Deletion of Light Polarisation]
\textbf{NMC REVISION WARNING FOR NEET UG:}
Under the revised NEET UG syllabus issued by the National Medical Commission:
\begin{itemize}
  \item \textbf{STRICTLY DELETED:} The entire topic of \textbf{Polarisation of Light} has been \textbf{completely eliminated}.
  \item Deleted concepts include: Brewster's Law of polarisation by reflection ($\tan i_p = \mu$), Brewster's angle, Malus' Law of transmitted intensity ($I = I_0 \cos^2\theta$), Polaroid sheets, polarisation by scattering, and optical activity.
  \item \textbf{EXAM ORIENTATION:} Candidates must focus solely on Wavefronts, Huygens' Principle, YDSE Interference, and Single-Slit Fraunhofer Diffraction.
\end{itemize}
\end{trapbox}

\section{Dual Nature of Matter \& Radiation: Photoelectric Effect \& de Broglie Waves}

\begin{conceptbox}[Photoelectric Effect: Experimental Laws \& Einstein's Quantum Derivation]
\textbf{1. Key Experimental Observations (Hertz, Hallwachs, Lenard):}
\begin{itemize}
  \item \textbf{Instantaneous Emission:} Electron emission occurs without any measurable time lag ($< 10^{-9}\,\mathrm{s}$), contradicting classical wave continuous energy buildup.
  \item \textbf{Threshold Frequency ($\nu_0$):} For each photosensitive metal, there exists a minimum characteristic threshold frequency $\nu_0$ below which \textbf{no photoelectrons are emitted}, regardless of incident light intensity or duration.
  \item \textbf{Kinetic Energy vs. Frequency:} Maximum kinetic energy $K_{\max}$ is strictly linear with incident frequency $\nu$ and completely independent of light intensity.
  \item \textbf{Saturation Current vs. Intensity:} Photoelectric saturation current is strictly proportional to incident light intensity (photon flux).
\end{itemize}
\textbf{2. Einstein's Photoelectric Equation Derivation:}
Light consists of discrete energy quanta called \textbf{photons}, each carrying energy $E = h\nu$.
When a photon of energy $h\nu$ collides with a conduction electron, its entire energy is transferred to that single electron. A portion equal to the metal's \textbf{work function} ($\Phi_0 = h\nu_0$) is consumed to overcome surface binding; the remainder appears as maximum electron kinetic energy:
\begin{equation}
  h\nu = \Phi_0 + K_{\max} \implies \mathbf{K_{\max} = h\nu - \Phi_0 = h(\nu - \nu_0) = \frac{hc}{\lambda} - \frac{hc}{\lambda_0}}
\end{equation}
where $h = 6.626 \times 10^{-34}\,\mathrm{J}\cdot\mathrm{s}$ is Planck's constant ($hc \approx 12400\,\mathrm{eV}\cdot\text{\AA} \approx 1240\,\mathrm{eV}\cdot\mathrm{nm}$).
\textbf{3. Stopping Potential ($V_0$):}
The negative retarding collector potential required to halt the most energetic photoelectrons:
\begin{equation}
  K_{\max} = e V_0 \implies \mathbf{V_0 = \left(\frac{h}{e}\right)\nu - \frac{\Phi_0}{e} = \left(\frac{h}{e}\right)(\nu - \nu_0)}
\end{equation}
\textit{Universal Graph Property:} In the $V_0$ vs. $\nu$ straight-line graph, the slope is a universal constant:
\begin{equation}
  \mathbf{\text{Slope} = \frac{h}{e} \approx 4.14 \times 10^{-15}\,\mathrm{V}\cdot\mathrm{s}}
\end{equation}
The horizontal axis intercept is threshold frequency $\nu_0$; vertical axis intercept is $-\Phi_0/e$.
\end{conceptbox}

\begin{formulabox}[de Broglie Matter Waves: Particle-Wave Duality Expressions]
Louis de Broglie postulated that any moving material particle of mass $m$ and velocity $v$ possesses an associated matter wave:
\begin{equation}
  \mathbf{\lambda = \frac{h}{p} = \frac{h}{m v} = \frac{h}{\sqrt{2 m K}}}
\end{equation}
where $K$ is kinetic energy.
\textbf{1. De Broglie Wavelength of Charged Particles Accelerated through Potential $V$:}
$K = q V \implies \lambda = \frac{h}{\sqrt{2 m q V}}$.
\begin{itemize}
  \item \textbf{Electron ($m = 9.109 \times 10^{-31}\,\mathrm{kg}, q = e$):}
  \begin{equation}
    \mathbf{\lambda_e = \frac{1.227}{\sqrt{V}}\,\mathrm{nm} = \frac{12.27}{\sqrt{V}}\,\text{\AA} = \sqrt{\frac{150}{V}}\,\text{\AA}}
  \end{equation}
  \item \textbf{Proton ($m_p = 1.673 \times 10^{-27}\,\mathrm{kg}, q = e$):}
  \begin{equation}
    \mathbf{\lambda_p = \frac{0.286}{\sqrt{V}}\,\text{\AA}}
  \end{equation}
  \item \textbf{Deuteron ($m_d = 2m_p, q = e$):}
  \begin{equation}
    \mathbf{\lambda_d = \frac{0.202}{\sqrt{V}}\,\text{\AA}}
  \end{equation}
  \item \textbf{Alpha Particle ($m_\alpha = 4m_p, q = 2e$):}
  \begin{equation}
    \mathbf{\lambda_\alpha = \frac{0.101}{\sqrt{V}}\,\text{\AA}}
  \end{equation}
\end{itemize}
\textbf{2. Thermal Neutrons in Equilibrium at Temperature $T$:}
Mean translational kinetic energy $K = \frac{3}{2}k_B T$:
\begin{equation}
  \mathbf{\lambda_{\text{thermal}} = \frac{h}{\sqrt{3 m k_B T}}}
\end{equation}
\end{formulabox}

\begin{trapbox}[Syllabus Boundary Alert: Deletion of Davisson-Germer Experiment]
\textbf{NMC REVISION WARNING FOR NEET UG:}
Under the revised NEET UG curriculum:
\begin{itemize}
  \item \textbf{STRICTLY DELETED:} The \textbf{Davisson-Germer Experiment} (nickel crystal electron diffraction, Bragg peak at $54\,\mathrm{V}$ and $\theta = 50^\circ$, accelerating electron gun details) has been \textbf{formally removed}.
  \item \textbf{EXAM ORIENTATION:} Candidates must be proficient in Einstein's photoelectric equation, photon momentum calculations, and de Broglie wavelength formulas for accelerated particles without studying the experimental apparatus of Davisson and Germer.
\end{itemize}
\end{trapbox}

\section{Atomic Physics: Bohr's Postulates, Spectral Series \& Energy Levels}

\begin{conceptbox}[Rutherford Scattering \& Bohr's Atomic Postulates]
\textbf{1. Rutherford Scattering (Geiger-Marsden Experiment):}
Distance of closest approach $r_0$ for an $\alpha$-particle ($q_\alpha = 2e$) directed head-on with kinetic energy $K$ toward a nucleus of atomic number $Z$:
\begin{equation}
  K = \frac{1}{4\pi\varepsilon_0}\frac{(2e)(Ze)}{r_0} \implies \mathbf{r_0 = \frac{1}{4\pi\varepsilon_0}\frac{2 Z e^2}{K}}
\end{equation}
Impact parameter: $b = \frac{1}{4\pi\varepsilon_0}\frac{Z e^2 \cot(\theta/2)}{K}$.
\textbf{2. Bohr's Foundational Postulates for Hydrogen-Like Species:}
\begin{enumerate}
  \item \textbf{Stable Stationary Orbits:} An orbital electron circles the positive nucleus solely under the Coulomb attraction without radiating electromagnetic energy, despite continuous centripetal acceleration.
  \item \textbf{Quantization of Orbital Angular Momentum:} Stable orbits are restricted strictly to those where orbital angular momentum is an integer multiple of $\hbar = h / 2\pi$:
  \begin{equation}
    \mathbf{L = m_e v r = \frac{n h}{2\pi}} \quad (n = 1, 2, 3, \dots)
  \end{equation}
  \item \textbf{Frequency Condition (Radiative Transitions):} Photons are emitted or absorbed exclusively when an electron undergoes a discrete transition between stationary energy levels $E_i$ and $E_f$:
  \begin{equation}
    \mathbf{\Delta E = E_i - E_f = h\nu = \frac{hc}{\lambda}}
  \end{equation}
\end{enumerate}
\end{conceptbox}

\begin{formulabox}[Bohr Atom Formulas: Rigorous Analytical Derivations for Hydrogen-Like Species]
For a hydrogen-like atom of nuclear charge $+Ze$ and electron mass $m$:
\begin{enumerate}
  \item \textbf{Electrostatic Centripetal Equilibrium:}
  \begin{equation}
    \frac{m v^2}{r} = \frac{1}{4\pi\varepsilon_0} \frac{Z e^2}{r^2} \implies m v^2 r = \frac{Z e^2}{4\pi\varepsilon_0} \tag{Eq. 1}
  \end{equation}
  \item \textbf{Angular Momentum Quantization:}
  \begin{equation}
    m v r = \frac{n h}{2\pi} \implies v = \frac{n h}{2\pi m r} \tag{Eq. 2}
  \end{equation}
  \item \textbf{Orbital Radius Derivation:} Substituting Eq. 2 into Eq. 1:
  \begin{equation}
    m \left(\frac{n^2 h^2}{4\pi^2 m^2 r^2}\right) r = \frac{Z e^2}{4\pi\varepsilon_0} \implies \mathbf{r_n = \frac{\varepsilon_0 n^2 h^2}{\pi m Z e^2} = r_0 \frac{n^2}{Z} = 0.529 \frac{n^2}{Z}\,\text{\AA}}
  \end{equation}
  where $r_0 = 0.529177\,\text{\AA}$ is the first Bohr radius of hydrogen ($n=1, Z=1$).
  \item \textbf{Orbital Velocity Derivation:}
  \begin{equation}
    \mathbf{v_n = \frac{Z e^2}{2\varepsilon_0 n h} = v_0 \frac{Z}{n} = \left(2.18 \times 10^6\,\mathrm{m/s}\right)\frac{Z}{n} = c \alpha \frac{Z}{n}}
  \end{equation}
  where $\alpha = \frac{e^2}{2\varepsilon_0 h c} \approx \frac{1}{137}$ is the fine-structure constant.
  \item \textbf{Energies in the $n$-th Stationary Orbit:}
  \begin{itemize}
    \item \textbf{Kinetic Energy:} $\mathbf{K_n = \frac{1}{2}m v_n^2 = \frac{m Z^2 e^4}{8\varepsilon_0^2 n^2 h^2} = +13.6 \frac{Z^2}{n^2}\,\mathrm{eV}}$.
    \item \textbf{Potential Energy:} $\mathbf{U_n = -\frac{1}{4\pi\varepsilon_0}\frac{Z e^2}{r_n} = -\frac{m Z^2 e^4}{4\varepsilon_0^2 n^2 h^2} = -27.2 \frac{Z^2}{n^2}\,\mathrm{eV}}$.
    \item \textbf{Total Mechanical Energy:}
    \begin{equation}
      \mathbf{E_n = K_n + U_n = -\frac{m Z^2 e^4}{8\varepsilon_0^2 n^2 h^2} = -13.6 \frac{Z^2}{n^2}\,\mathrm{eV}}
    \end{equation}
    \item \textbf{Golden Energy Ratios:}
    \begin{equation}
      \mathbf{E_n = -K_n = \frac{U_n}{2}}
    \end{equation}
  \end{itemize}
\end{enumerate}
\end{formulabox}

\begin{formulabox}[Hydrogen Spectral Emission Series \& Rydberg Equation]
When an electron drops from an upper stationary orbit $n_2$ to a lower orbit $n_1$ ($n_2 > n_1$):
\begin{equation}
  \Delta E = E_{n_2} - E_{n_1} = 13.6 Z^2 \left(\frac{1}{n_1^2} - \frac{1}{n_2^2}\right)\,\mathrm{eV}
\end{equation}
\begin{equation}
  \mathbf{\frac{1}{\lambda} = \overline{\nu} = R_\infty Z^2 \left(\frac{1}{n_1^2} - \frac{1}{n_2^2}\right)}
\end{equation}
where $R_\infty = \frac{m e^4}{8\varepsilon_0^2 c h^3} \approx 1.09737 \times 10^7\,\mathrm{m}^{-1}$ is the Rydberg constant ($\frac{1}{R_\infty} \approx 912\,\text{\AA}$).
\begin{center}
\begin{tabularx}{\linewidth}{|l|c|c|c|Y|}
\hline
\textbf{Spectral Series} & \textbf{Lower Level ($n_1$)} & \textbf{Upper Level ($n_2$)} & \textbf{Spectral Domain} & \textbf{Wavelength Range} \\
\hline
\textbf{Lyman} & $1$ & $2, 3, 4, \dots$ & \textbf{Ultraviolet (UV)} & $\lambda_{\min} = 91.2\,\mathrm{nm}, \lambda_{\max} = 121.6\,\mathrm{nm}$ \\
\hline
\textbf{Balmer} & $2$ & $3, 4, 5, \dots$ & \textbf{Visible} (part UV) & $\lambda_{\min} = 364.6\,\mathrm{nm}, \lambda_{\max} = 656.3\,\mathrm{nm}$ \\
\hline
\textbf{Paschen} & $3$ & $4, 5, 6, \dots$ & \textbf{Near Infrared} & $\lambda_{\min} = 820.4\,\mathrm{nm}, \lambda_{\max} = 1875\,\mathrm{nm}$ \\
\hline
\textbf{Brackett} & $4$ & $5, 6, 7, \dots$ & \textbf{Mid Infrared} & $\lambda_{\min} = 1458\,\mathrm{nm}, \lambda_{\max} = 4051\,\mathrm{nm}$ \\
\hline
\textbf{Pfund} & $5$ & $6, 7, 8, \dots$ & \textbf{Far Infrared} & $\lambda_{\min} = 2279\,\mathrm{nm}, \lambda_{\max} = 7458\,\mathrm{nm}$ \\
\hline
\end{tabularx}
\end{center}
\textbf{Total Spectral Lines Emitted from Excited Level $n$ to Ground State:}
\begin{equation}
  \mathbf{N = \frac{n(n - 1)}{2}}
\end{equation}
\end{formulabox}

\section{Nuclear Physics: Composition, Density, Mass Defect \& Binding Energy}

\begin{conceptbox}[Nuclear Dimensions, Universal Nuclear Density \& Strong Nuclear Force]
\textbf{1. Nuclear Constituents:}
Atomic nucleus contains $Z$ protons and $N = A - Z$ neutrons (nucleons).
\textbf{2. Nuclear Radius Empirical Law:}
\begin{equation}
  \mathbf{R = R_0 A^{1/3}}
\end{equation}
where $R_0 \approx 1.2 \times 10^{-15}\,\mathrm{m} = 1.2\,\mathrm{fm}$.
\textbf{3. Proof of Constant Nuclear Density ($\rho_{\text{nuc}}$):}
Nuclear volume: $V = \frac{4}{3}\pi R^3 = \frac{4}{3}\pi (R_0 A^{1/3})^3 = \frac{4}{3}\pi R_0^3 A$.
Nuclear mass: $M \approx A \cdot m_N$ (where nucleon mass $m_N \approx 1.66 \times 10^{-27}\,\mathrm{kg}$).
\begin{equation}
  \mathbf{\rho_{\text{nuc}} = \frac{M}{V} = \frac{A m_N}{\frac{4}{3}\pi R_0^3 A} = \frac{3 m_N}{4\pi R_0^3} \approx 2.3 \times 10^{17}\,\mathrm{kg/m}^3}
\end{equation}
Nuclear density is an \textbf{absolute universal constant}, completely independent of mass number $A$.
\textbf{4. Properties of the Strong Nuclear Force:}
\begin{itemize}
  \item Strongest fundamental force in nature ($\sim 100$ times stronger than electrostatic repulsion).
  \item Strictly short-ranged: active solely across $1 - 2\,\mathrm{fm}$. For $r < 0.8\,\mathrm{fm}$, it becomes strongly repulsive.
  \item Charge-independent: $F_{p-p} = F_{p-n} = F_{n-n}$ under identical quantum states.
  \item Saturated: each nucleon interacts exclusively with its immediate neighbouring nucleons.
\end{itemize}
\end{conceptbox}

\begin{formulabox}[Mass Defect, Nuclear Binding Energy Curve \& Energy Reactions]
\textbf{1. Mass Defect ($\Delta m$):}
The rest mass of any stable bound nucleus is strictly less than the combined rest masses of its constituent free protons and neutrons:
\begin{equation}
  \mathbf{\Delta m = [Z m_p + (A - Z) m_n] - M_{\text{nucleus}}(A, Z)}
\end{equation}
\textbf{2. Nuclear Binding Energy ($E_b$):}
By Einstein's mass-energy equivalence ($E = \Delta m c^2$):
\begin{equation}
  \mathbf{E_b = \Delta m c^2 = \Delta m\,[\mathrm{u}] \times 931.5\,\mathrm{MeV}}
\end{equation}
where $1\,\mathrm{u} \equiv 931.494\,\mathrm{MeV}/c^2$.
\textbf{3. Binding Energy per Nucleon ($E_b/A$ Curve):}
The stability of a nucleus is dictated strictly by binding energy per nucleon ($E_b/A$):
\begin{itemize}
  \item Rapid increase from light nuclei ($^2\mathrm{H} \approx 1.1\,\mathrm{MeV/nucleon}$) with local peaks at tightly bound alpha-multiples ($^4\mathrm{He}, ^{12}\mathrm{C}, ^{16}\mathrm{O}$).
  \item Broad flat maximum in intermediate mass range $30 < A < 120$ with peak at $\mathbf{^{56}\mathrm{Fe} \approx 8.75\,\mathrm{MeV/nucleon}}$.
  \item Gradual drop toward heavy nuclei due to cumulative Coulomb repulsion between protons, dropping to $\approx 7.6\,\mathrm{MeV/nucleon}$ for $^{238}\mathrm{U}$.
  \item \textbf{Nuclear Fission:} Heavy nucleus splitting ($A \sim 240 \to A_1 + A_2 \sim 120$) gains $\approx 0.9\,\mathrm{MeV/nucleon} \implies Q \approx 240 \times 0.9 \approx \mathbf{200\,\mathrm{MeV}}$ per fission.
  \item \textbf{Nuclear Fusion:} Two very light nuclei fusing into a heavier nucleus ($^2\mathrm{H} + ^3\mathrm{H} \to ^4\mathrm{He} + n$) releases $\approx \mathbf{17.6\,\mathrm{MeV}}$ ($3.5\,\mathrm{MeV/nucleon}$).
\end{itemize}
\end{formulabox}

\begin{trapbox}[Syllabus Boundary Alert: Complete Deletion of Radioactivity]
\textbf{NMC REVISION WARNING FOR NEET UG:}
Under the revised NEET UG syllabus issued by the National Medical Commission:
\begin{itemize}
  \item \textbf{STRICTLY DELETED:} The entire subtopic of \textbf{Radioactivity} has been \textbf{completely eliminated}.
  \item Deleted topics include: $\alpha$-decay, $\beta$-decay ($\beta^-$ and $\beta^+$), $\gamma$-decay, Radioactive Decay Law ($N(t) = N_0 e^{-\lambda t}$), Half-life ($T_{1/2} = \frac{0.693}{\lambda}$), Mean Life ($\tau = \frac{1}{\lambda}$), and Activity ($A = \lambda N$).
  \item \textbf{EXAM ORIENTATION:} Candidates must focus on Mass Defect, Binding Energy per Nucleon, Nuclear Density, and Fission/Fusion $Q$-values without solving radioactive decay kinematics.
\end{itemize}
\end{trapbox}

\section{Semiconductor Physics, Diodes \& Digital Logic Gates}

\begin{conceptbox}[Energy Bands, Intrinsic \& Extrinsic Semiconductors]
\textbf{1. Energy Band Theory of Solids:}
\begin{itemize}
  \item \textbf{Conductors:} Conduction band and valence band overlap, or conduction band is partially filled ($E_g = 0$). High electron density $n \sim 10^{28}\,\mathrm{m}^{-3}$.
  \item \textbf{Insulators:} Large forbidden energy gap: $\mathbf{E_g > 3.0\,\mathrm{eV}}$ (e.g., Diamond $E_g \approx 5.4\,\mathrm{eV}$).
  \item \textbf{Semiconductors:} Narrow forbidden energy gap: $\mathbf{E_g \le 3.0\,\mathrm{eV}}$ ($\mathrm{Si}: E_g = 1.1\,\mathrm{eV}$; $\mathrm{Ge}: E_g = 0.72\,\mathrm{eV}$; $\mathrm{GaAs}: E_g = 1.42\,\mathrm{eV}$).
\end{itemize}
\textbf{2. Intrinsic Semiconductors (Pure Crystals):}
Thermal energy excites valence electrons across $E_g$ into the conduction band, leaving equal numbers of vacant states (\textbf{holes}) in the valence band:
\begin{equation}
  \mathbf{n_e = n_h = n_i}
\end{equation}
Total electrical conductivity: $\mathbf{\sigma = e(n_e \mu_e + n_h \mu_h)}$.
\textbf{3. Extrinsic Semiconductors (Doped Crystals):}
\begin{itemize}
  \item \textbf{n-Type:} Doped with pentavalent donor impurities (P, As, Sb). Introduces donor discrete energy level $E_d$ just below conduction band ($0.05\,\mathrm{eV}$ for Si).
  \begin{equation}
    \mathbf{n_e \approx N_D \gg n_h} \quad (\text{Electrons: majority carriers; Holes: minority carriers})
  \end{equation}
  \item \textbf{p-Type:} Doped with trivalent acceptor impurities (B, Al, In, Ga). Introduces acceptor discrete energy level $E_a$ just above valence band ($0.05\,\mathrm{eV}$ for Si).
  \begin{equation}
    \mathbf{n_h \approx N_A \gg n_e} \quad (\text{Holes: majority carriers; Electrons: minority carriers})
  \end{equation}
  \item \textbf{Law of Mass Action (Thermal Equilibrium):}
  \begin{equation}
    \mathbf{n_e \cdot n_h = n_i^2(T)}
  \end{equation}
  Both n-type and p-type semiconductors are strictly \textbf{electrically neutral} overall.
\end{itemize}
\end{conceptbox}

\begin{formulabox}[p-n Junction Diode Physics, Biasing \& Rectifiers]
\textbf{1. p-n Junction Formation:}
Majority holes from p-side diffuse into n-side, and majority electrons from n-side diffuse into p-side.
Uncompensated immobile donor ions ($+$) on n-side and acceptor ions ($-$) on p-side form a \textbf{depletion layer} and built-in \textbf{barrier potential} ($V_0 \approx 0.7\,\mathrm{V}$ for Si, $0.3\,\mathrm{V}$ for Ge) pointing from n to p.
\textbf{2. Biasing Modes:}
\begin{itemize}
  \item \textbf{Forward Bias (p to $+$, n to $-$):} External field opposes barrier field $\implies$ barrier potential lowers to $(V_0 - V)$, depletion layer narrows, forward current rises exponentially: $I \approx I_0 (e^{eV/\eta k_B T} - 1)$.
  \item \textbf{Reverse Bias (p to $-$, n to $+$):} External field aids barrier field $\implies$ barrier increases to $(V_0 + V)$, depletion layer widens, current is tiny minority saturation current $I_0 \sim \mu\mathrm{A}$ (Si: $\mathrm{nA}$).
\end{itemize}
\textbf{3. Diode Rectifiers:}
\begin{itemize}
  \item \textbf{Half-Wave Rectifier:} One diode conducting during positive AC half-cycles.
  \begin{equation}
    \mathbf{\eta_{\max} = \frac{0.406 R_L}{r_f + R_L} \approx 40.6\%}, \qquad \mathbf{\nu_{\text{ripple}} = \nu_{\text{AC}}}
  \end{equation}
  \item \textbf{Full-Wave Rectifier (Center-Tapped or Bridge):} Two or four diodes conducting alternately.
  \begin{equation}
    \mathbf{\eta_{\max} = \frac{0.812 R_L}{r_f + R_L} \approx 81.2\%}, \qquad \mathbf{\nu_{\text{ripple}} = 2\nu_{\text{AC}}}
  \end{equation}
\end{itemize}
\end{formulabox}

\begin{formulabox}[Special Purpose Diodes: Zener, Photodiode, LED \& Solar Cell]
\textbf{1. Zener Diode as a Voltage Regulator:}
Heavily doped p-n junction with an ultra-thin depletion layer ($\sim 10^{-6}\,\mathrm{m}$) that undergoes sharp reverse breakdown at a fixed Zener breakdown voltage $V_z$.
Connected in \textbf{reverse bias} parallel with fluctuating load $R_L$:
\begin{equation}
  V_{\text{out}} = V_z = \text{constant}, \qquad I_{\text{supply}} = I_z + I_L
\end{equation}
Series current-limiting resistor $R_s$:
\begin{equation}
  \mathbf{R_s = \frac{V_{\text{unregulated}} - V_z}{I_z + I_L}}
\end{equation}
\textbf{2. Photodiode:} Operates in \textbf{reverse bias}. Incident photons ($h\nu > E_g$) generate electron-hole pairs in the depletion region. Fractional change in minority reverse current is much larger than in forward current, making reverse bias ideal for optical detection.
\textbf{3. Light Emitting Diode (LED):} Operates in \textbf{forward bias}. Electrons recombine with holes at junction, releasing recombination energy as spontaneous photons across direct bandgap: $\mathbf{\lambda = \frac{hc}{E_g}}$.
\textbf{4. Solar Cell:} Unbiased p-n junction operating in the \textbf{fourth quadrant} of $I-V$ characteristic, converting optical solar flux into electrical power.
\end{formulabox}

\begin{formulabox}[Digital Logic Gates: Truth Tables \& De Morgan's Theorems]
\textbf{1. Fundamental Logic Gates:}
\begin{itemize}
  \item \textbf{AND Gate:} $Y = A \cdot B$ (Output is 1 strictly when both inputs are 1).
  \item \textbf{OR Gate:} $Y = A + B$ (Output is 1 when at least one input is 1).
  \item \textbf{NOT Gate (Inverter):} $Y = \overline{A}$ ($0 \to 1, 1 \to 0$).
\end{itemize}
\textbf{2. Universal Logic Gates (Can Synthesize Any Boolean Function):}
\begin{itemize}
  \item \textbf{NAND Gate:} $Y = \overline{A \cdot B}$.
  \item \textbf{NOR Gate:} $Y = \overline{A + B}$.
\end{itemize}
\textbf{3. De Morgan's Master Theorems:}
\begin{equation}
  \mathbf{\overline{A + B} = \overline{A} \cdot \overline{B}}, \qquad \mathbf{\overline{A \cdot B} = \overline{A} + \overline{B}}
\end{equation}
\end{formulabox}

\begin{trapbox}[Syllabus Boundary Alert: Complete Deletion of Transistors]
\textbf{NMC REVISION WARNING FOR NEET UG:}
Under the revised NEET UG syllabus issued by the National Medical Commission:
\begin{itemize}
  \item \textbf{STRICTLY DELETED:} The entire topic of \textbf{Transistors} (bipolar junction transistors npn and pnp, emitter, base, collector actions) has been \textbf{completely eliminated}.
  \item Deleted subtopics include: Common Emitter (CE) input and output characteristics, current gain factors ($\alpha = I_c/I_e, \beta = I_c/I_b, \beta = \frac{\alpha}{1-\alpha}$), transistor as a CE amplifier (voltage gain $A_v = -\beta \frac{R_L}{R_i}$), transistor as a digital switch, and transistor as an oscillator.
  \item \textbf{EXAM ORIENTATION:} Candidates must focus exclusively on Semiconductors, p-n junction diodes, rectifiers, Zener diode regulators, optoelectronic diodes (LED, photodiode, solar cell), and Digital Logic Gates.
\end{itemize}
\end{trapbox}

\begin{ncertbox}[NCERT High-Yield Extraction: Optoelectronic Materials \& Eye Sensitivity]
The semiconductor bandgap for visible LEDs must exceed $1.8\,\mathrm{eV}$ ($hc/\lambda_{\max} \approx 1240 / 700\,\mathrm{nm} \approx 1.8\,\mathrm{eV}$).
\begin{itemize}
  \item \textbf{Gallium Arsenide Phosphide ($\mathrm{GaAs}_{1-x}\mathrm{P}_x$):} Used for red, amber, and yellow LEDs ($E_g \approx 1.9\,\mathrm{eV}$).
  \item \textbf{Gallium Nitride ($\mathrm{GaN}$):} Direct bandgap of $3.4\,\mathrm{eV}$ enables bright blue and ultraviolet LEDs.
  \item \textbf{Photodiode Detection Limit:} To detect an optical signal of wavelength $\lambda$, the semiconductor bandgap must satisfy $\mathbf{E_g \le \frac{hc}{\lambda}}$.
\end{itemize}
\end{ncertbox}

\section{Master High-Yield Problem Archetypes \& Worked Solutions}

\begin{usecasebox}[Archetype 1: Silvered Plano-Convex Lens Equivalent Mirror Focal Length]
\textbf{Problem:} A plano-convex lens of refractive index $\mu = 1.5$ has a curved surface of radius of curvature $R = 20\,\mathrm{cm}$. Compute the equivalent focal length when: (i) the flat plane surface is silvered, and (ii) the curved surface is silvered.
\textbf{Step-by-Step Solution:}
\begin{enumerate}
  \item Power of the plano-convex lens in air:
  \begin{equation}
    P_L = \frac{1}{f_L} = (\mu - 1)\left(\frac{1}{R_1} - \frac{1}{R_2}\right) = (1.5 - 1)\left(\frac{1}{20} - 0\right) = \frac{0.5}{20} = \frac{1}{40}\,\mathrm{cm}^{-1}
  \end{equation}
  \item Any silvered lens system acts as an equivalent mirror of power:
  \begin{equation}
    P_{\text{eq}} = 2 P_L + P_M \implies \mathbf{F_{\text{eq}} = -\frac{1}{P_{\text{eq}}}}
  \end{equation}
  \item \textbf{Case (i): Flat Plane Surface is Silvered:}
  Plane mirror has $R_M = \infty \implies P_M = -\frac{1}{f_M} = 0$.
  \begin{equation}
    P_{\text{eq}} = 2 P_L + 0 = 2 \left(\frac{1}{40}\right) = \frac{1}{20}\,\mathrm{cm}^{-1} \implies \mathbf{F_{\text{eq}} = -20\,\mathrm{cm}}
  \end{equation}
  Acts as a \textbf{concave mirror of focal length $20\,\mathrm{cm}$}.
  \item \textbf{Case (ii): Curved Surface is Silvered:}
  Concave mirror formed by silvered curved face has $R_M = 20\,\mathrm{cm} \implies f_M = -10\,\mathrm{cm} \implies P_M = -\frac{1}{f_M} = +\frac{1}{10}\,\mathrm{cm}^{-1}$.
  \begin{equation}
    P_{\text{eq}} = 2 P_L + P_M = 2 \left(\frac{1}{40}\right) + \frac{1}{10} = \frac{1}{20} + \frac{1}{10} = \frac{3}{20}\,\mathrm{cm}^{-1} \implies \mathbf{F_{\text{eq}} = -\frac{20}{3}\,\mathrm{cm} \approx -6.67\,\mathrm{cm}}
  \end{equation}
  Acts as a \textbf{concave mirror of focal length $6.67\,\mathrm{cm}$}.
\end{enumerate}
\end{usecasebox}

\begin{usecasebox}[Archetype 2: Multi-Wavelength Bright Fringe Coincidence in YDSE]
\textbf{Problem:} In a YDSE setup with slit separation $d = 1.0\,\mathrm{mm}$ and screen distance $D = 1.0\,\mathrm{m}$, light containing two wavelengths $\lambda_1 = 600\,\mathrm{nm}$ and $\lambda_2 = 480\,\mathrm{nm}$ is used. Determine the minimum non-zero distance from the central maximum where bright fringes of both wavelengths coincide.
\textbf{Step-by-Step Solution:}
\begin{enumerate}
  \item For bright fringes to coincide at position $y$:
  \begin{equation}
    y = \frac{n_1 \lambda_1 D}{d} = \frac{n_2 \lambda_2 D}{d} \implies n_1 \lambda_1 = n_2 \lambda_2
  \end{equation}
  \item Taking the ratio of integers:
  \begin{equation}
    \frac{n_1}{n_2} = \frac{\lambda_2}{\lambda_1} = \frac{480\,\mathrm{nm}}{600\,\mathrm{nm}} = \frac{4}{5}
  \end{equation}
  \item Smallest non-zero integers satisfying this condition: $n_1 = 4$ and $n_2 = 5$.
  \item Minimum distance $y_{\min}$:
  \begin{equation}
    y_{\min} = \frac{4 \times (600 \times 10^{-9}\,\mathrm{m}) \times 1.0\,\mathrm{m}}{1.0 \times 10^{-3}\,\mathrm{m}} = 2.4 \times 10^{-3}\,\mathrm{m} = \mathbf{2.4\,\mathrm{mm}}
  \end{equation}
  The 4th bright fringe of $\lambda_1$ coincides with the 5th bright fringe of $\lambda_2$ at $y = 2.4\,\mathrm{mm}$.
\end{enumerate}
\end{usecasebox}

\begin{usecasebox}[Archetype 3: Zener Diode Voltage Regulator Design Calculations]
\textbf{Problem:} A Zener diode with breakdown voltage $V_z = 10\,\mathrm{V}$ and maximum power dissipation $P_{z,\max} = 1.0\,\mathrm{W}$ is used to regulate a DC supply that fluctuates between $15\,\mathrm{V}$ and $25\,\mathrm{V}$. The load resistance is $R_L = 1.0\,\mathrm{k}\Omega$. Calculate: (i) the load current $I_L$, and (ii) the suitable series resistance $R_s$ to ensure the Zener current remains within safe limits.
\textbf{Step-by-Step Solution:}
\begin{enumerate}
  \item The regulated voltage across the load is fixed at $V_L = V_z = 10\,\mathrm{V}$.
  \item Constant load current:
  \begin{equation}
    I_L = \frac{V_z}{R_L} = \frac{10\,\mathrm{V}}{1000\,\Omega} = 10\,\mathrm{mA} = 0.01\,\mathrm{A}
  \end{equation}
  \item Maximum safe Zener diode current:
  \begin{equation}
    I_{z,\max} = \frac{P_{z,\max}}{V_z} = \frac{1.0\,\mathrm{W}}{10\,\mathrm{V}} = 100\,\mathrm{mA} = 0.10\,\mathrm{A}
  \end{equation}
  \item Total maximum supply current: $I_{\text{total},\max} = I_{z,\max} + I_L = 100\,\mathrm{mA} + 10\,\mathrm{mA} = 110\,\mathrm{mA}$.
  \item At maximum input voltage $V_{\text{in}} = 25\,\mathrm{V}$:
  \begin{equation}
    R_{s,\min} = \frac{V_{\text{in},\max} - V_z}{I_{\text{total},\max}} = \frac{25 - 10}{0.11} = \frac{15}{0.11} \approx \mathbf{136.4\,\Omega}
  \end{equation}
  \item At minimum input voltage $V_{\text{in}} = 15\,\mathrm{V}$ with minimum holding current $I_{z,\min} \approx 5\,\mathrm{mA}$:
  \begin{equation}
    R_{s,\max} = \frac{V_{\text{in},\min} - V_z}{I_{z,\min} + I_L} = \frac{15 - 10}{0.005 + 0.010} = \frac{5}{0.015} \approx \mathbf{333.3\,\Omega}
  \end{equation}
  A standard resistance of $R_s \approx 200\,\Omega$ provides robust regulation across the entire input range.
\end{enumerate}
\end{usecasebox}

% =========================================================================
% PART V: EXPERIMENTAL SKILLS & PRACTICAL PHYSICS
% =========================================================================
\part{Experimental Skills \& Practical Physics}

\chapter{Practical Physics \& Prescribed Experimental Lab Skills (NMC Unit 20)}

Under the revised NMC guidelines, experimental physics forms a dedicated unit tested directly in both Section A and Section B. This chapter provides a rigorous, step-by-step analytical manual for all 18 prescribed laboratory experiments, focusing on least count derivations, zero errors, reading formulas, experimental precautions, and high-frequency exam traps.

\section{Vernier Calipers \& Screw Gauge}

\begin{formulabox}[Vernier Calipers: Least Count, Zero Error \& Measurement]
\textbf{1. Least Count (Vernier Constant, VC):}
Smallest distance measurable by the instrument. If $n$ Vernier Scale Divisions (VSD) equal $m$ Main Scale Divisions (MSD):
\begin{equation}
  \mathbf{\mathrm{LC} = 1\,\mathrm{MSD} - 1\,\mathrm{VSD} = \left(1 - \frac{m}{n}\right)\mathrm{MSD}}
\end{equation}
\textit{Standard Metric Caliper:} $10\,\mathrm{VSD} = 9\,\mathrm{MSD}$ with $1\,\mathrm{MSD} = 1\,\mathrm{mm} = 0.1\,\mathrm{cm}$:
\begin{equation}
  \mathrm{LC} = 1\,\mathrm{mm} - \frac{9}{10}\,\mathrm{mm} = \mathbf{0.1\,\mathrm{mm} = 0.01\,\mathrm{cm}}
\end{equation}
\textbf{2. True Reading Formulation \& Zero Error Corrections:}
\begin{equation}
  \mathbf{\text{Observed Reading} = \mathrm{MSR} + (\mathrm{VSR} \times \mathrm{LC})} \qquad \mathbf{\text{True Reading} = \text{Observed Reading} - (\text{Zero Error})}
\end{equation}
\begin{itemize}
  \item \textbf{Positive Zero Error ($+e$):} Vernier zero lies to the \textbf{right of main scale zero} when jaws are closed.
  \begin{equation}
    e = +(n \times \mathrm{LC}) \implies \text{True Reading} = \text{Observed} - e \quad (\text{Must be subtracted})
  \end{equation}
  \item \textbf{Negative Zero Error ($-e$):} Vernier zero lies to the \textbf{left of main scale zero} when jaws are closed.
  \begin{equation}
    e = -(N - n)\times \mathrm{LC} \implies \text{True Reading} = \text{Observed} - (-e) = \text{Observed} + |e| \quad (\text{Must be added})
  \end{equation}
\end{itemize}
\end{formulabox}

\begin{formulabox}[Screw Gauge (Micrometer): Pitch \& Backlash Error]
\textbf{1. Pitch \& Least Count (LC):}
\begin{equation}
  \mathbf{\text{Pitch}} = \frac{\text{Distance moved on linear pitch scale}}{\text{Number of complete rotations of circular cap}} \qquad (\text{Usually } 1\,\mathrm{mm} \text{ or } 0.5\,\mathrm{mm})
\end{equation}
\begin{equation}
  \mathbf{\mathrm{LC}} = \frac{\text{Pitch}}{\text{Total number of circular scale divisions } (N)} = \frac{1\,\mathrm{mm}}{100} = \mathbf{0.01\,\mathrm{mm} = 0.001\,\mathrm{cm} = 10\,\mu\mathrm{m}}
\end{equation}
\textbf{2. True Thickness / Diameter Calculation:}
\begin{equation}
  \mathbf{\text{True Reading} = \mathrm{PSR} + (\mathrm{CSR} \times \mathrm{LC}) - (\text{Zero Error})}
\end{equation}
\begin{itemize}
  \item \textbf{Positive Zero Error:} Circular zero is \textbf{below the reference line} when studs touch $\implies e = +(n \times \mathrm{LC})$.
  \item \textbf{Negative Zero Error:} Circular zero is \textbf{above the reference line} when studs touch $\implies e = -(N - n) \times \mathrm{LC}$.
  \item \textbf{Backlash Error:} Due to loose thread fitting between screw and nut. Avoided by turning screw \textbf{strictly in one direction only}.
\end{itemize}
\end{formulabox}

\section{Mechanics \& Hydrodynamics Lab Experiments}

\begin{conceptbox}[Simple Pendulum, Metre Scale \& Elasticity (Searle's Apparatus)]
\textbf{1. Simple Pendulum ($L \text{ vs } T^2$ Graph):}
Time period: $T = 2\pi\sqrt{\frac{L}{g}} \implies T^2 = \left(\frac{4\pi^2}{g}\right)L$.
\begin{itemize}
  \item Graph of $L$ versus $T^2$ is a \textbf{straight line passing through the origin}.
  \begin{equation}
    \mathbf{g = 4\pi^2 \left(\frac{L}{T^2}\right) = \frac{4\pi^2}{\text{Slope of } T^2 \text{ vs } L \text{ graph}}}
  \end{equation}
  \item \textbf{Effective Length:} $L = l + h + r$ ($l$ = wire length, $h$ = hook length, $r$ = bob radius).
  \item Angular amplitude must be small ($\theta \le 4^\circ$) so that $\sin\theta \approx \theta$.
\end{itemize}
\textbf{2. Metre Scale (Principle of Moments / Determination of Mass):}
A balanced uniform metre scale obeys rotational equilibrium: $\sum \tau_{\mathrm{ext}} = 0$.
\begin{equation}
  m_1 g (x_{\mathrm{knife}} - x_1) = M g (x_{\mathrm{knife}} - x_{\mathrm{cg}}) + m_2 g (x_2 - x_{\mathrm{knife}})
\end{equation}
Where $M$ is mass of metre scale, $x_{\mathrm{cg}}$ is center of gravity (usually $50\,\mathrm{cm}$ mark for uniform scale).

\textbf{3. Young's Modulus of Elasticity (Searle's Apparatus):}
Two identical wires (reference and experimental) avoid thermal expansion errors:
\begin{equation}
  \mathbf{Y = \frac{\text{Stress}}{\text{Strain}} = \frac{M g L}{\pi r^2 l}} \implies \text{Elongation } l = \left(\frac{g L}{\pi r^2 Y}\right) M
\end{equation}
Graph of elongation $l$ versus suspended mass $M$ is a straight line. $Y = \frac{g L}{\pi r^2 (\text{Slope})}$.
\end{conceptbox}

\begin{conceptbox}[Surface Tension, Viscosity \& Speed of Sound in Air]
\textbf{1. Surface Tension of Water by Capillary Rise Method:}
\begin{equation}
  \mathbf{T = \frac{r (h + r/3) \rho g}{2 \cos\theta}} \approx \mathbf{\frac{r h \rho g}{2}} \quad (\text{For pure water in clean glass, contact angle } \theta \approx 0^\circ, \cos\theta = 1)
\end{equation}
Meniscus correction adds $r/3$ to measured height $h$. Internal radius $r$ measured by travelling microscope.

\textbf{2. Viscosity by Terminal Velocity (Stokes' Law):}
Spherical lead/steel ball falling through viscous medium (glycerine/castor oil) reaches constant terminal velocity $v_t$:
\begin{equation}
  F_{\mathrm{viscous}} = 6\pi \eta r v_t = \frac{4}{3}\pi r^3 (\rho - \sigma) g \implies \mathbf{\eta = \frac{2 r^2 (\rho - \sigma) g}{9 v_t}}
\end{equation}
Where $\rho$ is density of sphere, $\sigma$ is density of liquid. Graph of $v_t$ versus $r^2$ is a straight line through the origin.

\textbf{3. Speed of Sound in Air using Resonance Tube (End Correction):}
Water level provides closed end. For first two resonant air columns of lengths $l_1$ and $l_2$ with tuning fork of frequency $f$:
\begin{equation}
  l_1 + e = \frac{\lambda}{4} \qquad l_2 + e = \frac{3\lambda}{4} \implies l_2 - l_1 = \frac{\lambda}{2}
\end{equation}
\begin{equation}
  \mathbf{v = f \lambda = 2 f (l_2 - l_1)} \qquad \mathbf{\text{End Correction } e = \frac{l_2 - 3l_1}{2} \approx 0.6\,r} \ (r = \text{internal tube radius})
\end{equation}
\end{conceptbox}

\begin{formulabox}[Calorimetry \& Newton's Law of Cooling]
\textbf{1. Specific Heat Capacity of Solid by Method of Mixtures:}
Heat lost by hot solid = Heat gained by water and calorimeter (at thermal equilibrium):
\begin{equation}
  m_s s_s (T_s - T) = (m_w s_w + m_c s_c)(T - T_0) \implies \mathbf{s_s = \frac{(m_w s_w + m_c s_c)(T - T_0)}{m_s (T_s - T)}}
\end{equation}
Where $m_c s_c = W$ is water equivalent of calorimeter.
\textbf{2. Cooling Curve (Newton's Law of Cooling):}
\begin{equation}
  \frac{dT}{dt} = -K (T - T_0) \implies \mathbf{\ln(T - T_0) = -K t + C}
\end{equation}
Graph of $\ln(T - T_0)$ versus time $t$ is a straight line with negative slope ($-K$). Graph of temperature $T$ vs time $t$ is an exponential decay curve flattening at ambient temperature $T_0$.
\end{formulabox}

\section{Electricity \& Magnetism Lab Experiments}

\begin{formulabox}[Ohm's Law \& Metre Bridge (Wheatstone Bridge)]
\textbf{1. Verification of Ohm's Law ($V-I$ Graph):}
\begin{equation}
  V = I R \implies R = \text{Slope of } V \text{ vs } I \text{ graph} \qquad \text{Resistivity } \mathbf{\rho = \frac{R A}{L} = \frac{\pi r^2 R}{L}}
\end{equation}
\textbf{2. Metre Bridge (Null Deflection Method):}
Uniform high-resistance manganin/constantan wire of length $100\,\mathrm{cm}$:
\begin{equation}
  \frac{P}{Q} = \frac{R}{S} \implies \frac{l_1}{100 - l_1} = \frac{R}{X} \implies \mathbf{X = R \left(\frac{100 - l_1}{l_1}\right)}
\end{equation}
\textbf{End Errors \& Precautions:}
\begin{itemize}
  \item Resistance of thick copper strips at ends introduces \textbf{End Errors} ($\alpha, \beta$); corrected by swapping $R$ and $X$.
  \item Maximum sensitivity occurs when balance point is near the center: $\mathbf{l_1 \approx 50\,\mathrm{cm}}$ ($40\,\mathrm{cm} \le l_1 \le 60\,\mathrm{cm}$).
\end{itemize}
\end{formulabox}

\begin{formulabox}[Galvanometer Resistance, Figure of Merit \& Conversions]
\textbf{1. Half-Deflection Method for Galvanometer Resistance ($G$):}
\begin{equation}
  \text{Initial deflection } \theta \text{ with series } R: \quad I_1 = \frac{E}{R + G} = k\theta
\end{equation}
\begin{equation}
  \text{Shunting galvanometer with } S \text{ reduces deflection to } \frac{\theta}{2}: \quad \mathbf{G = \frac{R \cdot S}{R - S}}
\end{equation}
\textbf{2. Figure of Merit ($k$):} Current required to produce unit division deflection ($1\,\text{div}$):
\begin{equation}
  \mathbf{k = \frac{I}{\theta} = \frac{E}{(R + G)\theta}} \implies \text{Full Scale Current } I_g = n \cdot k \quad (n = \text{total scale divisions})
\end{equation}
\textbf{3. Conversion of Galvanometer into Ammeter and Voltmeter:}
\begin{itemize}
  \item \textbf{Ammeter (Range $0 \to I$):} Connect small resistance (\textbf{Shunt $S$}) in \textbf{parallel} with galvanometer:
  \begin{equation}
    (I - I_g)S = I_g G \implies \mathbf{S = \frac{I_g G}{I - I_g}} \approx \frac{I_g G}{I}
  \end{equation}
  Resistance of ideal ammeter = $0$. Equivalent ammeter resistance: $R_A = \frac{G S}{G + S} < S$.
  \item \textbf{Voltmeter (Range $0 \to V$):} Connect high resistance ($R_V$) in \textbf{series} with galvanometer:
  \begin{equation}
    V = I_g (G + R_V) \implies \mathbf{R_V = \frac{V}{I_g} - G}
  \end{equation}
  Resistance of ideal voltmeter = $\infty$. Equivalent voltmeter resistance: $R_{\text{eff}} = G + R_V$.
\end{itemize}
\end{formulabox}

\section{Optics \& Semiconductor Lab Experiments}

\begin{formulabox}[Focal Length of Mirror/Lens, Prism Deviation \& Glass Slab]
\textbf{1. Focal Length by $u-v$ Method (Convex Lens \& Concave Mirror):}
\begin{equation}
  \text{Convex Lens: } \frac{1}{f} = \frac{1}{v} - \frac{1}{u} \implies \text{Graph of } \frac{1}{v} \text{ vs } \frac{1}{u} \text{ has slope } -1 \text{ and intercepts } \frac{1}{f}
\end{equation}
\begin{equation}
  \text{Concave Mirror: } \frac{1}{f} = \frac{1}{v} + \frac{1}{u} \implies \text{At } u = 2f, \ v = 2f \quad (\text{Object at center of curvature})
\end{equation}
\textbf{2. Refractive Index of Prism ($i-\delta$ Curve):}
Graph of angle of incidence $i$ versus deviation $\delta$ is U-shaped with a unique minimum deviation $\delta_m$:
\begin{equation}
  \mathbf{\mu = \frac{\sin\left(\frac{A + \delta_m}{2}\right)}{\sin\left(\frac{A}{2}\right)}} \qquad (\text{At minimum deviation: } i_1 = i_2, \ r_1 = r_2 = A/2, \text{ ray inside is parallel to base})
\end{equation}
\textbf{3. Refractive Index of Glass Slab using Travelling Microscope:}
\begin{equation}
  \mathbf{\mu = \frac{\text{Real Thickness}}{\text{Apparent Thickness}} = \frac{R_3 - R_1}{R_3 - R_2}}
\end{equation}
Where $R_1$ = reading on ink mark on paper, $R_2$ = reading through glass slab, $R_3$ = reading of lycopodium powder on top surface.
\end{formulabox}

\begin{conceptbox}[Semiconductor Diode Characteristics (Forward \& Reverse Bias)]
\begin{itemize}
  \item \textbf{Forward Bias Characteristics:} $p$-side connected to positive, $n$-side to negative terminal.
  \begin{itemize}
    \item Depletion layer width decreases; barrier potential opposes applied voltage.
    \item \textbf{Knee / Cut-in Voltage ($V_k$):} Current is negligible until threshold voltage is reached ($\mathbf{0.7\,\mathrm{V}}$ for Silicon, $\mathbf{0.3\,\mathrm{V}}$ for Germanium), beyond which current rises exponentially.
    \item \textbf{Dynamic Forward Resistance:} $\mathbf{r_f = \frac{\Delta V_f}{\Delta I_f}}$ (typically $10-100\,\Omega$).
  \end{itemize}
  \item \textbf{Reverse Bias Characteristics:} $p$-side connected to negative, $n$-side to positive terminal.
  \begin{itemize}
    \item Depletion layer width widens; tiny constant \textbf{Reverse Saturation Current} ($I_0$) flows due to minority carriers (microamperes $\mu\mathrm{A}$ for Ge, nanoamperes $\mathrm{nA}$ for Si).
    \item \textbf{Reverse Breakdown Voltage ($V_{\mathrm{br}}$):} Sudden sharp increase in current due to avalanche/Zener breakdown.
    \item \textbf{Dynamic Reverse Resistance:} $\mathbf{r_r = \frac{\Delta V_r}{\Delta I_r}}$ (typically $> 1\,\mathrm{M}\Omega$).
  \end{itemize}
\end{itemize}
\end{conceptbox}

% =========================================================================
% PART VI: MASTER FORMULA COMPENDIUMS & QUICK REVISION HANDBOOK
% =========================================================================
\part{Master Formula Compendiums \& Quick Revision Handbook}

\chapter{High-Yield Formula Compendium: Mechanics, Gravitation, Thermal \& SHM}

This chapter brings together every governing equation, parameter definition, validity boundary condition, and NEET trap across Class 11 Physics (Units 2 through 10). Use this compendium for rapid revision before full-length tests and final preparation cycles.

\section{Kinematics (Motion in a Straight Line \& Plane)}

\subsection{Rectilinear Motion under Uniform Acceleration}

\begin{formulabox}[Master Kinematic Equations (Constant Acceleration $a = \text{const}$)]
\begin{align}
  v &= u + at \\
  s &= ut + \frac{1}{2}at^2 \\
  v^2 &= u^2 + 2as \\
  s_n &= u + \frac{a}{2}(2n - 1) \quad \text{(Displacement in the $n$-th second)} \\
  s &= \left(\frac{u + v}{2}\right)t
\end{align}
\textbf{Parameter Specification:}
\begin{itemize}
  \item $u$: Initial velocity vector at $t=0$ $[\si{m.s^{-1}}]$
  \item $v$: Instantaneous velocity vector at time $t$ $[\si{m.s^{-1}}]$
  \item $a$: Constant acceleration vector $[\si{m.s^{-2}}]$
  \item $s$: Net displacement vector from origin $[\si{m}]$
  \item $s_n$: Displacement accumulated strictly between $t = (n-1)$ and $t = n$ $[\si{m}]$
\end{itemize}
\end{formulabox}

\begin{usecasebox}[Applicability Conditions \& Shortcuts]
\begin{itemize}
  \item \textbf{Criterion of Validity:} Strictly applicable \textbf{only} when acceleration $a$ is constant in both magnitude and direction. If $a = f(t)$ or $a = f(x)$, kinematic equations fail; one must integrate:
  \begin{equation*}
    v = \int a\,dt, \quad s = \int v\,dt, \quad \int v\,dv = \int a\,dx
  \end{equation*}
  \item \textbf{Galileo's Law of Odd Numbers:} For a body starting from rest ($u=0$) under constant acceleration, the distances traversed in equal successive intervals of time $\Delta t$ follow the ratio:
  \begin{equation*}
    s_1 : s_2 : s_3 : \dots : s_n = 1 : 3 : 5 : \dots : (2n - 1)
  \end{equation*}
  \item \textbf{Stopping Distance \& Stopping Time:}
  \begin{equation*}
    d_{\text{stop}} = \frac{u^2}{2a} \implies d_{\text{stop}} \propto u^2, \qquad t_{\text{stop}} = \frac{u}{a} \implies t_{\text{stop}} \propto u
  \end{equation*}
\end{itemize}
\end{usecasebox}

\begin{trapbox}[NEET Traps: Distance vs. Displacement \& Turning Points]
\begin{itemize}
  \item \textbf{Trap 1 (Distance vs Displacement):} The formula $s = ut + \frac{1}{2}at^2$ calculates \textbf{displacement}, NOT total distance traveled. If velocity changes sign ($v=0$ at some $t_{\text{turn}} < t_{\text{total}}$), distance must be calculated piecewise:
  \begin{equation*}
    \text{Distance} = |s(0 \to t_{\text{turn}})| + |s(t_{\text{turn}} \to t_{\text{total}})|
  \end{equation*}
  \item \textbf{Trap 2 ($n$-th Second Dimensionality):} In $s_n = u + \frac{a}{2}(2n-1)$, $u$ is multiplied by an implicit $(1\text{ s})$, preserving dimensional homogeneity ($[L] = [LT^{-1} \cdot T] = [L]$). Do not drop signs for $u$ and $a$.
\end{itemize}
\end{trapbox}

\subsection{Projectile Motion in a Uniform Gravitational Field}

\begin{formulabox}[Oblique Projectile Motion (Launch angle $\theta$ with horizontal)]
\begin{align}
  T &= \frac{2u\sin\theta}{g} = \frac{2u_y}{g} && \text{[Total Time of Flight]} \\
  H_{\max} &= \frac{u^2\sin^2\theta}{2g} = \frac{u_y^2}{2g} && \text{[Maximum Height]} \\
  R &= \frac{u^2\sin 2\theta}{g} = \frac{2u_x u_y}{g} && \text{[Horizontal Range]} \\
  y &= x\tan\theta - \frac{g x^2}{2u^2\cos^2\theta} = x\tan\theta\left(1 - \frac{x}{R}\right) && \text{[Trajectory Equation]}
\end{align}
\end{formulabox}

\begin{usecasebox}[High-Yield Variations \& Relations]
\begin{itemize}
  \item \textbf{Relation Between $H_{\max}$ and $R$:}
  \begin{equation*}
    \frac{H_{\max}}{R} = \frac{\tan\theta}{4} \implies R\tan\theta = 4H_{\max}
  \end{equation*}
  When launch angle $\theta = 45^\circ$, $R_{\max} = 4H_{\max}$.
  \item \textbf{Complementary Angles of Projection:} For identical launch speed $u$, angles $\theta$ and $(90^\circ - \theta)$ yield identical horizontal ranges ($R_1 = R_2$). The corresponding flight times $T_1, T_2$ and heights $H_1, H_2$ satisfy:
  \begin{equation*}
    R = \frac{1}{2}g T_1 T_2 = 4\sqrt{H_1 H_2}
  \end{equation*}
  \item \textbf{Velocity at Highest Point:} $v_{\text{top}} = u_x = u\cos\theta$ (purely horizontal, acceleration $g$ is perpendicular: $\mathbf{a} \perp \mathbf{v}$).
  \item \textbf{Radius of Curvature:} At apex: $\rho_{\text{apex}} = \frac{u^2\cos^2\theta}{g}$; At launch: $\rho_{\text{launch}} = \frac{u^2}{g\cos\theta}$.
\end{itemize}
\end{usecasebox}

\begin{trapbox}[Common Projectile Traps]
\begin{itemize}
  \item \textbf{Kinetic Energy at Apex:} If initial KE is $K_0 = \frac{1}{2}mu^2$, KE at highest point is \textbf{not zero}; it is $K_{\text{apex}} = K_0 \cos^2\theta$. For $\theta=60^\circ$, $K_{\text{apex}} = \frac{1}{4}K_0$.
  \item \textbf{Angle Given with Vertical:} If problem states "projected at angle $\alpha$ to the vertical", replace $\theta$ with $(90^\circ - \alpha)$. Always double check the reference axis!
\end{itemize}
\end{trapbox}

\newpage

% =========================================================================
% SECTION 2: LAWS OF MOTION \& FRICTION
% =========================================================================
\section{Laws of Motion \& Friction}

\subsection{Newton's Second Law, Momentum \& Impulse}

\begin{formulabox}[Momentum, Force \& Impulse Formulations]
\begin{align}
  \mathbf{F}_{\text{net}} &= \frac{d\mathbf{p}}{dt} = m\frac{d\mathbf{v}}{dt} + \mathbf{v}_{\text{rel}}\frac{dm}{dt} \\
  \mathbf{J} &= \int_{t_1}^{t_2} \mathbf{F}\,dt = \Delta\mathbf{p} = \mathbf{p}_2 - \mathbf{p}_1 \quad \text{[Impulse-Momentum Theorem]}
\end{align}
\textbf{Variable Mass System (Rocket Propulsion):}
\begin{equation}
  v = u + v_{\text{rel}}\ln\left(\frac{m_0}{m}\right) - gt, \qquad F_{\text{thrust}} = v_{\text{rel}}\left(-\frac{dm}{dt}\right)
\end{equation}
\end{formulabox}

\subsection{Friction Mechanics}

\begin{formulabox}[Static, Limiting \& Kinetic Friction Laws]
\begin{align}
  f_s &\le f_{s,\max} = \mu_s N && \text{[Static Friction: Self-Adjusting, $0 \le f_s \le \mu_s N$]} \\
  f_k &= \mu_k N && \text{[Kinetic Friction: Constant for given surfaces, $\mu_k \le \mu_s$]} \\
  \tan\lambda &= \mu_s && \text{[Angle of Friction $\lambda$]} \\
  \theta_r &= \lambda = \tan^{-1}(\mu_s) && \text{[Angle of Repose $\theta_r$]}
\end{align}
\end{formulabox}

\begin{usecasebox}[Motion on an Inclined Plane of Inclination $\theta$]
\begin{itemize}
  \item If $\theta \le \theta_r$: Block remains at rest on incline; static friction $f_s = mg\sin\theta$.
  \item If $\theta > \theta_r$: Block accelerates downward:
  \begin{equation*}
    a_{\text{down}} = g(\sin\theta - \mu_k\cos\theta)
  \end{equation*}
  \item If projected upward along incline:
  \begin{equation*}
    a_{\text{up}} = g(\sin\theta + \mu_k\cos\theta) \implies a_{\text{up}} > a_{\text{down}} \implies t_{\text{up}} < t_{\text{down}}
  \end{equation*}
  \item \textbf{Ratio of Times ($n$-times longer down a rough plane vs smooth plane):}
  \begin{equation*}
    \mu = \tan\theta\left(1 - \frac{1}{n^2}\right)
  \end{equation*}
\end{itemize}
\end{usecasebox}

\begin{trapbox}[Friction Calculation Traps]
\begin{itemize}
  \item \textbf{Normal Reaction Modification:} $N \ne mg$ in general! On incline: $N = mg\cos\theta$. When pulling at angle $\alpha$: $N = mg - F\sin\alpha$. When pushing at angle $\alpha$: $N = mg + F\sin\alpha$. (Hence pulling is always easier than pushing).
  \item \textbf{Static Friction is NOT Always $\mu_s N$:} $\mu_s N$ is the \textbf{maximum limit} (impending motion). If applied force $F_{\text{applied}} < \mu_s N$, the friction force is exactly equal to $F_{\text{applied}}$, NOT $\mu_s N$.
\end{itemize}
\end{trapbox}

\subsection{Circular Dynamics \& Highway Banking}

\begin{formulabox}[Vehicle Dynamics on Circular Tracks]
\begin{align}
  v_{\max,\text{flat}} &= \sqrt{\mu_s g R} && \text{[Unbanked Flat Road]} \\
  v_{\text{opt}} &= \sqrt{g R \tan\theta} && \text{[Banked Road, Zero Friction Needed]} \\
  v_{\max,\text{banked}} &= \sqrt{g R \left(\frac{\mu_s + \tan\theta}{1 - \mu_s\tan\theta}\right)} && \text{[Banked Road with Friction, Upper Speed Limit]} \\
  v_{\min,\text{banked}} &= \sqrt{g R \left(\frac{\tan\theta - \mu_s}{1 + \mu_s\tan\theta}\right)} && \text{[Banked Road with Friction, Lower Speed Limit]}
\end{align}
\end{formulabox}

\newpage

% =========================================================================
% SECTION 3: WORK, ENERGY, POWER \& COLLISIONS
% =========================================================================
\section{Work, Energy, Power \& Collisions}

\subsection{Work-Energy Theorem \& Potential Energy}

\begin{formulabox}[Fundamental Work \& Potential Relations]
\begin{align}
  W_{\text{net}} &= \Delta K = K_f - K_i = \frac{1}{2}m v_f^2 - \frac{1}{2}m v_i^2 && \text{[Universal Work-Energy Theorem]} \\
  W_{\text{cons}} &= -\Delta U = -(U_f - U_i) && \text{[Conservative Force Definition]} \\
  \mathbf{F} &= -\boldsymbol{\nabla} U = -\left(\frac{\partial U}{\partial x}\hat{\mathbf{i}} + \frac{\partial U}{\partial y}\hat{\mathbf{j}} + \frac{\partial U}{\partial z}\hat{\mathbf{k}}\right) && \text{[Force-Potential Gradient Relation]}
\end{align}
\textbf{Spring Elastic Potential Energy:}
\begin{equation}
  U_{\text{spring}} = \frac{1}{2}k x^2, \qquad W_{\text{spring}} = -\frac{1}{2}k(x_f^2 - x_i^2)
\end{equation}
\end{formulabox}

\begin{usecasebox}[Equilibrium Criteria from Potential Function $U(x)$]
\begin{itemize}
  \item \textbf{Equilibrium Condition:} $\frac{dU}{dx} = 0$ (Force $F=0$).
  \item \textbf{Stable Equilibrium:} $\frac{d^2 U}{dx^2} > 0$ (Potential energy is minimum; restoring force acts on perturbation; SHM possible).
  \item \textbf{Unstable Equilibrium:} $\frac{d^2 U}{dx^2} < 0$ (Potential energy is maximum; destabilizing force acts on perturbation).
  \item \textbf{Neutral Equilibrium:} $\frac{d^2 U}{dx^2} = 0$ (Potential energy is locally constant).
\end{itemize}
\end{usecasebox}

\subsection{Vertical Circular Motion (Mass $m$ on Inextensible String Length $L$)}

\begin{formulabox}[Critical Velocity \& Tension Conditions]
\begin{align}
  v_{\text{bottom},\min} &= \sqrt{5gL} && \text{[Condition to complete full circle]} \\
  v_{\text{top},\min} &= \sqrt{gL} && \text{[Velocity at top for non-slack string ($T_{\text{top}} \ge 0$)]} \\
  v_{\text{horizontal}} &= \sqrt{3gL} && \text{[Velocity when string is horizontal]} \\
  T_{\text{bottom}} - T_{\text{top}} &= 6mg && \text{[Tension difference is independent of launch speed!]}
\end{align}
\end{formulabox}

\begin{trapbox}[String vs Light Rigid Rod in Vertical Circle]
\begin{itemize}
  \item \textbf{String (can slack):} Minimum velocity at bottom to complete circle is $\sqrt{5gL}$; minimum velocity at apex is $\sqrt{gL}$.
  \item \textbf{Light Rigid Rod (cannot slack):} Velocity at highest point can drop to \textbf{zero} ($v_{\text{top}} = 0$). Hence minimum velocity at bottom is:
  \begin{equation*}
    v_{\text{bottom},\min} = \sqrt{4gL} = 2\sqrt{gL}
  \end{equation*}
  Never confuse the string case with the rigid rod / sphere on track case!
\end{itemize}
\end{trapbox}

\subsection{1D \& 2D Collisions}

\begin{formulabox}[Coefficient of Restitution \& Collision Velocities]
\begin{equation}
  e = \frac{v_{\text{sep}}}{v_{\text{app}}} = \frac{v_2 - v_1}{u_1 - u_2} \quad (0 \le e \le 1)
\end{equation}
\textbf{1D Head-On Collision Velocities ($m_1$ moving with $u_1$, $m_2$ moving with $u_2$):}
\begin{align}
  v_1 &= \left(\frac{m_1 - e m_2}{m_1 + m_2}\right)u_1 + \frac{(1 + e)m_2}{m_1 + m_2}u_2 \\
  v_2 &= \frac{(1 + e)m_1}{m_1 + m_2}u_1 + \left(\frac{m_2 - e m_1}{m_1 + m_2}\right)u_2
\end{align}
\textbf{Loss of Kinetic Energy in Inelastic Collision ($\Delta K$):}
\begin{equation}
  \Delta K_{\text{loss}} = \frac{1}{2}\left(\frac{m_1 m_2}{m_1 + m_2}\right)(u_1 - u_2)^2 (1 - e^2) = \frac{1}{2}\mu v_{\text{rel}}^2 (1 - e^2)
\end{equation}
\end{formulabox}

\begin{usecasebox}[Special Elastic Collision Limits ($e=1$)]
\begin{itemize}
  \item \textbf{Equal Masses ($m_1 = m_2$):} Velocities simply interchange after collision:
  \begin{equation*}
    v_1 = u_2, \qquad v_2 = u_1
  \end{equation*}
  \item \textbf{Heavy Target at Rest ($m_2 \gg m_1, u_2 = 0$):} $v_1 \approx -u_1$ (rebounds with same speed), $v_2 \approx 0$.
  \item \textbf{Heavy Projectile against Light Target ($m_1 \gg m_2, u_2 = 0$):} $v_1 \approx u_1$, $v_2 \approx 2u_1$ (target launched at twice projectile speed).
  \item \textbf{Rebound from Floor (Floor $m_2 \to \infty, u_2 = 0$):}
  \begin{equation*}
    v_n = e^n v_0, \qquad h_n = e^{2n} h_0, \qquad t_{\text{total}} = t_0\left(\frac{1 + e}{1 - e}\right), \qquad H_{\text{total}} = h_0\left(\frac{1 + e^2}{1 - e^2}\right)
  \end{equation*}
\end{itemize}
\end{usecasebox}

\vspace{15pt}
\begin{center}
  \rule{0.6\textwidth}{0.8pt}\\[6pt]
  {\small\textbf{\color{primarydark}End of Module 01 Formula Handbook}}\\[2pt]
  {\footnotesize\color{accentblue}Prepared by Satya Prakash \quad\textbar\quad NEET UG 2027 Series}
\end{center}

\section{Rotational Dynamics \& Centre of Mass (Unit 5)}

\subsection{Centre of Mass (CM) Formulations}

\begin{formulabox}[Discrete and Continuous Centre of Mass Coordinates]
\begin{equation}
  \mathbf{r}_{\text{cm}} = \frac{\sum_{i=1}^n m_i \mathbf{r}_i}{\sum_{i=1}^n m_i} = \frac{1}{M}\sum m_i \mathbf{r}_i, \qquad \mathbf{r}_{\text{cm}} = \frac{\int \mathbf{r}\,dm}{\int dm} = \frac{1}{M}\int \mathbf{r}\,dm
\end{equation}
\textbf{Centre of Mass Velocity, Acceleration \& Truncated Body Relation:}
\begin{equation}
  \mathbf{v}_{\text{cm}} = \frac{\sum m_i \mathbf{v}_i}{M}, \quad \mathbf{a}_{\text{cm}} = \frac{\mathbf{F}_{\text{ext}}}{M}, \quad \mathbf{r}_{\text{rem}} = \frac{M\mathbf{r}_{\text{orig}} - m\mathbf{r}_{\text{cut}}}{M - m}
\end{equation}
\end{formulabox}

\begin{table}[htbp]
\centering
\small
\renewcommand{\arraystretch}{1.2}
\begin{tabularx}{\textwidth}{L c L c}
\toprule
\textbf{Geometry / Symmetric Body} & \textbf{CM Location ($y_{\text{cm}}$)} & \textbf{Geometry / Symmetric Body} & \textbf{CM Location ($y_{\text{cm}}$)} \\
\midrule
Uniform Semicircular Ring & $\frac{2R}{\pi}$ from base & Uniform Semicircular Disc & $\frac{4R}{3\pi}$ from base \\
Hemispherical Shell (Hollow) & $\frac{R}{2}$ from center & Solid Hemisphere & $\frac{3R}{8}$ from center \\
Hollow Cone (without base) & $\frac{h}{3}$ from base & Solid Cone & $\frac{h}{4}$ from base \\
\bottomrule
\end{tabularx}
\caption{Exact Centre of Mass Coordinates from Base / Symmetric Center along Axis.}
\end{table}

\subsection{Rotational Kinematics, Dynamics \& Angular Momentum}

\begin{formulabox}[Governing Torque and Angular Momentum Relations]
\begin{align}
  \boldsymbol{\tau} &= \mathbf{r} \times \mathbf{F} = I\boldsymbol{\alpha} = \frac{d\mathbf{L}}{dt} && \text{[Torque Equation; $\tau = r F \sin\theta$]} \\
  \mathbf{L} &= \mathbf{r} \times \mathbf{p} = I\boldsymbol{\omega} && \text{[Angular Momentum of Rigid Body]} \\
  W_{\text{rot}} &= \int \tau\,d\theta = \frac{1}{2}I\omega_f^2 - \frac{1}{2}I\omega_i^2 && \text{[Rotational Work-Energy Theorem]}
\end{align}
\textbf{Conservation of Angular Momentum (When $\boldsymbol{\tau}_{\text{ext}} = \mathbf{0}$):}
\begin{equation}
  \mathbf{L} = I_1 \boldsymbol{\omega}_1 = I_2 \boldsymbol{\omega}_2 = \text{constant} \implies \omega \propto \frac{1}{I}
\end{equation}
\end{formulabox}

\subsection{Moment of Inertia Theorems \& Standard Geometrical Values}

\begin{formulabox}[Moment of Inertia (MI) Fundamental Theorems]
\begin{align}
  I &= \sum m_i r_i^2 = \int r^2\,dm = M k^2 && \text{[$k$ = Radius of Gyration]} \\
  I &= I_{\text{cm}} + M d^2 && \text{[Parallel Axis Theorem: valid for ANY rigid body]} \\
  I_z &= I_x + I_y && \text{[Perpendicular Axis Theorem: strictly PLANAR laminar bodies only]}
\end{align}
\end{formulabox}

\begin{table}[htbp]
\centering
\small
\renewcommand{\arraystretch}{1.2}
\begin{tabularx}{\textwidth}{L c c c}
\toprule
\textbf{Body (Mass $M$, Radius $R$)} & \textbf{Axis of Rotation} & \textbf{Moment of Inertia ($I$)} & $k^2/R^2$ \\
\midrule
Thin Rod (Length $L$) & Perpendicular through center & $\frac{1}{12} M L^2$ & $L^2 / (12 R^2)$ \\
Thin Rod (Length $L$) & Perpendicular through one end & $\frac{1}{3} M L^2$ & $L^2 / (3 R^2)$ \\
Thin Circular Ring & Central perpendicular axis & $M R^2$ & $1$ \\
Thin Circular Ring & Across diameter & $\frac{1}{2} M R^2$ & $1/2$ \\
Uniform Circular Disc & Central perpendicular axis & $\frac{1}{2} M R^2$ & $1/2$ \\
Uniform Circular Disc & Across diameter & $\frac{1}{4} M R^2$ & $1/4$ \\
Solid Cylinder & Longitudinal central axis & $\frac{1}{2} M R^2$ & $1/2$ \\
Hollow Cylinder (Thin Pipe) & Longitudinal central axis & $M R^2$ & $1$ \\
Solid Sphere & Through any diameter & $\frac{2}{5} M R^2$ & $2/5 = 0.40$ \\
Spherical Shell (Hollow Sphere) & Through any diameter & $\frac{2}{3} M R^2$ & $2/3 \approx 0.67$ \\
\bottomrule
\end{tabularx}
\caption{Standard Moments of Inertia and Dimensionless Gyration Factor $k^2/R^2$.}
\end{table}

\subsection{Pure Rolling Dynamics (Rolling Without Slipping)}

\begin{formulabox}[Rolling Motion Governing Equations]
\begin{align}
  v_{\text{cm}} &= R\omega, \qquad a_{\text{cm}} = R\alpha \quad \text{[No-slip Kinematic Constraint]} \\
  K_{\text{total}} &= K_{\text{trans}} + K_{\text{rot}} = \frac{1}{2}M v_{\text{cm}}^2 + \frac{1}{2}I_{\text{cm}}\omega^2 = \frac{1}{2}M v_{\text{cm}}^2\left(1 + \frac{k^2}{R^2}\right) \\
  \frac{K_{\text{trans}}}{K_{\text{total}}} &= \frac{1}{1 + k^2/R^2}, \qquad \frac{K_{\text{rot}}}{K_{\text{total}}} = \frac{k^2/R^2}{1 + k^2/R^2}
\end{align}
\textbf{Rolling Down an Incline of Angle $\theta$ and Height $h$:}
\begin{equation}
  a = \frac{g\sin\theta}{1 + \frac{k^2}{R^2}}, \qquad v_{\text{bottom}} = \sqrt{\frac{2gh}{1 + \frac{k^2}{R^2}}}, \qquad \mu_{\min} = \frac{\tan\theta}{1 + \frac{R^2}{k^2}}
\end{equation}
\end{formulabox}

\begin{usecasebox}[Race Down an Incline (Identical $M, R, \theta, h$)]
\begin{itemize}
  \item Ordering of $k^2/R^2$: $\text{Solid Sphere } (0.40) < \text{Disc/Cylinder } (0.50) < \text{Hollow Sphere } (0.67) < \text{Ring/Hollow Cylinder } (1.00)$.
  \item Acceleration Order: $a_{\text{solid sphere}} > a_{\text{disc}} > a_{\text{hollow sphere}} > a_{\text{ring}}$.
  \item Time of Descent ($t = \sqrt{2s/a}$): $t_{\text{solid sphere}} < t_{\text{disc}} < t_{\text{hollow sphere}} < t_{\text{ring}}$. Solid sphere reaches first!
  \item Velocity at Bottom: $v_{\text{solid sphere}} > v_{\text{disc}} > v_{\text{hollow sphere}} > v_{\text{ring}}$.
\end{itemize}
\end{usecasebox}

\begin{trapbox}[Rotational Traps: Friction \& Perpendicular Axis Theorem]
\begin{itemize}
  \item \textbf{Friction in Pure Rolling:} Static friction acts at the instantaneous point of contact. Because the point of contact is instantaneously at rest ($v_{\text{contact}} = 0$), \textbf{work done by friction in pure rolling is identically zero} ($W_f = 0$). Mechanical energy is conserved.
  \item \textbf{Perpendicular Axis Theorem Trap:} Cannot be applied to 3D bodies like spheres, solid cylinders, or cones. It is strictly valid for 2D plane laminar sheets ($I_z = I_x + I_y$).
\end{itemize}
\end{trapbox}

\newpage

% =========================================================================
% SECTION 2: GRAVITATION
% =========================================================================
\section{Gravitation (Unit 6)}

\subsection{Universal Gravitation, Acceleration ($g$) \& Field Variations}

\begin{formulabox}[Newton's Law \& Surface Acceleration $g$]
\begin{align}
  F &= \frac{G M_1 M_2}{r^2} && [G = 6.674\times 10^{-11}\si{N.m^2.kg^{-2}}] \\
  g &= \frac{GM}{R^2} = \frac{4}{3}\pi G \rho R && [\text{Surface acceleration, Earth } g \approx 9.8\si{m.s^{-2}}]
\end{align}
\textbf{Variations of $g$ (Altitude, Depth \& Latitude):}
\begin{align}
  g_h &= g\left(\frac{R}{R+h}\right)^2 = g\left(1 + \frac{h}{R}\right)^{-2} \approx g\left(1 - \frac{2h}{R}\right) && [\text{Approximation valid strictly for } h \ll R] \\
  g_d &= g\left(1 - \frac{d}{R}\right) && [\text{Exact at all depths } 0 \le d \le R] \\
  g' &= g - \omega^2 R \cos^2\lambda && [\lambda = \text{latitude; Equator } \lambda=0, \text{Pole } \lambda=90^\circ]
\end{align}
\end{formulabox}

\begin{usecasebox}[Exact Comparison \& Extreme Earth Rotation Cases]
\begin{itemize}
  \item \textbf{Equal Reduction at Altitude and Depth:} For small $h \ll R$, $g_h = g_d \implies d = 2h$. A height of $1\si{km}$ causes the same reduction in $g$ as a depth of $2\si{km}$.
  \item \textbf{Weightlessness at Equator:} For bodies at equator to become weightless ($g'=0$):
  \begin{equation*}
    \omega = \sqrt{\frac{g}{R}} \approx \frac{1}{800}\si{rad.s^{-1}}, \qquad T = 2\pi\sqrt{\frac{R}{g}} \approx 84.6\text{ minutes} \approx 1.41\text{ hours}
  \end{equation*}
  This requires the Earth to rotate roughly \textbf{17 times faster} than its current angular velocity.
  \item \textbf{Earth Stops Rotating:} If $\omega \to 0$, $g$ increases by $\omega^2 R$ at the equator, but \textbf{remains unchanged at the poles}.
\end{itemize}
\end{usecasebox}

\subsection{Gravitational Potential, Escape \& Orbital Velocities}

\begin{formulabox}[Potential, Escape Speed, Orbital Speed \& Satellite Energies]
\begin{align}
  V(r) &= -\frac{GM}{r}, \qquad U(r) = -\frac{GMm}{r} && [E_g = -\frac{dV}{dr} = -\frac{GM}{r^2}] \\
  v_e &= \sqrt{\frac{2GM}{R}} = \sqrt{2gR} \approx 11.2\si{km.s^{-1}} && [\text{Escape Velocity from Earth Surface}] \\
  v_o &= \sqrt{\frac{GM}{r}} = \sqrt{\frac{g R^2}{R+h}} && [\text{Orbital Velocity at orbit radius } r = R+h] \\
  v_e &= \sqrt{2} v_o \approx 1.414 v_o && [\text{Near-Earth orbit } h \ll R: v_o \approx 7.92\si{km.s^{-1}}] \\
  T^2 &= \left(\frac{4\pi^2}{GM}\right)r^3 \implies T = 2\pi\sqrt{\frac{r^3}{GM}} && [\text{Kepler's Third Law of Periods}]
\end{align}
\textbf{Satellite Energy Spectrum (Orbit radius $r = R+h$):}
\begin{equation}
  K = \frac{GMm}{2r}, \quad U = -\frac{GMm}{r}, \quad E_{\text{total}} = -\frac{GMm}{2r} = -K = \frac{U}{2}, \quad \text{B.E.} = -E_{\text{total}} = \frac{GMm}{2r}
\end{equation}
\end{formulabox}

\begin{trapbox}[Gravitation Traps: Launch Angle \& Solid Sphere Potential]
\begin{itemize}
  \item \textbf{Escape Velocity Independence:} Escape velocity $v_e = \sqrt{2gR}$ is \textbf{completely independent} of the mass of the projectile $m$ and the angle of launch $\theta$. A body launched at $30^\circ, 60^\circ$ or $90^\circ$ requires the identical escape speed $11.2\si{km.s^{-1}}$.
  \item \textbf{Gravitational Potential Inside Uniform Solid Sphere (Earth):}
  \begin{equation*}
    V_{\text{in}}(r) = -\frac{GM}{2R^3}\left(3R^2 - r^2\right) \implies V_{\text{center}} = -\frac{3}{2}\frac{GM}{R} = 1.5\,V_{\text{surface}}
  \end{equation*}
  Potential is NOT zero at the centre; its magnitude is 1.5 times the surface value.
\end{itemize}
\end{trapbox}

\newpage

% =========================================================================
% SECTION 3: PROPERTIES OF SOLIDS \& FLUIDS
% =========================================================================
\section{Properties of Solids \& Fluids (Unit 7)}

\subsection{Elasticity \& Moduli of Elasticity}

\begin{formulabox}[Hooke's Law, Elastic Moduli \& Energy Density]
\begin{align}
  Y &= \frac{\text{Longitudinal Stress}}{\text{Longitudinal Strain}} = \frac{F/A}{\Delta L/L} = \frac{F L}{A \Delta L} && \text{[Young's Modulus, $\Delta L = \frac{F L}{A Y}$]} \\
  B &= -\frac{\Delta P}{\Delta V / V} = -V\frac{\Delta P}{\Delta V}, \qquad K = \frac{1}{B} && \text{[Bulk Modulus $B$ \& Compressibility $K$]} \\
  \eta &= \frac{\text{Shear Stress}}{\text{Shear Strain}} = \frac{F/A}{\theta} && \text{[Shear Modulus / Modulus of Rigidity]} \\
  \sigma &= -\frac{\text{Lateral Strain}}{\text{Longitudinal Strain}} = -\frac{\Delta d / d}{\Delta L / L} && \text{[Poisson's Ratio: Theoretical: $-1 \le \sigma \le 0.5$; Actual: $0 < \sigma < 0.5$]}
\end{align}
\textbf{Elastic Potential Energy Stored in Stretched Wire:}
\begin{equation}
  U = \frac{1}{2} F \Delta L = \frac{1}{2} k (\Delta L)^2, \qquad u = \frac{U}{V} = \frac{1}{2}\times\text{stress}\times\text{strain} = \frac{(\text{stress})^2}{2Y} = \frac{1}{2}Y(\text{strain})^2
\end{equation}
\textbf{Thermal Stress Developed in Clamped Rod:}
\begin{equation}
  \text{Stress} = Y\alpha\Delta T, \qquad F_{\text{tension}} = Y A \alpha \Delta T
\end{equation}
\end{formulabox}

\subsection{Fluid Statics, Pascal's Principle \& Buoyancy}

\begin{formulabox}[Pressure, Pascal's Law \& Archimedes' Principle]
\begin{align}
  P &= P_0 + \rho g h && \text{[Gauge Pressure } P_{\text{gauge}} = \rho g h\text{]} \\
  \frac{F_1}{A_1} &= \frac{F_2}{A_2} \implies F_2 = F_1\left(\frac{A_2}{A_1}\right) && \text{[Hydraulic Lift Principle]} \\
  F_B &= V_{\text{submerged}}\,\rho_{\text{liquid}}\,g && \text{[Archimedes' Buoyant Force]} \\
  \frac{V_{\text{submerged}}}{V_{\text{total}}} &= \frac{\rho_{\text{body}}}{\rho_{\text{liquid}}}, \qquad W_{\text{apparent}} = W_{\text{actual}}\left(1 - \frac{\rho_{\text{liquid}}}{\rho_{\text{body}}}\right)
\end{align}
\end{formulabox}

\subsection{Fluid Dynamics, Viscosity \& Surface Tension}

\begin{formulabox}[Continuity, Bernoulli, Viscosity \& Capillarity]
\begin{align}
  A_1 v_1 &= A_2 v_2 = \text{constant} && \text{[Equation of Continuity for Incompressible Fluid]} \\
  P + \frac{1}{2}\rho v^2 + \rho g h &= \text{constant} && \text{[Bernoulli's Theorem]} \\
  v_{\text{efflux}} &= \sqrt{2gh}, \quad R_{\text{horizontal}} = 2\sqrt{h(H-h)} && \text{[Torricelli's Law; $R_{\max} = H$ at $h = H/2$]} \\
  F_v &= 6\pi \eta r v && \text{[Stokes' Viscous Drag Law]} \\
  v_t &= \frac{2}{9}\frac{r^2(\rho - \sigma)g}{\eta} \implies v_t \propto r^2 && \text{[Terminal Velocity; $\rho=\text{body density}, \sigma=\text{fluid density}$]} \\
  R_e &= \frac{\rho v d}{\eta} && \text{[Reynolds Number: $R_e < 2000$ Laminar; $R_e > 3000$ Turbulent]}
\end{align}
\textbf{Surface Tension, Excess Pressure \& Capillary Rise:}
\begin{align}
  T &= \frac{F}{L} = \frac{\Delta W}{\Delta A}, \qquad W_{\text{splitting}} = 4\pi R^2 T(n^{1/3} - 1) && \text{[$1$ big drop $\to n$ identical droplets]} \\
  \Delta P_{\text{liquid drop}} &= \frac{2T}{R}, \qquad \Delta P_{\text{soap bubble}} = \frac{4T}{R} && \text{[Soap bubble has TWO free air-liquid interfaces]} \\
  h &= \frac{2T\cos\theta}{r\rho g} = \frac{2T}{R_c \rho g} && \text{[Jurin's Law: $h r = \text{constant}$; $R_c = r/\cos\theta$ is meniscus radius]}
\end{align}
\end{formulabox}

\begin{trapbox}[Fluid Traps: Terminal Velocity \& Insufficient Capillary Length]
\begin{itemize}
  \item \textbf{Terminal Velocity ($r^2$ vs $r$):} Viscous force $F_v \propto r$, but terminal velocity $v_t \propto r^2$. If two drops of radii $1:2$ coalesce, the radius becomes $R = 2^{1/3}r$; new terminal velocity increases by a factor of $2^{2/3}$.
  \item \textbf{Insufficient Capillary Tube Length ($l < h$):} Liquid \textbf{never overflows} from a capillary tube of insufficient height $l$. Instead, the liquid reaches the upper edge and increases its meniscus radius of curvature $R_c'$ such that $h r = l R_c'$ ($R_c' > R_c$).
\end{itemize}
\end{trapbox}

\newpage

% =========================================================================
% SECTION 4: THERMAL PHYSICS \& THERMODYNAMICS
% =========================================================================
\section{Thermal Physics \& Thermodynamics (Units 7, 8 \& 9)}

\subsection{Thermal Expansion, Calorimetry \& Heat Transfer}

\begin{formulabox}[Expansion, Calorimetry, Conduction \& Radiation]
\begin{align}
  \Delta L &= L_0 \alpha \Delta T, \quad \Delta A = A_0 \beta \Delta T, \quad \Delta V = V_0 \gamma \Delta T && [\alpha : \beta : \gamma = 1 : 2 : 3] \\
  Q &= m c \Delta T, \qquad Q = m L && [c = \text{sp. heat capacity}, L = \text{latent heat}] \\
  \frac{dQ}{dt} &= \frac{kA(T_1 - T_2)}{L} = \frac{\Delta T}{R_{\text{th}}} && [R_{\text{th}} = \frac{L}{kA} = \text{Thermal Resistance}] \\
  E &= \sigma e A T^4, \qquad P_{\text{net}} = \sigma e A(T^4 - T_0^4) && [\text{Stefan-Boltzmann Law, } \sigma = 5.67\times 10^{-8}\si{W.m^{-2}.K^{-4}}] \\
  \lambda_m T &= b = 2.898\times 10^{-3}\si{m.K} && [\text{Wien's Displacement Law}] \\
  \frac{T_1 - T_2}{t} &= K\left(\frac{T_1 + T_2}{2} - T_0\right) && [\text{Newton's Law of Cooling (Average Form)}]
\end{align}
\end{formulabox}

\subsection{First Law of Thermodynamics \& Process Formulations}

\begin{formulabox}[First Law \& Thermodynamic Processes]
\begin{equation}
  \Delta Q = \Delta U + \Delta W \quad [\text{Ideal Gas: } \Delta U = n C_v \Delta T = \frac{P_2V_2 - P_1V_1}{\gamma - 1} \text{ for ANY process}]
\end{equation}
\textbf{Work Done in Standard Thermodynamic Processes:}
\begin{align}
  W_{\text{isochoric}} &= 0, \qquad \Delta Q = \Delta U = n C_v \Delta T && [V = \text{constant}] \\
  W_{\text{isobaric}} &= P(V_2 - V_1) = n R(T_2 - T_1), \quad \Delta Q = n C_p \Delta T && [P = \text{constant}] \\
  W_{\text{isothermal}} &= n R T \ln\left(\frac{V_2}{V_1}\right) = 2.303 n R T \log_{10}\left(\frac{P_1}{P_2}\right) && [T = \text{const}, \Delta U = 0, \Delta Q = W] \\
  W_{\text{adiabatic}} &= \frac{P_1 V_1 - P_2 V_2}{\gamma - 1} = \frac{n R(T_1 - T_2)}{\gamma - 1} = -\Delta U && [\Delta Q = 0, P V^\gamma = \text{const}]
\end{align}
\end{formulabox}

\begin{table}[htbp]
\centering
\small
\renewcommand{\arraystretch}{1.2}
\begin{tabularx}{\textwidth}{L c c c c}
\toprule
\textbf{Gas Atomicity} & \textbf{Deg. of Freedom ($f$)} & $C_v = \frac{f}{2}R$ & $C_p = C_v + R$ & $\gamma = C_p/C_v = 1 + \frac{2}{f}$ \\
\midrule
Monoatomic ($\text{He}, \text{Ne}, \text{Ar}$) & $3$ (translational) & $\frac{3}{2}R$ & $\frac{5}{2}R$ & $\frac{5}{3} \approx 1.67$ \\
Diatomic (Rigid: $\text{O}_2, \text{N}_2$) & $3\text{ trans} + 2\text{ rot} = 5$ & $\frac{5}{2}R$ & $\frac{7}{2}R$ & $\frac{7}{5} = 1.40$ \\
Non-Linear Polyatomic ($\text{H}_2\text{O}, \text{NH}_3$) & $3\text{ trans} + 3\text{ rot} = 6$ & $3R$ & $4R$ & $\frac{4}{3} \approx 1.33$ \\
\bottomrule
\end{tabularx}
\caption{Molar Heat Capacities and Adiabatic Exponents for Standard Gases.}
\end{table}

\begin{usecasebox}[High-Yield Thermodynamic Relations \& Polytropic Processes]
\begin{itemize}
  \item \textbf{Energy Partition in Isobaric Process:}
  \begin{equation*}
    \Delta W : \Delta U : \Delta Q = nR\Delta T : nC_v\Delta T : nC_p\Delta T = (\gamma - 1) : 1 : \gamma
  \end{equation*}
  For a diatomic gas ($\gamma = 7/5$), fraction of heat converted to work is $\frac{\Delta W}{\Delta Q} = \frac{\gamma - 1}{\gamma} = \frac{2}{7} \approx 28.6\%$.
  \item \textbf{General Polytropic Process ($P V^n = \text{const}$):}
  \begin{equation*}
    W = \frac{nR(T_1 - T_2)}{n - 1}, \qquad C = C_v + \frac{R}{1 - n}
  \end{equation*}
  \item \textbf{Comparison of $P-V$ Curves:} $\left(\frac{dP}{dV}\right)_{\text{adiabatic}} = \gamma\left(\frac{dP}{dV}\right)_{\text{isothermal}}$. Adiabatic slope is $\gamma$-times steeper!
\end{itemize}
\end{usecasebox}

\subsection{Kinetic Theory of Gases (KTG)}

\begin{formulabox}[Pressure, Speeds \& Mean Free Path]
\begin{align}
  P &= \frac{1}{3}\rho v_{\text{rms}}^2 = \frac{1}{3}\frac{M}{V}v_{\text{rms}}^2, \qquad E_{\text{total, translation}} = \frac{3}{2}n R T = \frac{3}{2}N k_B T \\
  v_{\text{rms}} &= \sqrt{\frac{3RT}{M}}, \quad v_{\text{avg}} = \sqrt{\frac{8RT}{\pi M}}, \quad v_{\text{mp}} = \sqrt{\frac{2RT}{M}} \implies v_{\text{mp}} : v_{\text{avg}} : v_{\text{rms}} = 1 : 1.128 : 1.224 \\
  \lambda &= \frac{1}{\sqrt{2}\pi n d^2} = \frac{k_B T}{\sqrt{2}\pi d^2 P} \quad [\lambda \propto \frac{T}{P}; \text{if } V=\text{const}, \lambda = \text{constant}]
\end{align}
\end{formulabox}

\newpage

% =========================================================================
% SECTION 5: OSCILLATIONS \& WAVES
% =========================================================================
\section{Oscillations \& Waves (Unit 10)}

\subsection{Simple Harmonic Motion (SHM) Kinematics \& Energy}

\begin{formulabox}[Governing Differential Equation, Velocity \& Acceleration in SHM]
\begin{align}
  \frac{d^2 x}{dt^2} + \omega^2 x &= 0 \implies x(t) = A\sin(\omega t + \phi) \\
  v(t) &= \frac{dx}{dt} = \omega\sqrt{A^2 - x^2}, \qquad a(t) = \frac{dv}{dt} = -\omega^2 x \\
  v_{\max} &= \omega A \quad (\text{at mean } x=0), \qquad a_{\max} = \omega^2 A \quad (\text{at extreme } x=\pm A)
\end{align}
\textbf{Energy Spectrum of Particle Executing SHM:}
\begin{align}
  K &= \frac{1}{2}m v^2 = \frac{1}{2}m\omega^2(A^2 - x^2), \qquad U = \frac{1}{2}m\omega^2 x^2 \\
  E_{\text{total}} &= K + U = \frac{1}{2}m\omega^2 A^2 = \text{constant for undamped SHM} \\
  \langle K \rangle_{\text{time}} &= \langle U \rangle_{\text{time}} = \frac{1}{4}m\omega^2 A^2 = \frac{1}{2}E_{\text{total}}, \qquad \langle K \rangle_{\text{position}} = \frac{2}{3}E_{\text{total}}, \quad \langle U \rangle_{\text{position}} = \frac{1}{3}E_{\text{total}}
\end{align}
\end{formulabox}

\subsection{Standard Oscillating Systems \& Combinations}

\begin{formulabox}[Time Period of Springs, Pendulums \& Liquid Columns]
\begin{align}
  T_{\text{spring}} &= 2\pi\sqrt{\frac{m}{k}}, \quad k_s = \frac{k_1 k_2}{k_1 + k_2} \text{ [Series]}, \quad k_p = k_1 + k_2 \text{ [Parallel]} \\
  T_{\text{cut spring}} &\implies k' = n k \implies T' = \frac{T}{\sqrt{n}} \quad [\text{Spring cut into } n \text{ equal pieces}] \\
  T_{\text{simple pendulum}} &= 2\pi\sqrt{\frac{L}{g_{\text{eff}}}}, \qquad g_{\text{eff}} = |\mathbf{g} - \mathbf{a}| \quad [\text{Lift accelerating upward: } g_{\text{eff}} = g+a] \\
  T_{\text{two block}} &= 2\pi\sqrt{\frac{\mu}{k}}, \qquad \mu = \frac{m_1 m_2}{m_1 + m_2} \quad [\text{Reduced Mass Formulation}] \\
  T_{\text{U-tube liquid}} &= 2\pi\sqrt{\frac{h}{g}} = 2\pi\sqrt{\frac{L}{2g}} \quad [h = \text{equilibrium column height}, L = \text{total liquid length}]
\end{align}
\end{formulabox}

\subsection{Wave Mechanics, Speed of Sound \& Organ Pipes}

\begin{formulabox}[Wave Speed, Harmonics in Strings \& Pipes]
\begin{align}
  v &= \frac{\omega}{k} = \nu\lambda, \qquad y(x,t) = A\sin(kx \mp \omega t + \phi) \quad [-\text{ for } +x\text{ direction}] \\
  v_{\text{string}} &= \sqrt{\frac{T}{\mu}}, \qquad v_{\text{gas}} = \sqrt{\frac{\gamma P}{\rho}} = \sqrt{\frac{\gamma R T}{M}} \quad [\text{Newton-Laplace Correction}] \\
  \nu_{\text{string, } n} &= \frac{n v}{2L} = \frac{n}{2L}\sqrt{\frac{T}{\mu}} \quad (n = 1, 2, 3, \dots) \quad [\text{All Harmonics Present}] \\
  \nu_{\text{closed pipe, } n} &= (2n - 1)\frac{v}{4L} \quad (n = 1, 2, 3, \dots) \implies \nu_1 : \nu_2 : \nu_3 = 1 : 3 : 5 \quad [\text{ODD Harmonics ONLY}] \\
  \nu_{\text{open pipe, } n} &= \frac{n v}{2L} \quad (n = 1, 2, 3, \dots) \implies \nu_1 : \nu_2 : \nu_3 = 1 : 2 : 3 \quad [\text{ALL Harmonics Present}]
\end{align}
\textbf{Resonance Tube Experiment \& Beat Frequency:}
\begin{align}
  \lambda &= 2(l_2 - l_1), \qquad e = \frac{l_2 - 3l_1}{2} \approx 0.6 r, \qquad v = 2\nu(l_2 - l_1) \\
  \nu_{\text{beat}} &= |\nu_1 - \nu_2| \quad [\text{Waxing decreases frequency: } \nu\downarrow; \text{Filing increases frequency: } \nu\uparrow]
\end{align}
\end{formulabox}

\begin{trapbox}[Wave Traps: Organ Pipes \& Speed of Sound vs Pressure]
\begin{itemize}
  \item \textbf{Sound Speed Independence from Pressure:} $v = \sqrt{\frac{\gamma P}{\rho}} = \sqrt{\frac{\gamma R T}{M}}$. At constant temperature, doubling pressure doubles density, leaving the ratio $P/\rho$ and sound speed $v$ \textbf{completely unchanged}.
  \item \textbf{Closed Pipe Odd Harmonics:} In a closed organ pipe of length $L$, the 3rd harmonic is $\nu = 3\left(\frac{v}{4L}\right)$ which is the \textbf{1st overtone}. The $n$-th overtone is the $(2n+1)$-th harmonic!
\end{itemize}
\end{trapbox}

\vspace{15pt}
\begin{center}
  \rule{0.6\textwidth}{0.8pt}\\[6pt]
  {\small\textbf{\color{primarydark}End of Module 02 Formula Handbook}}\\[2pt]
  {\footnotesize\color{accentblue}Prepared, Compiled \& Authored by Satya Prakash \quad\textbar\quad NEET UG 2027 Series}
\end{center}

\chapter{High-Yield Formula Compendium: Electrodynamics, Optics \& Modern Physics}

This chapter compiles all core formulas, vector relations, circuit equations, optical relations, quantum constants, and logic gate truth tables spanning Class 12 Physics (Units 11 through 19).

\section{Electrostatics \& Capacitance (Unit 11)}

\subsection{Coulomb's Law, Electric Field \& Dipoles}

\begin{formulabox}[Coulomb's Law, Medium Effect \& Dipole Dynamics]
\begin{align}
  \mathbf{F}_{12} &= \frac{1}{4\pi\varepsilon_0}\frac{q_1 q_2}{r^2}\hat{\mathbf{r}}_{12}, \qquad F_{\text{medium}} = \frac{F_{\text{vacuum}}}{\kappa} = \frac{F_0}{\varepsilon_r} && \left[\frac{1}{4\pi\varepsilon_0} = 8.988\times 10^9\,\si{N.m^2.C^{-2}}\right] \\
  \mathbf{p} &= q(2\mathbf{a}), \qquad \boldsymbol{\tau} = \mathbf{p}\times\mathbf{E}, \qquad U = -\mathbf{p}\cdot\mathbf{E} && [\mathbf{p} \text{ directed from } -q \text{ to } +q] \\
  E_{\text{axial}} &= \frac{1}{4\pi\varepsilon_0}\frac{2pr}{(r^2-a^2)^2} \approx \frac{2kp}{r^3} && [r \gg a; \text{ direction parallel to } \mathbf{p}] \\
  E_{\text{equatorial}} &= \frac{1}{4\pi\varepsilon_0}\frac{p}{(r^2+a^2)^{3/2}} \approx \frac{kp}{r^3} && [r \gg a; \text{ direction antiparallel to } \mathbf{p}] \\
  E(r,\theta) &= \frac{kp}{r^3}\sqrt{1 + 3\cos^2\theta}, \qquad \tan\alpha = \frac{1}{2}\tan\theta && [\theta = \text{angle with dipole axis}]
\end{align}
\textbf{Work Done in Rotating an Electric Dipole from $\theta_1$ to $\theta_2$:}
\begin{equation}
  W = U_f - U_i = pE(\cos\theta_1 - \cos\theta_2) \implies W_{0 \to 180^\circ} = 2pE \quad \text{[Max work; unstable equilibrium]}
\end{equation}
\end{formulabox}

\subsection{Gauss's Law \& Symmetrical Charge Distributions}

\begin{formulabox}[Gauss's Theorem \& Canonical Field Formulations]
\begin{equation}
  \Phi_E = \oint \mathbf{E}\cdot d\mathbf{A} = \frac{q_{\text{enclosed}}}{\varepsilon_0}
\end{equation}
\textbf{Symmetrical Electric Field Solutions:}
\begin{align}
  E_{\text{line}} &= \frac{\lambda}{2\pi\varepsilon_0 r} = \frac{2k\lambda}{r} && \text{[Infinitely Long Uniform Line Charge]} \\
  E_{\text{sheet}} &= \frac{\sigma}{2\varepsilon_0} && \text{[Infinite Non-Conducting Plane Sheet (independent of $r$)]} \\
  E_{\text{conductor}} &= \frac{\sigma}{\varepsilon_0} && \text{[Just outside any charged conducting surface]} \\
  E_{\text{shell}}(r) &= \begin{cases} 0 & (r < R) \\ \frac{kq}{R^2} & (r = R) \\ \frac{kq}{r^2} & (r > R) \end{cases}, \qquad V_{\text{shell}}(r) = \begin{cases} \frac{kq}{R} & (r \le R) \\ \frac{kq}{r} & (r > R) \end{cases} && \text{[Thin Spherical Shell of Radius } R\text{]}
\end{align}
\textbf{Uniformly Charged Solid Non-Conducting Sphere (Total Charge $q$, Radius $R$):}
\begin{equation}
  E_{\text{inside}}(r) = \frac{k q r}{R^3} = \frac{\rho r}{3\varepsilon_0}, \qquad V_{\text{inside}}(r) = \frac{kq}{2R^3}(3R^2 - r^2) \implies V_{\text{center}} = \frac{3}{2}\frac{kq}{R} = 1.5\,V_{\text{surface}}
\end{equation}
\end{formulabox}

\subsection{Capacitance, Dielectric Insertion \& Stored Energy}

\begin{formulabox}[Capacitance, Energy Density \& Force between Plates]
\begin{align}
  C_0 &= \frac{\varepsilon_0 A}{d}, \qquad C_{\text{dielectric}} = \frac{\varepsilon_0 A}{d - t\left(1 - \frac{1}{\kappa}\right)}, \qquad C_{\text{conducting slab}} = \frac{\varepsilon_0 A}{d - t} \\
  U &= \frac{1}{2}CV^2 = \frac{Q^2}{2C} = \frac{1}{2}QV, \qquad u_E = \frac{U}{\text{Volume}} = \frac{1}{2}\varepsilon_0 E^2 \quad [\si{J.m^{-3}}] \\
  F_{\text{attraction}} &= \frac{Q^2}{2\varepsilon_0 A} = \frac{1}{2}QE = \frac{1}{2}\varepsilon_0 A E^2, \qquad P_{\text{electrostatic}} = \frac{F}{A} = \frac{1}{2}\varepsilon_0 E^2
\end{align}
\textbf{Common Potential \& Heat Dissipation on Connecting Two Capacitors:}
\begin{equation}
  V_{\text{common}} = \frac{C_1 V_1 + C_2 V_2}{C_1 + C_2}, \qquad \Delta U_{\text{loss}} = \frac{C_1 C_2}{2(C_1 + C_2)}(V_1 - V_2)^2
\end{equation}
\end{formulabox}

\begin{table}[htbp]
\centering
\small
\renewcommand{\arraystretch}{1.2}
\begin{tabularx}{\textwidth}{L c c}
\toprule
\textbf{Physical Parameter} & \textbf{Case A: Battery REMAINING CONNECTED} & \textbf{Case B: Battery DISCONNECTED} \\
\midrule
Capacitance ($C$) & Increases by $\kappa$ ($C = \kappa C_0$) & Increases by $\kappa$ ($C = \kappa C_0$) \\
Potential Difference ($V$) & Remains Constant ($V = V_0$) & Decreases by $\kappa$ ($V = V_0/\kappa$) \\
Electric Field ($E = V/d$) & Remains Constant ($E = E_0$) & Decreases by $\kappa$ ($E = E_0/\kappa$) \\
Stored Charge ($Q$) & Increases by $\kappa$ ($Q = \kappa Q_0$) & Remains Constant ($Q = Q_0$) \\
Stored Energy ($U$) & Increases by $\kappa$ ($U = \kappa U_0$) & Decreases by $\kappa$ ($U = U_0/\kappa$) \\
\bottomrule
\end{tabularx}
\caption{Effect of Inserting a Dielectric Slab ($\kappa$) into a Parallel Plate Capacitor.}
\end{table}

\begin{trapbox}[Capacitor Traps: Force between Plates \& Common Potential]
\begin{itemize}
  \item \textbf{Factor of $1/2$ in Plate Force:} $F = \frac{1}{2}QE$, NOT $QE$. One plate experiences the field produced strictly by the \textbf{other} plate ($E/2 = \frac{\sigma}{2\varepsilon_0}$), not the total field $E$.
  \item \textbf{Opposite Polarity Connection:} If positive plate of $C_1$ is connected to negative plate of $C_2$, the numerator of common potential becomes $(C_1 V_1 - C_2 V_2)$ and energy loss is $\frac{C_1 C_2}{2(C_1 + C_2)}(V_1 + V_2)^2$.
\end{itemize}
\end{trapbox}

\newpage

% =========================================================================
% SECTION 2: CURRENT ELECTRICITY
% =========================================================================
\section{Current Electricity (Unit 12)}

\subsection{Microscopic Conduction, Drift Velocity \& Resistance}

\begin{formulabox}[Microscopic Drift, Ohm's Law \& Temperature Dependence]
\begin{align}
  I &= n e A v_d, \qquad v_d = \frac{e E \tau}{m} = \frac{e V \tau}{m L}, \qquad J = \frac{I}{A} = n e v_d = \sigma E = \frac{E}{\rho} \\
  \mu &= \frac{v_d}{E} = \frac{e\tau}{m} \quad [\si{m^2.V^{-1}.s^{-1}}], \qquad \rho = \frac{m}{n e^2 \tau}, \qquad R = \rho\frac{L}{A} \\
  R_T &= R_0(1 + \alpha\Delta T), \qquad \rho_T = \rho_0(1 + \alpha\Delta T) \quad [\alpha > 0 \text{ metals; } \alpha < 0 \text{ semiconductors}]
\end{align}
\textbf{Stretching of Wires (Constant Volume $V = A\cdot L = \text{const}$):}
\begin{equation}
  R \propto L^2 \propto \frac{1}{A^2} \propto \frac{1}{r^4} \implies \text{If } L \to nL, \quad R' = n^2 R; \quad \text{For small } \%\Delta L < 5\%, \quad \frac{\Delta R}{R} \approx 2\frac{\Delta L}{L}
\end{equation}
\end{formulabox}

\subsection{Cells, Internal Resistance \& Power Transfer}

\begin{formulabox}[Terminal EMF, Cell Combinations \& Max Power Theorem]
\begin{align}
  V &= \mathcal{E} - Ir \quad \text{[Discharging; } V < \mathcal{E}\text{]}, \qquad V = \mathcal{E} + Ir \quad \text{[Charging; } V > \mathcal{E}\text{]} \\
  r &= R\left(\frac{\mathcal{E}}{V} - 1\right) = R\left(\frac{l_1 - l_2}{l_2}\right) && \text{[Internal Resistance Formulation]} \\
  \mathcal{E}_{\text{eq, parallel}} &= \frac{\frac{\mathcal{E}_1}{r_1} + \frac{\mathcal{E}_2}{r_2}}{\frac{1}{r_1} + \frac{1}{r_2}} = \frac{\mathcal{E}_1 r_2 + \mathcal{E}_2 r_1}{r_1 + r_2}, \qquad r_{\text{eq, parallel}} = \frac{r_1 r_2}{r_1 + r_2}
\end{align}
\textbf{Maximum Power Transfer Theorem (Load $R$ connected to cell $\mathcal{E}, r$):}
\begin{equation}
  P_{\text{load}} = I^2 R = \frac{\mathcal{E}^2 R}{(R + r)^2} \implies P = P_{\max} \iff R = r \implies P_{\max} = \frac{\mathcal{E}^2}{4r}
\end{equation}
\end{formulabox}

\subsection{Bridge Circuits \& Electrical Instruments}

\begin{formulabox}[Wheatstone Bridge \& Metre Bridge Balance Equations]
\begin{equation}
  \frac{P}{Q} = \frac{R}{S} \iff I_G = 0 \quad \text{[Wheatstone Balance Condition: Galv. current is null]}
\end{equation}
\textbf{Metre Bridge Unknown Resistance $S$ (Balance length $l$ from left):}
\begin{equation}
  \frac{R}{S} = \frac{l}{100 - l} \implies S = R\left(\frac{100 - l}{l}\right)
\end{equation}
\end{formulabox}

\begin{trapbox}[Current Electricity Traps: Cell Grouping \& Terminal Voltage]
\begin{itemize}
  \item \textbf{Reversed Cell in Series Battery:} In a combination of $n$ identical cells each of EMF $\mathcal{E}$ and internal resistance $r$, if $m$ cells are connected with reversed polarity:
  \begin{equation*}
    \mathcal{E}_{\text{net}} = (n - 2m)\mathcal{E}, \qquad r_{\text{net}} = n r \quad (\text{Internal resistances always add!})
  \end{equation*}
  \item \textbf{Short Circuit Terminal Potential:} When cell is short-circuited ($R=0$), current is maximum $I_{\text{sc}} = \mathcal{E}/r$, and terminal voltage drops to \textbf{zero} ($V = \mathcal{E} - I_{\text{sc}}r = 0$).
\end{itemize}
\end{trapbox}

\newpage

% =========================================================================
% SECTION 3: MAGNETIC EFFECTS OF CURRENT \& MAGNETISM
% =========================================================================
\section{Magnetic Effects of Current \& Magnetism (Unit 13)}

\subsection{Biot-Savart Law \& Ampere's Circuital Law}

\begin{formulabox}[Biot-Savart Law \& Symmetrical Field Equations]
\begin{align}
  d\mathbf{B} &= \frac{\mu_0}{4\pi}\frac{I d\mathbf{l}\times\hat{\mathbf{r}}}{r^2} && \left[\frac{\mu_0}{4\pi} = 10^{-7}\,\si{T.m.A^{-1}}\right] \\
  B_{\text{straight, finite}} &= \frac{\mu_0 I}{4\pi d}(\sin\phi_1 + \sin\phi_2) \implies B_{\text{infinite}} = \frac{\mu_0 I}{2\pi d}, \quad B_{\text{semi-infinite}} = \frac{\mu_0 I}{4\pi d} \\
  B_{\text{circular coil, axis}} &= \frac{\mu_0 N I R^2}{2(R^2 + x^2)^{3/2}} \implies B_{\text{center}} = \frac{\mu_0 N I}{2R} \quad (x = 0) \\
  B_{\text{arc}} &= \frac{\mu_0 I \theta}{4\pi R} \quad [\theta \text{ in radians}], \qquad B_{\text{solenoid}} = \mu_0 n I = \mu_0 \left(\frac{N}{L}\right)I \\
  \oint \mathbf{B}\cdot d\mathbf{l} &= \mu_0 I_{\text{enclosed}} && \text{[Ampere's Circuital Law]}
\end{align}
\end{formulabox}

\subsection{Lorentz Force, Particle Trajectories \& Force on Conductors}

\begin{formulabox}[Lorentz Equation \& Charged Particle Dynamics in Magnetic Fields]
\begin{align}
  \mathbf{F} &= q(\mathbf{E} + \mathbf{v}\times\mathbf{B}) \implies \mathbf{F}_B = q(\mathbf{v}\times\mathbf{B}) \implies W_B = \int \mathbf{F}_B\cdot d\mathbf{r} = 0 \quad [\text{Speed } v = \text{const}] \\
  r &= \frac{mv}{qB} = \frac{p}{qB} = \frac{\sqrt{2mK}}{qB} = \frac{\sqrt{2mqV}}{qB} && \text{[Circular Orbit Radius for } \mathbf{v}\perp\mathbf{B}\text{]} \\
  T &= \frac{2\pi m}{qB}, \qquad \nu = \frac{qB}{2\pi m} && \text{[Period and Frequency: INDEPENDENT of speed } v\text{]} \\
  p_{\text{pitch}} &= v_{\parallel} T = \frac{2\pi m v\cos\theta}{qB} && \text{[Helical Path Pitch for launch angle } \theta\text{]} \\
  \mathbf{F} &= I(\mathbf{L}\times\mathbf{B}), \qquad \frac{F}{L} = \frac{\mu_0 I_1 I_2}{2\pi d} && \text{[Parallel attract; antiparallel repel]}
\end{align}
\end{formulabox}

\subsection{Galvanometer Conversions \& Magnetic Materials}

\begin{formulabox}[Moving Coil Galvanometer, Ammeter, Voltmeter \& Magnetic Matter]
\begin{align}
  \tau &= N I A B = C\theta \implies I = \left(\frac{C}{NAB}\right)\theta, \quad I_s = \frac{\theta}{I} = \frac{NAB}{C}, \quad V_s = \frac{\theta}{V} = \frac{I_s}{R_G} \\
  S &= \frac{I_g R_g}{I - I_g} \quad \text{[Ammeter Shunt in Parallel]}, \qquad R_s = \frac{V}{I_g} - R_g \quad \text{[Voltmeter Multiplier in Series]} \\
  \mathbf{B} &= \mu_0(\mathbf{H} + \mathbf{M}) = \mu H, \qquad \chi = \frac{M}{H}, \qquad \mu_r = 1 + \chi = \frac{\mu}{\mu_0}
\end{align}
\textbf{Curie's Law for Paramagnetic Materials \& Curie-Weiss Law for Ferromagnets:}
\begin{equation}
  \chi_{\text{para}} = \frac{C}{T} \implies \chi \propto \frac{1}{T}, \qquad \chi_{\text{ferro}} = \frac{C'}{T - T_c} \quad (T > T_c)
\end{equation}
\end{formulabox}

\begin{trapbox}[Magnetism Traps: Work Done by Magnetic Force \& Galvanometer Sensitivity]
\begin{itemize}
  \item \textbf{Work Done by Magnetic Force is Zero:} Because $\mathbf{F}_B \perp \mathbf{v}$ at every instant, $\mathbf{F}_B\cdot\mathbf{v} = 0$. \textbf{A static magnetic field does zero work on a charged particle}. It changes the direction of motion, never the kinetic energy or speed.
  \item \textbf{Current Sensitivity vs Voltage Sensitivity:} Doubling the number of turns $N$ doubles current sensitivity ($I_s \propto N$), but resistance also doubles ($R \propto N$), leaving voltage sensitivity $V_s = I_s / R_G$ \textbf{completely unchanged}!
\end{itemize}
\end{trapbox}

\newpage

% =========================================================================
% SECTION 4: ELECTROMAGNETIC INDUCTION \& ALTERNATING CURRENT
% =========================================================================
\section{Electromagnetic Induction \& Alternating Current (Unit 14)}

\subsection{Faraday's Laws, Motional EMF \& Inductance}

\begin{formulabox}[Faraday-Lenz Law, Motional EMF \& Self/Mutual Inductance]
\begin{align}
  \mathcal{E} &= -\frac{d\Phi_B}{dt} = -N\frac{d}{dt}(B A\cos\theta), \qquad q_{\text{induced}} = \frac{\Delta\Phi_B}{R} \quad [\text{Independent of time!}] \\
  \mathcal{E}_{\text{motional}} &= B v l \quad [\text{Translational Rod}], \qquad \mathcal{E}_{\text{rotational}} = \frac{1}{2}B\omega L^2 \quad [\text{Rotating Rod about one end}] \\
  \mathcal{E} &= -L\frac{dI}{dt}, \qquad U = \frac{1}{2}L I^2, \qquad L_{\text{solenoid}} = \mu_0 n^2 A l = \frac{\mu_0 N^2 A}{l} \\
  \mathcal{E}_2 &= -M\frac{dI_1}{dt}, \qquad M = k\sqrt{L_1 L_2} \quad (0 \le k \le 1)
\end{align}
\end{formulabox}

\subsection{Alternating Current Circuits, Impedance \& Resonance}

\begin{formulabox}[AC Values, Reactance, Series LCR Circuit \& Resonance]
\begin{align}
  I_{\text{rms}} &= \frac{I_0}{\sqrt{2}} \approx 0.707 I_0, \qquad V_{\text{rms}} = \frac{V_0}{\sqrt{2}} \approx 0.707 V_0, \qquad I_{\text{avg, half}} = \frac{2I_0}{\pi} \approx 0.637 I_0 \\
  X_L &= \omega L = 2\pi\nu L, \qquad X_C = \frac{1}{\omega C} = \frac{1}{2\pi\nu C} \quad [X_L \text{ leads by } \pi/2; X_C \text{ lags by } \pi/2] \\
  Z &= \sqrt{R^2 + (X_L - X_C)^2}, \qquad \tan\phi = \frac{X_L - X_C}{R}, \qquad V_0 = \sqrt{V_R^2 + (V_L - V_C)^2} \\
  \omega_0 &= \frac{1}{\sqrt{LC}} \iff \nu_0 = \frac{1}{2\pi\sqrt{LC}} \quad [\text{Series Resonance Condition: } X_L = X_C] \\
  Z_{\min} &= R, \quad I_{\max} = \frac{V_{\text{rms}}}{R}, \quad \phi = 0, \quad Q = \frac{\omega_0 L}{R} = \frac{1}{\omega_0 C R} = \frac{1}{R}\sqrt{\frac{L}{C}} = \frac{\nu_0}{\Delta\nu}
\end{align}
\textbf{Average Power, Power Factor \& Ideal Transformer:}
\begin{align}
  P_{\text{avg}} &= V_{\text{rms}} I_{\text{rms}} \cos\phi = I_{\text{rms}}^2 R, \qquad \cos\phi = \frac{R}{Z} \quad [\text{Wattless Current } I_w = I_{\text{rms}}\sin\phi] \\
  \frac{V_s}{V_p} &= \frac{N_s}{N_p} = \frac{I_p}{I_s} = K, \qquad \eta = \frac{P_{\text{out}}}{P_{\text{in}}} = \frac{V_s I_s}{V_p I_p}
\end{align}
\end{formulabox}

\newpage

% =========================================================================
% SECTION 5: ELECTROMAGNETIC WAVES
% =========================================================================
\section{Electromagnetic Waves (Unit 15)}

\begin{formulabox}[Displacement Current, Wave Speed, Energy \& Radiation Pressure]
\begin{align}
  I_d &= \varepsilon_0 \frac{d\Phi_E}{dt}, \qquad \oint \mathbf{B}\cdot d\mathbf{l} = \mu_0\left(I_c + \varepsilon_0\frac{d\Phi_E}{dt}\right) && \text{[Ampere-Maxwell Law]} \\
  c &= \frac{1}{\sqrt{\mu_0\varepsilon_0}} = \frac{E_0}{B_0} \approx 3\times 10^8\,\si{m.s^{-1}}, \qquad v_{\text{medium}} = \frac{1}{\sqrt{\mu\varepsilon}} = \frac{c}{\sqrt{\mu_r \varepsilon_r}} \\
  u_{\text{total}} &= u_E + u_B = \frac{1}{2}\varepsilon_0 E_{\text{rms}}^2 + \frac{B_{\text{rms}}^2}{2\mu_0} = \varepsilon_0 E_{\text{rms}}^2 = \frac{B_{\text{rms}}^2}{\mu_0} && \text{[Equal Electric \& Magnetic Energy]} \\
  I &= \langle S \rangle = \frac{1}{2}\varepsilon_0 c E_0^2 = \varepsilon_0 c E_{\text{rms}}^2, \qquad P_{\text{rad}} = \begin{cases} \frac{I}{c} & \text{(Total Absorption)} \\ \frac{2I}{c} & \text{(Total Reflection)} \end{cases}
\end{align}
\end{formulabox}

\begin{table}[htbp]
\centering
\small
\renewcommand{\arraystretch}{1.2}
\begin{tabularx}{\textwidth}{L c c L}
\toprule
\textbf{Wave Type} & \textbf{Wavelength ($\lambda$)} & \textbf{Frequency ($\nu$)} & \textbf{Primary Source \& High-Yield Application} \\
\midrule
$\gamma$-Rays & $< 10^{-3}\,\si{nm}$ & $> 10^{20}\,\si{Hz}$ & Radioactive decay; cancer radiotherapy \\
X-Rays & $1\,\si{nm} - 10^{-3}\,\si{nm}$ & $10^{18} - 10^{20}\,\si{Hz}$ & Fast electrons hitting heavy metal target; bone imaging \\
UV Rays & $400\,\si{nm} - 1\,\si{nm}$ & $10^{15} - 10^{17}\,\si{Hz}$ & Very hot bodies (Sun), arcs; LASIK eye surgery, germicidal \\
Visible & $700\,\si{nm} - 400\,\si{nm}$ & $4\times 10^{14} - 7.5\times 10^{14}\,\si{Hz}$ & Atomic electron transitions; human optical vision \\
Infrared & $1\,\si{mm} - 700\,\si{nm}$ & $10^{11} - 4\times 10^{14}\,\si{Hz}$ & Hot bodies, molecules; thermal therapy, remote controls \\
Microwaves & $0.1\,\si{m} - 1\,\si{mm}$ & $10^9 - 10^{11}\,\si{Hz}$ & Klystrons, magnetrons; RADAR systems, microwave ovens \\
Radio Waves & $> 0.1\,\si{m}$ & $< 10^9\,\si{Hz}$ & Accelerated charges in antennas; cellular/radio broadcasting \\
\bottomrule
\end{tabularx}
\caption{Electromagnetic Spectrum in Decreasing Order of Frequency / Increasing Wavelength.}
\end{table}

\vspace{15pt}
\begin{center}
  \rule{0.6\textwidth}{0.8pt}\\[6pt]
  {\small\textbf{\color{primarydark}End of Module 03 Formula Handbook}}\\[2pt]
  {\footnotesize\color{accentblue}Prepared, Compiled \& Authored by Satya Prakash \quad\textbar\quad NEET UG 2027 Series}
\end{center}

\section{Ray Optics \& Optical Instruments (Unit 16)}

\subsection{Spherical Mirrors, Snell's Law \& Critical Angle}

\begin{formulabox}[Mirror Formula, Magnification \& Normal Shift]
\begin{align}
  \frac{1}{v} + \frac{1}{u} &= \frac{1}{f} = \frac{2}{R} && \text{[Cartesian Sign Convention: Pole is origin; light direction is $+x$]} \\
  m &= -\frac{v}{u} = \frac{f}{f - u} = \frac{f - v}{f} && \text{[Transverse / Linear Magnification]} \\
  m_L &= \frac{dv}{du} = -m^2 = -\left(\frac{v}{u}\right)^2 && \text{[Longitudinal Magnification for small object along axis]} \\
  \mu_1 \sin i &= \mu_2 \sin r \implies \frac{\sin i}{\sin r} = \frac{\mu_2}{\mu_1} = \frac{v_1}{v_2} = \frac{\lambda_1}{\lambda_2} && \text{[Snell's Law of Refraction; Frequency $\nu$ remains invariant]} \\
  d_{\text{apparent}} &= \frac{d_{\text{real}}}{\mu_{\text{relative}}} = d\left(\frac{\mu_{\text{observer}}}{\mu_{\text{object}}}\right), \qquad \Delta t = t\left(1 - \frac{1}{\mu}\right) && \text{[Apparent Depth \& Normal Shift through slab]}
\end{align}
\textbf{Total Internal Reflection (TIR) \& Underwater Illuminance Cone:}
\begin{equation}
  \sin i_c = \frac{\mu_{\text{rarer}}}{\mu_{\text{denser}}} = \frac{1}{\mu}, \qquad r_{\text{circle}} = \frac{h}{\sqrt{\mu^2 - 1}}, \qquad \text{Area} = \frac{\pi h^2}{\mu^2 - 1} \quad [\text{Point source at depth } h]
\end{equation}
\end{formulabox}

\subsection{Prisms, Dispersion \& Spherical Refracting Surfaces}

\begin{formulabox}[Prism Deviation, Dispersive Power \& Curved Surface Refraction]
\begin{align}
  \delta &= (i_1 + i_2) - A, \qquad r_1 + r_2 = A && \text{[General Triangular Prism of Angle } A\text{]} \\
  \mu &= \frac{\sin\left(\frac{A + \delta_m}{2}\right)}{\sin\left(\frac{A}{2}\right)} \iff i_1 = i_2 = \frac{A + \delta_m}{2}, \quad r_1 = r_2 = \frac{A}{2} && \text{[Minimum Deviation Condition; ray parallel to base]} \\
  \delta_{\text{thin}} &= (\mu - 1)A, \qquad \theta_{\text{dispersion}} = (\mu_V - \mu_R)A, \qquad \omega = \frac{\mu_V - \mu_R}{\mu_Y - 1} && \text{[Thin Prism ($A \le 10^\circ$) \& Dispersive Power $\omega$]} \\
  \frac{\mu_2}{v} - \frac{\mu_1}{u} &= \frac{\mu_2 - \mu_1}{R}, \qquad m = \frac{\mu_1 v}{\mu_2 u} && \text{[Refraction at Single Spherical Interface]}
\end{align}
\end{formulabox}

\subsection{Thin Lenses, Combinations \& Silvered Systems}

\begin{formulabox}[Lens Maker's Equation, Power \& Silvering Formulations]
\begin{align}
  \frac{1}{f} &= \left(\frac{\mu_{\text{lens}}}{\mu_{\text{medium}}} - 1\right)\left(\frac{1}{R_1} - \frac{1}{R_2}\right) && \text{[Lens Maker's Formula]} \\
  \frac{1}{v} - \frac{1}{u} &= \frac{1}{f}, \qquad m = \frac{v}{u} = \frac{f}{f + u} = \frac{f - v}{f} && \text{[Thin Lens Equation \& Magnification]} \\
  P &= \frac{1}{f\text{ (in m)}} = \frac{100}{f\text{ (in cm)}} \quad [\si{D = m^{-1}}], \qquad P_{\text{contact}} = \sum P_i = P_1 + P_2 \\
  \frac{1}{F} &= \frac{1}{f_1} + \frac{1}{f_2} - \frac{d}{f_1 f_2} \implies P = P_1 + P_2 - d P_1 P_2 && \text{[Two Lenses Separated by Distance } d\text{]}
\end{align}
\textbf{Silvering of a Lens Surface (Acts as an Equivalent Spherical Mirror):}
\begin{equation}
  P_{\text{eq}} = 2P_{\text{lens}} + P_{\text{mirror}} = \frac{2}{f_L} + \frac{1}{f_M} \implies F_{\text{eq}} = -\frac{1}{P_{\text{eq}}} \quad [\text{System always acts as a mirror}]
\end{equation}
\end{formulabox}

\begin{usecasebox}[Lens Cutting Rules \& Medium Immersion Effects]
\begin{itemize}
  \item \textbf{Vertical Cut (Perpendicular to Principal Axis):} Cutting a convex lens of focal length $f$ into two halves produces two plano-convex lenses: $f' = 2f, P' = P/2$. Intensity remains \textbf{unchanged}.
  \item \textbf{Horizontal Cut (Along Principal Axis):} Each half retains identical curvature: $f' = f, P' = P$. Intensity is halved because light-gathering aperture area is reduced by $50\%$.
  \item \textbf{Immersion in Liquid ($\mu_{\text{med}}$):}
  \begin{equation*}
    \frac{f_{\text{liquid}}}{f_{\text{air}}} = \frac{(\mu_g - 1)}{\left(\frac{\mu_g}{\mu_l} - 1\right)} \implies \text{For glass } (\mu=1.5) \text{ in water } (\mu=4/3): \quad f_{\text{water}} = 4 f_{\text{air}}
  \end{equation*}
  If $\mu_l > \mu_g$, focal length reverses sign: a converging convex lens behaves as a diverging concave lens!
\end{itemize}
\end{usecasebox}

\subsection{Optical Instruments}

\begin{table}[htbp]
\centering
\small
\renewcommand{\arraystretch}{1.2}
\begin{tabularx}{\textwidth}{L c c c}
\toprule
\textbf{Instrument} & \textbf{Normal Adjustment ($v = \infty$)} & \textbf{Near Point Adjustment ($v = D = 25\si{cm}$)} & \textbf{Tube Length ($L$)} \\
\midrule
Simple Microscope & $m = \frac{D}{f}$ & $m = 1 + \frac{D}{f}$ & Eye at lens \\
Compound Microscope & $m = -\left(\frac{v_o}{u_o}\right)\left(\frac{D}{f_e}\right) \approx -\frac{L}{f_o}\frac{D}{f_e}$ & $m = -\left(\frac{v_o}{u_o}\right)\left(1 + \frac{D}{f_e}\right) \approx -\frac{L}{f_o}\left(1 + \frac{D}{f_e}\right)$ & $L = v_o + u_e \approx v_o + f_e$ \\
Astronomical Telescope & $m = -\frac{f_o}{f_e}$ & $m = -\frac{f_o}{f_e}\left(1 + \frac{f_e}{D}\right)$ & $L = f_o + f_e \quad \text{or} \quad f_o + u_e$ \\
\bottomrule
\end{tabularx}
\caption{Magnifying Powers and Tube Lengths of Canonical Optical Instruments.}
\end{table}

\newpage

% =========================================================================
% SECTION 2: WAVE OPTICS
% =========================================================================
\section{Wave Optics (Unit 16)}

\subsection{Interference \& Young's Double Slit Experiment (YDSE)}

\begin{formulabox}[Interference Intensity, Fringe Position \& Fringe Width]
\begin{align}
  I &= I_1 + I_2 + 2\sqrt{I_1 I_2}\cos\phi \implies I = 4I_0 \cos^2\left(\frac{\phi}{2}\right) \quad (\text{when } I_1 = I_2 = I_0) \\
  \frac{I_{\max}}{I_{\min}} &= \left(\frac{\sqrt{I_1} + \sqrt{I_2}}{\sqrt{I_1} - \sqrt{I_2}}\right)^2 = \left(\frac{a_1 + a_2}{a_1 - a_2}\right)^2, \qquad \Delta\phi = \frac{2\pi}{\lambda}\Delta x \\
  y_n &= n\frac{\lambda D}{d} \quad (n = 0, 1, 2, \dots) && \text{[$n$-th Bright Fringe / Maximum; } \Delta x = n\lambda\text{]} \\
  y_n' &= (2n - 1)\frac{\lambda D}{2d} \quad (n = 1, 2, 3, \dots) && \text{[$n$-th Dark Fringe / Minimum; } \Delta x = (2n-1)\frac{\lambda}{2}\text{]} \\
  \beta &= \frac{\lambda D}{d}, \qquad \theta = \frac{\beta}{D} = \frac{\lambda}{d} \quad [\si{rad}] && \text{[Linear Fringe Width $\beta$ \& Angular Fringe Width $\theta$]}
\end{align}
\textbf{Shift Produced by Thin Dielectric Slab of Thickness $t$ and Refractive Index $\mu$:}
\begin{equation}
  \Delta y = \frac{D}{d}(\mu - 1)t, \qquad N_{\text{shifted}} = \frac{\Delta y}{\beta} = \frac{(\mu - 1)t}{\lambda}
\end{equation}
\end{formulabox}

\subsection{Diffraction due to a Single Slit}

\begin{formulabox}[Single Slit Diffraction Maxima, Minima \& Fresnel Distance]
\begin{align}
  a\sin\theta &= n\lambda \implies y_n = n\frac{\lambda D}{a} \quad (n = 1, 2, 3, \dots) && \text{[Diffraction Minima Condition]} \\
  a\sin\theta &\approx (2n + 1)\frac{\lambda}{2} \implies y_n' = (2n + 1)\frac{\lambda D}{2a} && \text{[Secondary Maxima Condition ($n \ge 1$)]} \\
  2\theta_0 &= \frac{2\lambda}{a}, \qquad W_{\text{central}} = 2y_1 = \frac{2\lambda D}{a} = 2\beta_{\text{diffraction}} && \text{[Width of Central Diffraction Maximum]} \\
  Z_F &= \frac{a^2}{\lambda} && \text{[Fresnel Distance: Ray optics valid for } z \ll Z_F\text{]}
\end{align}
\end{formulabox}

\begin{trapbox}[Wave Optics Traps: Interference vs Diffraction Central Fringe]
\begin{itemize}
  \item \textbf{Central Maximum Width in Diffraction is Doubled:} In single-slit diffraction, the central maximum spreads between the first minima on both sides ($-\lambda/a$ to $+\lambda/a$). Its linear width is $\frac{2\lambda D}{a}$, exactly \textbf{twice} the width of secondary maxima ($\frac{\lambda D}{a}$). In YDSE, all fringes have identical width $\beta = \frac{\lambda D}{d}$.
  \item \textbf{YDSE Immersed in Liquid:} When apparatus is dipped in liquid of refractive index $\mu$, wavelength shrinks ($\lambda' = \lambda/\mu$). Hence fringe width decreases: $\beta' = \beta/\mu$. Angular width also shrinks: $\theta' = \theta/\mu$.
\end{itemize}
\end{trapbox}

\newpage

% =========================================================================
% SECTION 3: DUAL NATURE OF MATTER \& RADIATION
% =========================================================================
\section{Dual Nature of Matter \& Radiation (Unit 17)}

\subsection{Photoelectric Effect \& Einstein's Equation}

\begin{formulabox}[Einstein's Photoelectric Governing Equations]
\begin{align}
  E &= h\nu = \Phi_0 + K_{\max} = \Phi_0 + e V_0 && \text{[Einstein's Photoelectric Law]} \\
  \Phi_0 &= h\nu_0 = \frac{hc}{\lambda_0} \approx \frac{1240}{\lambda_0\text{ (in nm)}}\,\si{eV} && \text{[Work Function $\Phi_0$; Cutoff wavelength $\lambda_0$]} \\
  e V_0 &= h\nu - \Phi_0 \implies V_0 = \left(\frac{h}{e}\right)\nu - \frac{\Phi_0}{e} && \text{[Slope of $V_0-\nu$ graph is universal: $\text{Slope} = h/e$]} \\
  K_{\max} &= hc\left(\frac{1}{\lambda} - \frac{1}{\lambda_0}\right) = 1240\left(\frac{1}{\lambda\text{ [nm]}} - \frac{1}{\lambda_0\text{ [nm]}}\right)\,\si{eV}
\end{align}
\textbf{Photon Dynamics (Rest Mass $m_0 = 0$):}
\begin{equation}
  E = h\nu = \frac{hc}{\lambda}, \qquad p = \frac{E}{c} = \frac{h}{\lambda}, \qquad m_{\text{effective}} = \frac{h\nu}{c^2}, \qquad N_{\text{photons/sec}} = \frac{P}{h\nu} = \frac{P\lambda}{hc}
\end{equation}
\end{formulabox}

\subsection{de Broglie Matter Wave Formulations}

\begin{formulabox}[de Broglie Wavelength of Particles \& Accelerating Potentials]
\begin{align}
  \lambda &= \frac{h}{p} = \frac{h}{mv} = \frac{h}{\sqrt{2mK}} = \frac{h}{\sqrt{2mqV}} && \text{[Fundamental Matter Wave Equation]} \\
  \lambda_{\text{electron}} &= \frac{h}{\sqrt{2m_e e V}} = \sqrt{\frac{150}{V}}\,\text{\AA} = \frac{1.227}{\sqrt{V\text{ [Volts]}}}\,\si{nm} && [V = \text{accelerating potential in Volts}] \\
  \lambda_{\text{proton}} &= \frac{0.286}{\sqrt{V}}\,\text{\AA}, \qquad \lambda_{\alpha} = \frac{0.101}{\sqrt{V}}\,\text{\AA}, \qquad \lambda_{\text{deuteron}} = \frac{0.202}{\sqrt{V}}\,\text{\AA} \\
  \lambda_{\text{gas molecule}} &= \frac{h}{\sqrt{3m k_B T}} = \frac{h}{\sqrt{2m\left(\frac{3}{2}k_B T\right)}} && \text{[Thermal de Broglie Wavelength at Temp } T\text{]}
\end{align}
\end{formulabox}

\begin{trapbox}[Photoelectric Traps: Intensity vs Frequency Effects]
\begin{itemize}
  \item \textbf{Intensity Modifies Current, NOT Kinetic Energy:} Doubling the intensity of incident light doubles the number of photons per second, thereby doubling the \textbf{saturation photoelectric current} ($I_{\text{sat}} \propto \text{Intensity}$). The maximum kinetic energy $K_{\max}$ and stopping potential $V_0$ remain \textbf{strictly unchanged}.
  \item \textbf{Frequency Modifies Kinetic Energy, NOT Current:} Increasing frequency $\nu$ increases $K_{\max}$ and $V_0$ linearly ($V_0 = \frac{h}{e}\nu - \frac{\Phi_0}{e}$). Saturation current is independent of frequency above $\nu_0$.
\end{itemize}
\end{trapbox}

\newpage

% =========================================================================
% SECTION 4: ATOMS \& NUCLEI
% =========================================================================
\section{Atoms \& Nuclei (Unit 18)}

\subsection{Bohr's Hydrogen-Like Atom Model}

\begin{formulabox}[Bohr Postulates, Radius, Velocity, Energy \& Rydberg Formula]
\begin{align}
  L &= m v_n r_n = \frac{n h}{2\pi} \quad (n = 1, 2, 3, \dots) && \text{[Bohr Angular Momentum Quantization]} \\
  r_n &= \frac{n^2 h^2 \varepsilon_0}{\pi m Z e^2} = 0.529\left(\frac{n^2}{Z}\right)\,\text{\AA} = r_1\left(\frac{n^2}{Z}\right) && \text{[Orbital Radius: } r \propto \frac{n^2}{Z}\text{]} \\
  v_n &= \frac{Z e^2}{2\varepsilon_0 n h} = \left(\frac{c}{137}\right)\frac{Z}{n} \approx 2.18\times 10^6\left(\frac{Z}{n}\right)\,\si{m.s^{-1}} && \text{[Orbital Velocity: } v \propto \frac{Z}{n}\text{]} \\
  E_n &= -\frac{m Z^2 e^4}{8\varepsilon_0^2 n^2 h^2} = -13.6\left(\frac{Z^2}{n^2}\right)\,\si{eV} && \text{[Total Energy: } K = -E_n, \quad U = 2E_n\text{]} \\
  \frac{1}{\lambda} &= R Z^2\left(\frac{1}{n_1^2} - \frac{1}{n_2^2}\right) && \left[R = 1.097\times 10^7\,\si{m^{-1}}, \quad \frac{1}{R} \approx 912\,\text{\AA}\right]
\end{align}
\textbf{Total Number of Emission Lines from Excited State $n$ to Ground State ($n=1$):}
\begin{equation}
  N = \frac{n(n - 1)}{2}, \qquad N_{\text{between } n_2 \text{ and } n_1} = \frac{(n_2 - n_1)(n_2 - n_1 + 1)}{2}
\end{equation}
\end{formulabox}

\begin{table}[htbp]
\centering
\small
\renewcommand{\arraystretch}{1.2}
\begin{tabularx}{\textwidth}{L c c c c}
\toprule
\textbf{Spectral Series} & \textbf{Lower Level ($n_1$)} & \textbf{Upper Level ($n_2$)} & \textbf{Spectral Region} & \textbf{Shortest Wavelength ($\lambda_{\min}$)} \\
\midrule
Lyman Series & $n_1 = 1$ & $2, 3, 4, \dots$ & Ultraviolet (UV) & $1/R \approx 912\,\text{\AA}$ \\
Balmer Series & $n_1 = 2$ & $3, 4, 5, \dots$ & Visible \& Near UV & $4/R \approx 3646\,\text{\AA}$ \\
Paschen Series & $n_1 = 3$ & $4, 5, 6, \dots$ & Infrared (IR) & $9/R \approx 8205\,\text{\AA}$ \\
Brackett Series & $n_1 = 4$ & $5, 6, 7, \dots$ & Infrared (IR) & $16/R \approx 14586\,\text{\AA}$ \\
Pfund Series & $n_1 = 5$ & $6, 7, 8, \dots$ & Far Infrared & $25/R \approx 22790\,\text{\AA}$ \\
\bottomrule
\end{tabularx}
\caption{Hydrogen Spectral Series Classifications and Spectral Regimes.}
\end{table}

\subsection{Nuclear Size, Mass Defect \& Binding Energy}

\begin{formulabox}[Nuclear Radius, Density, Mass Defect \& Fission/Fusion Energies]
\begin{align}
  R &= R_0 A^{1/3} \quad [R_0 \approx 1.2\,\si{fm} = 1.2\times 10^{-15}\,\si{m}], \qquad \rho_{\text{nuclear}} = \frac{M}{\frac{4}{3}\pi R^3} \approx 2.3\times 10^{17}\,\si{kg.m^{-3}} \\
  \Delta m &= [Z m_p + (A - Z)m_n] - M_{\text{nucleus}} && \text{[Mass Defect]} \\
  E_b &= \Delta m\,c^2 = \Delta m\text{ (in u)}\times 931.5\,\si{MeV} && \text{[Total Binding Energy]} \\
  \Delta E_{\text{reaction}} &= Q = [m_{\text{reactants}} - m_{\text{products}}]c^2 = [BE_{\text{products}} - BE_{\text{reactants}}]
\end{align}
\end{formulabox}

\begin{trapbox}[Nuclear Traps: Invariance of Nuclear Density]
\begin{itemize}
  \item \textbf{Nuclear Density is Constant:} Because $R \propto A^{1/3}$, nuclear volume $V \propto R^3 \propto A$. Thus nuclear density $\rho = \frac{m A}{V} \propto \frac{A}{A} = \text{constant}$. A Uranium nucleus ($A=238$) has approximately the \textbf{same density} as a Helium nucleus ($A=4$).
  \item \textbf{Radioactivity Excluded by NMC:} Note that $\alpha, \beta, \gamma$ decay rate equations and half-life decay laws have been removed from NEET UG under current NMC guidelines; focus strictly on binding energy per nucleon curves, mass defects, and fission/fusion energy calculations.
\end{itemize}
\end{trapbox}

\newpage

% =========================================================================
% SECTION 5: ELECTRONIC DEVICES (SEMICONDUCTORS)
% =========================================================================
\section{Electronic Devices: Semiconductors (Unit 19)}

\subsection{Intrinsic \& Extrinsic Semiconductors}

\begin{formulabox}[Carrier Densities, Mass Action Law \& Electrical Conductivity]
\begin{align}
  n_e n_h &= n_i^2 && \text{[Law of Mass Action: valid for both intrinsic and extrinsic]} \\
  n_e &= n_h = n_i && \text{[Intrinsic Semiconductor]} \\
  n_e &\approx N_d \gg n_h && \text{[n-Type Semiconductor: Pentavalent Donor Doping (P, As, Sb)]} \\
  n_h &\approx N_a \gg n_e && \text{[p-Type Semiconductor: Trivalent Acceptor Doping (B, Al, In)]} \\
  \sigma &= e(n_e \mu_e + n_h \mu_h), \qquad I = I_e + I_h = e A E(n_e \mu_e + n_h \mu_h) \quad [\text{Total Drift Current}]
\end{align}
\end{formulabox}

\subsection{p-n Junction Diodes, Rectifiers \& Optoelectronic Devices}

\begin{formulabox}[Diode Dynamic Resistance \& Rectifier Efficiencies]
\begin{align}
  r_d &= \frac{\Delta V}{\Delta I} \quad [\text{Dynamic Resistance: Low in Forward Bias ($\sim 10-100\,\Omega$), High in Reverse Bias ($\sim 10^6\,\Omega$)}] \\
  \eta_{\text{half-wave}} &= \frac{P_{\text{dc}}}{P_{\text{ac}}} = \frac{0.406 R_L}{r_f + R_L} \implies \eta_{\max} = 40.6\%, \qquad \nu_{\text{ripple}} = \nu_{\text{input}} \\
  \eta_{\text{full-wave}} &= \frac{P_{\text{dc}}}{P_{\text{ac}}} = \frac{0.812 R_L}{r_f + R_L} \implies \eta_{\max} = 81.2\%, \qquad \nu_{\text{ripple}} = 2\nu_{\text{input}}
\end{align}
\end{formulabox}

\subsection{Digital Electronics: Logic Gates \& Boolean Algebra}

\begin{table}[htbp]
\centering
\small
\renewcommand{\arraystretch}{1.2}
\begin{tabularx}{\textwidth}{l c c c c c c c}
\toprule
\textbf{Input $A$} & \textbf{Input $B$} & \textbf{AND ($A\cdot B$)} & \textbf{OR ($A+B$)} & \textbf{NAND ($\overline{A\cdot B}$)} & \textbf{NOR ($\overline{A+B}$)} & \textbf{XOR ($A\oplus B$)} & \textbf{XNOR ($\overline{A\oplus B}$)} \\
\midrule
0 & 0 & 0 & 0 & 1 & 1 & 0 & 1 \\
0 & 1 & 0 & 1 & 1 & 0 & 1 & 0 \\
1 & 0 & 0 & 1 & 1 & 0 & 1 & 0 \\
1 & 1 & 1 & 1 & 0 & 0 & 0 & 1 \\
\bottomrule
\end{tabularx}
\caption{Comprehensive Truth Tables for Fundamental, Universal and Derived Logic Gates.}
\end{table}

\begin{usecasebox}[De Morgan's Theorems \& Universal Logic Construction]
\begin{itemize}
  \item \textbf{De Morgan's Laws:}
  \begin{equation*}
    \overline{A + B} = \overline{A}\cdot\overline{B}, \qquad \overline{A\cdot B} = \overline{A} + \overline{B}
  \end{equation*}
  \item \textbf{Universal Gates Minimal Realizations:}
  \begin{itemize}
    \item Number of NAND gates to make: $\text{NOT} = 1$, $\text{AND} = 2$, $\text{OR} = 3$, $\text{NOR} = 4$, $\text{XOR} = 4$.
    \item Number of NOR gates to make: $\text{NOT} = 1$, $\text{OR} = 2$, $\text{AND} = 3$, $\text{NAND} = 4$, $\text{XNOR} = 4$.
  \end{itemize}
\end{itemize}
\end{usecasebox}

\vspace{15pt}
\begin{center}
  \rule{0.6\textwidth}{0.8pt}\\[6pt]
  {\small\textbf{\color{primarydark}End of Module 04 Formula Handbook (Complete Physics Series)}}\\[2pt]
  {\footnotesize\color{accentblue}Prepared, Compiled \& Authored by Satya Prakash \quad\textbar\quad NEET UG 2027 Series}
\end{center}

% =========================================================================
% BACK MATTER
% =========================================================================
\backmatter

\chapter*{About the Author \& Errata Submissions}
\addcontentsline{toc}{chapter}{About the Author \& Errata Submissions}

\textbf{Satya Prakash} is an educator, curriculum researcher, and author dedicated to creating rigorous, student-first learning resources for competitive examinations in India. His focus is on eliminating cognitive overload, dismantling obsolete legacy syllabi, and restoring conceptual rigor and mathematical transparency to competitive pre-medical preparation.

\section*{Official Project Repository \& Continuous Improvement}

Science education thrives on continuous scrutiny and collective refinement. If you identify a typo, numerical discrepancy, or pedagogical ambiguity, or if you wish to suggest an enhancement for future editions, please submit your feedback:

\begin{itemize}
  \item \textbf{Direct Author Email:} \href{mailto:spdedsec@hotmail.com}{spdedsec@hotmail.com}
  \item \textbf{GitHub Open-Source Repository:} \url{https://github.com/spdedsec/neet-study-materials}
  \item \textbf{Errata Tracker:} Verified corrections and updates are posted directly to the project repository.
\end{itemize}

\vspace{1.5cm}
\begin{center}
  \rule{0.6\textwidth}{0.6pt}\\[6pt]
  {\small\textbf{\color{primarydark}NEET UG 2027 Master Reference Series}}\\[2pt]
  {\footnotesize\color{gray}Authored by Satya Prakash $\cdot$ All Rights Reserved}
\end{center}

\end{document}
